% Generated human-review packet template.  Source records, Lean statements,
% and raw-source-to-Spec screening fields are substituted by
% scripts/review_dashboard_packet.py.
% Do not hand-edit generated packets; annotate the PDF or reconcile notes into
% the paper-local human-review log.
\documentclass[10pt]{article}
\usepackage[margin=0.43in,footskip=0.2in]{geometry}
\usepackage{fontspec}
% Use a Unicode-complete main face as well as a Unicode-complete code face.
% Reviewer reasons and source labels can contain exact Greek symbols outside
% verbatim blocks; silently dropping one would corrupt the review packet.
\setmainfont{DejaVu Serif}
% The broad DejaVu Sans Unicode coverage preserves Lean's mathematical glyphs
% (including U+2216 set difference and superscript identifiers) in verbatim
% review blocks.  Its proportional metrics are preferable to silently missing
% exact semantic text from a narrower monospace font.
\setmonofont{DejaVu Sans}
\usepackage{hyperref}
\usepackage{amssymb}
\usepackage{fvextra}
\usepackage{enumitem}
\setlength{\parindent}{0pt}
\setlength{\parskip}{0.15em}
\linespread{0.92}
\hypersetup{colorlinks=true,linkcolor=blue,urlcolor=blue,breaklinks=true}
\DefineVerbatimEnvironment{ReviewVerbatim}{Verbatim}{breaklines=true,breakanywhere=true,fontsize=\scriptsize}
\newcommand{\reviewerbox}{%
  \par\noindent\fbox{\parbox{0.96\linewidth}{\rule{0pt}{0.38in}}}\par
}
\newcommand{\reviewmatch}[1]{%
  \par\noindent\textbf{Reviewer check:} $\square$\; #1\par
  \noindent\emph{If left unchecked, explain the discrepancy or uncertainty below.}\par
}

\begin{document}
\fontsize{9}{10}\selectfont

\begin{center}
{\LARGE Human review packet: PZMH20PerformativePrediction}\\[0.35em]
\small Generated from byte-pinned source excerpts, the expanded PaperInterface
specifications, and hash-bound raw-source-to-Spec screening records.
\end{center}

\noindent\textbf{Paper:} PZMH20PerformativePrediction\\
\textbf{Paper id:} \texttt{PZMH20PerformativePrediction}\\
\textbf{Source version:} \parbox[t]{0.76\linewidth}{\raggedright PMLR 119 official paper text, byte-pinned in cited publication; arXiv v4 retained for documented alternate-proof provenance}\\
\textbf{Source record:} \texttt{audit/paper\_statement\_map.json}\\
\textbf{Rows:} 11\\
\textbf{Generated:} 2026-09-06\\
\textbf{Source URL:} \href{https://proceedings.mlr.press/v119/perdomo20a.html}{official source}





\section*{How to use this packet}
For each source claim, compare the exact source excerpts with the expanded
PaperInterface specification. Mark up this PDF or fill the blank reviewer box in the TeX source;
return the annotated file for reconciliation into the paper-local human-review
log.

\section*{Open the interactive dashboard}
The interactive dashboard is an optional alternative to using this PDF; it
presents the same review surface rather than an additional review requirement.
After forking and cloning (or pulling) AppliedModelingLib, run the following from the
repository root. The cache command is needed only the first time in a new clone.
\begin{ReviewVerbatim}
cd <your-repository-checkout>
lake exe cache get
python3 scripts/review_dashboard.py --paper PZMH20PerformativePrediction --serve
\end{ReviewVerbatim}
Open \url{http://127.0.0.1:8765/} in a browser and keep that terminal running
while reviewing. The dashboard records saved annotations in this paper's local
reviewer-owned review record.

\section*{Contents}
\begin{itemize}[leftmargin=1.2em]
\item \hyperlink{governing-declaration-section-5ef5ef0364b6939c}{Governing Lean declarations (136; supporting context)}
\item \hyperlink{library-section-858f18808007153a}{Semantic library prerequisites (14; not paper claims)}
\begin{itemize}[leftmargin=1.2em]
\item \hyperlink{library-prerequisite-2e05b739f5b71f23}{Applied\allowbreak{}Modeling\allowbreak{}Lib.\allowbreak{}Is\allowbreak{}Jointly\allowbreak{}Smooth\allowbreak{}Gradient}
\item \hyperlink{library-prerequisite-96553d4fbb82f1ac}{Applied\allowbreak{}Modeling\allowbreak{}Lib.\allowbreak{}Measure\allowbreak{}Performative\allowbreak{}Model}
\item \hyperlink{library-prerequisite-49496bb3a06c52e5}{Applied\allowbreak{}Modeling\allowbreak{}Lib.\allowbreak{}Is\allowbreak{}Measure\allowbreak{}Data\allowbreak{}Loss\allowbreak{}Lipschitz}
\item \hyperlink{library-prerequisite-b41cbb32daa3dce7}{Applied\allowbreak{}Modeling\allowbreak{}Lib.\allowbreak{}Is\allowbreak{}Measure\allowbreak{}Pointwise\allowbreak{}Gradient\allowbreak{}Strongly\allowbreak{}Convex}
\item \hyperlink{library-prerequisite-3c6dd4ad7ac116e3}{Applied\allowbreak{}Modeling\allowbreak{}Lib.\allowbreak{}Is\allowbreak{}Measure\allowbreak{}Wasserstein\allowbreak{}Sensitive}
\item \hyperlink{library-prerequisite-e14cc5e8b1e1460e}{Applied\allowbreak{}Modeling\allowbreak{}Lib.\allowbreak{}Measure\allowbreak{}Performative\allowbreak{}Model\allowbreak{}On}
\item \hyperlink{library-prerequisite-0de79b655cc05174}{Applied\allowbreak{}Modeling\allowbreak{}Lib.\allowbreak{}Measure\allowbreak{}Performative\allowbreak{}Model\allowbreak{}On.\allowbreak{}Is\allowbreak{}Data\allowbreak{}Loss\allowbreak{}Lipschitz}
\item \hyperlink{library-prerequisite-25221c4f931beacb}{Applied\allowbreak{}Modeling\allowbreak{}Lib.\allowbreak{}Measure\allowbreak{}Performative\allowbreak{}Model\allowbreak{}On.\allowbreak{}decoupled\allowbreak{}Performative\allowbreak{}Risk}
\item \hyperlink{library-prerequisite-8eef8d8ba05e7ebe}{Applied\allowbreak{}Modeling\allowbreak{}Lib.\allowbreak{}Measure\allowbreak{}Performative\allowbreak{}Model\allowbreak{}On.\allowbreak{}performative\allowbreak{}Risk}
\item \hyperlink{library-prerequisite-555f831d5bfa2a27}{Applied\allowbreak{}Modeling\allowbreak{}Lib.\allowbreak{}Measure\allowbreak{}Performative\allowbreak{}Model\allowbreak{}On.\allowbreak{}Is\allowbreak{}Performatively\allowbreak{}Optimal}
\item \hyperlink{library-prerequisite-4af6c70f954a80a0}{Applied\allowbreak{}Modeling\allowbreak{}Lib.\allowbreak{}Measure\allowbreak{}Performative\allowbreak{}Model\allowbreak{}On.\allowbreak{}Is\allowbreak{}Performatively\allowbreak{}Stable}
\item \hyperlink{library-prerequisite-eb9ab8e4235e1d28}{Applied\allowbreak{}Modeling\allowbreak{}Lib.\allowbreak{}Measure\allowbreak{}Performative\allowbreak{}Model\allowbreak{}On.\allowbreak{}Is\allowbreak{}Pointwise\allowbreak{}Gradient\allowbreak{}Strongly\allowbreak{}Convex}
\item \hyperlink{library-prerequisite-0e1e192eee739a4b}{Applied\allowbreak{}Modeling\allowbreak{}Lib.\allowbreak{}Measure\allowbreak{}Performative\allowbreak{}Model\allowbreak{}On.\allowbreak{}Is\allowbreak{}Wasserstein\allowbreak{}Sensitive}
\item \hyperlink{library-prerequisite-d4de05e57b5947f0}{Applied\allowbreak{}Modeling\allowbreak{}Lib.\allowbreak{}measure\allowbreak{}Decoupled\allowbreak{}Performative\allowbreak{}Risk}
\end{itemize}
\item \hyperlink{paper-prerequisite-section-5ef5ef0364b6939c}{Paper-specific semantic prerequisites (6; not paper claims)}
\begin{itemize}[leftmargin=1.2em]
\item \hyperlink{paper-prerequisite-48fa8b1c1ddbaa9c}{PZMH20\allowbreak{}Performative\allowbreak{}Prediction.\allowbreak{}Is\allowbreak{}Measure\allowbreak{}Performatively\allowbreak{}Optimal}
\item \hyperlink{paper-prerequisite-28cbc41dc9d12634}{PZMH20\allowbreak{}Performative\allowbreak{}Prediction.\allowbreak{}Is\allowbreak{}Measure\allowbreak{}RRM\allowbreak{}On}
\item \hyperlink{paper-prerequisite-f6b3bb0ef1cf7d32}{PZMH20\allowbreak{}Performative\allowbreak{}Prediction.\allowbreak{}Is\allowbreak{}Measure\allowbreak{}Stable}
\item \hyperlink{paper-prerequisite-2ea2f5769822c5c4}{PZMH20\allowbreak{}Performative\allowbreak{}Prediction.\allowbreak{}Is\allowbreak{}Measure\allowbreak{}Wasserstein\allowbreak{}Sensitive}
\item \hyperlink{paper-prerequisite-613020783cc01a15}{PZMH20\allowbreak{}Performative\allowbreak{}Prediction.\allowbreak{}Theorem310\allowbreak{}Corrected\allowbreak{}RERM\allowbreak{}All\allowbreak{}Round\allowbreak{}Trajectory\allowbreak{}Data}
\item \hyperlink{paper-prerequisite-5a443fd08b184a8b}{PZMH20\allowbreak{}Performative\allowbreak{}Prediction.\allowbreak{}Theorem310\allowbreak{}Corrected\allowbreak{}RGD\allowbreak{}All\allowbreak{}Round\allowbreak{}Trajectory\allowbreak{}Data}
\end{itemize}
\item \hyperlink{source-section-0c7af8265958a9ea}{Main-text source claims (11)}
\begin{itemize}[leftmargin=1.2em]
\item \hyperlink{source-claim-5b4f70cb8ea743cb}{Theorem 3.\allowbreak{}5}
\item \hyperlink{source-claim-d10117a376c0a20b}{Theorem 3.\allowbreak{}8}
\item \hyperlink{source-claim-bb0077763720ebbf}{Proposition 4.\allowbreak{}1}
\item \hyperlink{source-claim-a3115b33fe478ce6}{Proposition 4.\allowbreak{}2}
\item \hyperlink{source-claim-cef541336c7da1ef}{Theorem 4.\allowbreak{}3}
\item \hyperlink{source-claim-a22662c7824bd286}{Corollary 5.\allowbreak{}1}
\item \hyperlink{source-claim-1127efac5a38e232}{Proposition 3.\allowbreak{}6}
\item \hyperlink{source-claim-c4686335fcfce867}{Proposition 3.\allowbreak{}6}
\item \hyperlink{source-claim-e1c86b604c5fef5c}{Proposition 3.\allowbreak{}6}
\item \hyperlink{source-claim-a27351d86a4be6d9}{Theorem 3.\allowbreak{}10}
\item \hyperlink{source-claim-2e42d97b744932d5}{Theorem 3.\allowbreak{}10}
\end{itemize}
\end{itemize}

\clearpage
\hypertarget{governing-declaration-section-5ef5ef0364b6939c}{}
\section*{Governing Lean declarations}
These exact graph-bound declaration bodies are retained by the displayed specifications or reviewed prerequisites. They are supporting context and do not add source-result rows or semantic judgments.
\hypertarget{governing-declaration-72513ba85c27319f}{}
\subsection*{\small\ttfamily\raggedright Applied\allowbreak{}Modeling\allowbreak{}Lib.\allowbreak{}Finite\allowbreak{}Coupling}
\textbf{Declaration location:} reusable library
\paragraph{Lean declaration body}
\begin{ReviewVerbatim}
field law:
{α : Type u_1} → {β : Type u_2} → {μ : PMF α} → {ν : PMF β} → AppliedModelingLib.FiniteCoupling μ ν → PMF (α × β)

field fst_marginal:
∀ {α : Type u_1} {β : Type u_2} {μ : PMF α} {ν : PMF β} (self : AppliedModelingLib.FiniteCoupling μ ν),
  PMF.map Prod.fst self.law = μ

field snd_marginal:
∀ {α : Type u_1} {β : Type u_2} {μ : PMF α} {ν : PMF β} (self : AppliedModelingLib.FiniteCoupling μ ν),
  PMF.map Prod.snd self.law = ν
\end{ReviewVerbatim}


\hypertarget{governing-declaration-24472be40793a412}{}
\subsection*{\small\ttfamily\raggedright Applied\allowbreak{}Modeling\allowbreak{}Lib.\allowbreak{}Finite\allowbreak{}Performative\allowbreak{}Model}
\textbf{Declaration location:} reusable library
\paragraph{Lean declaration body}
\begin{ReviewVerbatim}
field dataLaw:
{Parameter : Type u_1} →
  {Data : Type u_2} →
    [inst : Fintype Data] →
      [inst_1 : DecidableEq Data] → AppliedModelingLib.FinitePerformativeModel Parameter Data → Parameter → PMF Data

field loss:
{Parameter : Type u_1} →
  {Data : Type u_2} →
    [inst : Fintype Data] →
      [inst_1 : DecidableEq Data] → AppliedModelingLib.FinitePerformativeModel Parameter Data → Data → Parameter → ℝ
\end{ReviewVerbatim}


\hypertarget{governing-declaration-307f711c4bf618c0}{}
\subsection*{\small\ttfamily\raggedright Applied\allowbreak{}Modeling\allowbreak{}Lib.\allowbreak{}Heterogeneous\allowbreak{}Batch\allowbreak{}Trace}
\textbf{Declaration location:} reusable library
\paragraph{Lean declaration body}
\begin{ReviewVerbatim}
type:
(ℕ → Type u_1) → Type u_2 → ℕ → Type (max u_1 u_2)

value:
fun (BatchIndex : ℕ → Type u_1) (Data : Type u_2) (iteration : ℕ) => (round : Fin iteration) → BatchIndex ↑round → Data
\end{ReviewVerbatim}


\hypertarget{governing-declaration-37d7ae049de215a7}{}
\subsection*{\small\ttfamily\raggedright Applied\allowbreak{}Modeling\allowbreak{}Lib.\allowbreak{}Heterogeneous\allowbreak{}Batch\allowbreak{}Trace\allowbreak{}Coordinate}
\textbf{Declaration location:} reusable library
\paragraph{Lean declaration body}
\begin{ReviewVerbatim}
type:
(ℕ → Type u_1) → Type u_2 → ℕ → Type (max u_1 u_2)

value:
fun (BatchIndex : ℕ → Type u_1) (Data : Type u_2) (a : ℕ) =>
  match a with
  | 0 => PUnit.{(max u_1 u_2) + 1}
  | round.succ => BatchIndex round → Data
\end{ReviewVerbatim}


\hypertarget{governing-declaration-698d9e6a0d26ec64}{}
\subsection*{\small\ttfamily\raggedright Applied\allowbreak{}Modeling\allowbreak{}Lib.\allowbreak{}Is\allowbreak{}Data\allowbreak{}Gradient\allowbreak{}Lipschitz}
\textbf{Declaration location:} reusable library
\paragraph{Lean declaration body}
\begin{ReviewVerbatim}
type:
{Parameter : Type u_1} →
  {Data : Type u_2} →
    [inst : NormedAddCommGroup Parameter] →
      [InnerProductSpace ℝ Parameter] → [MetricSpace Data] → (Data → Parameter → Parameter) → NNReal → Prop

value:
fun {Parameter : Type u_1} {Data : Type u_2} [NormedAddCommGroup Parameter] [InnerProductSpace ℝ Parameter]
    [MetricSpace Data] (gradient : Data → Parameter → Parameter) (smoothness : NNReal) =>
  ∀ (first second : Data) (parameter : Parameter),
    ‖gradient first parameter - gradient second parameter‖ ≤ ↑smoothness * dist first second
\end{ReviewVerbatim}


\hypertarget{governing-declaration-41e49a877b22dbe7}{}
\subsection*{\small\ttfamily\raggedright Applied\allowbreak{}Modeling\allowbreak{}Lib.\allowbreak{}Is\allowbreak{}Variational\allowbreak{}Euclidean\allowbreak{}Projection\allowbreak{}On}
\textbf{Declaration location:} reusable library
\paragraph{Lean declaration body}
\begin{ReviewVerbatim}
type:
{Parameter : Type u_1} →
  [inst : NormedAddCommGroup Parameter] →
    [InnerProductSpace ℝ Parameter] → Set Parameter → (Parameter → Parameter) → Prop

value:
fun {Parameter : Type u_1} [NormedAddCommGroup Parameter] [InnerProductSpace ℝ Parameter] (domain : Set Parameter)
    (project : Parameter → Parameter) =>
  (∀ (input : Parameter), project input ∈ domain) ∧
    ∀ (input candidate : Parameter),
      candidate ∈ domain → inner ℝ (input - project input) (candidate - project input) ≤ 0
\end{ReviewVerbatim}


\hypertarget{governing-declaration-b849531e0a75a336}{}
\subsection*{\small\ttfamily\raggedright Applied\allowbreak{}Modeling\allowbreak{}Lib.\allowbreak{}Math.\allowbreak{}dyadic\allowbreak{}Effective\allowbreak{}Count\allowbreak{}Depth}
\textbf{Declaration location:} reusable library
\paragraph{Lean declaration body}
\begin{ReviewVerbatim}
type:
ℕ → ℝ → ℕ

value:
fun (dimension : ℕ) (effectiveCount : ℝ) => ⌊Real.logb (↑(2 ^ dimension)) (effectiveCount / ↑(2 ^ dimension))⌋₊
\end{ReviewVerbatim}


\hypertarget{governing-declaration-4f196894fb58689f}{}
\subsection*{\small\ttfamily\raggedright Applied\allowbreak{}Modeling\allowbreak{}Lib.\allowbreak{}Math.\allowbreak{}dyadic\allowbreak{}Geometric\allowbreak{}Effective\allowbreak{}Depth}
\textbf{Declaration location:} reusable library
\paragraph{Lean declaration body}
\begin{ReviewVerbatim}
type:
ℕ → ℝ → ℝ → ℕ

value:
fun (dimension : ℕ) (effectiveCount confidence : ℝ) =>
  if 16 * ↑(2 ^ dimension) ≤ effectiveCount ∧ 8 * confidence ≤ effectiveCount then
    AppliedModelingLib.Math.dyadicEffectiveCountDepth dimension (effectiveCount / 16)
  else 0
\end{ReviewVerbatim}

\paragraph{Direct retained dependencies}
\begin{itemize}\raggedright
\item \hyperlink{governing-declaration-b849531e0a75a336}{Applied\allowbreak{}Modeling\allowbreak{}Lib.\allowbreak{}Math.\allowbreak{}dyadic\allowbreak{}Effective\allowbreak{}Count\allowbreak{}Depth}
\end{itemize}
\hypertarget{governing-declaration-77c24384a8319860}{}
\subsection*{\small\ttfamily\raggedright Applied\allowbreak{}Modeling\allowbreak{}Lib.\allowbreak{}Math.\allowbreak{}exp\allowbreak{}Confidence\allowbreak{}For\allowbreak{}Half\allowbreak{}Budget}
\textbf{Declaration location:} reusable library
\paragraph{Lean declaration body}
\begin{ReviewVerbatim}
type:
ℝ → ℝ

value:
fun (budget : ℝ) => max 0 (Real.log (2 / budget))
\end{ReviewVerbatim}


\hypertarget{governing-declaration-01ad3e76a1f28ec0}{}
\subsection*{\small\ttfamily\raggedright Applied\allowbreak{}Modeling\allowbreak{}Lib.\allowbreak{}Math.\allowbreak{}positive\allowbreak{}Nat\allowbreak{}Ceil}
\textbf{Declaration location:} reusable library
\paragraph{Lean declaration body}
\begin{ReviewVerbatim}
type:
ℝ → ℕ

value:
fun (threshold : ℝ) => ⌈max 1 threshold⌉₊
\end{ReviewVerbatim}


\hypertarget{governing-declaration-ef89f6befb052e48}{}
\subsection*{\small\ttfamily\raggedright Applied\allowbreak{}Modeling\allowbreak{}Lib.\allowbreak{}Is\allowbreak{}Measure\allowbreak{}Parameter\allowbreak{}Loss\allowbreak{}Lipschitz}
\textbf{Declaration location:} reusable library
\paragraph{Lean declaration body}
\begin{ReviewVerbatim}
type:
{Parameter : Type u_1} →
  {Data : Type u_2} →
    [NormedAddCommGroup Parameter] →
      [inst : MeasurableSpace Data] → AppliedModelingLib.MeasurePerformativeModel Parameter Data → NNReal → Prop

value:
fun {Parameter : Type u_1} {Data : Type u_2} [NormedAddCommGroup Parameter] [MeasurableSpace Data]
    (model : AppliedModelingLib.MeasurePerformativeModel Parameter Data) (constant : NNReal) =>
  ∀ (datum : Data), LipschitzWith constant (model.loss datum)
\end{ReviewVerbatim}

\paragraph{Direct retained dependencies}
\begin{itemize}\raggedright
\item \hyperlink{library-prerequisite-96553d4fbb82f1ac}{Applied\allowbreak{}Modeling\allowbreak{}Lib.\allowbreak{}Measure\allowbreak{}Performative\allowbreak{}Model}
\end{itemize}
\hypertarget{governing-declaration-527443e0fb077287}{}
\subsection*{\small\ttfamily\raggedright Applied\allowbreak{}Modeling\allowbreak{}Lib.\allowbreak{}Measure\allowbreak{}Performative\allowbreak{}Sampling\allowbreak{}Kernel}
\textbf{Declaration location:} reusable library
\paragraph{Lean declaration body}
\begin{ReviewVerbatim}
field kernel:
{Parameter : Type u_1} →
  {Data : Type u_2} →
    [inst : MeasurableSpace Parameter] →
      [inst_1 : MeasurableSpace Data] →
        {model : AppliedModelingLib.MeasurePerformativeModel Parameter Data} →
          AppliedModelingLib.MeasurePerformativeSamplingKernel model → ProbabilityTheory.Kernel Parameter Data

field isMarkov:
∀ {Parameter : Type u_1} {Data : Type u_2} [inst : MeasurableSpace Parameter] [inst_1 : MeasurableSpace Data]
  {model : AppliedModelingLib.MeasurePerformativeModel Parameter Data}
  (self : AppliedModelingLib.MeasurePerformativeSamplingKernel model), ProbabilityTheory.IsMarkovKernel self.kernel

field apply_eq_dataLaw:
∀ {Parameter : Type u_1} {Data : Type u_2} [inst : MeasurableSpace Parameter] [inst_1 : MeasurableSpace Data]
  {model : AppliedModelingLib.MeasurePerformativeModel Parameter Data}
  (self : AppliedModelingLib.MeasurePerformativeSamplingKernel model) (parameter : Parameter),
  self.kernel parameter = ↑(model.dataLaw parameter)
\end{ReviewVerbatim}

\paragraph{Direct retained dependencies}
\begin{itemize}\raggedright
\item \hyperlink{library-prerequisite-96553d4fbb82f1ac}{Applied\allowbreak{}Modeling\allowbreak{}Lib.\allowbreak{}Measure\allowbreak{}Performative\allowbreak{}Model}
\end{itemize}
\hypertarget{governing-declaration-c1d248907d4e893c}{}
\subsection*{\small\ttfamily\raggedright Applied\allowbreak{}Modeling\allowbreak{}Lib.\allowbreak{}Probability.\allowbreak{}Has\allowbreak{}Exponential\allowbreak{}Radial\allowbreak{}Moment}
\textbf{Declaration location:} reusable library
\paragraph{Lean declaration body}
\begin{ReviewVerbatim}
type:
{E : Type u_1} → [inst : MeasurableSpace E] → [NormedAddCommGroup E] → MeasureTheory.ProbabilityMeasure E → ℝ → ℝ → Prop

value:
fun {E : Type u_1} [MeasurableSpace E] [NormedAddCommGroup E] (law : MeasureTheory.ProbabilityMeasure E)
    (alpha gamma : ℝ) =>
  MeasureTheory.Integrable (fun (x : E) => Real.exp (gamma * ‖x‖.rpow alpha)) ↑law
\end{ReviewVerbatim}


\hypertarget{governing-declaration-9604ed56aea35d91}{}
\subsection*{\small\ttfamily\raggedright Applied\allowbreak{}Modeling\allowbreak{}Lib.\allowbreak{}Probability.\allowbreak{}Has\allowbreak{}Bounded\allowbreak{}Exponential\allowbreak{}Radial\allowbreak{}Moment}
\textbf{Declaration location:} reusable library
\paragraph{Lean declaration body}
\begin{ReviewVerbatim}
type:
{E : Type u_1} →
  [inst : MeasurableSpace E] → [NormedAddCommGroup E] → MeasureTheory.ProbabilityMeasure E → ℝ → ℝ → ℝ → Prop

value:
fun {E : Type u_1} [MeasurableSpace E] [NormedAddCommGroup E] (law : MeasureTheory.ProbabilityMeasure E)
    (alpha gamma momentBound : ℝ) =>
  AppliedModelingLib.Probability.HasExponentialRadialMoment law alpha gamma ∧
    ∫ (x : E), Real.exp (gamma * ‖x‖.rpow alpha) ∂↑law ≤ momentBound
\end{ReviewVerbatim}

\paragraph{Direct retained dependencies}
\begin{itemize}\raggedright
\item \hyperlink{governing-declaration-c1d248907d4e893c}{Applied\allowbreak{}Modeling\allowbreak{}Lib.\allowbreak{}Probability.\allowbreak{}Has\allowbreak{}Exponential\allowbreak{}Radial\allowbreak{}Moment}
\end{itemize}
\hypertarget{governing-declaration-d5b34777e35d2553}{}
\subsection*{\small\ttfamily\raggedright Applied\allowbreak{}Modeling\allowbreak{}Lib.\allowbreak{}Probability.\allowbreak{}dyadic\allowbreak{}Cell\allowbreak{}Index}
\textbf{Declaration location:} reusable library
\paragraph{Lean declaration body}
\begin{ReviewVerbatim}
type:
ℕ → Type

value:
fun (level : ℕ) => Fin (2 ^ level)
\end{ReviewVerbatim}


\hypertarget{governing-declaration-24e8f7da653baa4e}{}
\subsection*{\small\ttfamily\raggedright Applied\allowbreak{}Modeling\allowbreak{}Lib.\allowbreak{}Probability.\allowbreak{}dyadic\allowbreak{}Cube\allowbreak{}Index}
\textbf{Declaration location:} reusable library
\paragraph{Lean declaration body}
\begin{ReviewVerbatim}
type:
ℕ → ℕ → Type

value:
fun (dimension level : ℕ) => Fin dimension → AppliedModelingLib.Probability.dyadicCellIndex level
\end{ReviewVerbatim}

\paragraph{Direct retained dependencies}
\begin{itemize}\raggedright
\item \hyperlink{governing-declaration-d5b34777e35d2553}{Applied\allowbreak{}Modeling\allowbreak{}Lib.\allowbreak{}Probability.\allowbreak{}dyadic\allowbreak{}Cell\allowbreak{}Index}
\end{itemize}
\hypertarget{governing-declaration-67088ec7429e6ab5}{}
\subsection*{\small\ttfamily\raggedright Applied\allowbreak{}Modeling\allowbreak{}Lib.\allowbreak{}Probability.\allowbreak{}euclidean\allowbreak{}Dyadic\allowbreak{}Cube\allowbreak{}Radius}
\textbf{Declaration location:} reusable library
\paragraph{Lean declaration body}
\begin{ReviewVerbatim}
type:
ℕ → ℝ

value:
fun (scale : ℕ) => ↑(2 ^ scale)
\end{ReviewVerbatim}


\hypertarget{governing-declaration-1e6521903c15d24c}{}
\subsection*{\small\ttfamily\raggedright Applied\allowbreak{}Modeling\allowbreak{}Lib.\allowbreak{}Probability.\allowbreak{}euclidean\allowbreak{}Dyadic\allowbreak{}Cube}
\textbf{Declaration location:} reusable library
\paragraph{Lean declaration body}
\begin{ReviewVerbatim}
type:
(dimension : ℕ) → ℕ → EuclideanSpace ℝ (Fin dimension) → Prop

value:
fun (dimension scale : ℕ) (a : EuclideanSpace ℝ (Fin dimension)) =>
  {x : EuclideanSpace ℝ (Fin dimension) |
      ∀ (coordinate : Fin dimension),
        x.ofLp coordinate ∈
          Set.Ioc (-AppliedModelingLib.Probability.euclideanDyadicCubeRadius scale)
            (AppliedModelingLib.Probability.euclideanDyadicCubeRadius scale)}
    a
\end{ReviewVerbatim}

\paragraph{Direct retained dependencies}
\begin{itemize}\raggedright
\item \hyperlink{governing-declaration-67088ec7429e6ab5}{Applied\allowbreak{}Modeling\allowbreak{}Lib.\allowbreak{}Probability.\allowbreak{}euclidean\allowbreak{}Dyadic\allowbreak{}Cube\allowbreak{}Radius}
\end{itemize}
\hypertarget{governing-declaration-13cdf9a43259c30a}{}
\subsection*{\small\ttfamily\raggedright Applied\allowbreak{}Modeling\allowbreak{}Lib.\allowbreak{}Probability.\allowbreak{}euclidean\allowbreak{}Dyadic\allowbreak{}Cube\allowbreak{}Shell}
\textbf{Declaration location:} reusable library
\paragraph{Lean declaration body}
\begin{ReviewVerbatim}
type:
(dimension : ℕ) → ℕ → EuclideanSpace ℝ (Fin dimension) → Prop

value:
fun (dimension shell : ℕ) (a : EuclideanSpace ℝ (Fin dimension)) =>
  (match shell with
    | 0 => AppliedModelingLib.Probability.euclideanDyadicCube dimension 0
    | shell.succ =>
      AppliedModelingLib.Probability.euclideanDyadicCube dimension (shell + 1) \
        AppliedModelingLib.Probability.euclideanDyadicCube dimension shell)
    a
\end{ReviewVerbatim}

\paragraph{Direct retained dependencies}
\begin{itemize}\raggedright
\item \hyperlink{governing-declaration-1e6521903c15d24c}{Applied\allowbreak{}Modeling\allowbreak{}Lib.\allowbreak{}Probability.\allowbreak{}euclidean\allowbreak{}Dyadic\allowbreak{}Cube}
\end{itemize}
\hypertarget{governing-declaration-93c495b865c4b878}{}
\subsection*{\small\ttfamily\raggedright Applied\allowbreak{}Modeling\allowbreak{}Lib.\allowbreak{}Probability.\allowbreak{}fifth\allowbreak{}Order\allowbreak{}Exponential\allowbreak{}Radial\allowbreak{}Tail\allowbreak{}Constant}
\textbf{Declaration location:} reusable library
\paragraph{Lean declaration body}
\begin{ReviewVerbatim}
type:
{E : Type u_1} → [inst : MeasurableSpace E] → [NormedAddCommGroup E] → MeasureTheory.ProbabilityMeasure E → ℝ → ℝ → ℝ

value:
fun {E : Type u_1} [MeasurableSpace E] [NormedAddCommGroup E] (law : MeasureTheory.ProbabilityMeasure E)
    (alpha gamma : ℝ) =>
  ↑(Nat.factorial 5) * gamma⁻¹ ^ 5 * ∫ (x : E), Real.exp (gamma * ‖x‖.rpow alpha) ∂↑law
\end{ReviewVerbatim}


\hypertarget{governing-declaration-4518f5a970874c84}{}
\subsection*{\small\ttfamily\raggedright Applied\allowbreak{}Modeling\allowbreak{}Lib.\allowbreak{}Probability.\allowbreak{}fin\allowbreak{}Snoc\allowbreak{}Measurable\allowbreak{}Equiv}
\textbf{Declaration location:} reusable library
\paragraph{Lean declaration body}
\begin{ReviewVerbatim}
type:
{beta : Type u_2} → [inst : MeasurableSpace beta] → (n : ℕ) → (Fin n → beta) × beta ≃ᵐ (Fin (n + 1) → beta)

value:
fun {beta : Type u_2} [MeasurableSpace beta] (n : ℕ) =>
  MeasurableEquiv.prodComm.trans (MeasurableEquiv.piFinSuccAbove (fun (x : Fin (n + 1)) => beta) (Fin.last n)).symm
\end{ReviewVerbatim}


\hypertarget{governing-declaration-aa3875e0dd44ed11}{}
\subsection*{\small\ttfamily\raggedright Applied\allowbreak{}Modeling\allowbreak{}Lib.\allowbreak{}Probability.\allowbreak{}finite\allowbreak{}IID\allowbreak{}Cell\allowbreak{}Frequency}
\textbf{Declaration location:} reusable library
\paragraph{Lean declaration body}
\begin{ReviewVerbatim}
type:
{E : Type u_1} → [MeasurableSpace E] → (count : ℕ) → Set E → (Fin count → E) → ℝ

value:
fun {E : Type u_1} [MeasurableSpace E] (count : ℕ) (cell : Set E) (sample : Fin count → E) =>
  (∑ index : Fin count, cell.indicator (fun (x : E) => 1) (sample index)) / ↑count
\end{ReviewVerbatim}


\hypertarget{governing-declaration-b8cc820488d73247}{}
\subsection*{\small\ttfamily\raggedright Applied\allowbreak{}Modeling\allowbreak{}Lib.\allowbreak{}Probability.\allowbreak{}finite\allowbreak{}IID\allowbreak{}Euclidean\allowbreak{}Dyadic\allowbreak{}Cube\allowbreak{}Shell\allowbreak{}Mass\allowbreak{}Relative\allowbreak{}Tail}
\textbf{Declaration location:} reusable library
\paragraph{Lean declaration body}
\begin{ReviewVerbatim}
type:
(dimension : ℕ) →
  MeasureTheory.ProbabilityMeasure (EuclideanSpace ℝ (Fin dimension)) →
    (count : ℕ) → ℕ → ℝ → (Fin count → EuclideanSpace ℝ (Fin dimension)) → Prop

value:
fun (dimension : ℕ) (law : MeasureTheory.ProbabilityMeasure (EuclideanSpace ℝ (Fin dimension))) (count shell : ℕ)
    (z : ℝ) (a : Fin count → EuclideanSpace ℝ (Fin dimension)) =>
  {sample : Fin count → EuclideanSpace ℝ (Fin dimension) |
      (↑law).real (AppliedModelingLib.Probability.euclideanDyadicCubeShell dimension shell) * z ≤
        |AppliedModelingLib.Probability.finiteIIDCellFrequency count
              (AppliedModelingLib.Probability.euclideanDyadicCubeShell dimension shell) sample -
            (↑law).real (AppliedModelingLib.Probability.euclideanDyadicCubeShell dimension shell)|}
    a
\end{ReviewVerbatim}

\paragraph{Direct retained dependencies}
\begin{itemize}\raggedright
\item \hyperlink{governing-declaration-13cdf9a43259c30a}{Applied\allowbreak{}Modeling\allowbreak{}Lib.\allowbreak{}Probability.\allowbreak{}euclidean\allowbreak{}Dyadic\allowbreak{}Cube\allowbreak{}Shell}
\item \hyperlink{governing-declaration-aa3875e0dd44ed11}{Applied\allowbreak{}Modeling\allowbreak{}Lib.\allowbreak{}Probability.\allowbreak{}finite\allowbreak{}IID\allowbreak{}Cell\allowbreak{}Frequency}
\end{itemize}
\hypertarget{governing-declaration-dbd328d5a293e61c}{}
\subsection*{\small\ttfamily\raggedright Applied\allowbreak{}Modeling\allowbreak{}Lib.\allowbreak{}Probability.\allowbreak{}finite\allowbreak{}IID\allowbreak{}Euclidean\allowbreak{}Dyadic\allowbreak{}Cube\allowbreak{}Shell\allowbreak{}Mass\allowbreak{}Relative\allowbreak{}Tail\allowbreak{}Totalized}
\textbf{Declaration location:} reusable library
\paragraph{Lean declaration body}
\begin{ReviewVerbatim}
type:
(dimension : ℕ) →
  MeasureTheory.ProbabilityMeasure (EuclideanSpace ℝ (Fin dimension)) →
    (count : ℕ) → ℕ → ℝ → (Fin count → EuclideanSpace ℝ (Fin dimension)) → Prop

value:
fun (dimension : ℕ) (law : MeasureTheory.ProbabilityMeasure (EuclideanSpace ℝ (Fin dimension))) (count shell : ℕ)
    (z : ℝ) (a : Fin count → EuclideanSpace ℝ (Fin dimension)) =>
  (if law.toFiniteMeasure.restrict (AppliedModelingLib.Probability.euclideanDyadicCubeShell dimension shell) = 0 then
      {sample : Fin count → EuclideanSpace ℝ (Fin dimension) |
        ∃ (index : Fin count), sample index ∈ AppliedModelingLib.Probability.euclideanDyadicCubeShell dimension shell}
    else AppliedModelingLib.Probability.finiteIIDEuclideanDyadicCubeShellMassRelativeTail dimension law count shell z)
    a
\end{ReviewVerbatim}

\paragraph{Direct retained dependencies}
\begin{itemize}\raggedright
\item \hyperlink{governing-declaration-13cdf9a43259c30a}{Applied\allowbreak{}Modeling\allowbreak{}Lib.\allowbreak{}Probability.\allowbreak{}euclidean\allowbreak{}Dyadic\allowbreak{}Cube\allowbreak{}Shell}
\item \hyperlink{governing-declaration-b8cc820488d73247}{Applied\allowbreak{}Modeling\allowbreak{}Lib.\allowbreak{}Probability.\allowbreak{}finite\allowbreak{}IID\allowbreak{}Euclidean\allowbreak{}Dyadic\allowbreak{}Cube\allowbreak{}Shell\allowbreak{}Mass\allowbreak{}Relative\allowbreak{}Tail}
\end{itemize}
\hypertarget{governing-declaration-91a777e8ba845265}{}
\subsection*{\small\ttfamily\raggedright Applied\allowbreak{}Modeling\allowbreak{}Lib.\allowbreak{}Probability.\allowbreak{}finite\allowbreak{}IID\allowbreak{}Sample\allowbreak{}Law}
\textbf{Declaration location:} reusable library
\paragraph{Lean declaration body}
\begin{ReviewVerbatim}
type:
{X : Type u_1} →
  [inst : MeasurableSpace X] → MeasureTheory.Measure X → (count : ℕ) → MeasureTheory.Measure (Fin count → X)

value:
fun {X : Type u_1} [MeasurableSpace X] (law : MeasureTheory.Measure X) (count : ℕ) =>
  MeasureTheory.Measure.pi fun (x : Fin count) => law
\end{ReviewVerbatim}


\hypertarget{governing-declaration-548cac74b010e6ce}{}
\subsection*{\small\ttfamily\raggedright Applied\allowbreak{}Modeling\allowbreak{}Lib.\allowbreak{}Probability.\allowbreak{}finite\allowbreak{}IID\allowbreak{}Sample\allowbreak{}Map}
\textbf{Declaration location:} reusable library
\paragraph{Lean declaration body}
\begin{ReviewVerbatim}
type:
{X : Type u_1} → {Y : Type u_2} → {count : ℕ} → (X → Y) → (Fin count → X) → Fin count → Y

value:
fun {X : Type u_1} {Y : Type u_2} {count : ℕ} (f : X → Y) (a : Fin count → X) (a_1 : Fin count) => f (a a_1)
\end{ReviewVerbatim}


\hypertarget{governing-declaration-185ae8e741246e89}{}
\subsection*{\small\ttfamily\raggedright Applied\allowbreak{}Modeling\allowbreak{}Lib.\allowbreak{}Probability.\allowbreak{}finite\allowbreak{}IID\allowbreak{}Sample\allowbreak{}On\allowbreak{}Finset}
\textbf{Declaration location:} reusable library
\paragraph{Lean declaration body}
\begin{ReviewVerbatim}
type:
{X : Type u_1} → {count : ℕ} → (inside : Finset (Fin count)) → (Fin count → X) → ↥inside → X

value:
fun {X : Type u_1} {count : ℕ} (inside : Finset (Fin count)) (a : Fin count → X) (a_1 : ↥inside) => a ↑a_1
\end{ReviewVerbatim}


\hypertarget{governing-declaration-dd7cf8954dcd7cb1}{}
\subsection*{\small\ttfamily\raggedright Applied\allowbreak{}Modeling\allowbreak{}Lib.\allowbreak{}Probability.\allowbreak{}finite\allowbreak{}IID\allowbreak{}Sample\allowbreak{}On\allowbreak{}Finset\allowbreak{}To\allowbreak{}Fin}
\textbf{Declaration location:} reusable library
\paragraph{Lean declaration body}
\begin{ReviewVerbatim}
type:
{X : Type u_1} → {count : ℕ} → (inside : Finset (Fin count)) → (Fin count → X) → Fin inside.card → X

value:
fun {X : Type u_1} {count : ℕ} (inside : Finset (Fin count)) (a : Fin count → X) (a_1 : Fin inside.card) =>
  AppliedModelingLib.Probability.finiteIIDSampleOnFinset inside a
    (((finCongr (AppliedModelingLib.Probability.finiteIIDSampleOnFinsetToFin._proof_1 inside)).trans
        (Fintype.equivFin ↥inside).symm)
      a_1)
\end{ReviewVerbatim}

\paragraph{Direct retained dependencies}
\begin{itemize}\raggedright
\item \hyperlink{governing-declaration-185ae8e741246e89}{Applied\allowbreak{}Modeling\allowbreak{}Lib.\allowbreak{}Probability.\allowbreak{}finite\allowbreak{}IID\allowbreak{}Sample\allowbreak{}On\allowbreak{}Finset}
\end{itemize}
\hypertarget{governing-declaration-8c9ed2e122e54528}{}
\subsection*{\small\ttfamily\raggedright Applied\allowbreak{}Modeling\allowbreak{}Lib.\allowbreak{}Probability.\allowbreak{}finite\allowbreak{}Kernel\allowbreak{}Product}
\textbf{Declaration location:} reusable library
\paragraph{Lean declaration body}
\begin{ReviewVerbatim}
type:
{alpha : Type u_1} →
  {beta : Type u_2} →
    [inst : MeasurableSpace alpha] →
      [inst_1 : MeasurableSpace beta] →
        (n : ℕ) → (Fin n → ProbabilityTheory.Kernel alpha beta) → ProbabilityTheory.Kernel alpha (Fin n → beta)

value:
fun {alpha : Type u_1} {beta : Type u_2} [MeasurableSpace alpha] [MeasurableSpace beta] (n : ℕ)
    (a : Fin n → ProbabilityTheory.Kernel alpha beta) =>
  Nat.brecOn (motive := fun (x : ℕ) =>
    (Fin x → ProbabilityTheory.Kernel alpha beta) → ProbabilityTheory.Kernel alpha (Fin x → beta)) n
    AppliedModelingLib.Probability.finiteKernelProduct._f a
\end{ReviewVerbatim}

\paragraph{Direct retained dependencies}
\begin{itemize}\raggedright
\item \hyperlink{governing-declaration-4518f5a970874c84}{Applied\allowbreak{}Modeling\allowbreak{}Lib.\allowbreak{}Probability.\allowbreak{}fin\allowbreak{}Snoc\allowbreak{}Measurable\allowbreak{}Equiv}
\end{itemize}
\hypertarget{governing-declaration-8ef9c28a526f8794}{}
\subsection*{\small\ttfamily\raggedright Applied\allowbreak{}Modeling\allowbreak{}Lib.\allowbreak{}Probability.\allowbreak{}finite\allowbreak{}IID\allowbreak{}Batch\allowbreak{}Kernel}
\textbf{Declaration location:} reusable library
\paragraph{Lean declaration body}
\begin{ReviewVerbatim}
type:
{alpha : Type u_1} →
  {beta : Type u_2} →
    [inst : MeasurableSpace alpha] →
      [inst_1 : MeasurableSpace beta] →
        (n : ℕ) → ProbabilityTheory.Kernel alpha beta → ProbabilityTheory.Kernel alpha (Fin n → beta)

value:
fun {alpha : Type u_1} {beta : Type u_2} [MeasurableSpace alpha] [MeasurableSpace beta] (n : ℕ)
    (kappa : ProbabilityTheory.Kernel alpha beta) =>
  AppliedModelingLib.Probability.finiteKernelProduct n fun (x : Fin n) => kappa
\end{ReviewVerbatim}

\paragraph{Direct retained dependencies}
\begin{itemize}\raggedright
\item \hyperlink{governing-declaration-8c9ed2e122e54528}{Applied\allowbreak{}Modeling\allowbreak{}Lib.\allowbreak{}Probability.\allowbreak{}finite\allowbreak{}Kernel\allowbreak{}Product}
\end{itemize}
\hypertarget{governing-declaration-5eb127d4a6b534fc}{}
\subsection*{\small\ttfamily\raggedright Applied\allowbreak{}Modeling\allowbreak{}Lib.\allowbreak{}Probability.\allowbreak{}half\allowbreak{}Open\allowbreak{}Unit\allowbreak{}Interval}
\textbf{Declaration location:} reusable library
\paragraph{Lean declaration body}
\begin{ReviewVerbatim}
type:
ℝ → Prop

value:
fun (a : ℝ) => Set.Ioc (-1) 1 a
\end{ReviewVerbatim}


\hypertarget{governing-declaration-68be765167e02884}{}
\subsection*{\small\ttfamily\raggedright Applied\allowbreak{}Modeling\allowbreak{}Lib.\allowbreak{}Probability.\allowbreak{}Half\allowbreak{}Open\allowbreak{}Unit\allowbreak{}Cube}
\textbf{Declaration location:} reusable library
\paragraph{Lean declaration body}
\begin{ReviewVerbatim}
type:
ℕ → Type

value:
fun (dimension : ℕ) =>
  { x : EuclideanSpace ℝ (Fin dimension) //
    ∀ (i : Fin dimension), x.ofLp i ∈ AppliedModelingLib.Probability.halfOpenUnitInterval }
\end{ReviewVerbatim}

\paragraph{Direct retained dependencies}
\begin{itemize}\raggedright
\item \hyperlink{governing-declaration-5eb127d4a6b534fc}{Applied\allowbreak{}Modeling\allowbreak{}Lib.\allowbreak{}Probability.\allowbreak{}half\allowbreak{}Open\allowbreak{}Unit\allowbreak{}Interval}
\end{itemize}
\hypertarget{governing-declaration-149e3d31533a6cd4}{}
\subsection*{\small\ttfamily\raggedright Applied\allowbreak{}Modeling\allowbreak{}Lib.\allowbreak{}Probability.\allowbreak{}euclidean\allowbreak{}Dyadic\allowbreak{}Cube\allowbreak{}Shell\allowbreak{}Scale}
\textbf{Declaration location:} reusable library
\paragraph{Lean declaration body}
\begin{ReviewVerbatim}
type:
(dimension shell : ℕ) →
  ↑(AppliedModelingLib.Probability.euclideanDyadicCubeShell dimension shell) →
    AppliedModelingLib.Probability.HalfOpenUnitCube dimension

value:
fun (dimension shell : ℕ) (a : ↑(AppliedModelingLib.Probability.euclideanDyadicCubeShell dimension shell)) =>
  ⟨(AppliedModelingLib.Probability.euclideanDyadicCubeRadius shell)⁻¹ • ↑a,
    AppliedModelingLib.Probability.euclideanDyadicCubeShellScale._proof_1 dimension shell a⟩
\end{ReviewVerbatim}

\paragraph{Direct retained dependencies}
\begin{itemize}\raggedright
\item \hyperlink{governing-declaration-68be765167e02884}{Applied\allowbreak{}Modeling\allowbreak{}Lib.\allowbreak{}Probability.\allowbreak{}Half\allowbreak{}Open\allowbreak{}Unit\allowbreak{}Cube}
\item \hyperlink{governing-declaration-67088ec7429e6ab5}{Applied\allowbreak{}Modeling\allowbreak{}Lib.\allowbreak{}Probability.\allowbreak{}euclidean\allowbreak{}Dyadic\allowbreak{}Cube\allowbreak{}Radius}
\item \hyperlink{governing-declaration-13cdf9a43259c30a}{Applied\allowbreak{}Modeling\allowbreak{}Lib.\allowbreak{}Probability.\allowbreak{}euclidean\allowbreak{}Dyadic\allowbreak{}Cube\allowbreak{}Shell}
\item \hyperlink{governing-declaration-5eb127d4a6b534fc}{Applied\allowbreak{}Modeling\allowbreak{}Lib.\allowbreak{}Probability.\allowbreak{}half\allowbreak{}Open\allowbreak{}Unit\allowbreak{}Interval}
\end{itemize}
\hypertarget{governing-declaration-4bcc855202b49a28}{}
\subsection*{\small\ttfamily\raggedright Applied\allowbreak{}Modeling\allowbreak{}Lib.\allowbreak{}Probability.\allowbreak{}euclidean\allowbreak{}Dyadic\allowbreak{}Cube\allowbreak{}Shell\allowbreak{}Scale\allowbreak{}Extension}
\textbf{Declaration location:} reusable library
\paragraph{Lean declaration body}
\begin{ReviewVerbatim}
type:
(dimension : ℕ) → ℕ → EuclideanSpace ℝ (Fin dimension) → AppliedModelingLib.Probability.HalfOpenUnitCube dimension

value:
fun (dimension shell : ℕ) (a : EuclideanSpace ℝ (Fin dimension)) =>
  Classical.choose
    (AppliedModelingLib.Probability.exists_measurable_euclideanDyadicCubeShellScaleExtension dimension shell) a
\end{ReviewVerbatim}

\paragraph{Direct retained dependencies}
\begin{itemize}\raggedright
\item \hyperlink{governing-declaration-68be765167e02884}{Applied\allowbreak{}Modeling\allowbreak{}Lib.\allowbreak{}Probability.\allowbreak{}Half\allowbreak{}Open\allowbreak{}Unit\allowbreak{}Cube}
\item \hyperlink{governing-declaration-13cdf9a43259c30a}{Applied\allowbreak{}Modeling\allowbreak{}Lib.\allowbreak{}Probability.\allowbreak{}euclidean\allowbreak{}Dyadic\allowbreak{}Cube\allowbreak{}Shell}
\item \hyperlink{governing-declaration-149e3d31533a6cd4}{Applied\allowbreak{}Modeling\allowbreak{}Lib.\allowbreak{}Probability.\allowbreak{}euclidean\allowbreak{}Dyadic\allowbreak{}Cube\allowbreak{}Shell\allowbreak{}Scale}
\item \hyperlink{governing-declaration-5eb127d4a6b534fc}{Applied\allowbreak{}Modeling\allowbreak{}Lib.\allowbreak{}Probability.\allowbreak{}half\allowbreak{}Open\allowbreak{}Unit\allowbreak{}Interval}
\end{itemize}
\hypertarget{governing-declaration-a45967b04b4e794e}{}
\subsection*{\small\ttfamily\raggedright Applied\allowbreak{}Modeling\allowbreak{}Lib.\allowbreak{}Probability.\allowbreak{}euclidean\allowbreak{}Dyadic\allowbreak{}Cube\allowbreak{}Shell\allowbreak{}Compact\allowbreak{}Law}
\textbf{Declaration location:} reusable library
\paragraph{Lean declaration body}
\begin{ReviewVerbatim}
type:
(dimension : ℕ) →
  ℕ →
    MeasureTheory.ProbabilityMeasure (EuclideanSpace ℝ (Fin dimension)) →
      MeasureTheory.ProbabilityMeasure (AppliedModelingLib.Probability.HalfOpenUnitCube dimension)

value:
fun (dimension shell : ℕ) (law : MeasureTheory.ProbabilityMeasure (EuclideanSpace ℝ (Fin dimension))) =>
  ((law.toFiniteMeasure.restrict (AppliedModelingLib.Probability.euclideanDyadicCubeShell dimension shell)).map
      (AppliedModelingLib.Probability.euclideanDyadicCubeShellScaleExtension dimension shell)).normalize
\end{ReviewVerbatim}

\paragraph{Direct retained dependencies}
\begin{itemize}\raggedright
\item \hyperlink{governing-declaration-68be765167e02884}{Applied\allowbreak{}Modeling\allowbreak{}Lib.\allowbreak{}Probability.\allowbreak{}Half\allowbreak{}Open\allowbreak{}Unit\allowbreak{}Cube}
\item \hyperlink{governing-declaration-13cdf9a43259c30a}{Applied\allowbreak{}Modeling\allowbreak{}Lib.\allowbreak{}Probability.\allowbreak{}euclidean\allowbreak{}Dyadic\allowbreak{}Cube\allowbreak{}Shell}
\item \hyperlink{governing-declaration-4bcc855202b49a28}{Applied\allowbreak{}Modeling\allowbreak{}Lib.\allowbreak{}Probability.\allowbreak{}euclidean\allowbreak{}Dyadic\allowbreak{}Cube\allowbreak{}Shell\allowbreak{}Scale\allowbreak{}Extension}
\item \hyperlink{governing-declaration-5eb127d4a6b534fc}{Applied\allowbreak{}Modeling\allowbreak{}Lib.\allowbreak{}Probability.\allowbreak{}half\allowbreak{}Open\allowbreak{}Unit\allowbreak{}Interval}
\end{itemize}
\hypertarget{governing-declaration-7436601cf8c8c162}{}
\subsection*{\small\ttfamily\raggedright Applied\allowbreak{}Modeling\allowbreak{}Lib.\allowbreak{}Probability.\allowbreak{}normalized\allowbreak{}Unit\allowbreak{}Coordinate}
\textbf{Declaration location:} reusable library
\paragraph{Lean declaration body}
\begin{ReviewVerbatim}
type:
ℝ → ℝ

value:
fun (x : ℝ) => (x + 1) / 2
\end{ReviewVerbatim}


\hypertarget{governing-declaration-f30eabfddc4020a6}{}
\subsection*{\small\ttfamily\raggedright Applied\allowbreak{}Modeling\allowbreak{}Lib.\allowbreak{}Probability.\allowbreak{}dyadic\allowbreak{}Coordinate\allowbreak{}Number}
\textbf{Declaration location:} reusable library
\paragraph{Lean declaration body}
\begin{ReviewVerbatim}
type:
ℕ → ℝ → ℕ

value:
fun (level : ℕ) (x : ℝ) => ⌈↑(2 ^ level) * AppliedModelingLib.Probability.normalizedUnitCoordinate x⌉₊ - 1
\end{ReviewVerbatim}

\paragraph{Direct retained dependencies}
\begin{itemize}\raggedright
\item \hyperlink{governing-declaration-7436601cf8c8c162}{Applied\allowbreak{}Modeling\allowbreak{}Lib.\allowbreak{}Probability.\allowbreak{}normalized\allowbreak{}Unit\allowbreak{}Coordinate}
\end{itemize}
\hypertarget{governing-declaration-48c904a903f54ed7}{}
\subsection*{\small\ttfamily\raggedright Applied\allowbreak{}Modeling\allowbreak{}Lib.\allowbreak{}Probability.\allowbreak{}dyadic\allowbreak{}Coordinate\allowbreak{}Index}
\textbf{Declaration location:} reusable library
\paragraph{Lean declaration body}
\begin{ReviewVerbatim}
type:
(level : ℕ) →
  (x : ℝ) →
    x ∈ AppliedModelingLib.Probability.halfOpenUnitInterval → AppliedModelingLib.Probability.dyadicCellIndex level

value:
fun (level : ℕ) (x : ℝ) (hx : x ∈ AppliedModelingLib.Probability.halfOpenUnitInterval) =>
  ⟨AppliedModelingLib.Probability.dyadicCoordinateNumber level x,
    AppliedModelingLib.Probability.dyadicCoordinateNumber_lt level x hx⟩
\end{ReviewVerbatim}

\paragraph{Direct retained dependencies}
\begin{itemize}\raggedright
\item \hyperlink{governing-declaration-d5b34777e35d2553}{Applied\allowbreak{}Modeling\allowbreak{}Lib.\allowbreak{}Probability.\allowbreak{}dyadic\allowbreak{}Cell\allowbreak{}Index}
\item \hyperlink{governing-declaration-f30eabfddc4020a6}{Applied\allowbreak{}Modeling\allowbreak{}Lib.\allowbreak{}Probability.\allowbreak{}dyadic\allowbreak{}Coordinate\allowbreak{}Number}
\item \hyperlink{governing-declaration-5eb127d4a6b534fc}{Applied\allowbreak{}Modeling\allowbreak{}Lib.\allowbreak{}Probability.\allowbreak{}half\allowbreak{}Open\allowbreak{}Unit\allowbreak{}Interval}
\end{itemize}
\hypertarget{governing-declaration-26c60b71e183df86}{}
\subsection*{\small\ttfamily\raggedright Applied\allowbreak{}Modeling\allowbreak{}Lib.\allowbreak{}Probability.\allowbreak{}dyadic\allowbreak{}Cube\allowbreak{}Locate}
\textbf{Declaration location:} reusable library
\paragraph{Lean declaration body}
\begin{ReviewVerbatim}
type:
(dimension level : ℕ) →
  AppliedModelingLib.Probability.HalfOpenUnitCube dimension →
    Fin dimension → AppliedModelingLib.Probability.dyadicCellIndex level

value:
fun (dimension level : ℕ) (x : AppliedModelingLib.Probability.HalfOpenUnitCube dimension) (a : Fin dimension) =>
  AppliedModelingLib.Probability.dyadicCoordinateIndex level ((↑x).ofLp a)
    (AppliedModelingLib.Probability.dyadicCubeLocate._proof_1 dimension x a)
\end{ReviewVerbatim}

\paragraph{Direct retained dependencies}
\begin{itemize}\raggedright
\item \hyperlink{governing-declaration-68be765167e02884}{Applied\allowbreak{}Modeling\allowbreak{}Lib.\allowbreak{}Probability.\allowbreak{}Half\allowbreak{}Open\allowbreak{}Unit\allowbreak{}Cube}
\item \hyperlink{governing-declaration-d5b34777e35d2553}{Applied\allowbreak{}Modeling\allowbreak{}Lib.\allowbreak{}Probability.\allowbreak{}dyadic\allowbreak{}Cell\allowbreak{}Index}
\item \hyperlink{governing-declaration-48c904a903f54ed7}{Applied\allowbreak{}Modeling\allowbreak{}Lib.\allowbreak{}Probability.\allowbreak{}dyadic\allowbreak{}Coordinate\allowbreak{}Index}
\item \hyperlink{governing-declaration-5eb127d4a6b534fc}{Applied\allowbreak{}Modeling\allowbreak{}Lib.\allowbreak{}Probability.\allowbreak{}half\allowbreak{}Open\allowbreak{}Unit\allowbreak{}Interval}
\end{itemize}
\hypertarget{governing-declaration-6be421a0d1e9b169}{}
\subsection*{\small\ttfamily\raggedright Applied\allowbreak{}Modeling\allowbreak{}Lib.\allowbreak{}Probability.\allowbreak{}dyadic\allowbreak{}Partition\allowbreak{}L1\allowbreak{}Deviation}
\textbf{Declaration location:} reusable library
\paragraph{Lean declaration body}
\begin{ReviewVerbatim}
type:
(dimension : ℕ) →
  MeasureTheory.ProbabilityMeasure (AppliedModelingLib.Probability.HalfOpenUnitCube dimension) →
    (count : ℕ) → ℕ → (Fin count → AppliedModelingLib.Probability.HalfOpenUnitCube dimension) → ℝ

value:
fun (dimension : ℕ) (law : MeasureTheory.ProbabilityMeasure (AppliedModelingLib.Probability.HalfOpenUnitCube dimension))
    (count level : ℕ) (sample : Fin count → AppliedModelingLib.Probability.HalfOpenUnitCube dimension) =>
  ∑ index : AppliedModelingLib.Probability.dyadicCubeIndex dimension level,
    |AppliedModelingLib.Probability.finiteIIDCellFrequency count
          {x : AppliedModelingLib.Probability.HalfOpenUnitCube dimension |
            AppliedModelingLib.Probability.dyadicCubeLocate dimension level x = index}
          sample -
        (↑law).real
          {x : AppliedModelingLib.Probability.HalfOpenUnitCube dimension |
            AppliedModelingLib.Probability.dyadicCubeLocate dimension level x = index}|
\end{ReviewVerbatim}

\paragraph{Direct retained dependencies}
\begin{itemize}\raggedright
\item \hyperlink{governing-declaration-68be765167e02884}{Applied\allowbreak{}Modeling\allowbreak{}Lib.\allowbreak{}Probability.\allowbreak{}Half\allowbreak{}Open\allowbreak{}Unit\allowbreak{}Cube}
\item \hyperlink{governing-declaration-d5b34777e35d2553}{Applied\allowbreak{}Modeling\allowbreak{}Lib.\allowbreak{}Probability.\allowbreak{}dyadic\allowbreak{}Cell\allowbreak{}Index}
\item \hyperlink{governing-declaration-24e8f7da653baa4e}{Applied\allowbreak{}Modeling\allowbreak{}Lib.\allowbreak{}Probability.\allowbreak{}dyadic\allowbreak{}Cube\allowbreak{}Index}
\item \hyperlink{governing-declaration-26c60b71e183df86}{Applied\allowbreak{}Modeling\allowbreak{}Lib.\allowbreak{}Probability.\allowbreak{}dyadic\allowbreak{}Cube\allowbreak{}Locate}
\item \hyperlink{governing-declaration-aa3875e0dd44ed11}{Applied\allowbreak{}Modeling\allowbreak{}Lib.\allowbreak{}Probability.\allowbreak{}finite\allowbreak{}IID\allowbreak{}Cell\allowbreak{}Frequency}
\item \hyperlink{governing-declaration-5eb127d4a6b534fc}{Applied\allowbreak{}Modeling\allowbreak{}Lib.\allowbreak{}Probability.\allowbreak{}half\allowbreak{}Open\allowbreak{}Unit\allowbreak{}Interval}
\end{itemize}
\hypertarget{governing-declaration-c6ed5828a32d7abc}{}
\subsection*{\small\ttfamily\raggedright Applied\allowbreak{}Modeling\allowbreak{}Lib.\allowbreak{}Probability.\allowbreak{}p\allowbreak{}One\allowbreak{}Fournier\allowbreak{}Guillin\allowbreak{}Head\allowbreak{}Gate\allowbreak{}Radius}
\textbf{Declaration location:} reusable library
\paragraph{Lean declaration body}
\begin{ReviewVerbatim}
type:
ℝ → ℕ → ℝ

value:
fun (eta : ℝ) (gateScale : ℕ) => (2 ^ gateScale).rpow (-(2 * (1 + eta)))
\end{ReviewVerbatim}


\hypertarget{governing-declaration-adeba479719ad06f}{}
\subsection*{\small\ttfamily\raggedright Applied\allowbreak{}Modeling\allowbreak{}Lib.\allowbreak{}Probability.\allowbreak{}p\allowbreak{}One\allowbreak{}Fournier\allowbreak{}Guillin\allowbreak{}Capped\allowbreak{}Scaled\allowbreak{}Deviation}
\textbf{Declaration location:} reusable library
\paragraph{Lean declaration body}
\begin{ReviewVerbatim}
type:
ℝ → ℕ → ℝ → ℝ

value:
fun (eta : ℝ) (gateScale : ℕ) (scaledDeviation : ℝ) =>
  min scaledDeviation (AppliedModelingLib.Probability.pOneFournierGuillinHeadGateRadius eta gateScale)
\end{ReviewVerbatim}

\paragraph{Direct retained dependencies}
\begin{itemize}\raggedright
\item \hyperlink{governing-declaration-c6ed5828a32d7abc}{Applied\allowbreak{}Modeling\allowbreak{}Lib.\allowbreak{}Probability.\allowbreak{}p\allowbreak{}One\allowbreak{}Fournier\allowbreak{}Guillin\allowbreak{}Head\allowbreak{}Gate\allowbreak{}Radius}
\end{itemize}
\hypertarget{governing-declaration-dc945d4458127dc6}{}
\subsection*{\small\ttfamily\raggedright Applied\allowbreak{}Modeling\allowbreak{}Lib.\allowbreak{}Probability.\allowbreak{}p\allowbreak{}One\allowbreak{}Fournier\allowbreak{}Guillin\allowbreak{}Capped\allowbreak{}Deviation}
\textbf{Declaration location:} reusable library
\paragraph{Lean declaration body}
\begin{ReviewVerbatim}
type:
ℝ → ℕ → ℝ → ℝ

value:
fun (eta : ℝ) (gateScale : ℕ) (scaledDeviation : ℝ) =>
  AppliedModelingLib.Probability.pOneFournierGuillinCappedScaledDeviation eta gateScale scaledDeviation /
    Real.rpow 2 (-(1 + eta))
\end{ReviewVerbatim}

\paragraph{Direct retained dependencies}
\begin{itemize}\raggedright
\item \hyperlink{governing-declaration-adeba479719ad06f}{Applied\allowbreak{}Modeling\allowbreak{}Lib.\allowbreak{}Probability.\allowbreak{}p\allowbreak{}One\allowbreak{}Fournier\allowbreak{}Guillin\allowbreak{}Capped\allowbreak{}Scaled\allowbreak{}Deviation}
\end{itemize}
\hypertarget{governing-declaration-1dd6dc54de4dcd4f}{}
\subsection*{\small\ttfamily\raggedright Applied\allowbreak{}Modeling\allowbreak{}Lib.\allowbreak{}Probability.\allowbreak{}p\allowbreak{}One\allowbreak{}Fournier\allowbreak{}Guillin\allowbreak{}Shell\allowbreak{}Coefficient}
\textbf{Declaration location:} reusable library
\paragraph{Lean declaration body}
\begin{ReviewVerbatim}
type:
ℝ → ℝ

value:
fun (eta : ℝ) => 1 - Real.rpow 2 (-eta)
\end{ReviewVerbatim}


\hypertarget{governing-declaration-831786d365674bcf}{}
\subsection*{\small\ttfamily\raggedright Applied\allowbreak{}Modeling\allowbreak{}Lib.\allowbreak{}Probability.\allowbreak{}p\allowbreak{}One\allowbreak{}Fournier\allowbreak{}Guillin\allowbreak{}Shell\allowbreak{}Weight}
\textbf{Declaration location:} reusable library
\paragraph{Lean declaration body}
\begin{ReviewVerbatim}
type:
ℝ → ℝ → ℕ → ℝ

value:
fun (eta deviation : ℝ) (shell : ℕ) =>
  AppliedModelingLib.Probability.pOneFournierGuillinShellCoefficient eta * deviation *
    Real.rpow 2 (-((1 + eta) * ↑shell))
\end{ReviewVerbatim}

\paragraph{Direct retained dependencies}
\begin{itemize}\raggedright
\item \hyperlink{governing-declaration-1dd6dc54de4dcd4f}{Applied\allowbreak{}Modeling\allowbreak{}Lib.\allowbreak{}Probability.\allowbreak{}p\allowbreak{}One\allowbreak{}Fournier\allowbreak{}Guillin\allowbreak{}Shell\allowbreak{}Coefficient}
\end{itemize}
\hypertarget{governing-declaration-1b744d1c0d378e63}{}
\subsection*{\small\ttfamily\raggedright Applied\allowbreak{}Modeling\allowbreak{}Lib.\allowbreak{}Probability.\allowbreak{}p\allowbreak{}One\allowbreak{}Fournier\allowbreak{}Guillin\allowbreak{}Shell\allowbreak{}Tolerance\allowbreak{}With\allowbreak{}Eta}
\textbf{Declaration location:} reusable library
\paragraph{Lean declaration body}
\begin{ReviewVerbatim}
type:
ℝ → ℝ → ℝ → ℕ → ℝ

value:
fun (eta deviation shellMass : ℝ) (shell : ℕ) =>
  AppliedModelingLib.Probability.pOneFournierGuillinShellWeight eta deviation shell / shellMass
\end{ReviewVerbatim}

\paragraph{Direct retained dependencies}
\begin{itemize}\raggedright
\item \hyperlink{governing-declaration-831786d365674bcf}{Applied\allowbreak{}Modeling\allowbreak{}Lib.\allowbreak{}Probability.\allowbreak{}p\allowbreak{}One\allowbreak{}Fournier\allowbreak{}Guillin\allowbreak{}Shell\allowbreak{}Weight}
\end{itemize}
\hypertarget{governing-declaration-2d779803f9998f61}{}
\subsection*{\small\ttfamily\raggedright Applied\allowbreak{}Modeling\allowbreak{}Lib.\allowbreak{}Probability.\allowbreak{}Is\allowbreak{}P\allowbreak{}One\allowbreak{}Fournier\allowbreak{}Guillin\allowbreak{}Uniform\allowbreak{}All\allowbreak{}Shell\allowbreak{}Mass\allowbreak{}Gate\allowbreak{}Scale}
\textbf{Declaration location:} reusable library
\paragraph{Lean declaration body}
\begin{ReviewVerbatim}
type:
ℕ → ℝ → ℝ → ℝ → ℝ → ℕ → Prop

value:
fun (dimension : ℕ) (eta alpha gamma momentBound : ℝ) (gateScale : ℕ) =>
  ∀ (count : ℕ),
    0 < count →
      ∀ (law : MeasureTheory.ProbabilityMeasure (EuclideanSpace ℝ (Fin dimension))),
        AppliedModelingLib.Probability.HasBoundedExponentialRadialMoment law alpha gamma momentBound →
          ∀ {deviation scaledDeviation : ℝ},
            scaledDeviation = deviation * Real.rpow 2 (-(1 + eta)) →
              0 < scaledDeviation →
                scaledDeviation ≤ 1 →
                  scaledDeviation ≤ AppliedModelingLib.Probability.pOneFournierGuillinHeadGateRadius eta gateScale →
                    1 ≤
                        ↑count * AppliedModelingLib.Probability.pOneFournierGuillinShellCoefficient eta ^ 2 *
                            scaledDeviation ^ 2 /
                          (18 *
                            AppliedModelingLib.Probability.fifthOrderExponentialRadialTailConstant law alpha gamma) →
                      1 ≤
                          ↑count * AppliedModelingLib.Probability.pOneFournierGuillinShellCoefficient eta *
                                  Real.rpow 2 (-(1 + eta)) *
                                scaledDeviation ^ 2 *
                              Real.log 9 /
                            8 →
                        1 ≤
                            ↑count * AppliedModelingLib.Probability.pOneFournierGuillinShellCoefficient eta *
                                    scaledDeviation ^ 2 *
                                  gamma /
                                4 *
                              (Real.rpow 2 (alpha - 1 - eta) - 1) →
                          (AppliedModelingLib.Probability.finiteIIDSampleLaw (↑law) count).real
                              (⋃ (shell : ℕ),
                                AppliedModelingLib.Probability.finiteIIDEuclideanDyadicCubeShellMassRelativeTailTotalized
                                  dimension law count shell
                                  (AppliedModelingLib.Probability.pOneFournierGuillinShellToleranceWithEta eta deviation
                                    ((↑law).real
                                      (AppliedModelingLib.Probability.euclideanDyadicCubeShell dimension shell))
                                    shell)) ≤
                            2 *
                                Real.exp
                                  (-(↑count *
                                          AppliedModelingLib.Probability.pOneFournierGuillinShellWeight eta deviation
                                            0) ^
                                        2 /
                                    (2 * ↑count * (1 / 4))) +
                              (4 *
                                    Real.exp
                                      (-(↑count *
                                              AppliedModelingLib.Probability.pOneFournierGuillinShellCoefficient eta ^
                                                2 *
                                            scaledDeviation ^ 2 /
                                          (18 *
                                            AppliedModelingLib.Probability.fifthOrderExponentialRadialTailConstant law
                                              alpha gamma))) +
                                  Real.exp
                                    (-(↑count * AppliedModelingLib.Probability.pOneFournierGuillinShellCoefficient eta *
                                              Real.rpow 2 (-(1 + eta)) *
                                            scaledDeviation ^ 2 *
                                          Real.log 9 /
                                        4)) +
                                Real.exp
                                    (-(↑count * AppliedModelingLib.Probability.pOneFournierGuillinShellCoefficient eta *
                                            scaledDeviation ^ 2 *
                                          gamma /
                                        4)) /
                                  (1 - Real.exp (-1)))
\end{ReviewVerbatim}

\paragraph{Direct retained dependencies}
\begin{itemize}\raggedright
\item \hyperlink{governing-declaration-9604ed56aea35d91}{Applied\allowbreak{}Modeling\allowbreak{}Lib.\allowbreak{}Probability.\allowbreak{}Has\allowbreak{}Bounded\allowbreak{}Exponential\allowbreak{}Radial\allowbreak{}Moment}
\item \hyperlink{governing-declaration-13cdf9a43259c30a}{Applied\allowbreak{}Modeling\allowbreak{}Lib.\allowbreak{}Probability.\allowbreak{}euclidean\allowbreak{}Dyadic\allowbreak{}Cube\allowbreak{}Shell}
\item \hyperlink{governing-declaration-93c495b865c4b878}{Applied\allowbreak{}Modeling\allowbreak{}Lib.\allowbreak{}Probability.\allowbreak{}fifth\allowbreak{}Order\allowbreak{}Exponential\allowbreak{}Radial\allowbreak{}Tail\allowbreak{}Constant}
\item \hyperlink{governing-declaration-dbd328d5a293e61c}{Applied\allowbreak{}Modeling\allowbreak{}Lib.\allowbreak{}Probability.\allowbreak{}finite\allowbreak{}IID\allowbreak{}Euclidean\allowbreak{}Dyadic\allowbreak{}Cube\allowbreak{}Shell\allowbreak{}Mass\allowbreak{}Relative\allowbreak{}Tail\allowbreak{}Totalized}
\item \hyperlink{governing-declaration-91a777e8ba845265}{Applied\allowbreak{}Modeling\allowbreak{}Lib.\allowbreak{}Probability.\allowbreak{}finite\allowbreak{}IID\allowbreak{}Sample\allowbreak{}Law}
\item \hyperlink{governing-declaration-c6ed5828a32d7abc}{Applied\allowbreak{}Modeling\allowbreak{}Lib.\allowbreak{}Probability.\allowbreak{}p\allowbreak{}One\allowbreak{}Fournier\allowbreak{}Guillin\allowbreak{}Head\allowbreak{}Gate\allowbreak{}Radius}
\item \hyperlink{governing-declaration-1dd6dc54de4dcd4f}{Applied\allowbreak{}Modeling\allowbreak{}Lib.\allowbreak{}Probability.\allowbreak{}p\allowbreak{}One\allowbreak{}Fournier\allowbreak{}Guillin\allowbreak{}Shell\allowbreak{}Coefficient}
\item \hyperlink{governing-declaration-1b744d1c0d378e63}{Applied\allowbreak{}Modeling\allowbreak{}Lib.\allowbreak{}Probability.\allowbreak{}p\allowbreak{}One\allowbreak{}Fournier\allowbreak{}Guillin\allowbreak{}Shell\allowbreak{}Tolerance\allowbreak{}With\allowbreak{}Eta}
\item \hyperlink{governing-declaration-831786d365674bcf}{Applied\allowbreak{}Modeling\allowbreak{}Lib.\allowbreak{}Probability.\allowbreak{}p\allowbreak{}One\allowbreak{}Fournier\allowbreak{}Guillin\allowbreak{}Shell\allowbreak{}Weight}
\end{itemize}
\hypertarget{governing-declaration-beaaebb5e90384d4}{}
\subsection*{\small\ttfamily\raggedright Applied\allowbreak{}Modeling\allowbreak{}Lib.\allowbreak{}Probability.\allowbreak{}p\allowbreak{}One\allowbreak{}Fournier\allowbreak{}Guillin\allowbreak{}Uniform\allowbreak{}All\allowbreak{}Shell\allowbreak{}Mass\allowbreak{}Effective\allowbreak{}Count\allowbreak{}Requirement}
\textbf{Declaration location:} reusable library
\paragraph{Lean declaration body}
\begin{ReviewVerbatim}
type:
ℝ → ℝ → ℝ → ℝ → ℝ → ℝ → ℝ

value:
fun (eta deviation scaledDeviation gamma tailBound tolerance : ℝ) =>
  max
    (Real.log (8 / tolerance) / (2 * AppliedModelingLib.Probability.pOneFournierGuillinShellWeight eta deviation 0 ^ 2))
    (max
      (Real.log (16 / tolerance) /
        (AppliedModelingLib.Probability.pOneFournierGuillinShellCoefficient eta ^ 2 * scaledDeviation ^ 2 /
          (18 * tailBound)))
      (max
        (Real.log (4 / tolerance) /
          (AppliedModelingLib.Probability.pOneFournierGuillinShellCoefficient eta * Real.rpow 2 (-(1 + eta)) *
                scaledDeviation ^ 2 *
              Real.log 9 /
            4))
        (Real.log (4 / (tolerance * (1 - Real.exp (-1)))) /
          (AppliedModelingLib.Probability.pOneFournierGuillinShellCoefficient eta * scaledDeviation ^ 2 * gamma / 4))))
\end{ReviewVerbatim}

\paragraph{Direct retained dependencies}
\begin{itemize}\raggedright
\item \hyperlink{governing-declaration-1dd6dc54de4dcd4f}{Applied\allowbreak{}Modeling\allowbreak{}Lib.\allowbreak{}Probability.\allowbreak{}p\allowbreak{}One\allowbreak{}Fournier\allowbreak{}Guillin\allowbreak{}Shell\allowbreak{}Coefficient}
\item \hyperlink{governing-declaration-831786d365674bcf}{Applied\allowbreak{}Modeling\allowbreak{}Lib.\allowbreak{}Probability.\allowbreak{}p\allowbreak{}One\allowbreak{}Fournier\allowbreak{}Guillin\allowbreak{}Shell\allowbreak{}Weight}
\end{itemize}
\hypertarget{governing-declaration-b2e74788febba02c}{}
\subsection*{\small\ttfamily\raggedright Applied\allowbreak{}Modeling\allowbreak{}Lib.\allowbreak{}Probability.\allowbreak{}p\allowbreak{}One\allowbreak{}Fournier\allowbreak{}Guillin\allowbreak{}Uniform\allowbreak{}All\allowbreak{}Shell\allowbreak{}Mass\allowbreak{}Gate\allowbreak{}Scale}
\textbf{Declaration location:} reusable library
\paragraph{Lean declaration body}
\begin{ReviewVerbatim}
type:
ℕ → ℝ → ℝ → ℝ → ℝ → ℕ

value:
fun (dimension : ℕ) (eta alpha gamma momentBound : ℝ) =>
  Classical.epsilon
    (AppliedModelingLib.Probability.IsPOneFournierGuillinUniformAllShellMassGateScale dimension eta alpha gamma
      momentBound)
\end{ReviewVerbatim}

\paragraph{Direct retained dependencies}
\begin{itemize}\raggedright
\item \hyperlink{governing-declaration-2d779803f9998f61}{Applied\allowbreak{}Modeling\allowbreak{}Lib.\allowbreak{}Probability.\allowbreak{}Is\allowbreak{}P\allowbreak{}One\allowbreak{}Fournier\allowbreak{}Guillin\allowbreak{}Uniform\allowbreak{}All\allowbreak{}Shell\allowbreak{}Mass\allowbreak{}Gate\allowbreak{}Scale}
\end{itemize}
\hypertarget{governing-declaration-d3b37af7ee378c37}{}
\subsection*{\small\ttfamily\raggedright Applied\allowbreak{}Modeling\allowbreak{}Lib.\allowbreak{}Probability.\allowbreak{}p\allowbreak{}One\allowbreak{}Fournier\allowbreak{}Guillin\allowbreak{}Uniform\allowbreak{}Capped\allowbreak{}Deviation}
\textbf{Declaration location:} reusable library
\paragraph{Lean declaration body}
\begin{ReviewVerbatim}
type:
ℕ → ℝ → ℝ → ℝ → ℝ → ℝ → ℝ

value:
fun (dimension : ℕ) (eta alpha gamma momentBound scaledDeviation : ℝ) =>
  AppliedModelingLib.Probability.pOneFournierGuillinCappedDeviation eta
    (AppliedModelingLib.Probability.pOneFournierGuillinUniformAllShellMassGateScale dimension eta alpha gamma
      momentBound)
    scaledDeviation
\end{ReviewVerbatim}

\paragraph{Direct retained dependencies}
\begin{itemize}\raggedright
\item \hyperlink{governing-declaration-dc945d4458127dc6}{Applied\allowbreak{}Modeling\allowbreak{}Lib.\allowbreak{}Probability.\allowbreak{}p\allowbreak{}One\allowbreak{}Fournier\allowbreak{}Guillin\allowbreak{}Capped\allowbreak{}Deviation}
\item \hyperlink{governing-declaration-b2e74788febba02c}{Applied\allowbreak{}Modeling\allowbreak{}Lib.\allowbreak{}Probability.\allowbreak{}p\allowbreak{}One\allowbreak{}Fournier\allowbreak{}Guillin\allowbreak{}Uniform\allowbreak{}All\allowbreak{}Shell\allowbreak{}Mass\allowbreak{}Gate\allowbreak{}Scale}
\end{itemize}
\hypertarget{governing-declaration-284d6604fe65e6f0}{}
\subsection*{\small\ttfamily\raggedright Applied\allowbreak{}Modeling\allowbreak{}Lib.\allowbreak{}Probability.\allowbreak{}p\allowbreak{}One\allowbreak{}Fournier\allowbreak{}Guillin\allowbreak{}Uniform\allowbreak{}Capped\allowbreak{}Scaled\allowbreak{}Deviation}
\textbf{Declaration location:} reusable library
\paragraph{Lean declaration body}
\begin{ReviewVerbatim}
type:
ℕ → ℝ → ℝ → ℝ → ℝ → ℝ → ℝ

value:
fun (dimension : ℕ) (eta alpha gamma momentBound scaledDeviation : ℝ) =>
  AppliedModelingLib.Probability.pOneFournierGuillinCappedScaledDeviation eta
    (AppliedModelingLib.Probability.pOneFournierGuillinUniformAllShellMassGateScale dimension eta alpha gamma
      momentBound)
    scaledDeviation
\end{ReviewVerbatim}

\paragraph{Direct retained dependencies}
\begin{itemize}\raggedright
\item \hyperlink{governing-declaration-adeba479719ad06f}{Applied\allowbreak{}Modeling\allowbreak{}Lib.\allowbreak{}Probability.\allowbreak{}p\allowbreak{}One\allowbreak{}Fournier\allowbreak{}Guillin\allowbreak{}Capped\allowbreak{}Scaled\allowbreak{}Deviation}
\item \hyperlink{governing-declaration-b2e74788febba02c}{Applied\allowbreak{}Modeling\allowbreak{}Lib.\allowbreak{}Probability.\allowbreak{}p\allowbreak{}One\allowbreak{}Fournier\allowbreak{}Guillin\allowbreak{}Uniform\allowbreak{}All\allowbreak{}Shell\allowbreak{}Mass\allowbreak{}Gate\allowbreak{}Scale}
\end{itemize}
\hypertarget{governing-declaration-01d48769dfdad19d}{}
\subsection*{\small\ttfamily\raggedright Applied\allowbreak{}Modeling\allowbreak{}Lib.\allowbreak{}Probability.\allowbreak{}p\allowbreak{}One\allowbreak{}Fournier\allowbreak{}Guillin\allowbreak{}Uniform\allowbreak{}Square\allowbreak{}Rate\allowbreak{}Effective\allowbreak{}Count\allowbreak{}Requirement}
\textbf{Declaration location:} reusable library
\paragraph{Lean declaration body}
\begin{ReviewVerbatim}
type:
ℝ → ℝ → ℝ → ℝ → ℝ → ℝ

value:
fun (eta alpha gamma tailBound scaledDeviation : ℝ) =>
  max
    (1 /
      (AppliedModelingLib.Probability.pOneFournierGuillinShellCoefficient eta ^ 2 * scaledDeviation ^ 2 /
        (18 * tailBound)))
    (max
      (1 /
        (AppliedModelingLib.Probability.pOneFournierGuillinShellCoefficient eta * Real.rpow 2 (-(1 + eta)) *
              scaledDeviation ^ 2 *
            Real.log 9 /
          8))
      (1 /
        (AppliedModelingLib.Probability.pOneFournierGuillinShellCoefficient eta * scaledDeviation ^ 2 * gamma / 4 *
          (Real.rpow 2 (alpha - 1 - eta) - 1))))
\end{ReviewVerbatim}

\paragraph{Direct retained dependencies}
\begin{itemize}\raggedright
\item \hyperlink{governing-declaration-1dd6dc54de4dcd4f}{Applied\allowbreak{}Modeling\allowbreak{}Lib.\allowbreak{}Probability.\allowbreak{}p\allowbreak{}One\allowbreak{}Fournier\allowbreak{}Guillin\allowbreak{}Shell\allowbreak{}Coefficient}
\end{itemize}
\hypertarget{governing-declaration-c803849fcd834479}{}
\subsection*{\small\ttfamily\raggedright Applied\allowbreak{}Modeling\allowbreak{}Lib.\allowbreak{}Probability.\allowbreak{}selected\allowbreak{}Shell\allowbreak{}Dyadic\allowbreak{}Geometric\allowbreak{}Confidence\allowbreak{}Threshold}
\textbf{Declaration location:} reusable library
\paragraph{Lean declaration body}
\begin{ReviewVerbatim}
type:
ℕ → ℕ → ℕ → ℝ → ℝ → ℝ

value:
fun (dimension count level : ℕ) (shellMass confidence : ℝ) =>
  √(4 * (2 * ↑(2 ^ (dimension * level)) + confidence) / (↑count * shellMass))
\end{ReviewVerbatim}


\hypertarget{governing-declaration-d481e01804396f20}{}
\subsection*{\small\ttfamily\raggedright Applied\allowbreak{}Modeling\allowbreak{}Lib.\allowbreak{}Probability.\allowbreak{}selected\allowbreak{}Shell\allowbreak{}Geometric\allowbreak{}Confidence}
\textbf{Declaration location:} reusable library
\paragraph{Lean declaration body}
\begin{ReviewVerbatim}
type:
ℝ → ℕ → ℝ

value:
fun (confidence : ℝ) (shell : ℕ) => confidence + (↑shell + 1) * Real.log 2
\end{ReviewVerbatim}


\hypertarget{governing-declaration-e32669468be7611b}{}
\subsection*{\small\ttfamily\raggedright Applied\allowbreak{}Modeling\allowbreak{}Lib.\allowbreak{}Probability.\allowbreak{}selected\allowbreak{}Shell\allowbreak{}Geometric\allowbreak{}Effective\allowbreak{}Depth}
\textbf{Declaration location:} reusable library
\paragraph{Lean declaration body}
\begin{ReviewVerbatim}
type:
(dimension : ℕ) → ℕ → MeasureTheory.ProbabilityMeasure (EuclideanSpace ℝ (Fin dimension)) → ℝ → ℕ → ℕ

value:
fun (dimension count : ℕ) (law : MeasureTheory.ProbabilityMeasure (EuclideanSpace ℝ (Fin dimension))) (confidence : ℝ)
    (shell : ℕ) =>
  AppliedModelingLib.Math.dyadicGeometricEffectiveDepth dimension
    (↑count * (↑law).real (AppliedModelingLib.Probability.euclideanDyadicCubeShell dimension shell))
    (AppliedModelingLib.Probability.selectedShellGeometricConfidence confidence shell)
\end{ReviewVerbatim}

\paragraph{Direct retained dependencies}
\begin{itemize}\raggedright
\item \hyperlink{governing-declaration-4f196894fb58689f}{Applied\allowbreak{}Modeling\allowbreak{}Lib.\allowbreak{}Math.\allowbreak{}dyadic\allowbreak{}Geometric\allowbreak{}Effective\allowbreak{}Depth}
\item \hyperlink{governing-declaration-13cdf9a43259c30a}{Applied\allowbreak{}Modeling\allowbreak{}Lib.\allowbreak{}Probability.\allowbreak{}euclidean\allowbreak{}Dyadic\allowbreak{}Cube\allowbreak{}Shell}
\item \hyperlink{governing-declaration-d481e01804396f20}{Applied\allowbreak{}Modeling\allowbreak{}Lib.\allowbreak{}Probability.\allowbreak{}selected\allowbreak{}Shell\allowbreak{}Geometric\allowbreak{}Confidence}
\end{itemize}
\hypertarget{governing-declaration-f9f338884954e2a9}{}
\subsection*{\small\ttfamily\raggedright Applied\allowbreak{}Modeling\allowbreak{}Lib.\allowbreak{}Probability.\allowbreak{}selected\allowbreak{}Shell\allowbreak{}Geometric\allowbreak{}Effective\allowbreak{}Head\allowbreak{}Depth}
\textbf{Declaration location:} reusable library
\paragraph{Lean declaration body}
\begin{ReviewVerbatim}
type:
(dimension : ℕ) → ℕ → ℕ → MeasureTheory.ProbabilityMeasure (EuclideanSpace ℝ (Fin dimension)) → ℝ → ℕ → ℕ

value:
fun (dimension cutoff count : ℕ) (law : MeasureTheory.ProbabilityMeasure (EuclideanSpace ℝ (Fin dimension)))
    (confidence : ℝ) (shell : ℕ) =>
  if shell ≤ cutoff then
    AppliedModelingLib.Probability.selectedShellGeometricEffectiveDepth dimension count law confidence shell
  else 0
\end{ReviewVerbatim}

\paragraph{Direct retained dependencies}
\begin{itemize}\raggedright
\item \hyperlink{governing-declaration-e32669468be7611b}{Applied\allowbreak{}Modeling\allowbreak{}Lib.\allowbreak{}Probability.\allowbreak{}selected\allowbreak{}Shell\allowbreak{}Geometric\allowbreak{}Effective\allowbreak{}Depth}
\end{itemize}
\hypertarget{governing-declaration-6c9438e2d953743a}{}
\subsection*{\small\ttfamily\raggedright Applied\allowbreak{}Modeling\allowbreak{}Lib.\allowbreak{}Probability.\allowbreak{}selected\allowbreak{}Shell\allowbreak{}Geometric\allowbreak{}Supercritical\allowbreak{}Rate\allowbreak{}Coefficient}
\textbf{Declaration location:} reusable library
\paragraph{Lean declaration body}
\begin{ReviewVerbatim}
type:
ℕ → ℝ

value:
fun (dimension : ℕ) =>
  (2 * √↑dimension + √↑dimension * √2 / (4 * (√↑(2 ^ dimension) / 2 - 1))) *
    (16 * ↑(2 ^ dimension) ^ 2) ^ (↑dimension)⁻¹
\end{ReviewVerbatim}


\hypertarget{governing-declaration-f4ab54eb3fb61074}{}
\subsection*{\small\ttfamily\raggedright Applied\allowbreak{}Modeling\allowbreak{}Lib.\allowbreak{}Probability.\allowbreak{}selected\allowbreak{}Shell\allowbreak{}Geometric\allowbreak{}Tail\allowbreak{}Multiplier}
\textbf{Declaration location:} reusable library
\paragraph{Lean declaration body}
\begin{ReviewVerbatim}
type:
ℝ → ℝ

value:
fun (tailConstant : ℝ) => 4 * √tailConstant
\end{ReviewVerbatim}


\hypertarget{governing-declaration-ed2d27408e82fc8e}{}
\subsection*{\small\ttfamily\raggedright Applied\allowbreak{}Modeling\allowbreak{}Lib.\allowbreak{}Probability.\allowbreak{}selected\allowbreak{}Shell\allowbreak{}Geometric\allowbreak{}Supercritical\allowbreak{}Effective\allowbreak{}Count\allowbreak{}Requirement}
\textbf{Declaration location:} reusable library
\paragraph{Lean declaration body}
\begin{ReviewVerbatim}
type:
ℕ → ℕ → ℝ → ℝ → ℝ → ℝ

value:
fun (dimension cutoff : ℕ) (confidence tailConstant radius : ℝ) =>
  max
    (AppliedModelingLib.Probability.selectedShellGeometricSupercriticalRateCoefficient dimension ^ dimension /
      (radius / 6) ^ dimension)
    (max (6 ^ 2 * (2 * √↑dimension * √(confidence + Real.log 2)) ^ 2 / radius ^ 2)
      (max (6 * (2 * √↑dimension * (16 * ↑(2 ^ dimension) + 8 * (confidence + Real.log 2))) / radius)
        (max
          ((AppliedModelingLib.Probability.selectedShellGeometricSupercriticalRateCoefficient dimension *
                AppliedModelingLib.Probability.selectedShellGeometricTailMultiplier tailConstant) ^
              dimension /
            (radius / 6) ^ dimension)
          (max
            (6 ^ 2 * (2 * √↑dimension * √(confidence + (↑cutoff + 1) * Real.log 2) * (4 * √tailConstant)) ^ 2 /
              radius ^ 2)
            (6 ^ 2 *
                (2 * √↑dimension * √(16 * ↑(2 ^ dimension) + 8 * (confidence + (↑cutoff + 1) * Real.log 2)) *
                    (4 * √tailConstant)) ^
                  2 /
              radius ^ 2)))))
\end{ReviewVerbatim}

\paragraph{Direct retained dependencies}
\begin{itemize}\raggedright
\item \hyperlink{governing-declaration-6c9438e2d953743a}{Applied\allowbreak{}Modeling\allowbreak{}Lib.\allowbreak{}Probability.\allowbreak{}selected\allowbreak{}Shell\allowbreak{}Geometric\allowbreak{}Supercritical\allowbreak{}Rate\allowbreak{}Coefficient}
\item \hyperlink{governing-declaration-f4ab54eb3fb61074}{Applied\allowbreak{}Modeling\allowbreak{}Lib.\allowbreak{}Probability.\allowbreak{}selected\allowbreak{}Shell\allowbreak{}Geometric\allowbreak{}Tail\allowbreak{}Multiplier}
\end{itemize}
\hypertarget{governing-declaration-512ca853567bde9d}{}
\subsection*{\small\ttfamily\raggedright Applied\allowbreak{}Modeling\allowbreak{}Lib.\allowbreak{}Probability.\allowbreak{}p\allowbreak{}One\allowbreak{}Fournier\allowbreak{}Guillin\allowbreak{}Supercritical\allowbreak{}All\allowbreak{}Shell\allowbreak{}Effective\allowbreak{}Count\allowbreak{}Requirement}
\textbf{Declaration location:} reusable library
\paragraph{Lean declaration body}
\begin{ReviewVerbatim}
type:
ℕ → ℕ → ℝ → ℝ → ℝ → ℝ → ℝ → ℝ → ℝ → ℝ → ℝ → ℝ

value:
fun (dimension cutoff : ℕ) (eta alpha gamma tailBound deviation scaledDeviation confidence radius tolerance : ℝ) =>
  max
    (AppliedModelingLib.Probability.selectedShellGeometricSupercriticalEffectiveCountRequirement dimension cutoff
      confidence tailBound radius)
    (max
      (AppliedModelingLib.Probability.pOneFournierGuillinUniformAllShellMassEffectiveCountRequirement eta deviation
        scaledDeviation gamma tailBound tolerance)
      (AppliedModelingLib.Probability.pOneFournierGuillinUniformSquareRateEffectiveCountRequirement eta alpha gamma
        tailBound scaledDeviation))
\end{ReviewVerbatim}

\paragraph{Direct retained dependencies}
\begin{itemize}\raggedright
\item \hyperlink{governing-declaration-beaaebb5e90384d4}{Applied\allowbreak{}Modeling\allowbreak{}Lib.\allowbreak{}Probability.\allowbreak{}p\allowbreak{}One\allowbreak{}Fournier\allowbreak{}Guillin\allowbreak{}Uniform\allowbreak{}All\allowbreak{}Shell\allowbreak{}Mass\allowbreak{}Effective\allowbreak{}Count\allowbreak{}Requirement}
\item \hyperlink{governing-declaration-01d48769dfdad19d}{Applied\allowbreak{}Modeling\allowbreak{}Lib.\allowbreak{}Probability.\allowbreak{}p\allowbreak{}One\allowbreak{}Fournier\allowbreak{}Guillin\allowbreak{}Uniform\allowbreak{}Square\allowbreak{}Rate\allowbreak{}Effective\allowbreak{}Count\allowbreak{}Requirement}
\item \hyperlink{governing-declaration-ed2d27408e82fc8e}{Applied\allowbreak{}Modeling\allowbreak{}Lib.\allowbreak{}Probability.\allowbreak{}selected\allowbreak{}Shell\allowbreak{}Geometric\allowbreak{}Supercritical\allowbreak{}Effective\allowbreak{}Count\allowbreak{}Requirement}
\end{itemize}
\hypertarget{governing-declaration-2116c3c509ce31b4}{}
\subsection*{\small\ttfamily\raggedright Applied\allowbreak{}Modeling\allowbreak{}Lib.\allowbreak{}Probability.\allowbreak{}p\allowbreak{}One\allowbreak{}Fournier\allowbreak{}Guillin\allowbreak{}Supercritical\allowbreak{}All\allowbreak{}Shell\allowbreak{}Effective\allowbreak{}Count}
\textbf{Declaration location:} reusable library
\paragraph{Lean declaration body}
\begin{ReviewVerbatim}
type:
ℕ → ℕ → ℝ → ℝ → ℝ → ℝ → ℝ → ℝ → ℝ → ℝ → ℝ → ℕ

value:
fun (dimension cutoff : ℕ) (eta alpha gamma tailBound deviation scaledDeviation confidence radius tolerance : ℝ) =>
  AppliedModelingLib.Math.positiveNatCeil
    (AppliedModelingLib.Probability.pOneFournierGuillinSupercriticalAllShellEffectiveCountRequirement dimension cutoff
      eta alpha gamma tailBound deviation scaledDeviation confidence radius tolerance)
\end{ReviewVerbatim}

\paragraph{Direct retained dependencies}
\begin{itemize}\raggedright
\item \hyperlink{governing-declaration-01ad3e76a1f28ec0}{Applied\allowbreak{}Modeling\allowbreak{}Lib.\allowbreak{}Math.\allowbreak{}positive\allowbreak{}Nat\allowbreak{}Ceil}
\item \hyperlink{governing-declaration-512ca853567bde9d}{Applied\allowbreak{}Modeling\allowbreak{}Lib.\allowbreak{}Probability.\allowbreak{}p\allowbreak{}One\allowbreak{}Fournier\allowbreak{}Guillin\allowbreak{}Supercritical\allowbreak{}All\allowbreak{}Shell\allowbreak{}Effective\allowbreak{}Count\allowbreak{}Requirement}
\end{itemize}
\hypertarget{governing-declaration-3d6d95e1d691574b}{}
\subsection*{\small\ttfamily\raggedright Applied\allowbreak{}Modeling\allowbreak{}Lib.\allowbreak{}Probability.\allowbreak{}p\allowbreak{}One\allowbreak{}Fournier\allowbreak{}Guillin\allowbreak{}Supercritical\allowbreak{}Capped\allowbreak{}All\allowbreak{}Shell\allowbreak{}Concrete\allowbreak{}Count\allowbreak{}Schedule}
\textbf{Declaration location:} reusable library
\paragraph{Lean declaration body}
\begin{ReviewVerbatim}
type:
ℕ → ℕ → ℝ → ℝ → ℝ → ℝ → ℝ → ℝ → (ℕ → ℝ) → (ℕ → ℝ) → ℕ → ℕ

value:
fun (dimension cutoff : ℕ) (eta alpha gamma momentBound tailBound scaledDeviation : ℝ)
    (failureBudget radiusSchedule : ℕ → ℝ) (a : ℕ) =>
  AppliedModelingLib.Probability.pOneFournierGuillinSupercriticalAllShellEffectiveCount dimension cutoff eta alpha gamma
    tailBound
    (AppliedModelingLib.Probability.pOneFournierGuillinUniformCappedDeviation dimension eta alpha gamma momentBound
      scaledDeviation)
    (AppliedModelingLib.Probability.pOneFournierGuillinUniformCappedScaledDeviation dimension eta alpha gamma
      momentBound scaledDeviation)
    (AppliedModelingLib.Math.expConfidenceForHalfBudget (failureBudget a)) (radiusSchedule a) (failureBudget a / 2)
\end{ReviewVerbatim}

\paragraph{Direct retained dependencies}
\begin{itemize}\raggedright
\item \hyperlink{governing-declaration-77c24384a8319860}{Applied\allowbreak{}Modeling\allowbreak{}Lib.\allowbreak{}Math.\allowbreak{}exp\allowbreak{}Confidence\allowbreak{}For\allowbreak{}Half\allowbreak{}Budget}
\item \hyperlink{governing-declaration-2116c3c509ce31b4}{Applied\allowbreak{}Modeling\allowbreak{}Lib.\allowbreak{}Probability.\allowbreak{}p\allowbreak{}One\allowbreak{}Fournier\allowbreak{}Guillin\allowbreak{}Supercritical\allowbreak{}All\allowbreak{}Shell\allowbreak{}Effective\allowbreak{}Count}
\item \hyperlink{governing-declaration-d3b37af7ee378c37}{Applied\allowbreak{}Modeling\allowbreak{}Lib.\allowbreak{}Probability.\allowbreak{}p\allowbreak{}One\allowbreak{}Fournier\allowbreak{}Guillin\allowbreak{}Uniform\allowbreak{}Capped\allowbreak{}Deviation}
\item \hyperlink{governing-declaration-284d6604fe65e6f0}{Applied\allowbreak{}Modeling\allowbreak{}Lib.\allowbreak{}Probability.\allowbreak{}p\allowbreak{}One\allowbreak{}Fournier\allowbreak{}Guillin\allowbreak{}Uniform\allowbreak{}Capped\allowbreak{}Scaled\allowbreak{}Deviation}
\end{itemize}
\hypertarget{governing-declaration-e0719ef6c387553e}{}
\subsection*{\small\ttfamily\raggedright Applied\allowbreak{}Modeling\allowbreak{}Lib.\allowbreak{}Probability\allowbreak{}Coupling}
\textbf{Declaration location:} reusable library
\paragraph{Lean declaration body}
\begin{ReviewVerbatim}
field joint:
{α : Type u_1} →
  {β : Type u_2} →
    [inst : MeasurableSpace α] →
      [inst_1 : MeasurableSpace β] →
        {μ : MeasureTheory.ProbabilityMeasure α} →
          {ν : MeasureTheory.ProbabilityMeasure β} →
            AppliedModelingLib.ProbabilityCoupling μ ν → MeasureTheory.ProbabilityMeasure (α × β)

field map_fst:
∀ {α : Type u_1} {β : Type u_2} [inst : MeasurableSpace α] [inst_1 : MeasurableSpace β]
  {μ : MeasureTheory.ProbabilityMeasure α} {ν : MeasureTheory.ProbabilityMeasure β}
  (self : AppliedModelingLib.ProbabilityCoupling μ ν), self.joint.map (Measurable.aemeasurable measurable_fst) = μ

field map_snd:
∀ {α : Type u_1} {β : Type u_2} [inst : MeasurableSpace α] [inst_1 : MeasurableSpace β]
  {μ : MeasureTheory.ProbabilityMeasure α} {ν : MeasureTheory.ProbabilityMeasure β}
  (self : AppliedModelingLib.ProbabilityCoupling μ ν), self.joint.map (Measurable.aemeasurable measurable_snd) = ν
\end{ReviewVerbatim}


\hypertarget{governing-declaration-0507d24247d6815a}{}
\subsection*{\small\ttfamily\raggedright Applied\allowbreak{}Modeling\allowbreak{}Lib.\allowbreak{}Probability\allowbreak{}Coupling.\allowbreak{}expected\allowbreak{}Distance\allowbreak{}Costs}
\textbf{Declaration location:} reusable library
\paragraph{Lean declaration body}
\begin{ReviewVerbatim}
type:
{α : Type u_3} →
  [inst : MeasurableSpace α] →
    [PseudoMetricSpace α] → MeasureTheory.ProbabilityMeasure α → MeasureTheory.ProbabilityMeasure α → ℝ → Prop

value:
fun {α : Type u_3} [MeasurableSpace α] [PseudoMetricSpace α] (μ ν : MeasureTheory.ProbabilityMeasure α) (a : ℝ) =>
  {r : ℝ |
      ∃ (π : AppliedModelingLib.ProbabilityCoupling μ ν) (_ :
        MeasureTheory.Integrable (fun (z : α × α) => dist z.1 z.2) ↑π.joint), r = ∫ (z : α × α), dist z.1 z.2 ∂↑π.joint}
    a
\end{ReviewVerbatim}

\paragraph{Direct retained dependencies}
\begin{itemize}\raggedright
\item \hyperlink{governing-declaration-e0719ef6c387553e}{Applied\allowbreak{}Modeling\allowbreak{}Lib.\allowbreak{}Probability\allowbreak{}Coupling}
\end{itemize}
\hypertarget{governing-declaration-7fdf7bcf28dbe52b}{}
\subsection*{\small\ttfamily\raggedright Applied\allowbreak{}Modeling\allowbreak{}Lib.\allowbreak{}Probability\allowbreak{}Coupling.\allowbreak{}wasserstein\allowbreak{}One\allowbreak{}Real}
\textbf{Declaration location:} reusable library
\paragraph{Lean declaration body}
\begin{ReviewVerbatim}
type:
{α : Type u_3} →
  [inst : MeasurableSpace α] →
    [inst_1 : PseudoMetricSpace α] →
      (μ ν : MeasureTheory.ProbabilityMeasure α) →
        (AppliedModelingLib.ProbabilityCoupling.expectedDistanceCosts μ ν).Nonempty → ℝ

value:
fun {α : Type u_3} [MeasurableSpace α] [PseudoMetricSpace α] (μ ν : MeasureTheory.ProbabilityMeasure α)
    (x : (AppliedModelingLib.ProbabilityCoupling.expectedDistanceCosts μ ν).Nonempty) =>
  sInf (AppliedModelingLib.ProbabilityCoupling.expectedDistanceCosts μ ν)
\end{ReviewVerbatim}

\paragraph{Direct retained dependencies}
\begin{itemize}\raggedright
\item \hyperlink{governing-declaration-0507d24247d6815a}{Applied\allowbreak{}Modeling\allowbreak{}Lib.\allowbreak{}Probability\allowbreak{}Coupling.\allowbreak{}expected\allowbreak{}Distance\allowbreak{}Costs}
\end{itemize}
\hypertarget{governing-declaration-ac8867b927a8953c}{}
\subsection*{\small\ttfamily\raggedright Applied\allowbreak{}Modeling\allowbreak{}Lib.\allowbreak{}empirical\allowbreak{}Fintype\allowbreak{}Indices\allowbreak{}In\allowbreak{}Set}
\textbf{Declaration location:} reusable library
\paragraph{Lean declaration body}
\begin{ReviewVerbatim}
type:
{Index : Type u_1} → {Sample : Type u_2} → [Fintype Index] → (Index → Sample) → Set Sample → Finset Index

value:
fun {Index : Type u_1} {Sample : Type u_2} [Fintype Index] (sample : Index → Sample) (s : Set Sample) =>
  {index : Index | sample index ∈ s}
\end{ReviewVerbatim}


\hypertarget{governing-declaration-3e69c45b89f95dd2}{}
\subsection*{\small\ttfamily\raggedright Applied\allowbreak{}Modeling\allowbreak{}Lib.\allowbreak{}empirical\allowbreak{}Sample\allowbreak{}Indices\allowbreak{}In\allowbreak{}Set}
\textbf{Declaration location:} reusable library
\paragraph{Lean declaration body}
\begin{ReviewVerbatim}
type:
{Sample : Type u_1} → {count : ℕ} → (Fin count → Sample) → Set Sample → Finset (Fin count)

value:
fun {Sample : Type u_1} {count : ℕ} (sample : Fin count → Sample) (s : Set Sample) =>
  AppliedModelingLib.empiricalFintypeIndicesInSet sample s
\end{ReviewVerbatim}

\paragraph{Direct retained dependencies}
\begin{itemize}\raggedright
\item \hyperlink{governing-declaration-ac8867b927a8953c}{Applied\allowbreak{}Modeling\allowbreak{}Lib.\allowbreak{}empirical\allowbreak{}Fintype\allowbreak{}Indices\allowbreak{}In\allowbreak{}Set}
\end{itemize}
\hypertarget{governing-declaration-715d733278a03252}{}
\subsection*{\small\ttfamily\raggedright Applied\allowbreak{}Modeling\allowbreak{}Lib.\allowbreak{}Probability.\allowbreak{}finite\allowbreak{}IID\allowbreak{}Realized\allowbreak{}Compact\allowbreak{}Shell\allowbreak{}Dyadic\allowbreak{}Partition\allowbreak{}Tail}
\textbf{Declaration location:} reusable library
\paragraph{Lean declaration body}
\begin{ReviewVerbatim}
type:
(dimension : ℕ) →
  ℕ →
    MeasureTheory.ProbabilityMeasure (EuclideanSpace ℝ (Fin dimension)) →
      (count : ℕ) → ℕ → ℝ → (Fin count → EuclideanSpace ℝ (Fin dimension)) → Prop

value:
fun (dimension shell : ℕ) (law : MeasureTheory.ProbabilityMeasure (EuclideanSpace ℝ (Fin dimension))) (count level : ℕ)
    (error : ℝ) (a : Fin count → EuclideanSpace ℝ (Fin dimension)) =>
  {sample : Fin count → EuclideanSpace ℝ (Fin dimension) |
      ∃ (_ :
        (AppliedModelingLib.empiricalSampleIndicesInSet sample
            (AppliedModelingLib.Probability.euclideanDyadicCubeShell dimension shell)).Nonempty),
        error ≤
          AppliedModelingLib.Probability.dyadicPartitionL1Deviation dimension
            (AppliedModelingLib.Probability.euclideanDyadicCubeShellCompactLaw dimension shell law)
            (AppliedModelingLib.empiricalSampleIndicesInSet sample
                (AppliedModelingLib.Probability.euclideanDyadicCubeShell dimension shell)).card
            level
            (AppliedModelingLib.Probability.finiteIIDSampleMap
              (AppliedModelingLib.Probability.euclideanDyadicCubeShellScaleExtension dimension shell)
              (AppliedModelingLib.Probability.finiteIIDSampleOnFinsetToFin
                (AppliedModelingLib.empiricalSampleIndicesInSet sample
                  (AppliedModelingLib.Probability.euclideanDyadicCubeShell dimension shell))
                sample))}
    a
\end{ReviewVerbatim}

\paragraph{Direct retained dependencies}
\begin{itemize}\raggedright
\item \hyperlink{governing-declaration-68be765167e02884}{Applied\allowbreak{}Modeling\allowbreak{}Lib.\allowbreak{}Probability.\allowbreak{}Half\allowbreak{}Open\allowbreak{}Unit\allowbreak{}Cube}
\item \hyperlink{governing-declaration-6be421a0d1e9b169}{Applied\allowbreak{}Modeling\allowbreak{}Lib.\allowbreak{}Probability.\allowbreak{}dyadic\allowbreak{}Partition\allowbreak{}L1\allowbreak{}Deviation}
\item \hyperlink{governing-declaration-13cdf9a43259c30a}{Applied\allowbreak{}Modeling\allowbreak{}Lib.\allowbreak{}Probability.\allowbreak{}euclidean\allowbreak{}Dyadic\allowbreak{}Cube\allowbreak{}Shell}
\item \hyperlink{governing-declaration-a45967b04b4e794e}{Applied\allowbreak{}Modeling\allowbreak{}Lib.\allowbreak{}Probability.\allowbreak{}euclidean\allowbreak{}Dyadic\allowbreak{}Cube\allowbreak{}Shell\allowbreak{}Compact\allowbreak{}Law}
\item \hyperlink{governing-declaration-4bcc855202b49a28}{Applied\allowbreak{}Modeling\allowbreak{}Lib.\allowbreak{}Probability.\allowbreak{}euclidean\allowbreak{}Dyadic\allowbreak{}Cube\allowbreak{}Shell\allowbreak{}Scale\allowbreak{}Extension}
\item \hyperlink{governing-declaration-548cac74b010e6ce}{Applied\allowbreak{}Modeling\allowbreak{}Lib.\allowbreak{}Probability.\allowbreak{}finite\allowbreak{}IID\allowbreak{}Sample\allowbreak{}Map}
\item \hyperlink{governing-declaration-dd7cf8954dcd7cb1}{Applied\allowbreak{}Modeling\allowbreak{}Lib.\allowbreak{}Probability.\allowbreak{}finite\allowbreak{}IID\allowbreak{}Sample\allowbreak{}On\allowbreak{}Finset\allowbreak{}To\allowbreak{}Fin}
\item \hyperlink{governing-declaration-3e69c45b89f95dd2}{Applied\allowbreak{}Modeling\allowbreak{}Lib.\allowbreak{}empirical\allowbreak{}Sample\allowbreak{}Indices\allowbreak{}In\allowbreak{}Set}
\end{itemize}
\hypertarget{governing-declaration-e282345978931f95}{}
\subsection*{\small\ttfamily\raggedright Applied\allowbreak{}Modeling\allowbreak{}Lib.\allowbreak{}Probability.\allowbreak{}p\allowbreak{}One\allowbreak{}Fournier\allowbreak{}Guillin\allowbreak{}Selected\allowbreak{}Shell\allowbreak{}Partition\allowbreak{}Certificate}
\textbf{Declaration location:} reusable library
\paragraph{Lean declaration body}
\begin{ReviewVerbatim}
type:
(dimension : ℕ) →
  MeasureTheory.ProbabilityMeasure (EuclideanSpace ℝ (Fin dimension)) →
    (count : ℕ) → ℝ → ℝ → (ℕ → ℕ) → ℝ → (Fin count → EuclideanSpace ℝ (Fin dimension)) → Prop

value:
fun (dimension : ℕ) (law : MeasureTheory.ProbabilityMeasure (EuclideanSpace ℝ (Fin dimension))) (count : ℕ)
    (eta deviation : ℝ) (depth : ℕ → ℕ) (confidence : ℝ) (a : Fin count → EuclideanSpace ℝ (Fin dimension)) =>
  ((⋃ (shell : ℕ),
        AppliedModelingLib.Probability.finiteIIDEuclideanDyadicCubeShellMassRelativeTailTotalized dimension law count
          shell
          (AppliedModelingLib.Probability.pOneFournierGuillinShellToleranceWithEta eta deviation
            ((↑law).real (AppliedModelingLib.Probability.euclideanDyadicCubeShell dimension shell)) shell)) ∪
      ⋃ (shell : ℕ),
        ⋃ level ∈ Finset.range (depth shell),
          AppliedModelingLib.Probability.finiteIIDRealizedCompactShellDyadicPartitionTail dimension shell law count
            (level + 1)
            (AppliedModelingLib.Probability.selectedShellDyadicGeometricConfidenceThreshold dimension count (level + 1)
              ((↑law).real (AppliedModelingLib.Probability.euclideanDyadicCubeShell dimension shell))
              (AppliedModelingLib.Probability.selectedShellGeometricConfidence confidence shell)))
    a
\end{ReviewVerbatim}

\paragraph{Direct retained dependencies}
\begin{itemize}\raggedright
\item \hyperlink{governing-declaration-13cdf9a43259c30a}{Applied\allowbreak{}Modeling\allowbreak{}Lib.\allowbreak{}Probability.\allowbreak{}euclidean\allowbreak{}Dyadic\allowbreak{}Cube\allowbreak{}Shell}
\item \hyperlink{governing-declaration-dbd328d5a293e61c}{Applied\allowbreak{}Modeling\allowbreak{}Lib.\allowbreak{}Probability.\allowbreak{}finite\allowbreak{}IID\allowbreak{}Euclidean\allowbreak{}Dyadic\allowbreak{}Cube\allowbreak{}Shell\allowbreak{}Mass\allowbreak{}Relative\allowbreak{}Tail\allowbreak{}Totalized}
\item \hyperlink{governing-declaration-715d733278a03252}{Applied\allowbreak{}Modeling\allowbreak{}Lib.\allowbreak{}Probability.\allowbreak{}finite\allowbreak{}IID\allowbreak{}Realized\allowbreak{}Compact\allowbreak{}Shell\allowbreak{}Dyadic\allowbreak{}Partition\allowbreak{}Tail}
\item \hyperlink{governing-declaration-1b744d1c0d378e63}{Applied\allowbreak{}Modeling\allowbreak{}Lib.\allowbreak{}Probability.\allowbreak{}p\allowbreak{}One\allowbreak{}Fournier\allowbreak{}Guillin\allowbreak{}Shell\allowbreak{}Tolerance\allowbreak{}With\allowbreak{}Eta}
\item \hyperlink{governing-declaration-c803849fcd834479}{Applied\allowbreak{}Modeling\allowbreak{}Lib.\allowbreak{}Probability.\allowbreak{}selected\allowbreak{}Shell\allowbreak{}Dyadic\allowbreak{}Geometric\allowbreak{}Confidence\allowbreak{}Threshold}
\item \hyperlink{governing-declaration-d481e01804396f20}{Applied\allowbreak{}Modeling\allowbreak{}Lib.\allowbreak{}Probability.\allowbreak{}selected\allowbreak{}Shell\allowbreak{}Geometric\allowbreak{}Confidence}
\end{itemize}
\hypertarget{governing-declaration-11c3706493ab6c46}{}
\subsection*{\small\ttfamily\raggedright Applied\allowbreak{}Modeling\allowbreak{}Lib.\allowbreak{}empirical\allowbreak{}Sample\allowbreak{}Law}
\textbf{Declaration location:} reusable library
\paragraph{Lean declaration body}
\begin{ReviewVerbatim}
type:
{Sample : Type u_1} → {count : ℕ} → [Nonempty (Fin count)] → (Fin count → Sample) → PMF Sample

value:
fun {Sample : Type u_1} {count : ℕ} [Nonempty (Fin count)] (sample : Fin count → Sample) =>
  PMF.map sample (PMF.uniformOfFintype (Fin count))
\end{ReviewVerbatim}


\hypertarget{governing-declaration-d849e10aed4c8dc0}{}
\subsection*{\small\ttfamily\raggedright Applied\allowbreak{}Modeling\allowbreak{}Lib.\allowbreak{}empirical\allowbreak{}Sample\allowbreak{}Probability\allowbreak{}Measure}
\textbf{Declaration location:} reusable library
\paragraph{Lean declaration body}
\begin{ReviewVerbatim}
type:
{Sample : Type u_1} →
  [inst : MeasurableSpace Sample] →
    {count : ℕ} → [Nonempty (Fin count)] → (Fin count → Sample) → MeasureTheory.ProbabilityMeasure Sample

value:
fun {Sample : Type u_1} [MeasurableSpace Sample] {count : ℕ} [Nonempty (Fin count)] (sample : Fin count → Sample) =>
  ⟨(AppliedModelingLib.empiricalSampleLaw sample).toMeasure,
    AppliedModelingLib.empiricalSampleProbabilityMeasure._proof_1 sample⟩
\end{ReviewVerbatim}

\paragraph{Direct retained dependencies}
\begin{itemize}\raggedright
\item \hyperlink{governing-declaration-11c3706493ab6c46}{Applied\allowbreak{}Modeling\allowbreak{}Lib.\allowbreak{}empirical\allowbreak{}Sample\allowbreak{}Law}
\end{itemize}
\hypertarget{governing-declaration-cbd1f5f969456e5f}{}
\subsection*{\small\ttfamily\raggedright Applied\allowbreak{}Modeling\allowbreak{}Lib.\allowbreak{}empirical\allowbreak{}Sample\allowbreak{}Probability\allowbreak{}Measure\allowbreak{}Of\allowbreak{}Pos}
\textbf{Declaration location:} reusable library
\paragraph{Lean declaration body}
\begin{ReviewVerbatim}
type:
{Sample : Type u_1} →
  [inst : MeasurableSpace Sample] →
    {count : ℕ} → 0 < count → (Fin count → Sample) → MeasureTheory.ProbabilityMeasure Sample

value:
fun {Sample : Type u_1} [MeasurableSpace Sample] {count : ℕ} (hcount : 0 < count) (sample : Fin count → Sample) =>
  AppliedModelingLib.empiricalSampleProbabilityMeasure sample
\end{ReviewVerbatim}

\paragraph{Direct retained dependencies}
\begin{itemize}\raggedright
\item \hyperlink{governing-declaration-d849e10aed4c8dc0}{Applied\allowbreak{}Modeling\allowbreak{}Lib.\allowbreak{}empirical\allowbreak{}Sample\allowbreak{}Probability\allowbreak{}Measure}
\end{itemize}
\hypertarget{governing-declaration-92d1582063c434f3}{}
\subsection*{\small\ttfamily\raggedright Applied\allowbreak{}Modeling\allowbreak{}Lib.\allowbreak{}heterogeneous\allowbreak{}Batch\allowbreak{}Trace\allowbreak{}Coordinate\allowbreak{}Measurable\allowbreak{}Space}
\textbf{Declaration location:} reusable library
\paragraph{Lean declaration body}
\begin{ReviewVerbatim}
type:
(BatchIndex : ℕ → Type u_1) →
  (Data : Type u_2) →
    [MeasurableSpace Data] →
      (round : ℕ) → MeasurableSpace (AppliedModelingLib.HeterogeneousBatchTraceCoordinate BatchIndex Data round)

value:
fun (BatchIndex : ℕ → Type u_1) (Data : Type u_2) [MeasurableSpace Data] (round : ℕ) =>
  Nat.casesAuxOn (motive := fun (a : ℕ) =>
    round = a → MeasurableSpace (AppliedModelingLib.HeterogeneousBatchTraceCoordinate BatchIndex Data round)) round
    (fun (h : round = 0) =>
      AppliedModelingLib.heterogeneousBatchTraceCoordinateFintype._proof_1 round h ▸ id inferInstance)
    (fun (n : ℕ) (h : round = n + 1) =>
      AppliedModelingLib.heterogeneousBatchTraceCoordinateFintype._proof_2 round n h ▸ id inferInstance)
    (Eq.refl round)
\end{ReviewVerbatim}

\paragraph{Direct retained dependencies}
\begin{itemize}\raggedright
\item \hyperlink{governing-declaration-37d7ae049de215a7}{Applied\allowbreak{}Modeling\allowbreak{}Lib.\allowbreak{}Heterogeneous\allowbreak{}Batch\allowbreak{}Trace\allowbreak{}Coordinate}
\end{itemize}
\hypertarget{governing-declaration-5aee541574e15cfb}{}
\subsection*{\small\ttfamily\raggedright Applied\allowbreak{}Modeling\allowbreak{}Lib.\allowbreak{}heterogeneous\allowbreak{}Batch\allowbreak{}Trace\allowbreak{}Infinite\allowbreak{}Law}
\textbf{Declaration location:} reusable library
\paragraph{Lean declaration body}
\begin{ReviewVerbatim}
type:
{Data : Type u_1} →
  [inst : MeasurableSpace Data] →
    (BatchIndex : ℕ → Type u_2) →
      (kernel :
          (iteration : ℕ) →
            ProbabilityTheory.Kernel
              ((i : ↥(Finset.Iic iteration)) → AppliedModelingLib.HeterogeneousBatchTraceCoordinate BatchIndex Data ↑i)
              (AppliedModelingLib.HeterogeneousBatchTraceCoordinate BatchIndex Data (iteration + 1))) →
        [∀ (iteration : ℕ), ProbabilityTheory.IsMarkovKernel (kernel iteration)] →
          MeasureTheory.Measure ((i : ℕ) → AppliedModelingLib.HeterogeneousBatchTraceCoordinate BatchIndex Data i)

value:
fun {Data : Type u_1} [MeasurableSpace Data] (BatchIndex : ℕ → Type u_2)
    (kernel :
      (iteration : ℕ) →
        ProbabilityTheory.Kernel
          ((i : ↥(Finset.Iic iteration)) → AppliedModelingLib.HeterogeneousBatchTraceCoordinate BatchIndex Data ↑i)
          (AppliedModelingLib.HeterogeneousBatchTraceCoordinate BatchIndex Data (iteration + 1)))
    [∀ (iteration : ℕ), ProbabilityTheory.IsMarkovKernel (kernel iteration)] =>
  ProbabilityTheory.Kernel.trajMeasure (MeasureTheory.Measure.dirac PUnit.unit) kernel
\end{ReviewVerbatim}

\paragraph{Direct retained dependencies}
\begin{itemize}\raggedright
\item \hyperlink{governing-declaration-37d7ae049de215a7}{Applied\allowbreak{}Modeling\allowbreak{}Lib.\allowbreak{}Heterogeneous\allowbreak{}Batch\allowbreak{}Trace\allowbreak{}Coordinate}
\item \hyperlink{governing-declaration-92d1582063c434f3}{Applied\allowbreak{}Modeling\allowbreak{}Lib.\allowbreak{}heterogeneous\allowbreak{}Batch\allowbreak{}Trace\allowbreak{}Coordinate\allowbreak{}Measurable\allowbreak{}Space}
\end{itemize}
\hypertarget{governing-declaration-caf4e7d2e1526646}{}
\subsection*{\small\ttfamily\raggedright Applied\allowbreak{}Modeling\allowbreak{}Lib.\allowbreak{}heterogeneous\allowbreak{}Batch\allowbreak{}Trace\allowbreak{}Measurable\allowbreak{}Space}
\textbf{Declaration location:} reusable library
\paragraph{Lean declaration body}
\begin{ReviewVerbatim}
type:
(BatchIndex : ℕ → Type u_1) →
  (Data : Type u_2) →
    [MeasurableSpace Data] →
      (iteration : ℕ) → MeasurableSpace (AppliedModelingLib.HeterogeneousBatchTrace BatchIndex Data iteration)

value:
fun (BatchIndex : ℕ → Type u_1) (Data : Type u_2) [MeasurableSpace Data] (iteration : ℕ) => id inferInstance
\end{ReviewVerbatim}

\paragraph{Direct retained dependencies}
\begin{itemize}\raggedright
\item \hyperlink{governing-declaration-307f711c4bf618c0}{Applied\allowbreak{}Modeling\allowbreak{}Lib.\allowbreak{}Heterogeneous\allowbreak{}Batch\allowbreak{}Trace}
\end{itemize}
\hypertarget{governing-declaration-644c195e65179f85}{}
\subsection*{\small\ttfamily\raggedright Applied\allowbreak{}Modeling\allowbreak{}Lib.\allowbreak{}heterogeneous\allowbreak{}Batch\allowbreak{}Trace\allowbreak{}Of\allowbreak{}Infinite\allowbreak{}Trace}
\textbf{Declaration location:} reusable library
\paragraph{Lean declaration body}
\begin{ReviewVerbatim}
type:
{BatchIndex : ℕ → Type u_1} →
  {Data : Type u_2} →
    (iteration : ℕ) →
      ((i : ℕ) → AppliedModelingLib.HeterogeneousBatchTraceCoordinate BatchIndex Data i) →
        (round : Fin iteration) → BatchIndex ↑round → Data

value:
fun {BatchIndex : ℕ → Type u_1} {Data : Type u_2} (iteration : ℕ)
    (trace : (i : ℕ) → AppliedModelingLib.HeterogeneousBatchTraceCoordinate BatchIndex Data i) (round : Fin iteration)
    (a : BatchIndex ↑round) =>
  trace (↑round + 1) a
\end{ReviewVerbatim}

\paragraph{Direct retained dependencies}
\begin{itemize}\raggedright
\item \hyperlink{governing-declaration-37d7ae049de215a7}{Applied\allowbreak{}Modeling\allowbreak{}Lib.\allowbreak{}Heterogeneous\allowbreak{}Batch\allowbreak{}Trace\allowbreak{}Coordinate}
\end{itemize}
\hypertarget{governing-declaration-2dac435c3947cb3b}{}
\subsection*{\small\ttfamily\raggedright Applied\allowbreak{}Modeling\allowbreak{}Lib.\allowbreak{}heterogeneous\allowbreak{}Batch\allowbreak{}Trace\allowbreak{}Of\allowbreak{}Prefix}
\textbf{Declaration location:} reusable library
\paragraph{Lean declaration body}
\begin{ReviewVerbatim}
type:
{BatchIndex : ℕ → Type u_1} →
  {Data : Type u_2} →
    (iteration : ℕ) →
      ((i : ↥(Finset.Iic iteration)) → AppliedModelingLib.HeterogeneousBatchTraceCoordinate BatchIndex Data ↑i) →
        (round : Fin iteration) → BatchIndex ↑round → Data

value:
fun {BatchIndex : ℕ → Type u_1} {Data : Type u_2} (iteration : ℕ)
    (historyPrefix :
      (i : ↥(Finset.Iic iteration)) → AppliedModelingLib.HeterogeneousBatchTraceCoordinate BatchIndex Data ↑i)
    (round : Fin iteration) (a : BatchIndex ↑round) =>
  historyPrefix ⟨↑round + 1, AppliedModelingLib.heterogeneousBatchTraceOfPrefix._proof_1 iteration round⟩ a
\end{ReviewVerbatim}

\paragraph{Direct retained dependencies}
\begin{itemize}\raggedright
\item \hyperlink{governing-declaration-37d7ae049de215a7}{Applied\allowbreak{}Modeling\allowbreak{}Lib.\allowbreak{}Heterogeneous\allowbreak{}Batch\allowbreak{}Trace\allowbreak{}Coordinate}
\end{itemize}
\hypertarget{governing-declaration-68ef005103543934}{}
\subsection*{\small\ttfamily\raggedright Applied\allowbreak{}Modeling\allowbreak{}Lib.\allowbreak{}heterogeneous\allowbreak{}Batch\allowbreak{}Trace\allowbreak{}Bad\allowbreak{}Event\allowbreak{}Pair}
\textbf{Declaration location:} reusable library
\paragraph{Lean declaration body}
\begin{ReviewVerbatim}
type:
{Data : Type u_1} →
  (BatchIndex : ℕ → Type u_2) →
    ((iteration : ℕ) →
        AppliedModelingLib.HeterogeneousBatchTrace BatchIndex Data iteration → (BatchIndex iteration → Data) → Prop) →
      (iteration : ℕ) →
        ((i : ↥(Finset.Iic iteration)) → AppliedModelingLib.HeterogeneousBatchTraceCoordinate BatchIndex Data ↑i) ×
            AppliedModelingLib.HeterogeneousBatchTraceCoordinate BatchIndex Data (iteration + 1) →
          Prop

value:
fun {Data : Type u_1} (BatchIndex : ℕ → Type u_2)
    (bad :
      (iteration : ℕ) →
        AppliedModelingLib.HeterogeneousBatchTrace BatchIndex Data iteration → (BatchIndex iteration → Data) → Prop)
    (iteration : ℕ)
    (a :
      ((i : ↥(Finset.Iic iteration)) → AppliedModelingLib.HeterogeneousBatchTraceCoordinate BatchIndex Data ↑i) ×
        AppliedModelingLib.HeterogeneousBatchTraceCoordinate BatchIndex Data (iteration + 1)) =>
  {pair :
      ((i : ↥(Finset.Iic iteration)) → AppliedModelingLib.HeterogeneousBatchTraceCoordinate BatchIndex Data ↑i) ×
        AppliedModelingLib.HeterogeneousBatchTraceCoordinate BatchIndex Data (iteration + 1) |
      bad iteration (AppliedModelingLib.heterogeneousBatchTraceOfPrefix iteration pair.1) pair.2}
    a
\end{ReviewVerbatim}

\paragraph{Direct retained dependencies}
\begin{itemize}\raggedright
\item \hyperlink{governing-declaration-307f711c4bf618c0}{Applied\allowbreak{}Modeling\allowbreak{}Lib.\allowbreak{}Heterogeneous\allowbreak{}Batch\allowbreak{}Trace}
\item \hyperlink{governing-declaration-37d7ae049de215a7}{Applied\allowbreak{}Modeling\allowbreak{}Lib.\allowbreak{}Heterogeneous\allowbreak{}Batch\allowbreak{}Trace\allowbreak{}Coordinate}
\item \hyperlink{governing-declaration-2dac435c3947cb3b}{Applied\allowbreak{}Modeling\allowbreak{}Lib.\allowbreak{}heterogeneous\allowbreak{}Batch\allowbreak{}Trace\allowbreak{}Of\allowbreak{}Prefix}
\end{itemize}
\hypertarget{governing-declaration-cb6bd2e9984f0d33}{}
\subsection*{\small\ttfamily\raggedright Applied\allowbreak{}Modeling\allowbreak{}Lib.\allowbreak{}heterogeneous\allowbreak{}Batch\allowbreak{}Trace\allowbreak{}One\allowbreak{}Observation\allowbreak{}Kernel}
\textbf{Declaration location:} reusable library
\paragraph{Lean declaration body}
\begin{ReviewVerbatim}
type:
{State : Type u_1} →
  {Data : Type u_2} →
    [inst : MeasurableSpace State] →
      [inst_1 : MeasurableSpace Data] →
        (BatchIndex : ℕ → Type u_3) →
          ProbabilityTheory.Kernel State Data →
            (stateOfHistory :
                (iteration : ℕ) → AppliedModelingLib.HeterogeneousBatchTrace BatchIndex Data iteration → State) →
              (iteration : ℕ) →
                Measurable (stateOfHistory iteration) →
                  ProbabilityTheory.Kernel
                    ((i : ↥(Finset.Iic iteration)) →
                      AppliedModelingLib.HeterogeneousBatchTraceCoordinate BatchIndex Data ↑i)
                    Data

value:
fun {State : Type u_1} {Data : Type u_2} [MeasurableSpace State] [MeasurableSpace Data] (BatchIndex : ℕ → Type u_3)
    (dataKernel : ProbabilityTheory.Kernel State Data)
    (stateOfHistory : (iteration : ℕ) → AppliedModelingLib.HeterogeneousBatchTrace BatchIndex Data iteration → State)
    (iteration : ℕ) (hstate : Measurable (stateOfHistory iteration)) =>
  dataKernel.comap
    (fun
        (historyPrefix :
          (i : ↥(Finset.Iic iteration)) → AppliedModelingLib.HeterogeneousBatchTraceCoordinate BatchIndex Data ↑i) =>
      stateOfHistory iteration (AppliedModelingLib.heterogeneousBatchTraceOfPrefix iteration historyPrefix))
    (AppliedModelingLib.heterogeneousBatchTraceOneObservationKernel._proof_1 BatchIndex stateOfHistory iteration hstate)
\end{ReviewVerbatim}

\paragraph{Direct retained dependencies}
\begin{itemize}\raggedright
\item \hyperlink{governing-declaration-307f711c4bf618c0}{Applied\allowbreak{}Modeling\allowbreak{}Lib.\allowbreak{}Heterogeneous\allowbreak{}Batch\allowbreak{}Trace}
\item \hyperlink{governing-declaration-37d7ae049de215a7}{Applied\allowbreak{}Modeling\allowbreak{}Lib.\allowbreak{}Heterogeneous\allowbreak{}Batch\allowbreak{}Trace\allowbreak{}Coordinate}
\item \hyperlink{governing-declaration-92d1582063c434f3}{Applied\allowbreak{}Modeling\allowbreak{}Lib.\allowbreak{}heterogeneous\allowbreak{}Batch\allowbreak{}Trace\allowbreak{}Coordinate\allowbreak{}Measurable\allowbreak{}Space}
\item \hyperlink{governing-declaration-caf4e7d2e1526646}{Applied\allowbreak{}Modeling\allowbreak{}Lib.\allowbreak{}heterogeneous\allowbreak{}Batch\allowbreak{}Trace\allowbreak{}Measurable\allowbreak{}Space}
\item \hyperlink{governing-declaration-2dac435c3947cb3b}{Applied\allowbreak{}Modeling\allowbreak{}Lib.\allowbreak{}heterogeneous\allowbreak{}Batch\allowbreak{}Trace\allowbreak{}Of\allowbreak{}Prefix}
\end{itemize}
\hypertarget{governing-declaration-73f8e033eb247305}{}
\subsection*{\small\ttfamily\raggedright Applied\allowbreak{}Modeling\allowbreak{}Lib.\allowbreak{}heterogeneous\allowbreak{}Batch\allowbreak{}Trace\allowbreak{}Snoc}
\textbf{Declaration location:} reusable library
\paragraph{Lean declaration body}
\begin{ReviewVerbatim}
type:
{BatchIndex : ℕ → Type u_1} →
  {Data : Type u_2} →
    {iteration : ℕ} →
      AppliedModelingLib.HeterogeneousBatchTrace BatchIndex Data iteration →
        (BatchIndex iteration → Data) → (round : Fin (iteration + 1)) → BatchIndex ↑round → Data

value:
fun {BatchIndex : ℕ → Type u_1} {Data : Type u_2} {iteration : ℕ}
    (history : AppliedModelingLib.HeterogeneousBatchTrace BatchIndex Data iteration)
    (batch : BatchIndex iteration → Data) (round : Fin (iteration + 1)) (a : BatchIndex ↑round) =>
  Fin.lastCases (motive := fun (x : Fin (iteration + 1)) => BatchIndex ↑x → Data) batch
    (fun (previous : Fin iteration) => history previous) round a
\end{ReviewVerbatim}

\paragraph{Direct retained dependencies}
\begin{itemize}\raggedright
\item \hyperlink{governing-declaration-307f711c4bf618c0}{Applied\allowbreak{}Modeling\allowbreak{}Lib.\allowbreak{}Heterogeneous\allowbreak{}Batch\allowbreak{}Trace}
\end{itemize}
\hypertarget{governing-declaration-e97175bd0a346777}{}
\subsection*{\small\ttfamily\raggedright Applied\allowbreak{}Modeling\allowbreak{}Lib.\allowbreak{}Is\allowbreak{}Measure\allowbreak{}Performatively\allowbreak{}Stable}
\textbf{Declaration location:} reusable library
\paragraph{Lean declaration body}
\begin{ReviewVerbatim}
type:
{Parameter : Type u_1} →
  {Data : Type u_2} →
    [inst : MeasurableSpace Data] → AppliedModelingLib.MeasurePerformativeModel Parameter Data → Parameter → Prop

value:
fun {Parameter : Type u_1} {Data : Type u_2} [MeasurableSpace Data]
    (model : AppliedModelingLib.MeasurePerformativeModel Parameter Data) (parameter : Parameter) =>
  ∀ (candidate : Parameter),
    AppliedModelingLib.measureDecoupledPerformativeRisk model parameter parameter ≤
      AppliedModelingLib.measureDecoupledPerformativeRisk model parameter candidate
\end{ReviewVerbatim}

\paragraph{Direct retained dependencies}
\begin{itemize}\raggedright
\item \hyperlink{library-prerequisite-96553d4fbb82f1ac}{Applied\allowbreak{}Modeling\allowbreak{}Lib.\allowbreak{}Measure\allowbreak{}Performative\allowbreak{}Model}
\item \hyperlink{library-prerequisite-d4de05e57b5947f0}{Applied\allowbreak{}Modeling\allowbreak{}Lib.\allowbreak{}measure\allowbreak{}Decoupled\allowbreak{}Performative\allowbreak{}Risk}
\end{itemize}
\hypertarget{governing-declaration-522b06e211b1e40a}{}
\subsection*{\small\ttfamily\raggedright Applied\allowbreak{}Modeling\allowbreak{}Lib.\allowbreak{}Is\allowbreak{}Measure\allowbreak{}Performatively\allowbreak{}Stable\allowbreak{}On}
\textbf{Declaration location:} reusable library
\paragraph{Lean declaration body}
\begin{ReviewVerbatim}
type:
{Parameter : Type u_1} →
  {Data : Type u_2} →
    [inst : MeasurableSpace Data] →
      AppliedModelingLib.MeasurePerformativeModel Parameter Data → Set Parameter → Parameter → Prop

value:
fun {Parameter : Type u_1} {Data : Type u_2} [MeasurableSpace Data]
    (model : AppliedModelingLib.MeasurePerformativeModel Parameter Data) (domain : Set Parameter)
    (parameter : Parameter) =>
  parameter ∈ domain ∧
    ∀ candidate ∈ domain,
      AppliedModelingLib.measureDecoupledPerformativeRisk model parameter parameter ≤
        AppliedModelingLib.measureDecoupledPerformativeRisk model parameter candidate
\end{ReviewVerbatim}

\paragraph{Direct retained dependencies}
\begin{itemize}\raggedright
\item \hyperlink{library-prerequisite-96553d4fbb82f1ac}{Applied\allowbreak{}Modeling\allowbreak{}Lib.\allowbreak{}Measure\allowbreak{}Performative\allowbreak{}Model}
\item \hyperlink{library-prerequisite-d4de05e57b5947f0}{Applied\allowbreak{}Modeling\allowbreak{}Lib.\allowbreak{}measure\allowbreak{}Decoupled\allowbreak{}Performative\allowbreak{}Risk}
\end{itemize}
\hypertarget{governing-declaration-916f213e8c73a94a}{}
\subsection*{\small\ttfamily\raggedright Applied\allowbreak{}Modeling\allowbreak{}Lib.\allowbreak{}Is\allowbreak{}Measure\allowbreak{}Repeated\allowbreak{}Risk\allowbreak{}Minimization}
\textbf{Declaration location:} reusable library
\paragraph{Lean declaration body}
\begin{ReviewVerbatim}
type:
{Parameter : Type u_1} →
  {Data : Type u_2} →
    [inst : MeasurableSpace Data] →
      AppliedModelingLib.MeasurePerformativeModel Parameter Data → (Parameter → Parameter) → Prop

value:
fun {Parameter : Type u_1} {Data : Type u_2} [MeasurableSpace Data]
    (model : AppliedModelingLib.MeasurePerformativeModel Parameter Data) (update : Parameter → Parameter) =>
  ∀ (deployed candidate : Parameter),
    AppliedModelingLib.measureDecoupledPerformativeRisk model deployed (update deployed) ≤
      AppliedModelingLib.measureDecoupledPerformativeRisk model deployed candidate
\end{ReviewVerbatim}

\paragraph{Direct retained dependencies}
\begin{itemize}\raggedright
\item \hyperlink{library-prerequisite-96553d4fbb82f1ac}{Applied\allowbreak{}Modeling\allowbreak{}Lib.\allowbreak{}Measure\allowbreak{}Performative\allowbreak{}Model}
\item \hyperlink{library-prerequisite-d4de05e57b5947f0}{Applied\allowbreak{}Modeling\allowbreak{}Lib.\allowbreak{}measure\allowbreak{}Decoupled\allowbreak{}Performative\allowbreak{}Risk}
\end{itemize}
\hypertarget{governing-declaration-6ae86e5703739171}{}
\subsection*{\small\ttfamily\raggedright Applied\allowbreak{}Modeling\allowbreak{}Lib.\allowbreak{}Is\allowbreak{}Measure\allowbreak{}Repeated\allowbreak{}Risk\allowbreak{}Minimization\allowbreak{}On}
\textbf{Declaration location:} reusable library
\paragraph{Lean declaration body}
\begin{ReviewVerbatim}
type:
{Parameter : Type u_1} →
  {Data : Type u_2} →
    [inst : MeasurableSpace Data] →
      AppliedModelingLib.MeasurePerformativeModel Parameter Data → Set Parameter → (Parameter → Parameter) → Prop

value:
fun {Parameter : Type u_1} {Data : Type u_2} [MeasurableSpace Data]
    (model : AppliedModelingLib.MeasurePerformativeModel Parameter Data) (domain : Set Parameter)
    (update : Parameter → Parameter) =>
  ∀ deployed ∈ domain,
    update deployed ∈ domain ∧
      ∀ candidate ∈ domain,
        AppliedModelingLib.measureDecoupledPerformativeRisk model deployed (update deployed) ≤
          AppliedModelingLib.measureDecoupledPerformativeRisk model deployed candidate
\end{ReviewVerbatim}

\paragraph{Direct retained dependencies}
\begin{itemize}\raggedright
\item \hyperlink{library-prerequisite-96553d4fbb82f1ac}{Applied\allowbreak{}Modeling\allowbreak{}Lib.\allowbreak{}Measure\allowbreak{}Performative\allowbreak{}Model}
\item \hyperlink{library-prerequisite-d4de05e57b5947f0}{Applied\allowbreak{}Modeling\allowbreak{}Lib.\allowbreak{}measure\allowbreak{}Decoupled\allowbreak{}Performative\allowbreak{}Risk}
\end{itemize}
\hypertarget{governing-declaration-3e758dff81ee2431}{}
\subsection*{\small\ttfamily\raggedright Applied\allowbreak{}Modeling\allowbreak{}Lib.\allowbreak{}measure\allowbreak{}Empirical\allowbreak{}Decoupled\allowbreak{}Performative\allowbreak{}Risk}
\textbf{Declaration location:} reusable library
\paragraph{Lean declaration body}
\begin{ReviewVerbatim}
type:
{Parameter : Type u_1} →
  {Data : Type u_2} →
    [inst : MeasurableSpace Data] →
      AppliedModelingLib.MeasurePerformativeModel Parameter Data → MeasureTheory.ProbabilityMeasure Data → Parameter → ℝ

value:
fun {Parameter : Type u_1} {Data : Type u_2} [MeasurableSpace Data]
    (model : AppliedModelingLib.MeasurePerformativeModel Parameter Data)
    (empiricalLaw : MeasureTheory.ProbabilityMeasure Data) (evaluatedParameter : Parameter) =>
  ∫ (datum : Data), model.loss datum evaluatedParameter ∂↑empiricalLaw
\end{ReviewVerbatim}

\paragraph{Direct retained dependencies}
\begin{itemize}\raggedright
\item \hyperlink{library-prerequisite-96553d4fbb82f1ac}{Applied\allowbreak{}Modeling\allowbreak{}Lib.\allowbreak{}Measure\allowbreak{}Performative\allowbreak{}Model}
\end{itemize}
\hypertarget{governing-declaration-abd83945ceff04f5}{}
\subsection*{\small\ttfamily\raggedright Applied\allowbreak{}Modeling\allowbreak{}Lib.\allowbreak{}Is\allowbreak{}Measure\allowbreak{}Empirical\allowbreak{}Risk\allowbreak{}Minimizer\allowbreak{}On}
\textbf{Declaration location:} reusable library
\paragraph{Lean declaration body}
\begin{ReviewVerbatim}
type:
{Parameter : Type u_1} →
  {Data : Type u_2} →
    [inst : MeasurableSpace Data] →
      AppliedModelingLib.MeasurePerformativeModel Parameter Data →
        Set Parameter → MeasureTheory.ProbabilityMeasure Data → Parameter → Prop

value:
fun {Parameter : Type u_1} {Data : Type u_2} [MeasurableSpace Data]
    (model : AppliedModelingLib.MeasurePerformativeModel Parameter Data) (domain : Set Parameter)
    (empiricalLaw : MeasureTheory.ProbabilityMeasure Data) (minimizer : Parameter) =>
  minimizer ∈ domain ∧
    ∀ candidate ∈ domain,
      AppliedModelingLib.measureEmpiricalDecoupledPerformativeRisk model empiricalLaw minimizer ≤
        AppliedModelingLib.measureEmpiricalDecoupledPerformativeRisk model empiricalLaw candidate
\end{ReviewVerbatim}

\paragraph{Direct retained dependencies}
\begin{itemize}\raggedright
\item \hyperlink{library-prerequisite-96553d4fbb82f1ac}{Applied\allowbreak{}Modeling\allowbreak{}Lib.\allowbreak{}Measure\allowbreak{}Performative\allowbreak{}Model}
\item \hyperlink{governing-declaration-3e758dff81ee2431}{Applied\allowbreak{}Modeling\allowbreak{}Lib.\allowbreak{}measure\allowbreak{}Empirical\allowbreak{}Decoupled\allowbreak{}Performative\allowbreak{}Risk}
\end{itemize}
\hypertarget{governing-declaration-50fa794af5bef559}{}
\subsection*{\small\ttfamily\raggedright Applied\allowbreak{}Modeling\allowbreak{}Lib.\allowbreak{}Is\allowbreak{}Heterogeneous\allowbreak{}Measure\allowbreak{}Repeated\allowbreak{}Empirical\allowbreak{}Risk\allowbreak{}Minimization\allowbreak{}Trajectory\allowbreak{}On}
\textbf{Declaration location:} reusable library
\paragraph{Lean declaration body}
\begin{ReviewVerbatim}
type:
{Parameter : Type u_1} →
  {Data : Type u_2} →
    [inst : MeasurableSpace Data] →
      AppliedModelingLib.MeasurePerformativeModel Parameter Data →
        Set Parameter →
          (sampleCount : ℕ → ℕ) →
            (∀ (iteration : ℕ), 0 < sampleCount iteration) →
              ((iteration : ℕ) →
                  AppliedModelingLib.HeterogeneousBatchTrace (fun (round : ℕ) => Fin (sampleCount round)) Data
                      iteration →
                    Parameter) →
                Prop

value:
fun {Parameter : Type u_1} {Data : Type u_2} [MeasurableSpace Data]
    (model : AppliedModelingLib.MeasurePerformativeModel Parameter Data) (domain : Set Parameter) (sampleCount : ℕ → ℕ)
    (hcountPositive : ∀ (iteration : ℕ), 0 < sampleCount iteration)
    (deployedOfHistory :
      (iteration : ℕ) →
        AppliedModelingLib.HeterogeneousBatchTrace (fun (round : ℕ) => Fin (sampleCount round)) Data iteration →
          Parameter) =>
  ∀ (iteration : ℕ)
    (history : AppliedModelingLib.HeterogeneousBatchTrace (fun (round : ℕ) => Fin (sampleCount round)) Data iteration)
    (batch : Fin (sampleCount iteration) → Data),
    AppliedModelingLib.IsMeasureEmpiricalRiskMinimizerOn model domain
      (AppliedModelingLib.empiricalSampleProbabilityMeasureOfPos (hcountPositive iteration) batch)
      (deployedOfHistory (iteration + 1) (AppliedModelingLib.heterogeneousBatchTraceSnoc history batch))
\end{ReviewVerbatim}

\paragraph{Direct retained dependencies}
\begin{itemize}\raggedright
\item \hyperlink{governing-declaration-307f711c4bf618c0}{Applied\allowbreak{}Modeling\allowbreak{}Lib.\allowbreak{}Heterogeneous\allowbreak{}Batch\allowbreak{}Trace}
\item \hyperlink{governing-declaration-abd83945ceff04f5}{Applied\allowbreak{}Modeling\allowbreak{}Lib.\allowbreak{}Is\allowbreak{}Measure\allowbreak{}Empirical\allowbreak{}Risk\allowbreak{}Minimizer\allowbreak{}On}
\item \hyperlink{library-prerequisite-96553d4fbb82f1ac}{Applied\allowbreak{}Modeling\allowbreak{}Lib.\allowbreak{}Measure\allowbreak{}Performative\allowbreak{}Model}
\item \hyperlink{governing-declaration-cbd1f5f969456e5f}{Applied\allowbreak{}Modeling\allowbreak{}Lib.\allowbreak{}empirical\allowbreak{}Sample\allowbreak{}Probability\allowbreak{}Measure\allowbreak{}Of\allowbreak{}Pos}
\item \hyperlink{governing-declaration-73f8e033eb247305}{Applied\allowbreak{}Modeling\allowbreak{}Lib.\allowbreak{}heterogeneous\allowbreak{}Batch\allowbreak{}Trace\allowbreak{}Snoc}
\end{itemize}
\hypertarget{governing-declaration-0bb19c93c4765c7a}{}
\subsection*{\small\ttfamily\raggedright Applied\allowbreak{}Modeling\allowbreak{}Lib.\allowbreak{}measure\allowbreak{}Empirical\allowbreak{}Population\allowbreak{}Gradient}
\textbf{Declaration location:} reusable library
\paragraph{Lean declaration body}
\begin{ReviewVerbatim}
type:
{Parameter : Type u_1} →
  {Data : Type u_2} →
    [inst : NormedAddCommGroup Parameter] →
      [InnerProductSpace ℝ Parameter] →
        [inst : MeasurableSpace Data] →
          (Data → Parameter → Parameter) → MeasureTheory.ProbabilityMeasure Data → Parameter → Parameter

value:
fun {Parameter : Type u_1} {Data : Type u_2} [NormedAddCommGroup Parameter] [InnerProductSpace ℝ Parameter]
    [MeasurableSpace Data] (gradient : Data → Parameter → Parameter)
    (empiricalLaw : MeasureTheory.ProbabilityMeasure Data) (evaluatedParameter : Parameter) =>
  ∫ (datum : Data), gradient datum evaluatedParameter ∂↑empiricalLaw
\end{ReviewVerbatim}


\hypertarget{governing-declaration-d3dfc78bf5c69cfe}{}
\subsection*{\small\ttfamily\raggedright Applied\allowbreak{}Modeling\allowbreak{}Lib.\allowbreak{}measure\allowbreak{}Frozen\allowbreak{}Population\allowbreak{}Gradient}
\textbf{Declaration location:} reusable library
\paragraph{Lean declaration body}
\begin{ReviewVerbatim}
type:
{Parameter : Type u_1} →
  {Data : Type u_2} →
    [inst : MeasurableSpace Data] →
      [inst_1 : NormedAddCommGroup Parameter] →
        [InnerProductSpace ℝ Parameter] →
          AppliedModelingLib.MeasurePerformativeModel Parameter Data →
            (Data → Parameter → Parameter) → Parameter → Parameter → Parameter

value:
fun {Parameter : Type u_1} {Data : Type u_2} [MeasurableSpace Data] [NormedAddCommGroup Parameter]
    [InnerProductSpace ℝ Parameter] (model : AppliedModelingLib.MeasurePerformativeModel Parameter Data)
    (gradient : Data → Parameter → Parameter) (distributionParameter evaluatedParameter : Parameter) =>
  ∫ (datum : Data), gradient datum evaluatedParameter ∂↑(model.dataLaw distributionParameter)
\end{ReviewVerbatim}

\paragraph{Direct retained dependencies}
\begin{itemize}\raggedright
\item \hyperlink{library-prerequisite-96553d4fbb82f1ac}{Applied\allowbreak{}Modeling\allowbreak{}Lib.\allowbreak{}Measure\allowbreak{}Performative\allowbreak{}Model}
\end{itemize}
\hypertarget{governing-declaration-7bc25db7b9b1aaf4}{}
\subsection*{\small\ttfamily\raggedright Applied\allowbreak{}Modeling\allowbreak{}Lib.\allowbreak{}measure\allowbreak{}Performative\allowbreak{}Deployed\allowbreak{}Law}
\textbf{Declaration location:} reusable library
\paragraph{Lean declaration body}
\begin{ReviewVerbatim}
type:
{Parameter : Type u_1} →
  {Data : Type u_2} →
    [inst : MeasurableSpace Data] →
      AppliedModelingLib.MeasurePerformativeModel Parameter Data →
        (BatchIndex : ℕ → Type u_3) →
          ((iteration : ℕ) → AppliedModelingLib.HeterogeneousBatchTrace BatchIndex Data iteration → Parameter) →
            (iteration : ℕ) →
              AppliedModelingLib.HeterogeneousBatchTrace BatchIndex Data iteration →
                MeasureTheory.ProbabilityMeasure Data

value:
fun {Parameter : Type u_1} {Data : Type u_2} [MeasurableSpace Data]
    (model : AppliedModelingLib.MeasurePerformativeModel Parameter Data) (BatchIndex : ℕ → Type u_3)
    (deployedOfHistory :
      (iteration : ℕ) → AppliedModelingLib.HeterogeneousBatchTrace BatchIndex Data iteration → Parameter)
    (iteration : ℕ) (a : AppliedModelingLib.HeterogeneousBatchTrace BatchIndex Data iteration) =>
  model.dataLaw (deployedOfHistory iteration a)
\end{ReviewVerbatim}

\paragraph{Direct retained dependencies}
\begin{itemize}\raggedright
\item \hyperlink{governing-declaration-307f711c4bf618c0}{Applied\allowbreak{}Modeling\allowbreak{}Lib.\allowbreak{}Heterogeneous\allowbreak{}Batch\allowbreak{}Trace}
\item \hyperlink{library-prerequisite-96553d4fbb82f1ac}{Applied\allowbreak{}Modeling\allowbreak{}Lib.\allowbreak{}Measure\allowbreak{}Performative\allowbreak{}Model}
\end{itemize}
\hypertarget{governing-declaration-2833c0091dd3eed5}{}
\subsection*{\small\ttfamily\raggedright Applied\allowbreak{}Modeling\allowbreak{}Lib.\allowbreak{}measure\allowbreak{}Performative\allowbreak{}Heterogeneous\allowbreak{}Batch\allowbreak{}Trace\allowbreak{}Infinite\allowbreak{}Law}
\textbf{Declaration location:} reusable library
\paragraph{Lean declaration body}
\begin{ReviewVerbatim}
type:
{Parameter : Type u_1} →
  {Data : Type u_2} →
    [inst : MeasurableSpace Parameter] →
      [inst_1 : MeasurableSpace Data] →
        (model : AppliedModelingLib.MeasurePerformativeModel Parameter Data) →
          AppliedModelingLib.MeasurePerformativeSamplingKernel model →
            (sampleCount : ℕ → ℕ) →
              (deployedOfHistory :
                  (iteration : ℕ) →
                    AppliedModelingLib.HeterogeneousBatchTrace (fun (round : ℕ) => Fin (sampleCount round)) Data
                        iteration →
                      Parameter) →
                (∀ (iteration : ℕ), Measurable (deployedOfHistory iteration)) →
                  MeasureTheory.Measure
                    ((i : ℕ) →
                      AppliedModelingLib.HeterogeneousBatchTraceCoordinate (fun (round : ℕ) => Fin (sampleCount round))
                        Data i)

value:
fun {Parameter : Type u_1} {Data : Type u_2} [MeasurableSpace Parameter] [MeasurableSpace Data]
    (model : AppliedModelingLib.MeasurePerformativeModel Parameter Data)
    (sampling : AppliedModelingLib.MeasurePerformativeSamplingKernel model) (sampleCount : ℕ → ℕ)
    (deployedOfHistory :
      (iteration : ℕ) →
        AppliedModelingLib.HeterogeneousBatchTrace (fun (round : ℕ) => Fin (sampleCount round)) Data iteration →
          Parameter)
    (hmeasurableDeployed : ∀ (iteration : ℕ), Measurable (deployedOfHistory iteration)) =>
  AppliedModelingLib.heterogeneousBatchTraceInfiniteLaw (fun (round : ℕ) => Fin (sampleCount round))
    fun (iteration : ℕ) =>
    AppliedModelingLib.Probability.finiteIIDBatchKernel (sampleCount iteration)
      (AppliedModelingLib.heterogeneousBatchTraceOneObservationKernel (fun (round : ℕ) => Fin (sampleCount round))
        sampling.kernel deployedOfHistory iteration (hmeasurableDeployed iteration))
\end{ReviewVerbatim}

\paragraph{Direct retained dependencies}
\begin{itemize}\raggedright
\item \hyperlink{governing-declaration-307f711c4bf618c0}{Applied\allowbreak{}Modeling\allowbreak{}Lib.\allowbreak{}Heterogeneous\allowbreak{}Batch\allowbreak{}Trace}
\item \hyperlink{governing-declaration-37d7ae049de215a7}{Applied\allowbreak{}Modeling\allowbreak{}Lib.\allowbreak{}Heterogeneous\allowbreak{}Batch\allowbreak{}Trace\allowbreak{}Coordinate}
\item \hyperlink{library-prerequisite-96553d4fbb82f1ac}{Applied\allowbreak{}Modeling\allowbreak{}Lib.\allowbreak{}Measure\allowbreak{}Performative\allowbreak{}Model}
\item \hyperlink{governing-declaration-527443e0fb077287}{Applied\allowbreak{}Modeling\allowbreak{}Lib.\allowbreak{}Measure\allowbreak{}Performative\allowbreak{}Sampling\allowbreak{}Kernel}
\item \hyperlink{governing-declaration-8ef9c28a526f8794}{Applied\allowbreak{}Modeling\allowbreak{}Lib.\allowbreak{}Probability.\allowbreak{}finite\allowbreak{}IID\allowbreak{}Batch\allowbreak{}Kernel}
\item \hyperlink{governing-declaration-92d1582063c434f3}{Applied\allowbreak{}Modeling\allowbreak{}Lib.\allowbreak{}heterogeneous\allowbreak{}Batch\allowbreak{}Trace\allowbreak{}Coordinate\allowbreak{}Measurable\allowbreak{}Space}
\item \hyperlink{governing-declaration-5aee541574e15cfb}{Applied\allowbreak{}Modeling\allowbreak{}Lib.\allowbreak{}heterogeneous\allowbreak{}Batch\allowbreak{}Trace\allowbreak{}Infinite\allowbreak{}Law}
\item \hyperlink{governing-declaration-caf4e7d2e1526646}{Applied\allowbreak{}Modeling\allowbreak{}Lib.\allowbreak{}heterogeneous\allowbreak{}Batch\allowbreak{}Trace\allowbreak{}Measurable\allowbreak{}Space}
\item \hyperlink{governing-declaration-cb6bd2e9984f0d33}{Applied\allowbreak{}Modeling\allowbreak{}Lib.\allowbreak{}heterogeneous\allowbreak{}Batch\allowbreak{}Trace\allowbreak{}One\allowbreak{}Observation\allowbreak{}Kernel}
\end{itemize}
\hypertarget{governing-declaration-13fb93228218a38e}{}
\subsection*{\small\ttfamily\raggedright Applied\allowbreak{}Modeling\allowbreak{}Lib.\allowbreak{}measure\allowbreak{}Performative\allowbreak{}Risk}
\textbf{Declaration location:} reusable library
\paragraph{Lean declaration body}
\begin{ReviewVerbatim}
type:
{Parameter : Type u_1} →
  {Data : Type u_2} →
    [inst : MeasurableSpace Data] → AppliedModelingLib.MeasurePerformativeModel Parameter Data → Parameter → ℝ

value:
fun {Parameter : Type u_1} {Data : Type u_2} [MeasurableSpace Data]
    (model : AppliedModelingLib.MeasurePerformativeModel Parameter Data) (parameter : Parameter) =>
  AppliedModelingLib.measureDecoupledPerformativeRisk model parameter parameter
\end{ReviewVerbatim}

\paragraph{Direct retained dependencies}
\begin{itemize}\raggedright
\item \hyperlink{library-prerequisite-96553d4fbb82f1ac}{Applied\allowbreak{}Modeling\allowbreak{}Lib.\allowbreak{}Measure\allowbreak{}Performative\allowbreak{}Model}
\item \hyperlink{library-prerequisite-d4de05e57b5947f0}{Applied\allowbreak{}Modeling\allowbreak{}Lib.\allowbreak{}measure\allowbreak{}Decoupled\allowbreak{}Performative\allowbreak{}Risk}
\end{itemize}
\hypertarget{governing-declaration-bf41b084ebc9aa18}{}
\subsection*{\small\ttfamily\raggedright Applied\allowbreak{}Modeling\allowbreak{}Lib.\allowbreak{}Is\allowbreak{}Measure\allowbreak{}Performatively\allowbreak{}Optimal}
\textbf{Declaration location:} reusable library
\paragraph{Lean declaration body}
\begin{ReviewVerbatim}
type:
{Parameter : Type u_1} →
  {Data : Type u_2} →
    [inst : MeasurableSpace Data] → AppliedModelingLib.MeasurePerformativeModel Parameter Data → Parameter → Prop

value:
fun {Parameter : Type u_1} {Data : Type u_2} [MeasurableSpace Data]
    (model : AppliedModelingLib.MeasurePerformativeModel Parameter Data) (parameter : Parameter) =>
  ∀ (candidate : Parameter),
    AppliedModelingLib.measurePerformativeRisk model parameter ≤
      AppliedModelingLib.measurePerformativeRisk model candidate
\end{ReviewVerbatim}

\paragraph{Direct retained dependencies}
\begin{itemize}\raggedright
\item \hyperlink{library-prerequisite-96553d4fbb82f1ac}{Applied\allowbreak{}Modeling\allowbreak{}Lib.\allowbreak{}Measure\allowbreak{}Performative\allowbreak{}Model}
\item \hyperlink{governing-declaration-13fb93228218a38e}{Applied\allowbreak{}Modeling\allowbreak{}Lib.\allowbreak{}measure\allowbreak{}Performative\allowbreak{}Risk}
\end{itemize}
\hypertarget{governing-declaration-836bd13b5596fdaf}{}
\subsection*{\small\ttfamily\raggedright Applied\allowbreak{}Modeling\allowbreak{}Lib.\allowbreak{}measure\allowbreak{}Repeated\allowbreak{}Gradient\allowbreak{}Descent\allowbreak{}Update}
\textbf{Declaration location:} reusable library
\paragraph{Lean declaration body}
\begin{ReviewVerbatim}
type:
{Parameter : Type u_1} →
  {Data : Type u_2} →
    [inst : NormedAddCommGroup Parameter] →
      [InnerProductSpace ℝ Parameter] →
        [inst : MeasurableSpace Data] →
          AppliedModelingLib.MeasurePerformativeModel Parameter Data →
            (Data → Parameter → Parameter) → (Parameter → Parameter) → ℝ → Parameter → Parameter

value:
fun {Parameter : Type u_1} {Data : Type u_2} [NormedAddCommGroup Parameter] [InnerProductSpace ℝ Parameter]
    [MeasurableSpace Data] (model : AppliedModelingLib.MeasurePerformativeModel Parameter Data)
    (gradient : Data → Parameter → Parameter) (project : Parameter → Parameter) (stepSize : ℝ) (a : Parameter) =>
  project (a - stepSize • AppliedModelingLib.measureFrozenPopulationGradient model gradient a a)
\end{ReviewVerbatim}

\paragraph{Direct retained dependencies}
\begin{itemize}\raggedright
\item \hyperlink{library-prerequisite-96553d4fbb82f1ac}{Applied\allowbreak{}Modeling\allowbreak{}Lib.\allowbreak{}Measure\allowbreak{}Performative\allowbreak{}Model}
\item \hyperlink{governing-declaration-d3dfc78bf5c69cfe}{Applied\allowbreak{}Modeling\allowbreak{}Lib.\allowbreak{}measure\allowbreak{}Frozen\allowbreak{}Population\allowbreak{}Gradient}
\end{itemize}
\hypertarget{governing-declaration-46a8b190a3977782}{}
\subsection*{\small\ttfamily\raggedright Applied\allowbreak{}Modeling\allowbreak{}Lib.\allowbreak{}measure\allowbreak{}Sampled\allowbreak{}Gradient\allowbreak{}Descent\allowbreak{}Update}
\textbf{Declaration location:} reusable library
\paragraph{Lean declaration body}
\begin{ReviewVerbatim}
type:
{Parameter : Type u_1} →
  {Data : Type u_2} →
    [inst : NormedAddCommGroup Parameter] →
      [InnerProductSpace ℝ Parameter] →
        [MeasurableSpace Data] →
          (Data → Parameter → Parameter) →
            (Parameter → Parameter) →
              ℝ → {sampleCount : ℕ} → 0 < sampleCount → (Fin sampleCount → Data) → Parameter → Parameter

value:
fun {Parameter : Type u_1} {Data : Type u_2} [NormedAddCommGroup Parameter] [InnerProductSpace ℝ Parameter]
    [MeasurableSpace Data] (gradient : Data → Parameter → Parameter) (project : Parameter → Parameter) (stepSize : ℝ)
    {sampleCount : ℕ} (hsampleCount : 0 < sampleCount) (sample : Fin sampleCount → Data) (a : Parameter) =>
  project
    (a -
      stepSize •
        AppliedModelingLib.measureEmpiricalPopulationGradient gradient
          (AppliedModelingLib.empiricalSampleProbabilityMeasureOfPos hsampleCount sample) a)
\end{ReviewVerbatim}

\paragraph{Direct retained dependencies}
\begin{itemize}\raggedright
\item \hyperlink{governing-declaration-cbd1f5f969456e5f}{Applied\allowbreak{}Modeling\allowbreak{}Lib.\allowbreak{}empirical\allowbreak{}Sample\allowbreak{}Probability\allowbreak{}Measure\allowbreak{}Of\allowbreak{}Pos}
\item \hyperlink{governing-declaration-0bb19c93c4765c7a}{Applied\allowbreak{}Modeling\allowbreak{}Lib.\allowbreak{}measure\allowbreak{}Empirical\allowbreak{}Population\allowbreak{}Gradient}
\end{itemize}
\hypertarget{governing-declaration-30e7d69ebe63768c}{}
\subsection*{\small\ttfamily\raggedright Applied\allowbreak{}Modeling\allowbreak{}Lib.\allowbreak{}p\allowbreak{}One\allowbreak{}Fournier\allowbreak{}Guillin\allowbreak{}Adaptive\allowbreak{}Euclidean\allowbreak{}Selected\allowbreak{}Shell\allowbreak{}Partition\allowbreak{}Certificate\allowbreak{}Bad\allowbreak{}Event}
\textbf{Declaration location:} reusable library
\paragraph{Lean declaration body}
\begin{ReviewVerbatim}
type:
(dimension : ℕ) →
  ℕ →
    ℝ →
      ℝ →
        (ℕ → ℝ) →
          (countSchedule : ℕ → ℕ) →
            ((iteration : ℕ) →
                AppliedModelingLib.HeterogeneousBatchTrace (fun (round : ℕ) => Fin (countSchedule round))
                    (EuclideanSpace ℝ (Fin dimension)) iteration →
                  MeasureTheory.ProbabilityMeasure (EuclideanSpace ℝ (Fin dimension))) →
              (iteration : ℕ) →
                AppliedModelingLib.HeterogeneousBatchTrace (fun (round : ℕ) => Fin (countSchedule round))
                    (EuclideanSpace ℝ (Fin dimension)) iteration →
                  (Fin (countSchedule iteration) → EuclideanSpace ℝ (Fin dimension)) → Prop

value:
fun (dimension cutoff : ℕ) (eta deviation : ℝ) (confidence : ℕ → ℝ) (countSchedule : ℕ → ℕ)
    (law :
      (iteration : ℕ) →
        AppliedModelingLib.HeterogeneousBatchTrace (fun (round : ℕ) => Fin (countSchedule round))
            (EuclideanSpace ℝ (Fin dimension)) iteration →
          MeasureTheory.ProbabilityMeasure (EuclideanSpace ℝ (Fin dimension)))
    (iteration : ℕ)
    (a :
      AppliedModelingLib.HeterogeneousBatchTrace (fun (round : ℕ) => Fin (countSchedule round))
        (EuclideanSpace ℝ (Fin dimension)) iteration)
    (a_1 : Fin (countSchedule iteration) → EuclideanSpace ℝ (Fin dimension)) =>
  AppliedModelingLib.Probability.pOneFournierGuillinSelectedShellPartitionCertificate dimension (law iteration a)
    (countSchedule iteration) eta deviation
    (AppliedModelingLib.Probability.selectedShellGeometricEffectiveHeadDepth dimension cutoff (countSchedule iteration)
      (law iteration a) (confidence iteration))
    (confidence iteration) a_1
\end{ReviewVerbatim}

\paragraph{Direct retained dependencies}
\begin{itemize}\raggedright
\item \hyperlink{governing-declaration-307f711c4bf618c0}{Applied\allowbreak{}Modeling\allowbreak{}Lib.\allowbreak{}Heterogeneous\allowbreak{}Batch\allowbreak{}Trace}
\item \hyperlink{governing-declaration-e282345978931f95}{Applied\allowbreak{}Modeling\allowbreak{}Lib.\allowbreak{}Probability.\allowbreak{}p\allowbreak{}One\allowbreak{}Fournier\allowbreak{}Guillin\allowbreak{}Selected\allowbreak{}Shell\allowbreak{}Partition\allowbreak{}Certificate}
\item \hyperlink{governing-declaration-f9f338884954e2a9}{Applied\allowbreak{}Modeling\allowbreak{}Lib.\allowbreak{}Probability.\allowbreak{}selected\allowbreak{}Shell\allowbreak{}Geometric\allowbreak{}Effective\allowbreak{}Head\allowbreak{}Depth}
\end{itemize}
\hypertarget{governing-declaration-25a90c9551dbd2b1}{}
\subsection*{\small\ttfamily\raggedright Applied\allowbreak{}Modeling\allowbreak{}Lib.\allowbreak{}p\allowbreak{}One\allowbreak{}Fournier\allowbreak{}Guillin\allowbreak{}Theorem310\allowbreak{}Capped\allowbreak{}Deviation}
\textbf{Declaration location:} reusable library
\paragraph{Lean declaration body}
\begin{ReviewVerbatim}
type:
ℕ → ℝ → ℝ → ℝ → ℝ → ℝ → ℝ

value:
fun (dimension : ℕ) (eta alpha gamma momentBound scaledDeviation : ℝ) =>
  AppliedModelingLib.Probability.pOneFournierGuillinUniformCappedDeviation dimension eta alpha gamma momentBound
    scaledDeviation
\end{ReviewVerbatim}

\paragraph{Direct retained dependencies}
\begin{itemize}\raggedright
\item \hyperlink{governing-declaration-d3b37af7ee378c37}{Applied\allowbreak{}Modeling\allowbreak{}Lib.\allowbreak{}Probability.\allowbreak{}p\allowbreak{}One\allowbreak{}Fournier\allowbreak{}Guillin\allowbreak{}Uniform\allowbreak{}Capped\allowbreak{}Deviation}
\end{itemize}
\hypertarget{governing-declaration-8e4b4389a4f6454b}{}
\subsection*{\small\ttfamily\raggedright Applied\allowbreak{}Modeling\allowbreak{}Lib.\allowbreak{}pmf\allowbreak{}Exp}
\textbf{Declaration location:} reusable library
\paragraph{Lean declaration body}
\begin{ReviewVerbatim}
type:
{α : Type u_1} → [Fintype α] → PMF α → (α → ℝ) → ℝ

value:
fun {α : Type u_1} [Fintype α] (μ : PMF α) (f : α → ℝ) => ∑ a : α, (μ a).toReal * f a
\end{ReviewVerbatim}


\hypertarget{governing-declaration-5b91668525b299f9}{}
\subsection*{\small\ttfamily\raggedright Applied\allowbreak{}Modeling\allowbreak{}Lib.\allowbreak{}Finite\allowbreak{}Coupling.\allowbreak{}expected\allowbreak{}Cost}
\textbf{Declaration location:} reusable library
\paragraph{Lean declaration body}
\begin{ReviewVerbatim}
type:
{α : Type u_1} →
  {β : Type u_2} →
    [Fintype α] → [Fintype β] → {μ : PMF α} → {ν : PMF β} → AppliedModelingLib.FiniteCoupling μ ν → (α → β → ℝ) → ℝ

value:
fun {α : Type u_1} {β : Type u_2} [Fintype α] [Fintype β] {μ : PMF α} {ν : PMF β}
    (coupling : AppliedModelingLib.FiniteCoupling μ ν) (cost : α → β → ℝ) =>
  AppliedModelingLib.pmfExp coupling.law fun (pair : α × β) => cost pair.1 pair.2
\end{ReviewVerbatim}

\paragraph{Direct retained dependencies}
\begin{itemize}\raggedright
\item \hyperlink{governing-declaration-72513ba85c27319f}{Applied\allowbreak{}Modeling\allowbreak{}Lib.\allowbreak{}Finite\allowbreak{}Coupling}
\item \hyperlink{governing-declaration-8e4b4389a4f6454b}{Applied\allowbreak{}Modeling\allowbreak{}Lib.\allowbreak{}pmf\allowbreak{}Exp}
\end{itemize}
\hypertarget{governing-declaration-527bb09065d3c073}{}
\subsection*{\small\ttfamily\raggedright Applied\allowbreak{}Modeling\allowbreak{}Lib.\allowbreak{}Finite\allowbreak{}Coupling.\allowbreak{}transport\allowbreak{}Cost\allowbreak{}Set}
\textbf{Declaration location:} reusable library
\paragraph{Lean declaration body}
\begin{ReviewVerbatim}
type:
{α : Type u_1} → {β : Type u_2} → [Fintype α] → [Fintype β] → PMF α → PMF β → (α → β → ℝ) → ℝ → Prop

value:
fun {α : Type u_1} {β : Type u_2} [Fintype α] [Fintype β] (μ : PMF α) (ν : PMF β) (cost : α → β → ℝ) (a : ℝ) =>
  {bound : ℝ | ∃ (coupling : AppliedModelingLib.FiniteCoupling μ ν), coupling.expectedCost cost = bound} a
\end{ReviewVerbatim}

\paragraph{Direct retained dependencies}
\begin{itemize}\raggedright
\item \hyperlink{governing-declaration-72513ba85c27319f}{Applied\allowbreak{}Modeling\allowbreak{}Lib.\allowbreak{}Finite\allowbreak{}Coupling}
\item \hyperlink{governing-declaration-5b91668525b299f9}{Applied\allowbreak{}Modeling\allowbreak{}Lib.\allowbreak{}Finite\allowbreak{}Coupling.\allowbreak{}expected\allowbreak{}Cost}
\end{itemize}
\hypertarget{governing-declaration-93a96c0bd7f262fe}{}
\subsection*{\small\ttfamily\raggedright Applied\allowbreak{}Modeling\allowbreak{}Lib.\allowbreak{}Finite\allowbreak{}Coupling.\allowbreak{}finite\allowbreak{}Wasserstein\allowbreak{}One}
\textbf{Declaration location:} reusable library
\paragraph{Lean declaration body}
\begin{ReviewVerbatim}
type:
{Data : Type u_3} → [Fintype Data] → [DecidableEq Data] → [MetricSpace Data] → PMF Data → PMF Data → ℝ

value:
fun {Data : Type u_3} [Fintype Data] [DecidableEq Data] [MetricSpace Data] (μ ν : PMF Data) =>
  sInf (AppliedModelingLib.FiniteCoupling.transportCostSet μ ν fun (left right : Data) => dist left right)
\end{ReviewVerbatim}

\paragraph{Direct retained dependencies}
\begin{itemize}\raggedright
\item \hyperlink{governing-declaration-527bb09065d3c073}{Applied\allowbreak{}Modeling\allowbreak{}Lib.\allowbreak{}Finite\allowbreak{}Coupling.\allowbreak{}transport\allowbreak{}Cost\allowbreak{}Set}
\end{itemize}
\hypertarget{governing-declaration-889200fe2a9ff2ab}{}
\subsection*{\small\ttfamily\raggedright Applied\allowbreak{}Modeling\allowbreak{}Lib.\allowbreak{}Is\allowbreak{}Finite\allowbreak{}Wasserstein\allowbreak{}Sensitive}
\textbf{Declaration location:} reusable library
\paragraph{Lean declaration body}
\begin{ReviewVerbatim}
type:
{Parameter : Type u_1} →
  {Data : Type u_2} →
    [inst : Fintype Data] →
      [inst_1 : DecidableEq Data] →
        [MetricSpace Parameter] →
          [MetricSpace Data] → AppliedModelingLib.FinitePerformativeModel Parameter Data → ℝ → Prop

value:
fun {Parameter : Type u_1} {Data : Type u_2} [Fintype Data] [DecidableEq Data] [MetricSpace Parameter]
    [MetricSpace Data] (model : AppliedModelingLib.FinitePerformativeModel Parameter Data) (sensitivity : ℝ) =>
  ∀ (first second : Parameter),
    AppliedModelingLib.FiniteCoupling.finiteWassersteinOne (model.dataLaw first) (model.dataLaw second) ≤
      sensitivity * dist first second
\end{ReviewVerbatim}

\paragraph{Direct retained dependencies}
\begin{itemize}\raggedright
\item \hyperlink{governing-declaration-93a96c0bd7f262fe}{Applied\allowbreak{}Modeling\allowbreak{}Lib.\allowbreak{}Finite\allowbreak{}Coupling.\allowbreak{}finite\allowbreak{}Wasserstein\allowbreak{}One}
\item \hyperlink{governing-declaration-24472be40793a412}{Applied\allowbreak{}Modeling\allowbreak{}Lib.\allowbreak{}Finite\allowbreak{}Performative\allowbreak{}Model}
\end{itemize}
\hypertarget{governing-declaration-4a2c8294f0669302}{}
\subsection*{\small\ttfamily\raggedright Applied\allowbreak{}Modeling\allowbreak{}Lib.\allowbreak{}performative\allowbreak{}Risk}
\textbf{Declaration location:} reusable library
\paragraph{Lean declaration body}
\begin{ReviewVerbatim}
type:
{Parameter : Type u_1} →
  {Data : Type u_2} →
    [inst : Fintype Data] →
      [inst_1 : DecidableEq Data] → AppliedModelingLib.FinitePerformativeModel Parameter Data → Parameter → ℝ

value:
fun {Parameter : Type u_1} {Data : Type u_2} [Fintype Data] [DecidableEq Data]
    (model : AppliedModelingLib.FinitePerformativeModel Parameter Data) (parameter : Parameter) =>
  AppliedModelingLib.pmfExp (model.dataLaw parameter) fun (datum : Data) => model.loss datum parameter
\end{ReviewVerbatim}

\paragraph{Direct retained dependencies}
\begin{itemize}\raggedright
\item \hyperlink{governing-declaration-24472be40793a412}{Applied\allowbreak{}Modeling\allowbreak{}Lib.\allowbreak{}Finite\allowbreak{}Performative\allowbreak{}Model}
\item \hyperlink{governing-declaration-8e4b4389a4f6454b}{Applied\allowbreak{}Modeling\allowbreak{}Lib.\allowbreak{}pmf\allowbreak{}Exp}
\end{itemize}
\hypertarget{governing-declaration-965db1948b41d1e4}{}
\subsection*{\small\ttfamily\raggedright Applied\allowbreak{}Modeling\allowbreak{}Lib.\allowbreak{}proposition36a\allowbreak{}Domain}
\textbf{Declaration location:} reusable library
\paragraph{Lean declaration body}
\begin{ReviewVerbatim}
type:
ℝ → Prop

value:
fun (a : ℝ) => Set.Icc (-1) 1 a
\end{ReviewVerbatim}


\hypertarget{governing-declaration-6b5ac0170880188f}{}
\subsection*{\small\ttfamily\raggedright Applied\allowbreak{}Modeling\allowbreak{}Lib.\allowbreak{}proposition36a\allowbreak{}Measure\allowbreak{}Gradient}
\textbf{Declaration location:} reusable library
\paragraph{Lean declaration body}
\begin{ReviewVerbatim}
type:
ℝ → ℝ → ℝ → ℝ

value:
fun (beta datum _candidate : ℝ) => beta * datum
\end{ReviewVerbatim}


\hypertarget{governing-declaration-199ecccc69085253}{}
\subsection*{\small\ttfamily\raggedright Applied\allowbreak{}Modeling\allowbreak{}Lib.\allowbreak{}proposition36a\allowbreak{}Measure\allowbreak{}Law}
\textbf{Declaration location:} reusable library
\paragraph{Lean declaration body}
\begin{ReviewVerbatim}
type:
ℝ → ℝ → MeasureTheory.ProbabilityMeasure ℝ

value:
fun (epsilon theta : ℝ) => MeasureTheory.diracProba (epsilon * theta)
\end{ReviewVerbatim}


\hypertarget{governing-declaration-10c605398ed70a9b}{}
\subsection*{\small\ttfamily\raggedright Applied\allowbreak{}Modeling\allowbreak{}Lib.\allowbreak{}proposition36a\allowbreak{}Measure\allowbreak{}Loss}
\textbf{Declaration location:} reusable library
\paragraph{Lean declaration body}
\begin{ReviewVerbatim}
type:
ℝ → ℝ → ℝ → ℝ

value:
fun (beta datum candidate : ℝ) => beta * datum * candidate
\end{ReviewVerbatim}


\hypertarget{governing-declaration-3d85baf63239a428}{}
\subsection*{\small\ttfamily\raggedright Applied\allowbreak{}Modeling\allowbreak{}Lib.\allowbreak{}proposition36a\allowbreak{}Measure\allowbreak{}Model}
\textbf{Declaration location:} reusable library
\paragraph{Lean declaration body}
\begin{ReviewVerbatim}
type:
ℝ → ℝ → AppliedModelingLib.MeasurePerformativeModel ℝ ℝ

value:
fun (epsilon beta : ℝ) =>
  { dataLaw := AppliedModelingLib.proposition36aMeasureLaw epsilon,
    loss := AppliedModelingLib.proposition36aMeasureLoss beta,
    loss_integrable := AppliedModelingLib.proposition36aMeasureModel._proof_1 epsilon beta }
\end{ReviewVerbatim}

\paragraph{Direct retained dependencies}
\begin{itemize}\raggedright
\item \hyperlink{library-prerequisite-96553d4fbb82f1ac}{Applied\allowbreak{}Modeling\allowbreak{}Lib.\allowbreak{}Measure\allowbreak{}Performative\allowbreak{}Model}
\item \hyperlink{governing-declaration-199ecccc69085253}{Applied\allowbreak{}Modeling\allowbreak{}Lib.\allowbreak{}proposition36a\allowbreak{}Measure\allowbreak{}Law}
\item \hyperlink{governing-declaration-10c605398ed70a9b}{Applied\allowbreak{}Modeling\allowbreak{}Lib.\allowbreak{}proposition36a\allowbreak{}Measure\allowbreak{}Loss}
\end{itemize}
\hypertarget{governing-declaration-03f87e42b92f1bad}{}
\subsection*{\small\ttfamily\raggedright Applied\allowbreak{}Modeling\allowbreak{}Lib.\allowbreak{}proposition36a\allowbreak{}RRM\allowbreak{}Update}
\textbf{Declaration location:} reusable library
\paragraph{Lean declaration body}
\begin{ReviewVerbatim}
type:
ℝ → ℝ

value:
fun (theta : ℝ) => if theta < 0 then 1 else -1
\end{ReviewVerbatim}


\hypertarget{governing-declaration-ac13973b79856aff}{}
\subsection*{\small\ttfamily\raggedright Applied\allowbreak{}Modeling\allowbreak{}Lib.\allowbreak{}proposition36b\allowbreak{}Lower\allowbreak{}Endpoint}
\textbf{Declaration location:} reusable library
\paragraph{Lean declaration body}
\begin{ReviewVerbatim}
type:
ℝ → ℝ

value:
fun (epsilon : ℝ) => -(1 / (2 * epsilon))
\end{ReviewVerbatim}


\hypertarget{governing-declaration-ee9adf0c7d04f002}{}
\subsection*{\small\ttfamily\raggedright Applied\allowbreak{}Modeling\allowbreak{}Lib.\allowbreak{}Is\allowbreak{}Proposition36b\allowbreak{}Large\allowbreak{}Penalty}
\textbf{Declaration location:} reusable library
\paragraph{Lean declaration body}
\begin{ReviewVerbatim}
type:
ℝ → ℝ → ℝ → Prop

value:
fun (epsilon gamma penalty : ℝ) =>
  gamma * (1 - AppliedModelingLib.proposition36bLowerEndpoint epsilon) ≤ 2 * penalty * epsilon ∧ 2 * gamma ≤ penalty
\end{ReviewVerbatim}

\paragraph{Direct retained dependencies}
\begin{itemize}\raggedright
\item \hyperlink{governing-declaration-ac13973b79856aff}{Applied\allowbreak{}Modeling\allowbreak{}Lib.\allowbreak{}proposition36b\allowbreak{}Lower\allowbreak{}Endpoint}
\end{itemize}
\hypertarget{governing-declaration-ebb8a5e04e4e48ce}{}
\subsection*{\small\ttfamily\raggedright Applied\allowbreak{}Modeling\allowbreak{}Lib.\allowbreak{}proposition36b\allowbreak{}Domain}
\textbf{Declaration location:} reusable library
\paragraph{Lean declaration body}
\begin{ReviewVerbatim}
type:
ℝ → ℝ → Prop

value:
fun (epsilon a : ℝ) => Set.Icc (AppliedModelingLib.proposition36bLowerEndpoint epsilon) 2 a
\end{ReviewVerbatim}

\paragraph{Direct retained dependencies}
\begin{itemize}\raggedright
\item \hyperlink{governing-declaration-ac13973b79856aff}{Applied\allowbreak{}Modeling\allowbreak{}Lib.\allowbreak{}proposition36b\allowbreak{}Lower\allowbreak{}Endpoint}
\end{itemize}
\hypertarget{governing-declaration-ceacbb45eebd01a4}{}
\subsection*{\small\ttfamily\raggedright Applied\allowbreak{}Modeling\allowbreak{}Lib.\allowbreak{}proposition36b\allowbreak{}Measure\allowbreak{}Law}
\textbf{Declaration location:} reusable library
\paragraph{Lean declaration body}
\begin{ReviewVerbatim}
type:
ℝ → ℝ → MeasureTheory.ProbabilityMeasure ℝ

value:
fun (epsilon theta : ℝ) => MeasureTheory.diracProba (epsilon * theta)
\end{ReviewVerbatim}


\hypertarget{governing-declaration-9fe26c2f815241d3}{}
\subsection*{\small\ttfamily\raggedright Applied\allowbreak{}Modeling\allowbreak{}Lib.\allowbreak{}proposition36b\allowbreak{}Measure\allowbreak{}Loss}
\textbf{Declaration location:} reusable library
\paragraph{Lean declaration body}
\begin{ReviewVerbatim}
type:
ℝ → ℝ → ℝ → ℝ → ℝ

value:
fun (gamma penalty datum candidate : ℝ) => penalty * max (-1) (datum * candidate) + gamma / 2 * (candidate - 1) ^ 2
\end{ReviewVerbatim}


\hypertarget{governing-declaration-14367e2ad0182576}{}
\subsection*{\small\ttfamily\raggedright Applied\allowbreak{}Modeling\allowbreak{}Lib.\allowbreak{}proposition36b\allowbreak{}Measure\allowbreak{}Model}
\textbf{Declaration location:} reusable library
\paragraph{Lean declaration body}
\begin{ReviewVerbatim}
type:
ℝ → ℝ → ℝ → AppliedModelingLib.MeasurePerformativeModel ℝ ℝ

value:
fun (epsilon gamma penalty : ℝ) =>
  { dataLaw := AppliedModelingLib.proposition36bMeasureLaw epsilon,
    loss := AppliedModelingLib.proposition36bMeasureLoss gamma penalty,
    loss_integrable := AppliedModelingLib.proposition36bMeasureModel._proof_1 epsilon gamma penalty }
\end{ReviewVerbatim}

\paragraph{Direct retained dependencies}
\begin{itemize}\raggedright
\item \hyperlink{library-prerequisite-96553d4fbb82f1ac}{Applied\allowbreak{}Modeling\allowbreak{}Lib.\allowbreak{}Measure\allowbreak{}Performative\allowbreak{}Model}
\item \hyperlink{governing-declaration-ceacbb45eebd01a4}{Applied\allowbreak{}Modeling\allowbreak{}Lib.\allowbreak{}proposition36b\allowbreak{}Measure\allowbreak{}Law}
\item \hyperlink{governing-declaration-9fe26c2f815241d3}{Applied\allowbreak{}Modeling\allowbreak{}Lib.\allowbreak{}proposition36b\allowbreak{}Measure\allowbreak{}Loss}
\end{itemize}
\hypertarget{governing-declaration-10d57ee3741bad42}{}
\subsection*{\small\ttfamily\raggedright Applied\allowbreak{}Modeling\allowbreak{}Lib.\allowbreak{}proposition36b\allowbreak{}RRM\allowbreak{}Update}
\textbf{Declaration location:} reusable library
\paragraph{Lean declaration body}
\begin{ReviewVerbatim}
type:
ℝ → ℝ → ℝ

value:
fun (epsilon deployed : ℝ) => if deployed = 2 then AppliedModelingLib.proposition36bLowerEndpoint epsilon else 2
\end{ReviewVerbatim}

\paragraph{Direct retained dependencies}
\begin{itemize}\raggedright
\item \hyperlink{governing-declaration-ac13973b79856aff}{Applied\allowbreak{}Modeling\allowbreak{}Lib.\allowbreak{}proposition36b\allowbreak{}Lower\allowbreak{}Endpoint}
\end{itemize}
\hypertarget{governing-declaration-d690bb26f4cad28e}{}
\subsection*{\small\ttfamily\raggedright Applied\allowbreak{}Modeling\allowbreak{}Lib.\allowbreak{}proposition36c\allowbreak{}Measure\allowbreak{}Law}
\textbf{Declaration location:} reusable library
\paragraph{Lean declaration body}
\begin{ReviewVerbatim}
type:
ℝ → ℝ → MeasureTheory.ProbabilityMeasure ℝ

value:
fun (epsilon theta : ℝ) => MeasureTheory.diracProba (1 + epsilon * theta)
\end{ReviewVerbatim}


\hypertarget{governing-declaration-954407b64b7a9ac2}{}
\subsection*{\small\ttfamily\raggedright Applied\allowbreak{}Modeling\allowbreak{}Lib.\allowbreak{}proposition36c\allowbreak{}Regularity\allowbreak{}Gradient}
\textbf{Declaration location:} reusable library
\paragraph{Lean declaration body}
\begin{ReviewVerbatim}
type:
ℝ → ℝ → ℝ → ℝ → ℝ

value:
fun (beta gamma datum candidate : ℝ) => gamma * candidate - beta * datum
\end{ReviewVerbatim}


\hypertarget{governing-declaration-9a6823e6fdefa446}{}
\subsection*{\small\ttfamily\raggedright Applied\allowbreak{}Modeling\allowbreak{}Lib.\allowbreak{}proposition36c\allowbreak{}Regularity\allowbreak{}Loss}
\textbf{Declaration location:} reusable library
\paragraph{Lean declaration body}
\begin{ReviewVerbatim}
type:
ℝ → ℝ → ℝ → ℝ → ℝ

value:
fun (beta gamma datum candidate : ℝ) => gamma / 2 * candidate ^ 2 - beta * datum * candidate
\end{ReviewVerbatim}


\hypertarget{governing-declaration-9b198d5b04ab29c5}{}
\subsection*{\small\ttfamily\raggedright Applied\allowbreak{}Modeling\allowbreak{}Lib.\allowbreak{}proposition36c\allowbreak{}Regularity\allowbreak{}Measure\allowbreak{}Model}
\textbf{Declaration location:} reusable library
\paragraph{Lean declaration body}
\begin{ReviewVerbatim}
type:
ℝ → ℝ → ℝ → AppliedModelingLib.MeasurePerformativeModel ℝ ℝ

value:
fun (epsilon beta gamma : ℝ) =>
  { dataLaw := AppliedModelingLib.proposition36cMeasureLaw epsilon,
    loss := AppliedModelingLib.proposition36cRegularityLoss beta gamma,
    loss_integrable := AppliedModelingLib.proposition36cRegularityMeasureModel._proof_1 epsilon beta gamma }
\end{ReviewVerbatim}

\paragraph{Direct retained dependencies}
\begin{itemize}\raggedright
\item \hyperlink{library-prerequisite-96553d4fbb82f1ac}{Applied\allowbreak{}Modeling\allowbreak{}Lib.\allowbreak{}Measure\allowbreak{}Performative\allowbreak{}Model}
\item \hyperlink{governing-declaration-d690bb26f4cad28e}{Applied\allowbreak{}Modeling\allowbreak{}Lib.\allowbreak{}proposition36c\allowbreak{}Measure\allowbreak{}Law}
\item \hyperlink{governing-declaration-9a6823e6fdefa446}{Applied\allowbreak{}Modeling\allowbreak{}Lib.\allowbreak{}proposition36c\allowbreak{}Regularity\allowbreak{}Loss}
\end{itemize}
\hypertarget{governing-declaration-39c6c11ddc030996}{}
\subsection*{\small\ttfamily\raggedright Applied\allowbreak{}Modeling\allowbreak{}Lib.\allowbreak{}proposition36c\allowbreak{}Regularity\allowbreak{}RRM\allowbreak{}Update}
\textbf{Declaration location:} reusable library
\paragraph{Lean declaration body}
\begin{ReviewVerbatim}
type:
ℝ → ℝ → ℝ → ℝ → ℝ

value:
fun (epsilon beta gamma theta : ℝ) => beta / gamma * (1 + epsilon * theta)
\end{ReviewVerbatim}


\hypertarget{governing-declaration-4e38b2c8890e2130}{}
\subsection*{\small\ttfamily\raggedright Applied\allowbreak{}Modeling\allowbreak{}Lib.\allowbreak{}proposition42\allowbreak{}Bernoulli\allowbreak{}Parameter}
\textbf{Declaration location:} reusable library
\paragraph{Lean declaration body}
\begin{ReviewVerbatim}
type:
ℝ → Bool → NNReal

value:
fun (bias : ℝ) (feature : Bool) => (if feature = true then 1 / 2 + bias else 1 / 2 - bias).toNNReal
\end{ReviewVerbatim}


\hypertarget{governing-declaration-d4d35ede9cba1f94}{}
\subsection*{\small\ttfamily\raggedright Applied\allowbreak{}Modeling\allowbreak{}Lib.\allowbreak{}proposition42\allowbreak{}Bool\allowbreak{}Metric\allowbreak{}Space}
\textbf{Declaration location:} reusable library
\paragraph{Lean declaration body}
\begin{ReviewVerbatim}
type:
MetricSpace Bool

value:
MetricSpace.induced (fun (value : Bool) => if value = true then 1 else 0)
  AppliedModelingLib.proposition42BoolMetricSpace._proof_1 inferInstance
\end{ReviewVerbatim}


\hypertarget{governing-declaration-8c485a232d3e9b18}{}
\subsection*{\small\ttfamily\raggedright Applied\allowbreak{}Modeling\allowbreak{}Lib.\allowbreak{}proposition42\allowbreak{}Gradient}
\textbf{Declaration location:} reusable library
\paragraph{Lean declaration body}
\begin{ReviewVerbatim}
type:
Bool × Bool → ℝ → ℝ

value:
fun (datum : Bool × Bool) (theta : ℝ) =>
  (2 * if datum.1 = true then 1 else -1) *
    ((theta * if datum.1 = true then 1 else -1) + 1 / 2 - if datum.2 = true then 1 else 0)
\end{ReviewVerbatim}


\hypertarget{governing-declaration-ea650343c67da273}{}
\subsection*{\small\ttfamily\raggedright Applied\allowbreak{}Modeling\allowbreak{}Lib.\allowbreak{}proposition42\allowbreak{}Law}
\textbf{Declaration location:} reusable library
\paragraph{Lean declaration body}
\begin{ReviewVerbatim}
type:
(bias : ℝ) → |bias| ≤ 1 / 2 → PMF (Bool × Bool)

value:
fun (bias : ℝ) (hbias : |bias| ≤ 1 / 2) =>
  (PMF.uniformOfFintype Bool).bind fun (feature : Bool) =>
    PMF.map (fun (label : Bool) => (feature, label))
      (PMF.bernoulli (AppliedModelingLib.proposition42BernoulliParameter bias feature)
        (AppliedModelingLib.proposition42BernoulliParameter_le_one bias feature hbias))
\end{ReviewVerbatim}

\paragraph{Direct retained dependencies}
\begin{itemize}\raggedright
\item \hyperlink{governing-declaration-4e38b2c8890e2130}{Applied\allowbreak{}Modeling\allowbreak{}Lib.\allowbreak{}proposition42\allowbreak{}Bernoulli\allowbreak{}Parameter}
\end{itemize}
\hypertarget{governing-declaration-ded6eb207c916fe6}{}
\subsection*{\small\ttfamily\raggedright Applied\allowbreak{}Modeling\allowbreak{}Lib.\allowbreak{}proposition42\allowbreak{}Loss}
\textbf{Declaration location:} reusable library
\paragraph{Lean declaration body}
\begin{ReviewVerbatim}
type:
ℝ → Bool × Bool → ℝ

value:
fun (theta : ℝ) (datum : Bool × Bool) =>
  ((if datum.2 = true then 1 else 0) - ((theta * if datum.1 = true then 1 else -1) + 1 / 2)) ^ 2
\end{ReviewVerbatim}


\hypertarget{governing-declaration-a540759b18439129}{}
\subsection*{\small\ttfamily\raggedright Applied\allowbreak{}Modeling\allowbreak{}Lib.\allowbreak{}corrected\allowbreak{}Proposition42\allowbreak{}Model}
\textbf{Declaration location:} reusable library
\paragraph{Lean declaration body}
\begin{ReviewVerbatim}
type:
(mu epsilon : ℝ) →
  |mu| ≤ 1 / 2 →
    |mu + epsilon| ≤ 1 / 2 →
      AppliedModelingLib.FinitePerformativeModel { theta : ℝ // theta ∈ Set.Icc 0 1 } (Bool × Bool)

value:
fun (mu epsilon : ℝ) (hmu : |mu| ≤ 1 / 2) (hmuEpsilon : |mu + epsilon| ≤ 1 / 2) =>
  {
    dataLaw := fun (theta : { theta : ℝ // theta ∈ Set.Icc 0 1 }) =>
      AppliedModelingLib.proposition42Law (mu + epsilon * ↑theta)
        (AppliedModelingLib.correctedProposition42Model._proof_1 mu epsilon hmu hmuEpsilon theta),
    loss := fun (datum : Bool × Bool) (theta : { theta : ℝ // theta ∈ Set.Icc 0 1 }) =>
      AppliedModelingLib.proposition42Loss (↑theta) datum }
\end{ReviewVerbatim}

\paragraph{Direct retained dependencies}
\begin{itemize}\raggedright
\item \hyperlink{governing-declaration-24472be40793a412}{Applied\allowbreak{}Modeling\allowbreak{}Lib.\allowbreak{}Finite\allowbreak{}Performative\allowbreak{}Model}
\item \hyperlink{governing-declaration-ea650343c67da273}{Applied\allowbreak{}Modeling\allowbreak{}Lib.\allowbreak{}proposition42\allowbreak{}Law}
\item \hyperlink{governing-declaration-ded6eb207c916fe6}{Applied\allowbreak{}Modeling\allowbreak{}Lib.\allowbreak{}proposition42\allowbreak{}Loss}
\end{itemize}
\hypertarget{governing-declaration-6a6b18e647182cee}{}
\subsection*{\small\ttfamily\raggedright Applied\allowbreak{}Modeling\allowbreak{}Lib.\allowbreak{}theorem310\allowbreak{}Failure\allowbreak{}Budget}
\textbf{Declaration location:} reusable library
\paragraph{Lean declaration body}
\begin{ReviewVerbatim}
type:
ℝ → ℕ → ℝ

value:
fun (p : ℝ) (iteration : ℕ) => 6 * p / (Real.pi ^ 2 * ↑(iteration + 1) ^ 2)
\end{ReviewVerbatim}


\hypertarget{governing-declaration-3cf6b14c7bcca80c}{}
\subsection*{\small\ttfamily\raggedright Applied\allowbreak{}Modeling\allowbreak{}Lib.\allowbreak{}p\allowbreak{}One\allowbreak{}Fournier\allowbreak{}Guillin\allowbreak{}Theorem310\allowbreak{}Supercritical\allowbreak{}Capped\allowbreak{}Concrete\allowbreak{}Count\allowbreak{}Schedule}
\textbf{Declaration location:} reusable library
\paragraph{Lean declaration body}
\begin{ReviewVerbatim}
type:
ℕ → ℕ → ℝ → ℝ → ℝ → ℝ → ℝ → ℝ → ℝ → ℝ → ℕ → ℕ

value:
fun (dimension cutoff : ℕ) (eta alpha gamma momentBound tailBound scaledDeviation p headTolerance : ℝ) (a : ℕ) =>
  AppliedModelingLib.Probability.pOneFournierGuillinSupercriticalCappedAllShellConcreteCountSchedule dimension cutoff
    eta alpha gamma momentBound tailBound scaledDeviation (AppliedModelingLib.theorem310FailureBudget p)
    (fun (x : ℕ) => headTolerance) a
\end{ReviewVerbatim}

\paragraph{Direct retained dependencies}
\begin{itemize}\raggedright
\item \hyperlink{governing-declaration-3d6d95e1d691574b}{Applied\allowbreak{}Modeling\allowbreak{}Lib.\allowbreak{}Probability.\allowbreak{}p\allowbreak{}One\allowbreak{}Fournier\allowbreak{}Guillin\allowbreak{}Supercritical\allowbreak{}Capped\allowbreak{}All\allowbreak{}Shell\allowbreak{}Concrete\allowbreak{}Count\allowbreak{}Schedule}
\item \hyperlink{governing-declaration-6a6b18e647182cee}{Applied\allowbreak{}Modeling\allowbreak{}Lib.\allowbreak{}theorem310\allowbreak{}Failure\allowbreak{}Budget}
\end{itemize}
\hypertarget{governing-declaration-6a2f3c5b322d78da}{}
\subsection*{\small\ttfamily\raggedright Applied\allowbreak{}Modeling\allowbreak{}Lib.\allowbreak{}p\allowbreak{}One\allowbreak{}Fournier\allowbreak{}Guillin\allowbreak{}Theorem310\allowbreak{}Supercritical\allowbreak{}Confidence\allowbreak{}Schedule}
\textbf{Declaration location:} reusable library
\paragraph{Lean declaration body}
\begin{ReviewVerbatim}
type:
ℝ → ℕ → ℝ

value:
fun (p : ℝ) (a : ℕ) =>
  AppliedModelingLib.Math.expConfidenceForHalfBudget (AppliedModelingLib.theorem310FailureBudget p a)
\end{ReviewVerbatim}

\paragraph{Direct retained dependencies}
\begin{itemize}\raggedright
\item \hyperlink{governing-declaration-77c24384a8319860}{Applied\allowbreak{}Modeling\allowbreak{}Lib.\allowbreak{}Math.\allowbreak{}exp\allowbreak{}Confidence\allowbreak{}For\allowbreak{}Half\allowbreak{}Budget}
\item \hyperlink{governing-declaration-6a6b18e647182cee}{Applied\allowbreak{}Modeling\allowbreak{}Lib.\allowbreak{}theorem310\allowbreak{}Failure\allowbreak{}Budget}
\end{itemize}
\hypertarget{governing-declaration-12799a7660a5d8a2}{}
\subsection*{\small\ttfamily\raggedright PZMH20\allowbreak{}Performative\allowbreak{}Prediction.\allowbreak{}Is\allowbreak{}Finite\allowbreak{}Wasserstein\allowbreak{}Sensitive}
\textbf{Declaration location:} paper-local
\paragraph{Lean declaration body}
\begin{ReviewVerbatim}
type:
{Parameter : Type u_1} →
  {Data : Type u_2} →
    [inst : Fintype Data] →
      [inst_1 : DecidableEq Data] →
        [MetricSpace Parameter] →
          [MetricSpace Data] → AppliedModelingLib.FinitePerformativeModel Parameter Data → ℝ → Prop

value:
fun {Parameter : Type u_1} {Data : Type u_2} [Fintype Data] [DecidableEq Data] [MetricSpace Parameter]
    [MetricSpace Data] (model : AppliedModelingLib.FinitePerformativeModel Parameter Data) (sensitivity : ℝ) =>
  AppliedModelingLib.IsFiniteWassersteinSensitive model sensitivity
\end{ReviewVerbatim}

\paragraph{Direct retained dependencies}
\begin{itemize}\raggedright
\item \hyperlink{governing-declaration-24472be40793a412}{Applied\allowbreak{}Modeling\allowbreak{}Lib.\allowbreak{}Finite\allowbreak{}Performative\allowbreak{}Model}
\item \hyperlink{governing-declaration-889200fe2a9ff2ab}{Applied\allowbreak{}Modeling\allowbreak{}Lib.\allowbreak{}Is\allowbreak{}Finite\allowbreak{}Wasserstein\allowbreak{}Sensitive}
\end{itemize}
\hypertarget{governing-declaration-11c9ab6e3c6f528b}{}
\subsection*{\small\ttfamily\raggedright PZMH20\allowbreak{}Performative\allowbreak{}Prediction.\allowbreak{}Is\allowbreak{}Measure\allowbreak{}Pointwise\allowbreak{}Gradient\allowbreak{}Strongly\allowbreak{}Convex}
\textbf{Declaration location:} paper-local
\paragraph{Lean declaration body}
\begin{ReviewVerbatim}
type:
{Parameter : Type u_1} →
  {Data : Type u_2} →
    [inst : MeasurableSpace Data] →
      [inst_1 : NormedAddCommGroup Parameter] →
        [InnerProductSpace ℝ Parameter] →
          AppliedModelingLib.MeasurePerformativeModel Parameter Data → (Data → Parameter → Parameter) → ℝ → Prop

value:
fun {Parameter : Type u_1} {Data : Type u_2} [MeasurableSpace Data] [NormedAddCommGroup Parameter]
    [InnerProductSpace ℝ Parameter] (model : AppliedModelingLib.MeasurePerformativeModel Parameter Data)
    (gradient : Data → Parameter → Parameter) (modulus : ℝ) =>
  AppliedModelingLib.IsMeasurePointwiseGradientStronglyConvex model gradient modulus
\end{ReviewVerbatim}

\paragraph{Direct retained dependencies}
\begin{itemize}\raggedright
\item \hyperlink{library-prerequisite-b41cbb32daa3dce7}{Applied\allowbreak{}Modeling\allowbreak{}Lib.\allowbreak{}Is\allowbreak{}Measure\allowbreak{}Pointwise\allowbreak{}Gradient\allowbreak{}Strongly\allowbreak{}Convex}
\item \hyperlink{library-prerequisite-96553d4fbb82f1ac}{Applied\allowbreak{}Modeling\allowbreak{}Lib.\allowbreak{}Measure\allowbreak{}Performative\allowbreak{}Model}
\end{itemize}
\hypertarget{governing-declaration-32f609ab3d5aaf05}{}
\subsection*{\small\ttfamily\raggedright PZMH20\allowbreak{}Performative\allowbreak{}Prediction.\allowbreak{}Is\allowbreak{}Measure\allowbreak{}RRM}
\textbf{Declaration location:} paper-local
\paragraph{Lean declaration body}
\begin{ReviewVerbatim}
type:
{Parameter : Type u_1} →
  {Data : Type u_2} →
    [inst : MeasurableSpace Data] →
      AppliedModelingLib.MeasurePerformativeModel Parameter Data → (Parameter → Parameter) → Prop

value:
fun {Parameter : Type u_1} {Data : Type u_2} [MeasurableSpace Data]
    (model : AppliedModelingLib.MeasurePerformativeModel Parameter Data) (update : Parameter → Parameter) =>
  AppliedModelingLib.IsMeasureRepeatedRiskMinimization model update
\end{ReviewVerbatim}

\paragraph{Direct retained dependencies}
\begin{itemize}\raggedright
\item \hyperlink{governing-declaration-916f213e8c73a94a}{Applied\allowbreak{}Modeling\allowbreak{}Lib.\allowbreak{}Is\allowbreak{}Measure\allowbreak{}Repeated\allowbreak{}Risk\allowbreak{}Minimization}
\item \hyperlink{library-prerequisite-96553d4fbb82f1ac}{Applied\allowbreak{}Modeling\allowbreak{}Lib.\allowbreak{}Measure\allowbreak{}Performative\allowbreak{}Model}
\end{itemize}
\hypertarget{governing-declaration-d4ecd5e34f1d7939}{}
\subsection*{\small\ttfamily\raggedright PZMH20\allowbreak{}Performative\allowbreak{}Prediction.\allowbreak{}Is\allowbreak{}Measure\allowbreak{}Stable\allowbreak{}On}
\textbf{Declaration location:} paper-local
\paragraph{Lean declaration body}
\begin{ReviewVerbatim}
type:
{Parameter : Type u_1} →
  {Data : Type u_2} →
    [inst : MeasurableSpace Data] →
      AppliedModelingLib.MeasurePerformativeModel Parameter Data → Set Parameter → Parameter → Prop

value:
fun {Parameter : Type u_1} {Data : Type u_2} [MeasurableSpace Data]
    (model : AppliedModelingLib.MeasurePerformativeModel Parameter Data) (domain : Set Parameter)
    (parameter : Parameter) =>
  AppliedModelingLib.IsMeasurePerformativelyStableOn model domain parameter
\end{ReviewVerbatim}

\paragraph{Direct retained dependencies}
\begin{itemize}\raggedright
\item \hyperlink{governing-declaration-522b06e211b1e40a}{Applied\allowbreak{}Modeling\allowbreak{}Lib.\allowbreak{}Is\allowbreak{}Measure\allowbreak{}Performatively\allowbreak{}Stable\allowbreak{}On}
\item \hyperlink{library-prerequisite-96553d4fbb82f1ac}{Applied\allowbreak{}Modeling\allowbreak{}Lib.\allowbreak{}Measure\allowbreak{}Performative\allowbreak{}Model}
\end{itemize}
\hypertarget{governing-declaration-e0c24b576a9d5f99}{}
\subsection*{\small\ttfamily\raggedright PZMH20\allowbreak{}Performative\allowbreak{}Prediction.\allowbreak{}proposition36b\allowbreak{}Exact\allowbreak{}Dirac\allowbreak{}Endpoint\allowbreak{}Counterexample\allowbreak{}Spec}
\textbf{Declaration location:} paper-local
\paragraph{Lean declaration body}
\begin{ReviewVerbatim}
type:
ℝ → ℝ → ℝ → Prop

value:
fun (epsilon gamma penalty : ℝ) =>
  0 < epsilon →
    0 < gamma →
      AppliedModelingLib.IsProposition36bLargePenalty epsilon gamma penalty →
        PZMH20PerformativePrediction.IsMeasureWassersteinSensitive
            (AppliedModelingLib.proposition36bMeasureModel epsilon gamma penalty) epsilon ∧
          (∀ (datum : ℝ),
              StrongConvexOn Set.univ gamma fun (candidate : ℝ) =>
                AppliedModelingLib.proposition36bMeasureLoss gamma penalty datum candidate) ∧
            ¬DifferentiableAt ℝ
                  (fun (candidate : ℝ) =>
                    AppliedModelingLib.proposition36bMeasureLoss gamma penalty (epsilon * 2) candidate)
                  (AppliedModelingLib.proposition36bLowerEndpoint epsilon) ∧
              ((AppliedModelingLib.proposition36bRRMUpdate epsilon 2 ∈ AppliedModelingLib.proposition36bDomain epsilon ∧
                    ∀ candidate ∈ AppliedModelingLib.proposition36bDomain epsilon,
                      AppliedModelingLib.measureDecoupledPerformativeRisk
                          (AppliedModelingLib.proposition36bMeasureModel epsilon gamma penalty) 2
                          (AppliedModelingLib.proposition36bRRMUpdate epsilon 2) ≤
                        AppliedModelingLib.measureDecoupledPerformativeRisk
                          (AppliedModelingLib.proposition36bMeasureModel epsilon gamma penalty) 2 candidate) ∧
                  AppliedModelingLib.proposition36bRRMUpdate epsilon
                        (AppliedModelingLib.proposition36bLowerEndpoint epsilon) ∈
                      AppliedModelingLib.proposition36bDomain epsilon ∧
                    ∀ candidate ∈ AppliedModelingLib.proposition36bDomain epsilon,
                      AppliedModelingLib.measureDecoupledPerformativeRisk
                          (AppliedModelingLib.proposition36bMeasureModel epsilon gamma penalty)
                          (AppliedModelingLib.proposition36bLowerEndpoint epsilon)
                          (AppliedModelingLib.proposition36bRRMUpdate epsilon
                            (AppliedModelingLib.proposition36bLowerEndpoint epsilon)) ≤
                        AppliedModelingLib.measureDecoupledPerformativeRisk
                          (AppliedModelingLib.proposition36bMeasureModel epsilon gamma penalty)
                          (AppliedModelingLib.proposition36bLowerEndpoint epsilon) candidate) ∧
                ¬∃ (limit : ℝ),
                    Filter.Tendsto
                      (fun (iteration : ℕ) => (AppliedModelingLib.proposition36bRRMUpdate epsilon)^[iteration] 2)
                      Filter.atTop (nhds limit)
\end{ReviewVerbatim}

\paragraph{Direct retained dependencies}
\begin{itemize}\raggedright
\item \hyperlink{governing-declaration-ee9adf0c7d04f002}{Applied\allowbreak{}Modeling\allowbreak{}Lib.\allowbreak{}Is\allowbreak{}Proposition36b\allowbreak{}Large\allowbreak{}Penalty}
\item \hyperlink{library-prerequisite-d4de05e57b5947f0}{Applied\allowbreak{}Modeling\allowbreak{}Lib.\allowbreak{}measure\allowbreak{}Decoupled\allowbreak{}Performative\allowbreak{}Risk}
\item \hyperlink{governing-declaration-ebb8a5e04e4e48ce}{Applied\allowbreak{}Modeling\allowbreak{}Lib.\allowbreak{}proposition36b\allowbreak{}Domain}
\item \hyperlink{governing-declaration-ac13973b79856aff}{Applied\allowbreak{}Modeling\allowbreak{}Lib.\allowbreak{}proposition36b\allowbreak{}Lower\allowbreak{}Endpoint}
\item \hyperlink{governing-declaration-9fe26c2f815241d3}{Applied\allowbreak{}Modeling\allowbreak{}Lib.\allowbreak{}proposition36b\allowbreak{}Measure\allowbreak{}Loss}
\item \hyperlink{governing-declaration-14367e2ad0182576}{Applied\allowbreak{}Modeling\allowbreak{}Lib.\allowbreak{}proposition36b\allowbreak{}Measure\allowbreak{}Model}
\item \hyperlink{governing-declaration-10d57ee3741bad42}{Applied\allowbreak{}Modeling\allowbreak{}Lib.\allowbreak{}proposition36b\allowbreak{}RRM\allowbreak{}Update}
\item \hyperlink{paper-prerequisite-2ea2f5769822c5c4}{PZMH20\allowbreak{}Performative\allowbreak{}Prediction.\allowbreak{}Is\allowbreak{}Measure\allowbreak{}Wasserstein\allowbreak{}Sensitive}
\end{itemize}
\hypertarget{governing-declaration-f4e329b046167e0c}{}
\subsection*{\small\ttfamily\raggedright PZMH20\allowbreak{}Performative\allowbreak{}Prediction.\allowbreak{}theorem310\allowbreak{}Corrected\allowbreak{}RERM\allowbreak{}Count\allowbreak{}Schedule}
\textbf{Declaration location:} paper-local
\paragraph{Lean declaration body}
\begin{ReviewVerbatim}
type:
ℕ → ℕ → ℝ → ℝ → ℝ → ℝ → ℝ → ℝ → ℝ → ℝ → ℕ → ℕ

value:
fun (dimension cutoff : ℕ) (eta alpha gamma momentBound tailBound scaledDeviation p headTolerance : ℝ) (a : ℕ) =>
  AppliedModelingLib.pOneFournierGuillinTheorem310SupercriticalCappedConcreteCountSchedule dimension cutoff eta alpha
    gamma momentBound tailBound scaledDeviation p headTolerance a
\end{ReviewVerbatim}

\paragraph{Direct retained dependencies}
\begin{itemize}\raggedright
\item \hyperlink{governing-declaration-3cf6b14c7bcca80c}{Applied\allowbreak{}Modeling\allowbreak{}Lib.\allowbreak{}p\allowbreak{}One\allowbreak{}Fournier\allowbreak{}Guillin\allowbreak{}Theorem310\allowbreak{}Supercritical\allowbreak{}Capped\allowbreak{}Concrete\allowbreak{}Count\allowbreak{}Schedule}
\end{itemize}
\clearpage
\hypertarget{library-section-858f18808007153a}{}
\section*{Material library prerequisites}
\hypertarget{library-prerequisite-2e05b739f5b71f23}{}
\subsection*{\small\ttfamily\raggedright Applied\allowbreak{}Modeling\allowbreak{}Lib.\allowbreak{}Is\allowbreak{}Jointly\allowbreak{}Smooth\allowbreak{}Gradient}
\textbf{Source locator:} {\footnotesize\raggedright cited publication, lines 373-\allowbreak{}415\par}
\paragraph{Verbatim paper-source connection}
\begin{ReviewVerbatim}
Theorem 3.10. Suppose that the loss `(z; θ) is β-jointly
Theorem 3.8. Suppose that the loss `(z; θ) is β-jointly                       smooth (A1) and γ-strongly convex (A2), and that there exist
smooth (A1) and γ-strongly convex (A2). If the distribution                                                def R        α
                                        γ                                     α > 1, µ > 0 such that ξα,µ = Rm eµ|x| dD(θ) is finite
map D(·) is ε-sensitive with ε < (β+γ)(1+1.5ηβ) , then RGD
                                                                              ∀θ ∈ Θ. Let δ ∈ (0, 1) be a radius of convergence.     Con-
                     2
with step size η 6 β+γ satisfies the following:
                                                                                                                                    � t
                                                                                                                            1
                                                                              sider running RERM or RGD with nt = O (εδ)m log p
(a) kGgd (θ) − Ggd (θ0 )k2                                                    samples at time t.
                                                                                                                            γ
                                                                              (a) If the map D(·) is ε-sensitive with ε < 2β  , then with
                         βγ                                          0
      6       1−η           − ε(1.5ηβ 2 + β)             kθ − θ k2 .              probability 1 − p, RERM satisfies,
                        β+γ
                                                                                                                  log 1δ kθ0 − θPS k2
                                                                                                                     �
(b) The iterates θt of RGD converge to a unique performa-                           kθt − θPS k2 6 δ, for all t >                       .
    tively stable point θPS at a linear rate,                                                                           1 − 2εβ
                                                                                                                              γ

                                                                                                                                 γ
                                               �1                             (b) If the map D(·) is ε-sensitive with ε < (β+γ)(1+1.5ηβ)  ,
                                         log    δ kθ0 − θPS k2
     kθt − θPS k2 6 δ for t >                                             .       then with probability 1 − p, REGD with satisfies,
                                          βγ          2
                                 η       β+γ − ε(1.5ηβ + β)
                                                                                                                    log 1δ kθ0 − θPS k2
                                                                                                                        �
                                                                                  kθt −θPS k2 6 δ, for all t >                              ,
The conclusion of Theorem 3.8 is a strict generalization of a                                                        βγ
                                                                                                                 η β+γ    − ε(3ηβ 2 + 2β)
classical optimization result which considers a static objec-

                                                        βγ
tive, in which case the rate of contraction is 1 − η β+γ                                                                   2
                                                                                   for a constant choice of step size η 6 β+γ .
(see for example Theorem 2.1.15 in Nesterov (2013) or
Lemma 3.7 in Hardt et al. (2016b)). Our rate exactly                          Proof sketch. The basic idea behind these results is the fol-
matches this standard result in the case that ε = 0. The                      lowing. While kθt − θPS k2 > δ, the sample size nt is
proof of Theorem 3.8 can be found in Appendix E.3.                            sufficiently large to ensure a behavior similar to that on a
\end{ReviewVerbatim}


\paragraph{Formalized review target}
The archival source above is not asserted equivalent to this formalized target.
\begin{ReviewVerbatim}
The REGD premise bundle for Theorem 3.10(b) specifies data dimension greater than two, the concrete capped count and inverse-square confidence schedules, a shared bounded exponential moment and tail envelope, measurable adaptive sampling and bad events, joint-gradient smoothness and pointwise strong convexity with Wasserstein sensitivity, a convex domain with a variational nonexpansive projection, gradient integrability and pointwise loss differentiation, a checked empirical-gradient trajectory and a stable point, and the step-size, numerical W1 budget, contraction slack, absorbing-radius, and initial-distance conditions. The trajectory theorem derives population contraction and a finite entry iteration; neither a caller-supplied small-deviation gate nor a supplied burn-in bound is a premise. The model and expected-loss integrability are defined for every ambient parameter. Joint-gradient smoothness and pointwise strong convexity quantify over all ambient parameters and data points; Wasserstein sensitivity, gradient integrability, and loss differentiation quantify over all ambient parameters. These global hypotheses are stronger in scope than the source assumptions on the parameter domain and the union of its induced supports; their necessity for the source conclusion is unresolved.
\end{ReviewVerbatim}

\textbf{Archival source anchor:} {\footnotesize\raggedright cited publication, lines 373-\allowbreak{}415\par}
\paragraph{Lean-expanded library semantic target}
\begin{ReviewVerbatim}
type:
{Parameter : Type u_1} →
  {Data : Type u_2} →
    [inst : NormedAddCommGroup Parameter] →
      [InnerProductSpace ℝ Parameter] → [MetricSpace Data] → (Data → Parameter → Parameter) → NNReal → Prop

value:
fun {Parameter : Type u_1} {Data : Type u_2} [NormedAddCommGroup Parameter] [InnerProductSpace ℝ Parameter]
    [MetricSpace Data] (gradient : Data → Parameter → Parameter) (smoothness : NNReal) =>
  (∀ (datum : Data) (first second : Parameter),
      ‖gradient datum first - gradient datum second‖ ≤ ↑smoothness * ‖first - second‖) ∧
    AppliedModelingLib.IsDataGradientLipschitz gradient smoothness
\end{ReviewVerbatim}

\paragraph{Exact Lean library declaration}
\begin{ReviewVerbatim}
/-- The two-coordinate joint-smoothness condition used in the PZ paper. -/
def IsJointlySmoothGradient {Parameter Data : Type*}
    [NormedAddCommGroup Parameter] [InnerProductSpace ℝ Parameter] [MetricSpace Data]
    (gradient : Data → Parameter → Parameter) (smoothness : NNReal) : Prop :=
  (∀ (datum : Data) (first second : Parameter),
    ‖gradient datum first - gradient datum second‖ ≤
      (smoothness : ℝ) * ‖first - second‖) ∧
  IsDataGradientLipschitz gradient smoothness
\end{ReviewVerbatim}

\paragraph{Direct retained dependencies}
\begin{itemize}\raggedright
\item \hyperlink{governing-declaration-698d9e6a0d26ec64}{Applied\allowbreak{}Modeling\allowbreak{}Lib.\allowbreak{}Is\allowbreak{}Data\allowbreak{}Gradient\allowbreak{}Lipschitz}
\end{itemize}
\paragraph{Recorded source-to-library screening}
\textbf{Verdict:} \texttt{matches\_approved\_corrected\_target}
\\
{\raggedright
\textbf{Reason:} The expanded definition is exactly the two A1 Lipschitz inequalities for the supplied gradient:\allowbreak{} uniform smoothness in the parameter coordinate and, through IsDataGradientLipschitz, in the data coordinate.\allowbreak{} Its graph-\allowbreak{}bound uses separately require HasGradientAt to identify this field with the loss gradient, and the definition's global parameter/\allowbreak{}data quantifiers are the stronger scope expressly disclosed in the corrected Theorem 3.\allowbreak{}10 target.\allowbreak{}
\par}
\reviewmatch{Matches formalized target}
\noindent\textbf{Reviewer annotation}\par
\reviewerbox
\clearpage
\hypertarget{library-prerequisite-96553d4fbb82f1ac}{}
\subsection*{\small\ttfamily\raggedright Applied\allowbreak{}Modeling\allowbreak{}Lib.\allowbreak{}Measure\allowbreak{}Performative\allowbreak{}Model}
\textbf{Source locator:} {\footnotesize\raggedright cited publication, lines 373-\allowbreak{}415\par}
\paragraph{Verbatim paper-source connection}
\begin{ReviewVerbatim}
Theorem 3.10. Suppose that the loss `(z; θ) is β-jointly
Theorem 3.8. Suppose that the loss `(z; θ) is β-jointly                       smooth (A1) and γ-strongly convex (A2), and that there exist
smooth (A1) and γ-strongly convex (A2). If the distribution                                                def R        α
                                        γ                                     α > 1, µ > 0 such that ξα,µ = Rm eµ|x| dD(θ) is finite
map D(·) is ε-sensitive with ε < (β+γ)(1+1.5ηβ) , then RGD
                                                                              ∀θ ∈ Θ. Let δ ∈ (0, 1) be a radius of convergence.     Con-
                     2
with step size η 6 β+γ satisfies the following:
                                                                                                                                    � t
                                                                                                                            1
                                                                              sider running RERM or RGD with nt = O (εδ)m log p
(a) kGgd (θ) − Ggd (θ0 )k2                                                    samples at time t.
                                                                                                                            γ
                                                                              (a) If the map D(·) is ε-sensitive with ε < 2β  , then with
                         βγ                                          0
      6       1−η           − ε(1.5ηβ 2 + β)             kθ − θ k2 .              probability 1 − p, RERM satisfies,
                        β+γ
                                                                                                                  log 1δ kθ0 − θPS k2
                                                                                                                     �
(b) The iterates θt of RGD converge to a unique performa-                           kθt − θPS k2 6 δ, for all t >                       .
    tively stable point θPS at a linear rate,                                                                           1 − 2εβ
                                                                                                                              γ

                                                                                                                                 γ
                                               �1                             (b) If the map D(·) is ε-sensitive with ε < (β+γ)(1+1.5ηβ)  ,
                                         log    δ kθ0 − θPS k2
     kθt − θPS k2 6 δ for t >                                             .       then with probability 1 − p, REGD with satisfies,
                                          βγ          2
                                 η       β+γ − ε(1.5ηβ + β)
                                                                                                                    log 1δ kθ0 − θPS k2
                                                                                                                        �
                                                                                  kθt −θPS k2 6 δ, for all t >                              ,
The conclusion of Theorem 3.8 is a strict generalization of a                                                        βγ
                                                                                                                 η β+γ    − ε(3ηβ 2 + 2β)
classical optimization result which considers a static objec-

                                                        βγ
tive, in which case the rate of contraction is 1 − η β+γ                                                                   2
                                                                                   for a constant choice of step size η 6 β+γ .
(see for example Theorem 2.1.15 in Nesterov (2013) or
Lemma 3.7 in Hardt et al. (2016b)). Our rate exactly                          Proof sketch. The basic idea behind these results is the fol-
matches this standard result in the case that ε = 0. The                      lowing. While kθt − θPS k2 > δ, the sample size nt is
proof of Theorem 3.8 can be found in Appendix E.3.                            sufficiently large to ensure a behavior similar to that on a
\end{ReviewVerbatim}


\paragraph{Formalized review target}
The archival source above is not asserted equivalent to this formalized target.
\begin{ReviewVerbatim}
The REGD premise bundle for Theorem 3.10(b) specifies data dimension greater than two, the concrete capped count and inverse-square confidence schedules, a shared bounded exponential moment and tail envelope, measurable adaptive sampling and bad events, joint-gradient smoothness and pointwise strong convexity with Wasserstein sensitivity, a convex domain with a variational nonexpansive projection, gradient integrability and pointwise loss differentiation, a checked empirical-gradient trajectory and a stable point, and the step-size, numerical W1 budget, contraction slack, absorbing-radius, and initial-distance conditions. The trajectory theorem derives population contraction and a finite entry iteration; neither a caller-supplied small-deviation gate nor a supplied burn-in bound is a premise. The model and expected-loss integrability are defined for every ambient parameter. Joint-gradient smoothness and pointwise strong convexity quantify over all ambient parameters and data points; Wasserstein sensitivity, gradient integrability, and loss differentiation quantify over all ambient parameters. These global hypotheses are stronger in scope than the source assumptions on the parameter domain and the union of its induced supports; their necessity for the source conclusion is unresolved.
\end{ReviewVerbatim}

\textbf{Archival source anchor:} {\footnotesize\raggedright cited publication, lines 373-\allowbreak{}415\par}
\paragraph{Lean declaration type (metadata)}
\begin{ReviewVerbatim}
field dataLaw:
{Parameter : Type u_1} →
  {Data : Type u_2} →
    [inst : MeasurableSpace Data] →
      AppliedModelingLib.MeasurePerformativeModel Parameter Data → Parameter → MeasureTheory.ProbabilityMeasure Data

field loss:
{Parameter : Type u_1} →
  {Data : Type u_2} →
    [inst : MeasurableSpace Data] → AppliedModelingLib.MeasurePerformativeModel Parameter Data → Data → Parameter → ℝ

field loss_integrable:
∀ {Parameter : Type u_1} {Data : Type u_2} [inst : MeasurableSpace Data]
  (self : AppliedModelingLib.MeasurePerformativeModel Parameter Data)
  (distributionParameter evaluatedParameter : Parameter),
  MeasureTheory.Integrable (fun (datum : Data) => self.loss datum evaluatedParameter)
    ↑(self.dataLaw distributionParameter)
\end{ReviewVerbatim}

\paragraph{Exact Lean library declaration}
\begin{ReviewVerbatim}
/--
A performative model over arbitrary probability measures.  Integrability is
part of the model because the general-distribution population risks below are
real Bochner integrals rather than finite sums.
-/
structure MeasurePerformativeModel (Parameter Data : Type*) [MeasurableSpace Data] where
  dataLaw : Parameter → ProbabilityMeasure Data
  loss : Data → Parameter → ℝ
  loss_integrable : ∀ distributionParameter evaluatedParameter,
    Integrable (fun datum => loss datum evaluatedParameter)
      (dataLaw distributionParameter : Measure Data)
\end{ReviewVerbatim}

\paragraph{Recorded source-to-library screening}
\textbf{Verdict:} \texttt{matches\_approved\_corrected\_target}
\\
{\raggedright
\textbf{Reason:} The expanded structure consists of a probability data law for every deployed parameter, a loss for every datum and evaluated parameter, and integrability of that loss under every deployed/\allowbreak{}evaluated pair.\allowbreak{} Those are exactly the globally defined model and expected-\allowbreak{}loss integrability conditions disclosed in the corrected Theorem 3.\allowbreak{}10 target and assumed by every displayed material use.\allowbreak{}
\par}
\reviewmatch{Matches formalized target}
\noindent\textbf{Reviewer annotation}\par
\reviewerbox
\clearpage
\hypertarget{library-prerequisite-49496bb3a06c52e5}{}
\subsection*{\small\ttfamily\raggedright Applied\allowbreak{}Modeling\allowbreak{}Lib.\allowbreak{}Is\allowbreak{}Measure\allowbreak{}Data\allowbreak{}Loss\allowbreak{}Lipschitz}
\textbf{Source locator:} {\footnotesize\raggedright cited publication, lines 373-\allowbreak{}415\par}
\paragraph{Verbatim paper-source connection}
\begin{ReviewVerbatim}
Theorem 3.10. Suppose that the loss `(z; θ) is β-jointly
Theorem 3.8. Suppose that the loss `(z; θ) is β-jointly                       smooth (A1) and γ-strongly convex (A2), and that there exist
smooth (A1) and γ-strongly convex (A2). If the distribution                                                def R        α
                                        γ                                     α > 1, µ > 0 such that ξα,µ = Rm eµ|x| dD(θ) is finite
map D(·) is ε-sensitive with ε < (β+γ)(1+1.5ηβ) , then RGD
                                                                              ∀θ ∈ Θ. Let δ ∈ (0, 1) be a radius of convergence.     Con-
                     2
with step size η 6 β+γ satisfies the following:
                                                                                                                                    � t
                                                                                                                            1
                                                                              sider running RERM or RGD with nt = O (εδ)m log p
(a) kGgd (θ) − Ggd (θ0 )k2                                                    samples at time t.
                                                                                                                            γ
                                                                              (a) If the map D(·) is ε-sensitive with ε < 2β  , then with
                         βγ                                          0
      6       1−η           − ε(1.5ηβ 2 + β)             kθ − θ k2 .              probability 1 − p, RERM satisfies,
                        β+γ
                                                                                                                  log 1δ kθ0 − θPS k2
                                                                                                                     �
(b) The iterates θt of RGD converge to a unique performa-                           kθt − θPS k2 6 δ, for all t >                       .
    tively stable point θPS at a linear rate,                                                                           1 − 2εβ
                                                                                                                              γ

                                                                                                                                 γ
                                               �1                             (b) If the map D(·) is ε-sensitive with ε < (β+γ)(1+1.5ηβ)  ,
                                         log    δ kθ0 − θPS k2
     kθt − θPS k2 6 δ for t >                                             .       then with probability 1 − p, REGD with satisfies,
                                          βγ          2
                                 η       β+γ − ε(1.5ηβ + β)
                                                                                                                    log 1δ kθ0 − θPS k2
                                                                                                                        �
                                                                                  kθt −θPS k2 6 δ, for all t >                              ,
The conclusion of Theorem 3.8 is a strict generalization of a                                                        βγ
                                                                                                                 η β+γ    − ε(3ηβ 2 + 2β)
classical optimization result which considers a static objec-

                                                        βγ
tive, in which case the rate of contraction is 1 − η β+γ                                                                   2
                                                                                   for a constant choice of step size η 6 β+γ .
(see for example Theorem 2.1.15 in Nesterov (2013) or
Lemma 3.7 in Hardt et al. (2016b)). Our rate exactly                          Proof sketch. The basic idea behind these results is the fol-
matches this standard result in the case that ε = 0. The                      lowing. While kθt − θPS k2 > δ, the sample size nt is
proof of Theorem 3.8 can be found in Appendix E.3.                            sufficiently large to ensure a behavior similar to that on a
\end{ReviewVerbatim}


\paragraph{Formalized review target}
The archival source above is not asserted equivalent to this formalized target.
\begin{ReviewVerbatim}
The RERM premise bundle for Theorem 3.10(a) specifies data dimension greater than two, the concrete capped count and inverse-square confidence schedules, a shared bounded exponential moment and tail envelope, measurable adaptive sampling and bad events, data-Lipschitz loss and empirical-loss integrability, smoothness and strong convexity with Wasserstein sensitivity, a convex domain with empirical and population risk-minimizing updates and a stable point, and the numerical W1 budget, contraction slack, absorbing-radius, and initial-distance conditions. The trajectory theorem derives population contraction and a finite entry iteration; neither a caller-supplied small-deviation gate nor a supplied burn-in bound is a premise. The model and expected-loss integrability are defined for every ambient parameter. Joint-gradient smoothness and pointwise strong convexity quantify over all ambient parameters and data points; Wasserstein sensitivity, gradient integrability, and loss differentiation quantify over all ambient parameters. RERM data-loss Lipschitzness also quantifies over every ambient parameter and data point. These global hypotheses are stronger in scope than the source assumptions on the parameter domain and the union of its induced supports; their necessity for the source conclusion is unresolved.
\end{ReviewVerbatim}

\textbf{Archival source anchor:} {\footnotesize\raggedright cited publication, lines 373-\allowbreak{}415\par}
\paragraph{Lean-expanded library semantic target}
\begin{ReviewVerbatim}
type:
{Parameter : Type u_1} →
  {Data : Type u_2} →
    [inst : MeasurableSpace Data] →
      [MetricSpace Data] → AppliedModelingLib.MeasurePerformativeModel Parameter Data → NNReal → Prop

value:
fun {Parameter : Type u_1} {Data : Type u_2} [MeasurableSpace Data] [MetricSpace Data]
    (model : AppliedModelingLib.MeasurePerformativeModel Parameter Data) (constant : NNReal) =>
  ∀ (parameter : Parameter), LipschitzWith constant fun (datum : Data) => model.loss datum parameter
\end{ReviewVerbatim}

\paragraph{Exact Lean library declaration}
\begin{ReviewVerbatim}
/-- A general-law loss is uniformly Lipschitz in its data coordinate. -/
def IsMeasureDataLossLipschitz {Parameter Data : Type*}
    [MeasurableSpace Data] [MetricSpace Data]
    (model : MeasurePerformativeModel Parameter Data) (constant : NNReal) : Prop :=
  ∀ parameter, LipschitzWith constant (fun datum => model.loss datum parameter)
\end{ReviewVerbatim}

\paragraph{Direct retained dependencies}
\begin{itemize}\raggedright
\item \hyperlink{library-prerequisite-96553d4fbb82f1ac}{Applied\allowbreak{}Modeling\allowbreak{}Lib.\allowbreak{}Measure\allowbreak{}Performative\allowbreak{}Model}
\end{itemize}
\paragraph{Recorded source-to-library screening}
\textbf{Verdict:} \texttt{matches\_approved\_corrected\_target}
\\
{\raggedright
\textbf{Reason:} The expanded definition requires one NNReal constant to make the loss Lipschitz in the data coordinate for every ambient parameter.\allowbreak{} This is exactly the globally scoped RERM Wasserstein-\allowbreak{}to-\allowbreak{}loss bridge expressly included in the corrected Theorem 3.\allowbreak{}10(a) premise bundle and used with the same model and constant.\allowbreak{}
\par}
\reviewmatch{Matches formalized target}
\noindent\textbf{Reviewer annotation}\par
\reviewerbox
\clearpage
\hypertarget{library-prerequisite-b41cbb32daa3dce7}{}
\subsection*{\small\ttfamily\raggedright Applied\allowbreak{}Modeling\allowbreak{}Lib.\allowbreak{}Is\allowbreak{}Measure\allowbreak{}Pointwise\allowbreak{}Gradient\allowbreak{}Strongly\allowbreak{}Convex}
\textbf{Source locator:} {\footnotesize\raggedright cited publication, lines 373-\allowbreak{}415\par}
\paragraph{Verbatim paper-source connection}
\begin{ReviewVerbatim}
Theorem 3.10. Suppose that the loss `(z; θ) is β-jointly
Theorem 3.8. Suppose that the loss `(z; θ) is β-jointly                       smooth (A1) and γ-strongly convex (A2), and that there exist
smooth (A1) and γ-strongly convex (A2). If the distribution                                                def R        α
                                        γ                                     α > 1, µ > 0 such that ξα,µ = Rm eµ|x| dD(θ) is finite
map D(·) is ε-sensitive with ε < (β+γ)(1+1.5ηβ) , then RGD
                                                                              ∀θ ∈ Θ. Let δ ∈ (0, 1) be a radius of convergence.     Con-
                     2
with step size η 6 β+γ satisfies the following:
                                                                                                                                    � t
                                                                                                                            1
                                                                              sider running RERM or RGD with nt = O (εδ)m log p
(a) kGgd (θ) − Ggd (θ0 )k2                                                    samples at time t.
                                                                                                                            γ
                                                                              (a) If the map D(·) is ε-sensitive with ε < 2β  , then with
                         βγ                                          0
      6       1−η           − ε(1.5ηβ 2 + β)             kθ − θ k2 .              probability 1 − p, RERM satisfies,
                        β+γ
                                                                                                                  log 1δ kθ0 − θPS k2
                                                                                                                     �
(b) The iterates θt of RGD converge to a unique performa-                           kθt − θPS k2 6 δ, for all t >                       .
    tively stable point θPS at a linear rate,                                                                           1 − 2εβ
                                                                                                                              γ

                                                                                                                                 γ
                                               �1                             (b) If the map D(·) is ε-sensitive with ε < (β+γ)(1+1.5ηβ)  ,
                                         log    δ kθ0 − θPS k2
     kθt − θPS k2 6 δ for t >                                             .       then with probability 1 − p, REGD with satisfies,
                                          βγ          2
                                 η       β+γ − ε(1.5ηβ + β)
                                                                                                                    log 1δ kθ0 − θPS k2
                                                                                                                        �
                                                                                  kθt −θPS k2 6 δ, for all t >                              ,
The conclusion of Theorem 3.8 is a strict generalization of a                                                        βγ
                                                                                                                 η β+γ    − ε(3ηβ 2 + 2β)
classical optimization result which considers a static objec-

                                                        βγ
tive, in which case the rate of contraction is 1 − η β+γ                                                                   2
                                                                                   for a constant choice of step size η 6 β+γ .
(see for example Theorem 2.1.15 in Nesterov (2013) or
Lemma 3.7 in Hardt et al. (2016b)). Our rate exactly                          Proof sketch. The basic idea behind these results is the fol-
matches this standard result in the case that ε = 0. The                      lowing. While kθt − θPS k2 > δ, the sample size nt is
proof of Theorem 3.8 can be found in Appendix E.3.                            sufficiently large to ensure a behavior similar to that on a
\end{ReviewVerbatim}


\paragraph{Formalized review target}
The archival source above is not asserted equivalent to this formalized target.
\begin{ReviewVerbatim}
The REGD premise bundle for Theorem 3.10(b) specifies data dimension greater than two, the concrete capped count and inverse-square confidence schedules, a shared bounded exponential moment and tail envelope, measurable adaptive sampling and bad events, joint-gradient smoothness and pointwise strong convexity with Wasserstein sensitivity, a convex domain with a variational nonexpansive projection, gradient integrability and pointwise loss differentiation, a checked empirical-gradient trajectory and a stable point, and the step-size, numerical W1 budget, contraction slack, absorbing-radius, and initial-distance conditions. The trajectory theorem derives population contraction and a finite entry iteration; neither a caller-supplied small-deviation gate nor a supplied burn-in bound is a premise. The model and expected-loss integrability are defined for every ambient parameter. Joint-gradient smoothness and pointwise strong convexity quantify over all ambient parameters and data points; Wasserstein sensitivity, gradient integrability, and loss differentiation quantify over all ambient parameters. These global hypotheses are stronger in scope than the source assumptions on the parameter domain and the union of its induced supports; their necessity for the source conclusion is unresolved.
\end{ReviewVerbatim}

\textbf{Archival source anchor:} {\footnotesize\raggedright cited publication, lines 373-\allowbreak{}415\par}
\paragraph{Lean-expanded library semantic target}
\begin{ReviewVerbatim}
type:
{Parameter : Type u_1} →
  {Data : Type u_2} →
    [inst : MeasurableSpace Data] →
      [inst_1 : NormedAddCommGroup Parameter] →
        [InnerProductSpace ℝ Parameter] →
          AppliedModelingLib.MeasurePerformativeModel Parameter Data → (Data → Parameter → Parameter) → ℝ → Prop

value:
fun {Parameter : Type u_1} {Data : Type u_2} [MeasurableSpace Data] [NormedAddCommGroup Parameter]
    [InnerProductSpace ℝ Parameter] (model : AppliedModelingLib.MeasurePerformativeModel Parameter Data)
    (gradient : Data → Parameter → Parameter) (modulus : ℝ) =>
  ∀ (datum : Data) (center candidate : Parameter),
    model.loss datum candidate ≥
      model.loss datum center + inner ℝ (gradient datum center) (candidate - center) +
        modulus / 2 * ‖candidate - center‖ ^ 2
\end{ReviewVerbatim}

\paragraph{Exact Lean library declaration}
\begin{ReviewVerbatim}
/--
The pointwise form of the source's A2 strong-convexity assumption for an
arbitrary probability data law.  Together with integrability of the gradient,
it implies strong monotonicity of the frozen expected gradient below.
-/
def IsMeasurePointwiseGradientStronglyConvex {Parameter Data : Type*}
    [MeasurableSpace Data] [NormedAddCommGroup Parameter] [InnerProductSpace ℝ Parameter]
    (model : MeasurePerformativeModel Parameter Data)
    (gradient : Data → Parameter → Parameter) (modulus : ℝ) : Prop :=
  ∀ datum center candidate,
    model.loss datum candidate ≥
      model.loss datum center +
        ⟪gradient datum center, candidate - center⟫_ℝ +
          modulus / 2 * ‖candidate - center‖ ^ 2
\end{ReviewVerbatim}

\paragraph{Direct retained dependencies}
\begin{itemize}\raggedright
\item \hyperlink{library-prerequisite-96553d4fbb82f1ac}{Applied\allowbreak{}Modeling\allowbreak{}Lib.\allowbreak{}Measure\allowbreak{}Performative\allowbreak{}Model}
\end{itemize}
\paragraph{Recorded source-to-library screening}
\textbf{Verdict:} \texttt{matches\_approved\_corrected\_target}
\\
{\raggedright
\textbf{Reason:} The expanded definition states the A2 lower bound with modulus /\allowbreak{} 2 times squared parameter distance for every datum, center, and candidate, using the supplied gradient at the center.\allowbreak{} The graph-\allowbreak{}bound carriers separately identify that gradient with HasGradientAt;\allowbreak{} the all-\allowbreak{}ambient quantifier scope is exactly the stronger scope disclosed by the corrected target.\allowbreak{}
\par}
\reviewmatch{Matches formalized target}
\noindent\textbf{Reviewer annotation}\par
\reviewerbox
\clearpage
\hypertarget{library-prerequisite-3c6dd4ad7ac116e3}{}
\subsection*{\small\ttfamily\raggedright Applied\allowbreak{}Modeling\allowbreak{}Lib.\allowbreak{}Is\allowbreak{}Measure\allowbreak{}Wasserstein\allowbreak{}Sensitive}
\textbf{Source locator:} {\footnotesize\raggedright cited publication, lines 373-\allowbreak{}415\par}
\paragraph{Verbatim paper-source connection}
\begin{ReviewVerbatim}
Theorem 3.10. Suppose that the loss `(z; θ) is β-jointly
Theorem 3.8. Suppose that the loss `(z; θ) is β-jointly                       smooth (A1) and γ-strongly convex (A2), and that there exist
smooth (A1) and γ-strongly convex (A2). If the distribution                                                def R        α
                                        γ                                     α > 1, µ > 0 such that ξα,µ = Rm eµ|x| dD(θ) is finite
map D(·) is ε-sensitive with ε < (β+γ)(1+1.5ηβ) , then RGD
                                                                              ∀θ ∈ Θ. Let δ ∈ (0, 1) be a radius of convergence.     Con-
                     2
with step size η 6 β+γ satisfies the following:
                                                                                                                                    � t
                                                                                                                            1
                                                                              sider running RERM or RGD with nt = O (εδ)m log p
(a) kGgd (θ) − Ggd (θ0 )k2                                                    samples at time t.
                                                                                                                            γ
                                                                              (a) If the map D(·) is ε-sensitive with ε < 2β  , then with
                         βγ                                          0
      6       1−η           − ε(1.5ηβ 2 + β)             kθ − θ k2 .              probability 1 − p, RERM satisfies,
                        β+γ
                                                                                                                  log 1δ kθ0 − θPS k2
                                                                                                                     �
(b) The iterates θt of RGD converge to a unique performa-                           kθt − θPS k2 6 δ, for all t >                       .
    tively stable point θPS at a linear rate,                                                                           1 − 2εβ
                                                                                                                              γ

                                                                                                                                 γ
                                               �1                             (b) If the map D(·) is ε-sensitive with ε < (β+γ)(1+1.5ηβ)  ,
                                         log    δ kθ0 − θPS k2
     kθt − θPS k2 6 δ for t >                                             .       then with probability 1 − p, REGD with satisfies,
                                          βγ          2
                                 η       β+γ − ε(1.5ηβ + β)
                                                                                                                    log 1δ kθ0 − θPS k2
                                                                                                                        �
                                                                                  kθt −θPS k2 6 δ, for all t >                              ,
The conclusion of Theorem 3.8 is a strict generalization of a                                                        βγ
                                                                                                                 η β+γ    − ε(3ηβ 2 + 2β)
classical optimization result which considers a static objec-

                                                        βγ
tive, in which case the rate of contraction is 1 − η β+γ                                                                   2
                                                                                   for a constant choice of step size η 6 β+γ .
(see for example Theorem 2.1.15 in Nesterov (2013) or
Lemma 3.7 in Hardt et al. (2016b)). Our rate exactly                          Proof sketch. The basic idea behind these results is the fol-
matches this standard result in the case that ε = 0. The                      lowing. While kθt − θPS k2 > δ, the sample size nt is
proof of Theorem 3.8 can be found in Appendix E.3.                            sufficiently large to ensure a behavior similar to that on a
\end{ReviewVerbatim}


\paragraph{Formalized review target}
The archival source above is not asserted equivalent to this formalized target.
\begin{ReviewVerbatim}
The REGD premise bundle for Theorem 3.10(b) specifies data dimension greater than two, the concrete capped count and inverse-square confidence schedules, a shared bounded exponential moment and tail envelope, measurable adaptive sampling and bad events, joint-gradient smoothness and pointwise strong convexity with Wasserstein sensitivity, a convex domain with a variational nonexpansive projection, gradient integrability and pointwise loss differentiation, a checked empirical-gradient trajectory and a stable point, and the step-size, numerical W1 budget, contraction slack, absorbing-radius, and initial-distance conditions. The trajectory theorem derives population contraction and a finite entry iteration; neither a caller-supplied small-deviation gate nor a supplied burn-in bound is a premise. The model and expected-loss integrability are defined for every ambient parameter. Joint-gradient smoothness and pointwise strong convexity quantify over all ambient parameters and data points; Wasserstein sensitivity, gradient integrability, and loss differentiation quantify over all ambient parameters. These global hypotheses are stronger in scope than the source assumptions on the parameter domain and the union of its induced supports; their necessity for the source conclusion is unresolved.
\end{ReviewVerbatim}

\textbf{Archival source anchor:} {\footnotesize\raggedright cited publication, lines 373-\allowbreak{}415\par}
\paragraph{Lean-expanded library semantic target}
\begin{ReviewVerbatim}
type:
{Parameter : Type u_1} →
  {Data : Type u_2} →
    [MetricSpace Parameter] →
      [inst : MeasurableSpace Data] →
        [MetricSpace Data] → AppliedModelingLib.MeasurePerformativeModel Parameter Data → ℝ → Prop

value:
fun {Parameter : Type u_1} {Data : Type u_2} [MetricSpace Parameter] [MeasurableSpace Data] [MetricSpace Data]
    (model : AppliedModelingLib.MeasurePerformativeModel Parameter Data) (sensitivity : ℝ) =>
  ∀ (first second : Parameter),
    ∃ (hfinite :
      (AppliedModelingLib.ProbabilityCoupling.expectedDistanceCosts (model.dataLaw first)
          (model.dataLaw second)).Nonempty),
      AppliedModelingLib.ProbabilityCoupling.wassersteinOneReal (model.dataLaw first) (model.dataLaw second) hfinite ≤
        sensitivity * dist first second
\end{ReviewVerbatim}

\paragraph{Exact Lean library declaration}
\begin{ReviewVerbatim}
/--
The source paper's general-distribution Wasserstein-1 sensitivity condition.
For every pair of deployed parameters, the induced laws have a finite
expected-distance coupling and their real W₁ cost is bounded by the displayed
parameter-distance modulus.
-/
def IsMeasureWassersteinSensitive {Parameter Data : Type*}
    [MetricSpace Parameter] [MeasurableSpace Data] [MetricSpace Data]
    (model : MeasurePerformativeModel Parameter Data) (sensitivity : ℝ) : Prop :=
  ∀ first second,
    ∃ hfinite : (ProbabilityCoupling.expectedDistanceCosts
      (model.dataLaw first) (model.dataLaw second)).Nonempty,
      ProbabilityCoupling.wassersteinOneReal (model.dataLaw first) (model.dataLaw second) hfinite ≤
        sensitivity * dist first second
\end{ReviewVerbatim}

\paragraph{Direct retained dependencies}
\begin{itemize}\raggedright
\item \hyperlink{library-prerequisite-96553d4fbb82f1ac}{Applied\allowbreak{}Modeling\allowbreak{}Lib.\allowbreak{}Measure\allowbreak{}Performative\allowbreak{}Model}
\item \hyperlink{governing-declaration-0507d24247d6815a}{Applied\allowbreak{}Modeling\allowbreak{}Lib.\allowbreak{}Probability\allowbreak{}Coupling.\allowbreak{}expected\allowbreak{}Distance\allowbreak{}Costs}
\item \hyperlink{governing-declaration-7fdf7bcf28dbe52b}{Applied\allowbreak{}Modeling\allowbreak{}Lib.\allowbreak{}Probability\allowbreak{}Coupling.\allowbreak{}wasserstein\allowbreak{}One\allowbreak{}Real}
\end{itemize}
\paragraph{Recorded source-to-library screening}
\textbf{Verdict:} \texttt{matches\_approved\_corrected\_target}
\\
{\raggedright
\textbf{Reason:} For every pair of ambient parameters, the expanded definition supplies a finite expected-\allowbreak{}distance coupling and bounds the resulting real W1 infimum by sensitivity times parameter distance.\allowbreak{} This is the finite-\allowbreak{}first-\allowbreak{}moment Wasserstein-\allowbreak{}1 reading and global scope expressly disclosed in the corrected Theorem 3.\allowbreak{}10 premises and used unchanged by the carriers.\allowbreak{}
\par}
\reviewmatch{Matches formalized target}
\noindent\textbf{Reviewer annotation}\par
\reviewerbox
\clearpage
\hypertarget{library-prerequisite-e14cc5e8b1e1460e}{}
\subsection*{\small\ttfamily\raggedright Applied\allowbreak{}Modeling\allowbreak{}Lib.\allowbreak{}Measure\allowbreak{}Performative\allowbreak{}Model\allowbreak{}On}
\textbf{Source locator:} {\footnotesize\raggedright cited publication, lines 151-\allowbreak{}172\par}
\paragraph{Verbatim paper-source connection}
\begin{ReviewVerbatim}
Recently, there has been increased interest within the al-        We contrast this with the performative optimum. As in-
gorithmic fairness community in classification dynamics.           troduced previously, in settings where predictions support
See, for example, Liu et al. (2018), Hu & Chen (2018), and        decisions, the manifested distribution over features and out-
Hashimoto et al. (2018). The latter work considers repeated       comes is in part determined by the deployed classifier. In-
risk minimization, but from the perspective of what it does       stead of considering a fixed distribution D, each classifier
to a measure of disparity between groups.                         fθ induces a potentially different distribution D(θ) over
                                                                  instances z. A predictor must therefore be evaluated with
Causal inference. The reader familiar with causality can          regard to the expected loss over the distribution D(θ) it
think of D(θ) as the interventional distribution over in-         induces: its performative risk.
stances Z resulting from a do-intervention that sets the
                                                                  Definition 2.1 (performative optimality and risk). A classi-
model parameters to θ in some underlying causal graph.
                                                                  fier fθPO is performatively optimal if the following relation-
Importantly, this mapping D(·) remains fixed and does not
                                                                  ship holds:
change over time or by intervention: deploying the same
classifier at two different points in time must induce the                       θPO = arg min      E      `(Z; θ).
same distribution over observations Z. While causal in-                                    θ     Z∼D(θ)

ference focuses on estimating properties of interventional                          def
distributions such as treatment effects (Pearl, 2009; Imbens      We define PR(θ) = EZ∼D(θ) `(Z; θ) as the performative
& Rubin, 2015), our focus is on a new stability notion and        risk; then, θPO = arg minθ PR(θ).

[Next verbatim source excerpt]

Throughout our presentation, we focus on predictive models        Here, ε represents the performative aspect of the model. If
that are parametrized by a vector θ ∈ Θ, where the param-         ε = 0, outcomes are independent of the assigned scores
eter space Θ ⊆ Rd is a closed, convex set. We use capi-           and the problem reduces to standard supervised learning
tal letters to denote random variables and their lowercase        where the optimal predictor is the conditional expectation
counterparts to denote realizations of these variables. We        fθSL (x) = E[Y | X = x] = 12 + µx, with θSL = µ.
consider instances z = (x, y) defined as feature, outcome          In the performative setting with ε 6= 0, the optimal predic-
pairs, where x ∈ Rm−1 and y ∈ R. Whenever we define a              tion θPO balances between its predictive accuracy as well

[Next verbatim source excerpt]

Definition 2.3 (performative stability and decoupled risk).              Points
A classifier fθPS is performatively stable if the following           Having introduced our framework for performative predic-
relationship holds:                                                  tion, we now address some of the basic questions that arise
             θPS = arg min          E        `(Z; θ).                in this setting and examine the behavior of common ma-
                        θ        Z∼D(θPS )                           chine learning practices, such as retraining, through the lens
                                                                     of performativity. We begin by analyzing the behavior of
                                  def
We define DPR(θ, θ0 )     =    EZ∼D(θ) `(Z; θ0 ) as                   these procedures when they operate at a population level
the decoupled performative risk; then, θPS      =                    and then extend our analysis to finite samples.
arg minθ DPR(θPS , θ).
\end{ReviewVerbatim}



\paragraph{Lean declaration type (metadata)}
\begin{ReviewVerbatim}
field dataLaw:
{Parameter : Type u_1} →
  {Data : Type u_2} →
    [inst : MeasurableSpace Data] →
      {domain : Set Parameter} →
        AppliedModelingLib.MeasurePerformativeModelOn Parameter Data domain →
          ↑domain → MeasureTheory.ProbabilityMeasure Data

field loss:
{Parameter : Type u_1} →
  {Data : Type u_2} →
    [inst : MeasurableSpace Data] →
      {domain : Set Parameter} →
        AppliedModelingLib.MeasurePerformativeModelOn Parameter Data domain → Data → ↑domain → ℝ

field loss_integrable:
∀ {Parameter : Type u_1} {Data : Type u_2} [inst : MeasurableSpace Data] {domain : Set Parameter}
  (self : AppliedModelingLib.MeasurePerformativeModelOn Parameter Data domain)
  (distributionParameter evaluatedParameter : ↑domain),
  MeasureTheory.Integrable (fun (datum : Data) => self.loss datum evaluatedParameter)
    ↑(self.dataLaw distributionParameter)
\end{ReviewVerbatim}

\paragraph{Exact Lean library declaration}
\begin{ReviewVerbatim}
/-- A performative model defined only on a designated parameter domain. -/
structure MeasurePerformativeModelOn (Parameter Data : Type*) [MeasurableSpace Data]
    (domain : Set Parameter) where
  dataLaw : domain → ProbabilityMeasure Data
  loss : Data → domain → ℝ
  loss_integrable : ∀ distributionParameter evaluatedParameter,
    Integrable (fun datum => loss datum evaluatedParameter)
      (dataLaw distributionParameter : Measure Data)
\end{ReviewVerbatim}

\paragraph{Recorded source-to-library screening}
\textbf{Verdict:} \texttt{matches}
\\
{\raggedright
\textbf{Reason:} Compared domain\_\allowbreak{}relative\_\allowbreak{}performative\_\allowbreak{}model at cited publication:\allowbreak{}151-\allowbreak{}172,188-\allowbreak{}194,219-\allowbreak{}228 with the three record-\allowbreak{}field displays and every listed risk/\allowbreak{}predicate/\allowbreak{}use display.\allowbreak{} dataLaw is a fixed probability-\allowbreak{}law map on Theta, and loss has a datum and a feasible parameter as its arguments.\allowbreak{} The independently quantified distributionParameter and evaluatedParameter in loss\_\allowbreak{}integrable cover exactly the two parameters of the source DPR expectation, with PR as its diagonal.\allowbreak{} These fields give the source risks as real expectations without requiring a distribution or loss value outside Theta.\allowbreak{} The displayed paper use supplies closedness and convexity of the domain, a positive A2 modulus, the uniform data-\allowbreak{}Lipschitz and nonnegative sensitivity hypotheses, and the optimal/\allowbreak{}stable domain members.\allowbreak{} Its A2 prerequisite binds the loss derivative within this same domain;\allowbreak{} the explicit zero extension is read only within the domain and therefore adds no off-\allowbreak{}domain regularity condition.\allowbreak{}
\par}
\reviewmatch{Matches source input}
\noindent\textbf{Reviewer annotation}\par
\reviewerbox
\clearpage
\hypertarget{library-prerequisite-0de79b655cc05174}{}
\subsection*{\small\ttfamily\raggedright Applied\allowbreak{}Modeling\allowbreak{}Lib.\allowbreak{}Measure\allowbreak{}Performative\allowbreak{}Model\allowbreak{}On.\allowbreak{}Is\allowbreak{}Data\allowbreak{}Loss\allowbreak{}Lipschitz}
\textbf{Source locator:} {\footnotesize\raggedright cited publication, lines 468-\allowbreak{}480\par}
\paragraph{Verbatim paper-source connection}
\begin{ReviewVerbatim}
tee that the resulting classifier has low performative risk. In        said to have a Stackelberg structure since agents adapt their
particular, our next result demonstrates that if the loss func-       features only after the bank has deployed their classifier.
tion `(z; θ) is Lipschitz in z and γ-strongly convex, then
all performatively stable points and performative optima              The optimal strategy for the institution in a strategic classifi-
lie in a small neighborhood around each other. Moreover,              cation setting is to deploy the solution corresponding to the
the theorem holds for cases where performative optima and             Stackelberg equilibrium, defined as the classifier fθ which
performatively stable points are not necessarily unique.              achieves minimal loss over the induced distribution D(θ)
                                                                      in which agents have strategically adapted their features in
Theorem 4.3. Suppose that the loss `(z; θ) is Lz -Lipschitz           response to fθ . In fact, we see that this equilibrium notion
in z, γ-strongly convex (A2), and that the distribution map           exactly matches our definition of performative optimality:
    1
      In particular, they are guaranteed to exist over the extended
real line, i.e. we allow θ ∈ (R ∪ {±∞})d .                            fθSE is a Stackelberg equilibrium ⇔ θSE ∈ arg min PR(θ).
\end{ReviewVerbatim}



\paragraph{Lean-expanded library semantic target}
\begin{ReviewVerbatim}
type:
{Parameter : Type u_1} →
  {Data : Type u_2} →
    [inst : MeasurableSpace Data] →
      [MetricSpace Data] →
        {domain : Set Parameter} → AppliedModelingLib.MeasurePerformativeModelOn Parameter Data domain → NNReal → Prop

value:
fun {Parameter : Type u_1} {Data : Type u_2} [MeasurableSpace Data] [MetricSpace Data] {domain : Set Parameter}
    (model : AppliedModelingLib.MeasurePerformativeModelOn Parameter Data domain) (constant : NNReal) =>
  ∀ (parameter : ↑domain), LipschitzWith constant fun (datum : Data) => model.loss datum parameter
\end{ReviewVerbatim}

\paragraph{Exact Lean library declaration}
\begin{ReviewVerbatim}
/-- Data-coordinate loss Lipschitzness on a model's parameter domain. -/
def IsDataLossLipschitz
    {Parameter Data : Type*} [MeasurableSpace Data] [MetricSpace Data]
    {domain : Set Parameter} (model : MeasurePerformativeModelOn Parameter Data domain)
    (constant : NNReal) : Prop :=
  ∀ parameter, LipschitzWith constant (fun datum => model.loss datum parameter)
\end{ReviewVerbatim}

\paragraph{Direct retained dependencies}
\begin{itemize}\raggedright
\item \hyperlink{library-prerequisite-e14cc5e8b1e1460e}{Applied\allowbreak{}Modeling\allowbreak{}Lib.\allowbreak{}Measure\allowbreak{}Performative\allowbreak{}Model\allowbreak{}On}
\end{itemize}
\paragraph{Recorded source-to-library screening}
\textbf{Verdict:} \texttt{matches}
\\
{\raggedright
\textbf{Reason:} Compared theorem\_\allowbreak{}4\_\allowbreak{}3\_\allowbreak{}data\_\allowbreak{}lipschitz\_\allowbreak{}premise at cited publication:\allowbreak{}468-\allowbreak{}480 with the expanded predicate, model fields, and theorem-\allowbreak{}use display.\allowbreak{} For every parameter in Theta, the same constant bounds the Lipschitz modulus of datum -\allowbreak{}> loss datum parameter.\allowbreak{} Thus the target is the uniform data-\allowbreak{}coordinate inequality |loss(z,theta)-\allowbreak{}loss(z-\allowbreak{}prime,theta)| <= Lz dist(z,z-\allowbreak{}prime), as required by the source premise.\allowbreak{} The NNReal parameter permits Lz=0 and supplies its nonnegativity;\allowbreak{} no parameter-\allowbreak{}coordinate Lipschitz condition is introduced.\allowbreak{} The use retains this same constant and the domain membership of all parameters.\allowbreak{} Model integrability provides the real expectations required by the displayed risk comparisons.\allowbreak{}
\par}
\reviewmatch{Matches source input}
\noindent\textbf{Reviewer annotation}\par
\reviewerbox
\clearpage
\hypertarget{library-prerequisite-25221c4f931beacb}{}
\subsection*{\small\ttfamily\raggedright Applied\allowbreak{}Modeling\allowbreak{}Lib.\allowbreak{}Measure\allowbreak{}Performative\allowbreak{}Model\allowbreak{}On.\allowbreak{}decoupled\allowbreak{}Performative\allowbreak{}Risk}
\textbf{Source locator:} {\footnotesize\raggedright cited publication, lines 219-\allowbreak{}228\par}
\paragraph{Verbatim paper-source connection}
\begin{ReviewVerbatim}
Definition 2.3 (performative stability and decoupled risk).              Points
A classifier fθPS is performatively stable if the following           Having introduced our framework for performative predic-
relationship holds:                                                  tion, we now address some of the basic questions that arise
             θPS = arg min          E        `(Z; θ).                in this setting and examine the behavior of common ma-
                        θ        Z∼D(θPS )                           chine learning practices, such as retraining, through the lens
                                                                     of performativity. We begin by analyzing the behavior of
                                  def
We define DPR(θ, θ0 )     =    EZ∼D(θ) `(Z; θ0 ) as                   these procedures when they operate at a population level
the decoupled performative risk; then, θPS      =                    and then extend our analysis to finite samples.
arg minθ DPR(θPS , θ).
\end{ReviewVerbatim}



\paragraph{Lean-expanded library semantic target}
\begin{ReviewVerbatim}
type:
{Parameter : Type u_1} →
  {Data : Type u_2} →
    [inst : MeasurableSpace Data] →
      {domain : Set Parameter} →
        AppliedModelingLib.MeasurePerformativeModelOn Parameter Data domain → ↑domain → ↑domain → ℝ

value:
fun {Parameter : Type u_1} {Data : Type u_2} [MeasurableSpace Data] {domain : Set Parameter}
    (model : AppliedModelingLib.MeasurePerformativeModelOn Parameter Data domain) (deployed evaluated : ↑domain) =>
  ∫ (datum : Data), model.loss datum evaluated ∂↑(model.dataLaw deployed)
\end{ReviewVerbatim}

\paragraph{Exact Lean library declaration}
\begin{ReviewVerbatim}
/-- Decoupled performative risk for a model whose parameter carrier is a set. -/
noncomputable def decoupledPerformativeRisk
    {Parameter Data : Type*} [MeasurableSpace Data] {domain : Set Parameter}
    (model : MeasurePerformativeModelOn Parameter Data domain)
    (deployed evaluated : domain) : ℝ :=
  ∫ datum, model.loss datum evaluated ∂(model.dataLaw deployed : Measure Data)
\end{ReviewVerbatim}

\paragraph{Direct retained dependencies}
\begin{itemize}\raggedright
\item \hyperlink{library-prerequisite-e14cc5e8b1e1460e}{Applied\allowbreak{}Modeling\allowbreak{}Lib.\allowbreak{}Measure\allowbreak{}Performative\allowbreak{}Model\allowbreak{}On}
\end{itemize}
\paragraph{Recorded source-to-library screening}
\textbf{Verdict:} \texttt{matches}
\\
{\raggedright
\textbf{Reason:} Compared definition\_\allowbreak{}2\_\allowbreak{}3\_\allowbreak{}performative\_\allowbreak{}stability\_\allowbreak{}and\_\allowbreak{}decoupled\_\allowbreak{}risk at cited publication:\allowbreak{}219-\allowbreak{}228 with the expanded integral and every listed model, risk, optimality, stability, and theorem-\allowbreak{}use display.\allowbreak{} The deployed argument selects dataLaw, while the evaluated argument selects the loss parameter, giving DPR(theta,theta-\allowbreak{}prime)=E under D(theta) of loss(Z,theta-\allowbreak{}prime).\allowbreak{} The model ensures integrability for these two independently chosen feasible arguments.\allowbreak{} The displayed diagonal call gives PR;\allowbreak{} the stability calls keep the law argument fixed;\allowbreak{} and optimality compares diagonals.\allowbreak{} All these uses preserve the source argument order, candidate quantifier, and feasible parameter set.\allowbreak{}
\par}
\reviewmatch{Matches source input}
\noindent\textbf{Reviewer annotation}\par
\reviewerbox
\clearpage
\hypertarget{library-prerequisite-8eef8d8ba05e7ebe}{}
\subsection*{\small\ttfamily\raggedright Applied\allowbreak{}Modeling\allowbreak{}Lib.\allowbreak{}Measure\allowbreak{}Performative\allowbreak{}Model\allowbreak{}On.\allowbreak{}performative\allowbreak{}Risk}
\textbf{Source locator:} {\footnotesize\raggedright cited publication, lines 161-\allowbreak{}172\par}
\paragraph{Verbatim paper-source connection}
\begin{ReviewVerbatim}
Definition 2.1 (performative optimality and risk). A classi-
model parameters to θ in some underlying causal graph.
                                                                  fier fθPO is performatively optimal if the following relation-
Importantly, this mapping D(·) remains fixed and does not
                                                                  ship holds:
change over time or by intervention: deploying the same
classifier at two different points in time must induce the                       θPO = arg min      E      `(Z; θ).
same distribution over observations Z. While causal in-                                    θ     Z∼D(θ)

ference focuses on estimating properties of interventional                          def
distributions such as treatment effects (Pearl, 2009; Imbens      We define PR(θ) = EZ∼D(θ) `(Z; θ) as the performative
& Rubin, 2015), our focus is on a new stability notion and        risk; then, θPO = arg minθ PR(θ).
\end{ReviewVerbatim}



\paragraph{Lean-expanded library semantic target}
\begin{ReviewVerbatim}
type:
{Parameter : Type u_1} →
  {Data : Type u_2} →
    [inst : MeasurableSpace Data] →
      {domain : Set Parameter} → AppliedModelingLib.MeasurePerformativeModelOn Parameter Data domain → ↑domain → ℝ

value:
fun {Parameter : Type u_1} {Data : Type u_2} [MeasurableSpace Data] {domain : Set Parameter}
    (model : AppliedModelingLib.MeasurePerformativeModelOn Parameter Data domain) (parameter : ↑domain) =>
  model.decoupledPerformativeRisk parameter parameter
\end{ReviewVerbatim}

\paragraph{Exact Lean library declaration}
\begin{ReviewVerbatim}
/-- Performative risk for a model whose parameter carrier is a set. -/
noncomputable def performativeRisk
    {Parameter Data : Type*} [MeasurableSpace Data] {domain : Set Parameter}
    (model : MeasurePerformativeModelOn Parameter Data domain)
    (parameter : domain) : ℝ :=
  model.decoupledPerformativeRisk parameter parameter
\end{ReviewVerbatim}

\paragraph{Direct retained dependencies}
\begin{itemize}\raggedright
\item \hyperlink{library-prerequisite-e14cc5e8b1e1460e}{Applied\allowbreak{}Modeling\allowbreak{}Lib.\allowbreak{}Measure\allowbreak{}Performative\allowbreak{}Model\allowbreak{}On}
\item \hyperlink{library-prerequisite-25221c4f931beacb}{Applied\allowbreak{}Modeling\allowbreak{}Lib.\allowbreak{}Measure\allowbreak{}Performative\allowbreak{}Model\allowbreak{}On.\allowbreak{}decoupled\allowbreak{}Performative\allowbreak{}Risk}
\end{itemize}
\paragraph{Recorded source-to-library screening}
\textbf{Verdict:} \texttt{matches}
\\
{\raggedright
\textbf{Reason:} Compared definition\_\allowbreak{}2\_\allowbreak{}1\_\allowbreak{}performative\_\allowbreak{}optimality\_\allowbreak{}and\_\allowbreak{}risk at cited publication:\allowbreak{}161-\allowbreak{}172 with the expanded definition and every listed decoupled-\allowbreak{}risk, model, optimality, and theorem-\allowbreak{}use display.\allowbreak{} The same feasible parameter is passed as both arguments of decoupledPerformativeRisk, so PR(theta) is exactly E under D(theta) of loss(Z,theta).\allowbreak{} In optimality comparisons, the induced law and evaluated parameter therefore change together, preserving the source objective.\allowbreak{} The model supplies the needed diagonal integrability;\allowbreak{} neither this definition nor the displayed use introduces a unique minimizer or evaluates a parameter outside Theta.\allowbreak{}
\par}
\reviewmatch{Matches source input}
\noindent\textbf{Reviewer annotation}\par
\reviewerbox
\clearpage
\hypertarget{library-prerequisite-555f831d5bfa2a27}{}
\subsection*{\small\ttfamily\raggedright Applied\allowbreak{}Modeling\allowbreak{}Lib.\allowbreak{}Measure\allowbreak{}Performative\allowbreak{}Model\allowbreak{}On.\allowbreak{}Is\allowbreak{}Performatively\allowbreak{}Optimal}
\textbf{Source locator:} {\footnotesize\raggedright cited publication, lines 161-\allowbreak{}172\par}
\paragraph{Verbatim paper-source connection}
\begin{ReviewVerbatim}
Definition 2.1 (performative optimality and risk). A classi-
model parameters to θ in some underlying causal graph.
                                                                  fier fθPO is performatively optimal if the following relation-
Importantly, this mapping D(·) remains fixed and does not
                                                                  ship holds:
change over time or by intervention: deploying the same
classifier at two different points in time must induce the                       θPO = arg min      E      `(Z; θ).
same distribution over observations Z. While causal in-                                    θ     Z∼D(θ)

ference focuses on estimating properties of interventional                          def
distributions such as treatment effects (Pearl, 2009; Imbens      We define PR(θ) = EZ∼D(θ) `(Z; θ) as the performative
& Rubin, 2015), our focus is on a new stability notion and        risk; then, θPO = arg minθ PR(θ).
\end{ReviewVerbatim}



\paragraph{Lean-expanded library semantic target}
\begin{ReviewVerbatim}
type:
{Parameter : Type u_1} →
  {Data : Type u_2} →
    [inst : MeasurableSpace Data] →
      {domain : Set Parameter} → AppliedModelingLib.MeasurePerformativeModelOn Parameter Data domain → ↑domain → Prop

value:
fun {Parameter : Type u_1} {Data : Type u_2} [MeasurableSpace Data] {domain : Set Parameter}
    (model : AppliedModelingLib.MeasurePerformativeModelOn Parameter Data domain) (optimal : ↑domain) =>
  ∀ (candidate : ↑domain), model.performativeRisk optimal ≤ model.performativeRisk candidate
\end{ReviewVerbatim}

\paragraph{Exact Lean library declaration}
\begin{ReviewVerbatim}
/-- Performative optimality for a model defined only on its parameter domain. -/
def IsPerformativelyOptimal
    {Parameter Data : Type*} [MeasurableSpace Data] {domain : Set Parameter}
    (model : MeasurePerformativeModelOn Parameter Data domain)
    (optimal : domain) : Prop :=
  ∀ candidate, model.performativeRisk optimal ≤ model.performativeRisk candidate
\end{ReviewVerbatim}

\paragraph{Direct retained dependencies}
\begin{itemize}\raggedright
\item \hyperlink{library-prerequisite-e14cc5e8b1e1460e}{Applied\allowbreak{}Modeling\allowbreak{}Lib.\allowbreak{}Measure\allowbreak{}Performative\allowbreak{}Model\allowbreak{}On}
\item \hyperlink{library-prerequisite-8eef8d8ba05e7ebe}{Applied\allowbreak{}Modeling\allowbreak{}Lib.\allowbreak{}Measure\allowbreak{}Performative\allowbreak{}Model\allowbreak{}On.\allowbreak{}performative\allowbreak{}Risk}
\end{itemize}
\paragraph{Recorded source-to-library screening}
\textbf{Verdict:} \texttt{matches}
\\
{\raggedright
\textbf{Reason:} Compared definition\_\allowbreak{}2\_\allowbreak{}1\_\allowbreak{}performative\_\allowbreak{}optimality\_\allowbreak{}and\_\allowbreak{}risk at cited publication:\allowbreak{}161-\allowbreak{}172 with the expanded predicate, both risk definitions, the model fields, and the listed theorem use.\allowbreak{} Expanding the two risk definitions gives E under D(optimal) of loss(Z,optimal) <= E under D(candidate) of loss(Z,candidate) for every candidate in Theta.\allowbreak{} Both the law and evaluation argument vary together on the candidate side, exactly the source argmin of PR.\allowbreak{} The supplied optimal member satisfies all comparisons;\allowbreak{} neither this predicate nor its paper use asserts uniqueness or constructs an optimum.\allowbreak{} The model covers the integrability of every diagonal expectation, and all minimization arguments stay inside the source domain.\allowbreak{}
\par}
\reviewmatch{Matches source input}
\noindent\textbf{Reviewer annotation}\par
\reviewerbox
\clearpage
\hypertarget{library-prerequisite-4af6c70f954a80a0}{}
\subsection*{\small\ttfamily\raggedright Applied\allowbreak{}Modeling\allowbreak{}Lib.\allowbreak{}Measure\allowbreak{}Performative\allowbreak{}Model\allowbreak{}On.\allowbreak{}Is\allowbreak{}Performatively\allowbreak{}Stable}
\textbf{Source locator:} {\footnotesize\raggedright cited publication, lines 219-\allowbreak{}228\par}
\paragraph{Verbatim paper-source connection}
\begin{ReviewVerbatim}
Definition 2.3 (performative stability and decoupled risk).              Points
A classifier fθPS is performatively stable if the following           Having introduced our framework for performative predic-
relationship holds:                                                  tion, we now address some of the basic questions that arise
             θPS = arg min          E        `(Z; θ).                in this setting and examine the behavior of common ma-
                        θ        Z∼D(θPS )                           chine learning practices, such as retraining, through the lens
                                                                     of performativity. We begin by analyzing the behavior of
                                  def
We define DPR(θ, θ0 )     =    EZ∼D(θ) `(Z; θ0 ) as                   these procedures when they operate at a population level
the decoupled performative risk; then, θPS      =                    and then extend our analysis to finite samples.
arg minθ DPR(θPS , θ).
\end{ReviewVerbatim}



\paragraph{Lean-expanded library semantic target}
\begin{ReviewVerbatim}
type:
{Parameter : Type u_1} →
  {Data : Type u_2} →
    [inst : MeasurableSpace Data] →
      {domain : Set Parameter} → AppliedModelingLib.MeasurePerformativeModelOn Parameter Data domain → ↑domain → Prop

value:
fun {Parameter : Type u_1} {Data : Type u_2} [MeasurableSpace Data] {domain : Set Parameter}
    (model : AppliedModelingLib.MeasurePerformativeModelOn Parameter Data domain) (stable : ↑domain) =>
  ∀ (candidate : ↑domain),
    model.decoupledPerformativeRisk stable stable ≤ model.decoupledPerformativeRisk stable candidate
\end{ReviewVerbatim}

\paragraph{Exact Lean library declaration}
\begin{ReviewVerbatim}
/-- Performative stability for a model defined only on its parameter domain. -/
def IsPerformativelyStable
    {Parameter Data : Type*} [MeasurableSpace Data] {domain : Set Parameter}
    (model : MeasurePerformativeModelOn Parameter Data domain)
    (stable : domain) : Prop :=
  ∀ candidate, model.decoupledPerformativeRisk stable stable ≤
    model.decoupledPerformativeRisk stable candidate
\end{ReviewVerbatim}

\paragraph{Direct retained dependencies}
\begin{itemize}\raggedright
\item \hyperlink{library-prerequisite-e14cc5e8b1e1460e}{Applied\allowbreak{}Modeling\allowbreak{}Lib.\allowbreak{}Measure\allowbreak{}Performative\allowbreak{}Model\allowbreak{}On}
\item \hyperlink{library-prerequisite-25221c4f931beacb}{Applied\allowbreak{}Modeling\allowbreak{}Lib.\allowbreak{}Measure\allowbreak{}Performative\allowbreak{}Model\allowbreak{}On.\allowbreak{}decoupled\allowbreak{}Performative\allowbreak{}Risk}
\end{itemize}
\paragraph{Recorded source-to-library screening}
\textbf{Verdict:} \texttt{matches}
\\
{\raggedright
\textbf{Reason:} Compared definition\_\allowbreak{}2\_\allowbreak{}3\_\allowbreak{}performative\_\allowbreak{}stability\_\allowbreak{}and\_\allowbreak{}decoupled\_\allowbreak{}risk at cited publication:\allowbreak{}219-\allowbreak{}228 with the expanded predicate and all listed model, decoupled-\allowbreak{}risk, and theorem-\allowbreak{}use displays.\allowbreak{} The resulting inequality is E under D(stable) of loss(Z,stable) <= E under D(stable) of loss(Z,candidate) for every candidate in Theta.\allowbreak{} The induced distribution is fixed on both sides, and only the evaluated classifier changes, exactly the source argmin\_\allowbreak{}theta DPR(thetaPS,theta).\allowbreak{} The supplied stable domain member need not be unique, and no existence assertion is substituted for the predicate.\allowbreak{} Universal cross-\allowbreak{}law integrability in the model covers these fixed-\allowbreak{}law real expectations.\allowbreak{}
\par}
\reviewmatch{Matches source input}
\noindent\textbf{Reviewer annotation}\par
\reviewerbox
\clearpage
\hypertarget{library-prerequisite-eb9ab8e4235e1d28}{}
\subsection*{\small\ttfamily\raggedright Applied\allowbreak{}Modeling\allowbreak{}Lib.\allowbreak{}Measure\allowbreak{}Performative\allowbreak{}Model\allowbreak{}On.\allowbreak{}Is\allowbreak{}Pointwise\allowbreak{}Gradient\allowbreak{}Strongly\allowbreak{}Convex}
\textbf{Source locator:} {\footnotesize\raggedright cited publication, lines 270-\allowbreak{}282\par}
\paragraph{Verbatim paper-source connection}
\begin{ReviewVerbatim}
(A1) (joint smoothness) A loss function `(z; θ) is β-jointly       The main message of this theorem is that in performative
     smooth if ∀θ, θ0 ∈ Θ and z, z 0 ∈ Z:                          prediction, if the loss function is sufficiently “nice” and the
                                                                   distribution map is sufficiently (in)sensitive, then one need
           krθ `(z; θ) − rθ `(z; θ0 )k2 6 β kθ − θ0 k2 ,
                                                                   only retrain a classifier a small number of times before it
           krθ `(z; θ) − rθ `(z 0 ; θ)k2 6 β kz − z 0 k2 .         converges to a unique stable point. Here, we provide a proof
                                                                   sketch illustrating the main ideas behind the theorem and
(A2) (strong convexity) A loss function `(z; θ) is γ-
                                                                   defer the full proof to Appendix E.1.
     strongly convex if ∀ θ, θ0 ∈ Θ and z ∈ Z:
                                                γ         2        Proof Sketch. Fix θ, θ0 ∈ Θ. Let f (ϕ) = EZ∼D(θ) `(Z; ϕ)
      `(z; θ) > `(z; θ0 )+rθ `(z; θ0 )> (θ−θ0 )+ kθ − θ0 k2 .
                                                2                  and f 0 (ϕ) = EZ∼D(θ0 ) `(Z; ϕ). By applying standard prop-
\end{ReviewVerbatim}



\paragraph{Lean-expanded library semantic target}
\begin{ReviewVerbatim}
type:
{Parameter : Type u_1} →
  {Data : Type u_2} →
    [inst : MeasurableSpace Data] →
      [inst_1 : NormedAddCommGroup Parameter] →
        [InnerProductSpace ℝ Parameter] →
          {domain : Set Parameter} →
            AppliedModelingLib.MeasurePerformativeModelOn Parameter Data domain →
              (Data → ↑domain → Parameter) → ℝ → Prop

value:
fun {Parameter : Type u_1} {Data : Type u_2} [MeasurableSpace Data] [NormedAddCommGroup Parameter]
    [InnerProductSpace ℝ Parameter] {domain : Set Parameter}
    (model : AppliedModelingLib.MeasurePerformativeModelOn Parameter Data domain)
    (gradient : Data → ↑domain → Parameter) (modulus : ℝ) =>
  (∀ (datum : Data) (center : ↑domain),
      HasFDerivWithinAt
        (fun (parameter : Parameter) =>
          if hparameter : parameter ∈ domain then model.loss datum ⟨parameter, hparameter⟩ else 0)
        ((innerSL ℝ) (gradient datum center)) domain ↑center) ∧
    ∀ (datum : Data) (center candidate : ↑domain),
      model.loss datum candidate ≥
        model.loss datum center + inner ℝ (gradient datum center) (↑candidate - ↑center) +
          modulus / 2 * ‖↑candidate - ↑center‖ ^ 2
\end{ReviewVerbatim}

\paragraph{Exact Lean library declaration}
\begin{ReviewVerbatim}
/-- A2 uses the actual loss derivative on the feasible parameter domain.
The zero extension is read only along the domain; no off-domain loss value or
regularity is required. The lower-model inequality uses this same derivative. -/
def IsPointwiseGradientStronglyConvex
    {Parameter Data : Type*} [MeasurableSpace Data]
    [NormedAddCommGroup Parameter] [InnerProductSpace ℝ Parameter]
    {domain : Set Parameter} (model : MeasurePerformativeModelOn Parameter Data domain)
    (gradient : Data → domain → Parameter) (modulus : ℝ) : Prop := by
  classical
  exact
    (∀ datum (center : domain),
      HasFDerivWithinAt
        (fun parameter => if hparameter : parameter ∈ domain then
          model.loss datum ⟨parameter, hparameter⟩ else 0)
        (innerSL ℝ (gradient datum center)) domain (center : Parameter)) ∧
    ∀ datum center candidate,
      model.loss datum candidate ≥
        model.loss datum center +
          ⟪gradient datum center, (candidate : Parameter) - (center : Parameter)⟫_ℝ +
            modulus / 2 * ‖(candidate : Parameter) - (center : Parameter)‖ ^ 2
\end{ReviewVerbatim}

\paragraph{Direct retained dependencies}
\begin{itemize}\raggedright
\item \hyperlink{library-prerequisite-e14cc5e8b1e1460e}{Applied\allowbreak{}Modeling\allowbreak{}Lib.\allowbreak{}Measure\allowbreak{}Performative\allowbreak{}Model\allowbreak{}On}
\end{itemize}
\paragraph{Recorded source-to-library screening}
\textbf{Verdict:} \texttt{matches}
\\
{\raggedright
\textbf{Reason:} Compared the A2 formula in assumptions\_\allowbreak{}a1\_\allowbreak{}a2\_\allowbreak{}joint\_\allowbreak{}smoothness\_\allowbreak{}and\_\allowbreak{}strong\_\allowbreak{}convexity at cited publication:\allowbreak{}270-\allowbreak{}282 with the full expanded conjunctive predicate, its model prerequisite, and the theorem-\allowbreak{}4.\allowbreak{}3 use.\allowbreak{} For every datum and feasible center, HasFDerivWithinAt identifies the derivative of the loss on Theta with the linear functional v -\allowbreak{}> inner(gradient datum center,v).\allowbreak{} The subsequent inequality uses exactly this same vector, every feasible candidate, the displacement candidate-\allowbreak{}center, and modulus/\allowbreak{}2 times its squared norm.\allowbreak{} This gives the actual first-\allowbreak{}order quadratic lower model in source A2 rather than an unrestricted supporting-\allowbreak{}vector condition.\allowbreak{} The zero extension agrees with the loss at every point used by the within-\allowbreak{}domain derivative, so it imposes no differentiability outside Theta or discontinuity condition across its boundary.\allowbreak{} On lower-\allowbreak{}dimensional domains, any undetermined normal component has zero pairing with feasible displacements and does not change A2.\allowbreak{} The paper use requires this whole predicate and a strictly positive modulus on its closed convex domain;\allowbreak{} A1 smoothness is not inadvertently added.\allowbreak{} In particular, theta\textasciicircum{}2+|theta| on [-\allowbreak{}1,1] cannot satisfy the derivative conjunct at its interior point 0.\allowbreak{}
\par}
\reviewmatch{Matches source input}
\noindent\textbf{Reviewer annotation}\par
\reviewerbox
\clearpage
\hypertarget{library-prerequisite-0e1e192eee739a4b}{}
\subsection*{\small\ttfamily\raggedright Applied\allowbreak{}Modeling\allowbreak{}Lib.\allowbreak{}Measure\allowbreak{}Performative\allowbreak{}Model\allowbreak{}On.\allowbreak{}Is\allowbreak{}Wasserstein\allowbreak{}Sensitive}
\textbf{Source locator:} {\footnotesize\raggedright cited publication, lines 242-\allowbreak{}249\par}
\paragraph{Verbatim paper-source connection}
\begin{ReviewVerbatim}
formatively optimal classifiers need not be performatively            Definition 3.1 (ε-sensitivity). We say that a distribution
stable and performatively stable classifiers need not be per-         map D(·) is ε-sensitive if for all θ, θ0 ∈ Θ:
formatively optimal. We illustrate this point in the context
                                                                                   W1 D(θ), D(θ0 ) 6 εkθ − θ0 k2 ,
                                                                                     �
of our previous biased coin toss example.
Example 2.2 (continued). Consider again our model of
a biased coin toss. In order for a classifier fθ to be perfor-        where W1 denotes the earth mover’s distance.
\end{ReviewVerbatim}



\paragraph{Lean-expanded library semantic target}
\begin{ReviewVerbatim}
type:
{Parameter : Type u_1} →
  {Data : Type u_2} →
    [MetricSpace Parameter] →
      [inst : MeasurableSpace Data] →
        [MetricSpace Data] →
          {domain : Set Parameter} → AppliedModelingLib.MeasurePerformativeModelOn Parameter Data domain → ℝ → Prop

value:
fun {Parameter : Type u_1} {Data : Type u_2} [MetricSpace Parameter] [MeasurableSpace Data] [MetricSpace Data]
    {domain : Set Parameter} (model : AppliedModelingLib.MeasurePerformativeModelOn Parameter Data domain)
    (sensitivity : ℝ) =>
  ∀ (first second : ↑domain),
    ∃ (hfinite :
      (AppliedModelingLib.ProbabilityCoupling.expectedDistanceCosts (model.dataLaw first)
          (model.dataLaw second)).Nonempty),
      AppliedModelingLib.ProbabilityCoupling.wassersteinOneReal (model.dataLaw first) (model.dataLaw second) hfinite ≤
        sensitivity * dist ↑first ↑second
\end{ReviewVerbatim}

\paragraph{Exact Lean library declaration}
\begin{ReviewVerbatim}
/-- Wasserstein sensitivity for a model defined only on its parameter domain. -/
def IsWassersteinSensitive
    {Parameter Data : Type*} [MetricSpace Parameter] [MeasurableSpace Data]
    [MetricSpace Data] {domain : Set Parameter}
    (model : MeasurePerformativeModelOn Parameter Data domain)
    (sensitivity : ℝ) : Prop :=
  ∀ first second,
    ∃ hfinite : (ProbabilityCoupling.expectedDistanceCosts
      (model.dataLaw first) (model.dataLaw second)).Nonempty,
      ProbabilityCoupling.wassersteinOneReal
          (model.dataLaw first) (model.dataLaw second) hfinite ≤
        sensitivity * dist (first : Parameter) (second : Parameter)
\end{ReviewVerbatim}

\paragraph{Direct retained dependencies}
\begin{itemize}\raggedright
\item \hyperlink{library-prerequisite-e14cc5e8b1e1460e}{Applied\allowbreak{}Modeling\allowbreak{}Lib.\allowbreak{}Measure\allowbreak{}Performative\allowbreak{}Model\allowbreak{}On}
\item \hyperlink{governing-declaration-0507d24247d6815a}{Applied\allowbreak{}Modeling\allowbreak{}Lib.\allowbreak{}Probability\allowbreak{}Coupling.\allowbreak{}expected\allowbreak{}Distance\allowbreak{}Costs}
\item \hyperlink{governing-declaration-7fdf7bcf28dbe52b}{Applied\allowbreak{}Modeling\allowbreak{}Lib.\allowbreak{}Probability\allowbreak{}Coupling.\allowbreak{}wasserstein\allowbreak{}One\allowbreak{}Real}
\end{itemize}
\paragraph{Recorded source-to-library screening}
\textbf{Verdict:} \texttt{matches}
\\
{\raggedright
\textbf{Reason:} Compared definition\_\allowbreak{}3\_\allowbreak{}1\_\allowbreak{}epsilon\_\allowbreak{}sensitivity at cited publication:\allowbreak{}242-\allowbreak{}249 with the expanded predicate and every listed coupling/\allowbreak{}model/\allowbreak{}use display.\allowbreak{} The joint probability law has exactly the two prescribed marginals.\allowbreak{} expectedDistanceCosts consists of its finite expected distances, and wassersteinOneReal takes their infimum with an explicit nonempty finite-\allowbreak{}cost witness.\allowbreak{} The target is therefore the source earth-\allowbreak{}mover bound W1(D(first),D(second)) <= epsilon dist(first,second) for all feasible parameter pairs.\allowbreak{} A finite source bound corresponds to this finite-\allowbreak{}cost infimum;\allowbreak{} infinite-\allowbreak{}cost couplings cannot lower it.\allowbreak{} In the inner-\allowbreak{}product paper use, parameter distance equals the norm of the difference, and the explicit sensitivity premise permits precisely nonnegative epsilon, including 0.\allowbreak{} No attainment of the infimum, finite support, or off-\allowbreak{}domain distribution is required.\allowbreak{}
\par}
\reviewmatch{Matches source input}
\noindent\textbf{Reviewer annotation}\par
\reviewerbox
\clearpage
\hypertarget{library-prerequisite-d4de05e57b5947f0}{}
\subsection*{\small\ttfamily\raggedright Applied\allowbreak{}Modeling\allowbreak{}Lib.\allowbreak{}measure\allowbreak{}Decoupled\allowbreak{}Performative\allowbreak{}Risk}
\textbf{Source locator:} {\footnotesize\raggedright cited publication, lines 219-\allowbreak{}228\par}
\paragraph{Verbatim paper-source connection}
\begin{ReviewVerbatim}
Definition 2.3 (performative stability and decoupled risk).              Points
A classifier fθPS is performatively stable if the following           Having introduced our framework for performative predic-
relationship holds:                                                  tion, we now address some of the basic questions that arise
             θPS = arg min          E        `(Z; θ).                in this setting and examine the behavior of common ma-
                        θ        Z∼D(θPS )                           chine learning practices, such as retraining, through the lens
                                                                     of performativity. We begin by analyzing the behavior of
                                  def
We define DPR(θ, θ0 )     =    EZ∼D(θ) `(Z; θ0 ) as                   these procedures when they operate at a population level
the decoupled performative risk; then, θPS      =                    and then extend our analysis to finite samples.
arg minθ DPR(θPS , θ).
\end{ReviewVerbatim}



\paragraph{Lean-expanded library semantic target}
\begin{ReviewVerbatim}
type:
{Parameter : Type u_1} →
  {Data : Type u_2} →
    [inst : MeasurableSpace Data] →
      AppliedModelingLib.MeasurePerformativeModel Parameter Data → Parameter → Parameter → ℝ

value:
fun {Parameter : Type u_1} {Data : Type u_2} [MeasurableSpace Data]
    (model : AppliedModelingLib.MeasurePerformativeModel Parameter Data)
    (distributionParameter evaluatedParameter : Parameter) =>
  ∫ (datum : Data), model.loss datum evaluatedParameter ∂↑(model.dataLaw distributionParameter)
\end{ReviewVerbatim}

\paragraph{Exact Lean library declaration}
\begin{ReviewVerbatim}
/-- General-distribution frozen population risk. -/
noncomputable def measureDecoupledPerformativeRisk {Parameter Data : Type*} [MeasurableSpace Data]
    (model : MeasurePerformativeModel Parameter Data)
    (distributionParameter evaluatedParameter : Parameter) : ℝ :=
  ∫ datum, model.loss datum evaluatedParameter ∂(model.dataLaw distributionParameter : Measure Data)
\end{ReviewVerbatim}

\paragraph{Direct retained dependencies}
\begin{itemize}\raggedright
\item \hyperlink{library-prerequisite-96553d4fbb82f1ac}{Applied\allowbreak{}Modeling\allowbreak{}Lib.\allowbreak{}Measure\allowbreak{}Performative\allowbreak{}Model}
\end{itemize}
\paragraph{Recorded source-to-library screening}
\textbf{Verdict:} \texttt{matches}
\\
{\raggedright
\textbf{Reason:} Definition 2.\allowbreak{}3 defines DPR(θ, θ′) = E\_\allowbreak{}\{Z∼D(θ)\} {\fontspec{Latin Modern Math}ℓ}(Z;\allowbreak{} θ′).\allowbreak{} The expanded Lean declaration is exactly the corresponding integral of model.\allowbreak{}loss z evaluatedParameter under model.\allowbreak{}dataLaw distributionParameter, preserving the distribution/\allowbreak{}evaluation parameter order;\allowbreak{} the model integrability condition only makes the expectation real-\allowbreak{}valued and does not change the definition.\allowbreak{}
\par}
\reviewmatch{Matches source input}
\noindent\textbf{Reviewer annotation}\par
\reviewerbox
\clearpage
\hypertarget{paper-prerequisite-section-5ef5ef0364b6939c}{}
\section*{Paper-specific semantic prerequisites}
\hypertarget{paper-prerequisite-48fa8b1c1ddbaa9c}{}
\subsection*{\small\ttfamily\raggedright PZMH20\allowbreak{}Performative\allowbreak{}Prediction.\allowbreak{}Is\allowbreak{}Measure\allowbreak{}Performatively\allowbreak{}Optimal}
\noindent\textbf{Paper-local declaration:}\par
{\footnotesize\ttfamily\raggedright PZMH20\allowbreak{}Performative\allowbreak{}Prediction.\allowbreak{}Is\allowbreak{}Measure\allowbreak{}Performatively\allowbreak{}Optimal\par}
\textbf{Source locator:} {\footnotesize\raggedright cited publication, lines 161-\allowbreak{}172\par}
\paragraph{Verbatim paper-source connection}
\begin{ReviewVerbatim}
Definition 2.1 (performative optimality and risk). A classi-
model parameters to θ in some underlying causal graph.
                                                                  fier fθPO is performatively optimal if the following relation-
Importantly, this mapping D(·) remains fixed and does not
                                                                  ship holds:
change over time or by intervention: deploying the same
classifier at two different points in time must induce the                       θPO = arg min      E      `(Z; θ).
same distribution over observations Z. While causal in-                                    θ     Z∼D(θ)

ference focuses on estimating properties of interventional                          def
distributions such as treatment effects (Pearl, 2009; Imbens      We define PR(θ) = EZ∼D(θ) `(Z; θ) as the performative
& Rubin, 2015), our focus is on a new stability notion and        risk; then, θPO = arg minθ PR(θ).
\end{ReviewVerbatim}



\paragraph{Lean-expanded paper semantic target}
\begin{ReviewVerbatim}
type:
{Parameter : Type u_1} →
  {Data : Type u_2} →
    [inst : MeasurableSpace Data] → AppliedModelingLib.MeasurePerformativeModel Parameter Data → Parameter → Prop

value:
fun {Parameter : Type u_1} {Data : Type u_2} [MeasurableSpace Data]
    (model : AppliedModelingLib.MeasurePerformativeModel Parameter Data) (parameter : Parameter) =>
  AppliedModelingLib.IsMeasurePerformativelyOptimal model parameter
\end{ReviewVerbatim}

\paragraph{Direct retained dependencies}
\begin{itemize}\raggedright
\item \hyperlink{governing-declaration-bf41b084ebc9aa18}{Applied\allowbreak{}Modeling\allowbreak{}Lib.\allowbreak{}Is\allowbreak{}Measure\allowbreak{}Performatively\allowbreak{}Optimal}
\item \hyperlink{library-prerequisite-96553d4fbb82f1ac}{Applied\allowbreak{}Modeling\allowbreak{}Lib.\allowbreak{}Measure\allowbreak{}Performative\allowbreak{}Model}
\end{itemize}
\paragraph{Recorded source-to-declaration screening}
\textbf{Verdict:} \texttt{matches}
\\
{\raggedright
\textbf{Reason:} Definition 2.\allowbreak{}1's θPO = argminθ PR(θ) is exactly the expanded universal minimizer condition for measurePerformativeRisk.\allowbreak{} The explicit probability law and loss integrability merely realize the displayed expectation.\allowbreak{}
\par}
\reviewmatch{Matches source input}
\noindent\textbf{Reviewer annotation}\par
\reviewerbox
\clearpage
\hypertarget{paper-prerequisite-28cbc41dc9d12634}{}
\subsection*{\small\ttfamily\raggedright PZMH20\allowbreak{}Performative\allowbreak{}Prediction.\allowbreak{}Is\allowbreak{}Measure\allowbreak{}RRM\allowbreak{}On}
\noindent\textbf{Paper-local declaration:}\par
{\footnotesize\ttfamily\raggedright PZMH20\allowbreak{}Performative\allowbreak{}Prediction.\allowbreak{}Is\allowbreak{}Measure\allowbreak{}RRM\allowbreak{}On\par}
\textbf{Source locator:} {\footnotesize\raggedright cited publication, lines 296-\allowbreak{}303\par}
\paragraph{Verbatim paper-source connection}
\begin{ReviewVerbatim}
Definition 3.3 (RRM). Repeated risk minimization (RRM)              well as the first-order conditions of optimality for convex
refers to the procedure where, starting from an initial model      functions, a short calculation reveals that
fθ0 we perform the following sequence of updates for t > 0:
                                                                   (G(θ) − G(θ0 ))> rf (G(θ0 )) > −εβkG(θ) − G(θ0 )k2 kθ − θ0 k2 .
                        def
        θt+1 = G(θt ) = arg min          E      `(Z; θ).           Claim (a) then follows by combining the previous two inequalities
                              θ∈Θ    Z∼D(θt )
                                                                   and rearranging. Intuitively, strong convexity forces the iterates
\end{ReviewVerbatim}



\paragraph{Lean-expanded paper semantic target}
\begin{ReviewVerbatim}
type:
{Parameter : Type u_1} →
  {Data : Type u_2} →
    [inst : MeasurableSpace Data] →
      AppliedModelingLib.MeasurePerformativeModel Parameter Data → Set Parameter → (Parameter → Parameter) → Prop

value:
fun {Parameter : Type u_1} {Data : Type u_2} [MeasurableSpace Data]
    (model : AppliedModelingLib.MeasurePerformativeModel Parameter Data) (domain : Set Parameter)
    (update : Parameter → Parameter) =>
  AppliedModelingLib.IsMeasureRepeatedRiskMinimizationOn model domain update
\end{ReviewVerbatim}

\paragraph{Direct retained dependencies}
\begin{itemize}\raggedright
\item \hyperlink{governing-declaration-6ae86e5703739171}{Applied\allowbreak{}Modeling\allowbreak{}Lib.\allowbreak{}Is\allowbreak{}Measure\allowbreak{}Repeated\allowbreak{}Risk\allowbreak{}Minimization\allowbreak{}On}
\item \hyperlink{library-prerequisite-96553d4fbb82f1ac}{Applied\allowbreak{}Modeling\allowbreak{}Lib.\allowbreak{}Measure\allowbreak{}Performative\allowbreak{}Model}
\end{itemize}
\paragraph{Recorded source-to-declaration screening}
\textbf{Verdict:} \texttt{matches}
\\
{\raggedright
\textbf{Reason:} Definition 3.\allowbreak{}3's G(θt) = argminθ∈Θ E\_\allowbreak{}\{D(θt)\} {\fontspec{Latin Modern Math}ℓ}(Z;\allowbreak{} θ) matches the selected frozen-\allowbreak{}risk minimizer over domain, including the requirement that the update remains in the parameter domain.\allowbreak{}
\par}
\reviewmatch{Matches source input}
\noindent\textbf{Reviewer annotation}\par
\reviewerbox
\clearpage
\hypertarget{paper-prerequisite-f6b3bb0ef1cf7d32}{}
\subsection*{\small\ttfamily\raggedright PZMH20\allowbreak{}Performative\allowbreak{}Prediction.\allowbreak{}Is\allowbreak{}Measure\allowbreak{}Stable}
\noindent\textbf{Paper-local declaration:}\par
{\footnotesize\ttfamily\raggedright PZMH20\allowbreak{}Performative\allowbreak{}Prediction.\allowbreak{}Is\allowbreak{}Measure\allowbreak{}Stable\par}
\textbf{Source locator:} {\footnotesize\raggedright cited publication, lines 219-\allowbreak{}228\par}
\paragraph{Verbatim paper-source connection}
\begin{ReviewVerbatim}
Definition 2.3 (performative stability and decoupled risk).              Points
A classifier fθPS is performatively stable if the following           Having introduced our framework for performative predic-
relationship holds:                                                  tion, we now address some of the basic questions that arise
             θPS = arg min          E        `(Z; θ).                in this setting and examine the behavior of common ma-
                        θ        Z∼D(θPS )                           chine learning practices, such as retraining, through the lens
                                                                     of performativity. We begin by analyzing the behavior of
                                  def
We define DPR(θ, θ0 )     =    EZ∼D(θ) `(Z; θ0 ) as                   these procedures when they operate at a population level
the decoupled performative risk; then, θPS      =                    and then extend our analysis to finite samples.
arg minθ DPR(θPS , θ).
\end{ReviewVerbatim}



\paragraph{Lean-expanded paper semantic target}
\begin{ReviewVerbatim}
type:
{Parameter : Type u_1} →
  {Data : Type u_2} →
    [inst : MeasurableSpace Data] → AppliedModelingLib.MeasurePerformativeModel Parameter Data → Parameter → Prop

value:
fun {Parameter : Type u_1} {Data : Type u_2} [MeasurableSpace Data]
    (model : AppliedModelingLib.MeasurePerformativeModel Parameter Data) (parameter : Parameter) =>
  AppliedModelingLib.IsMeasurePerformativelyStable model parameter
\end{ReviewVerbatim}

\paragraph{Direct retained dependencies}
\begin{itemize}\raggedright
\item \hyperlink{governing-declaration-e97175bd0a346777}{Applied\allowbreak{}Modeling\allowbreak{}Lib.\allowbreak{}Is\allowbreak{}Measure\allowbreak{}Performatively\allowbreak{}Stable}
\item \hyperlink{library-prerequisite-96553d4fbb82f1ac}{Applied\allowbreak{}Modeling\allowbreak{}Lib.\allowbreak{}Measure\allowbreak{}Performative\allowbreak{}Model}
\end{itemize}
\paragraph{Recorded source-to-declaration screening}
\textbf{Verdict:} \texttt{matches}
\\
{\raggedright
\textbf{Reason:} Definition 2.\allowbreak{}3's θPS = argminθ DPR(θPS, θ) is the expanded condition that the frozen law at parameter has no lower-\allowbreak{}risk candidate.\allowbreak{}
\par}
\reviewmatch{Matches source input}
\noindent\textbf{Reviewer annotation}\par
\reviewerbox
\clearpage
\hypertarget{paper-prerequisite-2ea2f5769822c5c4}{}
\subsection*{\small\ttfamily\raggedright PZMH20\allowbreak{}Performative\allowbreak{}Prediction.\allowbreak{}Is\allowbreak{}Measure\allowbreak{}Wasserstein\allowbreak{}Sensitive}
\noindent\textbf{Paper-local declaration:}\par
{\footnotesize\ttfamily\raggedright PZMH20\allowbreak{}Performative\allowbreak{}Prediction.\allowbreak{}Is\allowbreak{}Measure\allowbreak{}Wasserstein\allowbreak{}Sensitive\par}
\textbf{Source locator:} {\footnotesize\raggedright cited publication, lines 242-\allowbreak{}249\par}
\paragraph{Verbatim paper-source connection}
\begin{ReviewVerbatim}
formatively optimal classifiers need not be performatively            Definition 3.1 (ε-sensitivity). We say that a distribution
stable and performatively stable classifiers need not be per-         map D(·) is ε-sensitive if for all θ, θ0 ∈ Θ:
formatively optimal. We illustrate this point in the context
                                                                                   W1 D(θ), D(θ0 ) 6 εkθ − θ0 k2 ,
                                                                                     �
of our previous biased coin toss example.
Example 2.2 (continued). Consider again our model of
a biased coin toss. In order for a classifier fθ to be perfor-        where W1 denotes the earth mover’s distance.
\end{ReviewVerbatim}



\paragraph{Lean-expanded paper semantic target}
\begin{ReviewVerbatim}
type:
{Parameter : Type u_1} →
  {Data : Type u_2} →
    [MetricSpace Parameter] →
      [inst : MeasurableSpace Data] →
        [MetricSpace Data] → AppliedModelingLib.MeasurePerformativeModel Parameter Data → ℝ → Prop

value:
fun {Parameter : Type u_1} {Data : Type u_2} [MetricSpace Parameter] [MeasurableSpace Data] [MetricSpace Data]
    (model : AppliedModelingLib.MeasurePerformativeModel Parameter Data) (sensitivity : ℝ) =>
  AppliedModelingLib.IsMeasureWassersteinSensitive model sensitivity
\end{ReviewVerbatim}

\paragraph{Direct retained dependencies}
\begin{itemize}\raggedright
\item \hyperlink{library-prerequisite-3c6dd4ad7ac116e3}{Applied\allowbreak{}Modeling\allowbreak{}Lib.\allowbreak{}Is\allowbreak{}Measure\allowbreak{}Wasserstein\allowbreak{}Sensitive}
\item \hyperlink{library-prerequisite-96553d4fbb82f1ac}{Applied\allowbreak{}Modeling\allowbreak{}Lib.\allowbreak{}Measure\allowbreak{}Performative\allowbreak{}Model}
\end{itemize}
\paragraph{Recorded source-to-declaration screening}
\textbf{Verdict:} \texttt{matches}
\\
{\raggedright
\textbf{Reason:} Definition 3.\allowbreak{}1's W₁(D(θ), D(θ′)) ≤ ε‖θ−θ′‖₂ is represented by a finite expected-\allowbreak{}distance coupling and its real W₁ infimum bounded by sensitivity * dist;\allowbreak{} finiteness is required for the source's finite displayed bound.\allowbreak{}
\par}
\reviewmatch{Matches source input}
\noindent\textbf{Reviewer annotation}\par
\reviewerbox
\clearpage
\hypertarget{paper-prerequisite-613020783cc01a15}{}
\subsection*{\small\ttfamily\raggedright PZMH20\allowbreak{}Performative\allowbreak{}Prediction.\allowbreak{}Theorem310\allowbreak{}Corrected\allowbreak{}RERM\allowbreak{}All\allowbreak{}Round\allowbreak{}Trajectory\allowbreak{}Data}
\noindent\textbf{Paper-local declaration:}\par
{\footnotesize\ttfamily\raggedright PZMH20\allowbreak{}Performative\allowbreak{}Prediction.\allowbreak{}Theorem310\allowbreak{}Corrected\allowbreak{}RERM\allowbreak{}All\allowbreak{}Round\allowbreak{}Trajectory\allowbreak{}Data\par}
\textbf{Source locator:} {\footnotesize\raggedright cited publication, lines 373-\allowbreak{}415\par}
\paragraph{Verbatim paper-source connection}
\begin{ReviewVerbatim}
Theorem 3.10. Suppose that the loss `(z; θ) is β-jointly
Theorem 3.8. Suppose that the loss `(z; θ) is β-jointly                       smooth (A1) and γ-strongly convex (A2), and that there exist
smooth (A1) and γ-strongly convex (A2). If the distribution                                                def R        α
                                        γ                                     α > 1, µ > 0 such that ξα,µ = Rm eµ|x| dD(θ) is finite
map D(·) is ε-sensitive with ε < (β+γ)(1+1.5ηβ) , then RGD
                                                                              ∀θ ∈ Θ. Let δ ∈ (0, 1) be a radius of convergence.     Con-
                     2
with step size η 6 β+γ satisfies the following:
                                                                                                                                    � t
                                                                                                                            1
                                                                              sider running RERM or RGD with nt = O (εδ)m log p
(a) kGgd (θ) − Ggd (θ0 )k2                                                    samples at time t.
                                                                                                                            γ
                                                                              (a) If the map D(·) is ε-sensitive with ε < 2β  , then with
                         βγ                                          0
      6       1−η           − ε(1.5ηβ 2 + β)             kθ − θ k2 .              probability 1 − p, RERM satisfies,
                        β+γ
                                                                                                                  log 1δ kθ0 − θPS k2
                                                                                                                     �
(b) The iterates θt of RGD converge to a unique performa-                           kθt − θPS k2 6 δ, for all t >                       .
    tively stable point θPS at a linear rate,                                                                           1 − 2εβ
                                                                                                                              γ

                                                                                                                                 γ
                                               �1                             (b) If the map D(·) is ε-sensitive with ε < (β+γ)(1+1.5ηβ)  ,
                                         log    δ kθ0 − θPS k2
     kθt − θPS k2 6 δ for t >                                             .       then with probability 1 − p, REGD with satisfies,
                                          βγ          2
                                 η       β+γ − ε(1.5ηβ + β)
                                                                                                                    log 1δ kθ0 − θPS k2
                                                                                                                        �
                                                                                  kθt −θPS k2 6 δ, for all t >                              ,
The conclusion of Theorem 3.8 is a strict generalization of a                                                        βγ
                                                                                                                 η β+γ    − ε(3ηβ 2 + 2β)
classical optimization result which considers a static objec-

                                                        βγ
tive, in which case the rate of contraction is 1 − η β+γ                                                                   2
                                                                                   for a constant choice of step size η 6 β+γ .
(see for example Theorem 2.1.15 in Nesterov (2013) or
Lemma 3.7 in Hardt et al. (2016b)). Our rate exactly                          Proof sketch. The basic idea behind these results is the fol-
matches this standard result in the case that ε = 0. The                      lowing. While kθt − θPS k2 > δ, the sample size nt is
proof of Theorem 3.8 can be found in Appendix E.3.                            sufficiently large to ensure a behavior similar to that on a
\end{ReviewVerbatim}


\paragraph{Formalized review target}
The archival source above is not asserted equivalent to this formalized target.
\begin{ReviewVerbatim}
The RERM premise bundle for Theorem 3.10(a) specifies data dimension greater than two, the concrete capped count and inverse-square confidence schedules, a shared bounded exponential moment and tail envelope, measurable adaptive sampling and bad events, data-Lipschitz loss and empirical-loss integrability, smoothness and strong convexity with Wasserstein sensitivity, a convex domain with empirical and population risk-minimizing updates and a stable point, and the numerical W1 budget, contraction slack, absorbing-radius, and initial-distance conditions. The trajectory theorem derives population contraction and a finite entry iteration; neither a caller-supplied small-deviation gate nor a supplied burn-in bound is a premise. The model and expected-loss integrability are defined for every ambient parameter. Joint-gradient smoothness and pointwise strong convexity quantify over all ambient parameters and data points; Wasserstein sensitivity, gradient integrability, and loss differentiation quantify over all ambient parameters. RERM data-loss Lipschitzness also quantifies over every ambient parameter and data point. These global hypotheses are stronger in scope than the source assumptions on the parameter domain and the union of its induced supports; their necessity for the source conclusion is unresolved.
\end{ReviewVerbatim}

\textbf{Archival source anchor:} {\footnotesize\raggedright cited publication, lines 373-\allowbreak{}415\par}
\paragraph{Lean-elaborated paper signature}
\begin{ReviewVerbatim}
field dimension:
{Parameter : Type u_1} →
  [inst : MeasurableSpace Parameter] →
    [inst_1 : NormedAddCommGroup Parameter] →
      [inst_2 : InnerProductSpace ℝ Parameter] →
        [inst_3 : CompleteSpace Parameter] →
          PZMH20PerformativePrediction.Theorem310CorrectedRERMAllRoundTrajectoryData Parameter → ℕ

field cutoff:
{Parameter : Type u_1} →
  [inst : MeasurableSpace Parameter] →
    [inst_1 : NormedAddCommGroup Parameter] →
      [inst_2 : InnerProductSpace ℝ Parameter] →
        [inst_3 : CompleteSpace Parameter] →
          PZMH20PerformativePrediction.Theorem310CorrectedRERMAllRoundTrajectoryData Parameter → ℕ

field hdimension:
∀ {Parameter : Type u_1} [inst : MeasurableSpace Parameter] [inst_1 : NormedAddCommGroup Parameter]
  [inst_2 : InnerProductSpace ℝ Parameter] [inst_3 : CompleteSpace Parameter]
  (self : PZMH20PerformativePrediction.Theorem310CorrectedRERMAllRoundTrajectoryData Parameter), 2 < self.dimension

field eta:
{Parameter : Type u_1} →
  [inst : MeasurableSpace Parameter] →
    [inst_1 : NormedAddCommGroup Parameter] →
      [inst_2 : InnerProductSpace ℝ Parameter] →
        [inst_3 : CompleteSpace Parameter] →
          PZMH20PerformativePrediction.Theorem310CorrectedRERMAllRoundTrajectoryData Parameter → ℝ

field alpha:
{Parameter : Type u_1} →
  [inst : MeasurableSpace Parameter] →
    [inst_1 : NormedAddCommGroup Parameter] →
      [inst_2 : InnerProductSpace ℝ Parameter] →
        [inst_3 : CompleteSpace Parameter] →
          PZMH20PerformativePrediction.Theorem310CorrectedRERMAllRoundTrajectoryData Parameter → ℝ

field gamma:
{Parameter : Type u_1} →
  [inst : MeasurableSpace Parameter] →
    [inst_1 : NormedAddCommGroup Parameter] →
      [inst_2 : InnerProductSpace ℝ Parameter] →
        [inst_3 : CompleteSpace Parameter] →
          PZMH20PerformativePrediction.Theorem310CorrectedRERMAllRoundTrajectoryData Parameter → ℝ

field momentBound:
{Parameter : Type u_1} →
  [inst : MeasurableSpace Parameter] →
    [inst_1 : NormedAddCommGroup Parameter] →
      [inst_2 : InnerProductSpace ℝ Parameter] →
        [inst_3 : CompleteSpace Parameter] →
          PZMH20PerformativePrediction.Theorem310CorrectedRERMAllRoundTrajectoryData Parameter → ℝ

field deviation:
{Parameter : Type u_1} →
  [inst : MeasurableSpace Parameter] →
    [inst_1 : NormedAddCommGroup Parameter] →
      [inst_2 : InnerProductSpace ℝ Parameter] →
        [inst_3 : CompleteSpace Parameter] →
          PZMH20PerformativePrediction.Theorem310CorrectedRERMAllRoundTrajectoryData Parameter → ℝ

field scaledDeviation:
{Parameter : Type u_1} →
  [inst : MeasurableSpace Parameter] →
    [inst_1 : NormedAddCommGroup Parameter] →
      [inst_2 : InnerProductSpace ℝ Parameter] →
        [inst_3 : CompleteSpace Parameter] →
          PZMH20PerformativePrediction.Theorem310CorrectedRERMAllRoundTrajectoryData Parameter → ℝ

field tailBound:
{Parameter : Type u_1} →
  [inst : MeasurableSpace Parameter] →
    [inst_1 : NormedAddCommGroup Parameter] →
      [inst_2 : InnerProductSpace ℝ Parameter] →
        [inst_3 : CompleteSpace Parameter] →
          PZMH20PerformativePrediction.Theorem310CorrectedRERMAllRoundTrajectoryData Parameter → ℝ

field heta_pos:
∀ {Parameter : Type u_1} [inst : MeasurableSpace Parameter] [inst_1 : NormedAddCommGroup Parameter]
  [inst_2 : InnerProductSpace ℝ Parameter] [inst_3 : CompleteSpace Parameter]
  (self : PZMH20PerformativePrediction.Theorem310CorrectedRERMAllRoundTrajectoryData Parameter), 0 < self.eta

field heta_le_one:
∀ {Parameter : Type u_1} [inst : MeasurableSpace Parameter] [inst_1 : NormedAddCommGroup Parameter]
  [inst_2 : InnerProductSpace ℝ Parameter] [inst_3 : CompleteSpace Parameter]
  (self : PZMH20PerformativePrediction.Theorem310CorrectedRERMAllRoundTrajectoryData Parameter), self.eta ≤ 1

field halpha_pos:
∀ {Parameter : Type u_1} [inst : MeasurableSpace Parameter] [inst_1 : NormedAddCommGroup Parameter]
  [inst_2 : InnerProductSpace ℝ Parameter] [inst_3 : CompleteSpace Parameter]
  (self : PZMH20PerformativePrediction.Theorem310CorrectedRERMAllRoundTrajectoryData Parameter), 0 < self.alpha

field hgamma_pos:
∀ {Parameter : Type u_1} [inst : MeasurableSpace Parameter] [inst_1 : NormedAddCommGroup Parameter]
  [inst_2 : InnerProductSpace ℝ Parameter] [inst_3 : CompleteSpace Parameter]
  (self : PZMH20PerformativePrediction.Theorem310CorrectedRERMAllRoundTrajectoryData Parameter), 0 < self.gamma

field halpha_gap:
∀ {Parameter : Type u_1} [inst : MeasurableSpace Parameter] [inst_1 : NormedAddCommGroup Parameter]
  [inst_2 : InnerProductSpace ℝ Parameter] [inst_3 : CompleteSpace Parameter]
  (self : PZMH20PerformativePrediction.Theorem310CorrectedRERMAllRoundTrajectoryData Parameter),
  1 + self.eta < self.alpha

field hscaled:
∀ {Parameter : Type u_1} [inst : MeasurableSpace Parameter] [inst_1 : NormedAddCommGroup Parameter]
  [inst_2 : InnerProductSpace ℝ Parameter] [inst_3 : CompleteSpace Parameter]
  (self : PZMH20PerformativePrediction.Theorem310CorrectedRERMAllRoundTrajectoryData Parameter),
  self.scaledDeviation = self.deviation * Real.rpow 2 (-(1 + self.eta))

field hscaled_pos:
∀ {Parameter : Type u_1} [inst : MeasurableSpace Parameter] [inst_1 : NormedAddCommGroup Parameter]
  [inst_2 : InnerProductSpace ℝ Parameter] [inst_3 : CompleteSpace Parameter]
  (self : PZMH20PerformativePrediction.Theorem310CorrectedRERMAllRoundTrajectoryData Parameter),
  0 < self.scaledDeviation

field hscaled_le_one:
∀ {Parameter : Type u_1} [inst : MeasurableSpace Parameter] [inst_1 : NormedAddCommGroup Parameter]
  [inst_2 : InnerProductSpace ℝ Parameter] [inst_3 : CompleteSpace Parameter]
  (self : PZMH20PerformativePrediction.Theorem310CorrectedRERMAllRoundTrajectoryData Parameter),
  self.scaledDeviation ≤ 1

field htailBound_pos:
∀ {Parameter : Type u_1} [inst : MeasurableSpace Parameter] [inst_1 : NormedAddCommGroup Parameter]
  [inst_2 : InnerProductSpace ℝ Parameter] [inst_3 : CompleteSpace Parameter]
  (self : PZMH20PerformativePrediction.Theorem310CorrectedRERMAllRoundTrajectoryData Parameter), 0 < self.tailBound

field htailEnvelope:
∀ {Parameter : Type u_1} [inst : MeasurableSpace Parameter] [inst_1 : NormedAddCommGroup Parameter]
  [inst_2 : InnerProductSpace ℝ Parameter] [inst_3 : CompleteSpace Parameter]
  (self : PZMH20PerformativePrediction.Theorem310CorrectedRERMAllRoundTrajectoryData Parameter)
  (law : MeasureTheory.ProbabilityMeasure (EuclideanSpace ℝ (Fin self.dimension))),
  AppliedModelingLib.Probability.HasBoundedExponentialRadialMoment law self.alpha self.gamma self.momentBound →
    AppliedModelingLib.Probability.fifthOrderExponentialRadialTailConstant law self.alpha self.gamma ≤ self.tailBound

field p:
{Parameter : Type u_1} →
  [inst : MeasurableSpace Parameter] →
    [inst_1 : NormedAddCommGroup Parameter] →
      [inst_2 : InnerProductSpace ℝ Parameter] →
        [inst_3 : CompleteSpace Parameter] →
          PZMH20PerformativePrediction.Theorem310CorrectedRERMAllRoundTrajectoryData Parameter → ℝ

field headTolerance:
{Parameter : Type u_1} →
  [inst : MeasurableSpace Parameter] →
    [inst_1 : NormedAddCommGroup Parameter] →
      [inst_2 : InnerProductSpace ℝ Parameter] →
        [inst_3 : CompleteSpace Parameter] →
          PZMH20PerformativePrediction.Theorem310CorrectedRERMAllRoundTrajectoryData Parameter → ℝ

field hp:
∀ {Parameter : Type u_1} [inst : MeasurableSpace Parameter] [inst_1 : NormedAddCommGroup Parameter]
  [inst_2 : InnerProductSpace ℝ Parameter] [inst_3 : CompleteSpace Parameter]
  (self : PZMH20PerformativePrediction.Theorem310CorrectedRERMAllRoundTrajectoryData Parameter), 0 < self.p

field hheadTolerance:
∀ {Parameter : Type u_1} [inst : MeasurableSpace Parameter] [inst_1 : NormedAddCommGroup Parameter]
  [inst_2 : InnerProductSpace ℝ Parameter] [inst_3 : CompleteSpace Parameter]
  (self : PZMH20PerformativePrediction.Theorem310CorrectedRERMAllRoundTrajectoryData Parameter), 0 < self.headTolerance

field htailBound_moment:
∀ {Parameter : Type u_1} [inst : MeasurableSpace Parameter] [inst_1 : NormedAddCommGroup Parameter]
  [inst_2 : InnerProductSpace ℝ Parameter] [inst_3 : CompleteSpace Parameter]
  (self : PZMH20PerformativePrediction.Theorem310CorrectedRERMAllRoundTrajectoryData Parameter),
  self.tailBound = ↑(Nat.factorial 5) * self.gamma⁻¹ ^ 5 * self.momentBound

field model:
{Parameter : Type u_1} →
  [inst : MeasurableSpace Parameter] →
    [inst_1 : NormedAddCommGroup Parameter] →
      [inst_2 : InnerProductSpace ℝ Parameter] →
        [inst_3 : CompleteSpace Parameter] →
          (self : PZMH20PerformativePrediction.Theorem310CorrectedRERMAllRoundTrajectoryData Parameter) →
            AppliedModelingLib.MeasurePerformativeModel Parameter (EuclideanSpace ℝ (Fin self.dimension))

field sampling:
{Parameter : Type u_1} →
  [inst : MeasurableSpace Parameter] →
    [inst_1 : NormedAddCommGroup Parameter] →
      [inst_2 : InnerProductSpace ℝ Parameter] →
        [inst_3 : CompleteSpace Parameter] →
          (self : PZMH20PerformativePrediction.Theorem310CorrectedRERMAllRoundTrajectoryData Parameter) →
            AppliedModelingLib.MeasurePerformativeSamplingKernel self.model

field constant:
{Parameter : Type u_1} →
  [inst : MeasurableSpace Parameter] →
    [inst_1 : NormedAddCommGroup Parameter] →
      [inst_2 : InnerProductSpace ℝ Parameter] →
        [inst_3 : CompleteSpace Parameter] →
          PZMH20PerformativePrediction.Theorem310CorrectedRERMAllRoundTrajectoryData Parameter → NNReal

field hloss:
∀ {Parameter : Type u_1} [inst : MeasurableSpace Parameter] [inst_1 : NormedAddCommGroup Parameter]
  [inst_2 : InnerProductSpace ℝ Parameter] [inst_3 : CompleteSpace Parameter]
  (self : PZMH20PerformativePrediction.Theorem310CorrectedRERMAllRoundTrajectoryData Parameter),
  AppliedModelingLib.IsMeasureDataLossLipschitz self.model self.constant

field hconstant:
∀ {Parameter : Type u_1} [inst : MeasurableSpace Parameter] [inst_1 : NormedAddCommGroup Parameter]
  [inst_2 : InnerProductSpace ℝ Parameter] [inst_3 : CompleteSpace Parameter]
  (self : PZMH20PerformativePrediction.Theorem310CorrectedRERMAllRoundTrajectoryData Parameter), self.constant ≠ 0

field gradient:
{Parameter : Type u_1} →
  [inst : MeasurableSpace Parameter] →
    [inst_1 : NormedAddCommGroup Parameter] →
      [inst_2 : InnerProductSpace ℝ Parameter] →
        [inst_3 : CompleteSpace Parameter] →
          (self : PZMH20PerformativePrediction.Theorem310CorrectedRERMAllRoundTrajectoryData Parameter) →
            EuclideanSpace ℝ (Fin self.dimension) → Parameter → Parameter

field smoothness:
{Parameter : Type u_1} →
  [inst : MeasurableSpace Parameter] →
    [inst_1 : NormedAddCommGroup Parameter] →
      [inst_2 : InnerProductSpace ℝ Parameter] →
        [inst_3 : CompleteSpace Parameter] →
          PZMH20PerformativePrediction.Theorem310CorrectedRERMAllRoundTrajectoryData Parameter → NNReal

field hjoint:
∀ {Parameter : Type u_1} [inst : MeasurableSpace Parameter] [inst_1 : NormedAddCommGroup Parameter]
  [inst_2 : InnerProductSpace ℝ Parameter] [inst_3 : CompleteSpace Parameter]
  (self : PZMH20PerformativePrediction.Theorem310CorrectedRERMAllRoundTrajectoryData Parameter),
  AppliedModelingLib.IsJointlySmoothGradient self.gradient self.smoothness

field sensitivity:
{Parameter : Type u_1} →
  [inst : MeasurableSpace Parameter] →
    [inst_1 : NormedAddCommGroup Parameter] →
      [inst_2 : InnerProductSpace ℝ Parameter] →
        [inst_3 : CompleteSpace Parameter] →
          PZMH20PerformativePrediction.Theorem310CorrectedRERMAllRoundTrajectoryData Parameter → ℝ

field modulus:
{Parameter : Type u_1} →
  [inst : MeasurableSpace Parameter] →
    [inst_1 : NormedAddCommGroup Parameter] →
      [inst_2 : InnerProductSpace ℝ Parameter] →
        [inst_3 : CompleteSpace Parameter] →
          PZMH20PerformativePrediction.Theorem310CorrectedRERMAllRoundTrajectoryData Parameter → ℝ

field hmodulus:
∀ {Parameter : Type u_1} [inst : MeasurableSpace Parameter] [inst_1 : NormedAddCommGroup Parameter]
  [inst_2 : InnerProductSpace ℝ Parameter] [inst_3 : CompleteSpace Parameter]
  (self : PZMH20PerformativePrediction.Theorem310CorrectedRERMAllRoundTrajectoryData Parameter), 0 < self.modulus

field hsensitivity_nonneg:
∀ {Parameter : Type u_1} [inst : MeasurableSpace Parameter] [inst_1 : NormedAddCommGroup Parameter]
  [inst_2 : InnerProductSpace ℝ Parameter] [inst_3 : CompleteSpace Parameter]
  (self : PZMH20PerformativePrediction.Theorem310CorrectedRERMAllRoundTrajectoryData Parameter), 0 ≤ self.sensitivity

field hsmoothness_pos:
∀ {Parameter : Type u_1} [inst : MeasurableSpace Parameter] [inst_1 : NormedAddCommGroup Parameter]
  [inst_2 : InnerProductSpace ℝ Parameter] [inst_3 : CompleteSpace Parameter]
  (self : PZMH20PerformativePrediction.Theorem310CorrectedRERMAllRoundTrajectoryData Parameter), 0 < ↑self.smoothness

field hsensitivity_lt:
∀ {Parameter : Type u_1} [inst : MeasurableSpace Parameter] [inst_1 : NormedAddCommGroup Parameter]
  [inst_2 : InnerProductSpace ℝ Parameter] [inst_3 : CompleteSpace Parameter]
  (self : PZMH20PerformativePrediction.Theorem310CorrectedRERMAllRoundTrajectoryData Parameter),
  self.sensitivity < self.modulus / (2 * ↑self.smoothness)

field hsensitive:
∀ {Parameter : Type u_1} [inst : MeasurableSpace Parameter] [inst_1 : NormedAddCommGroup Parameter]
  [inst_2 : InnerProductSpace ℝ Parameter] [inst_3 : CompleteSpace Parameter]
  (self : PZMH20PerformativePrediction.Theorem310CorrectedRERMAllRoundTrajectoryData Parameter),
  PZMH20PerformativePrediction.IsMeasureWassersteinSensitive self.model self.sensitivity

field hstrong:
∀ {Parameter : Type u_1} [inst : MeasurableSpace Parameter] [inst_1 : NormedAddCommGroup Parameter]
  [inst_2 : InnerProductSpace ℝ Parameter] [inst_3 : CompleteSpace Parameter]
  (self : PZMH20PerformativePrediction.Theorem310CorrectedRERMAllRoundTrajectoryData Parameter),
  PZMH20PerformativePrediction.IsMeasurePointwiseGradientStronglyConvex self.model self.gradient self.modulus

field hintegrable:
∀ {Parameter : Type u_1} [inst : MeasurableSpace Parameter] [inst_1 : NormedAddCommGroup Parameter]
  [inst_2 : InnerProductSpace ℝ Parameter] [inst_3 : CompleteSpace Parameter]
  (self : PZMH20PerformativePrediction.Theorem310CorrectedRERMAllRoundTrajectoryData Parameter)
  (distributionParameter evaluatedParameter : Parameter),
  MeasureTheory.Integrable
    (fun (datum : EuclideanSpace ℝ (Fin self.dimension)) => self.gradient datum evaluatedParameter)
    ↑(self.model.dataLaw distributionParameter)

field domain:
{Parameter : Type u_1} →
  [inst : MeasurableSpace Parameter] →
    [inst_1 : NormedAddCommGroup Parameter] →
      [inst_2 : InnerProductSpace ℝ Parameter] →
        [inst_3 : CompleteSpace Parameter] →
          PZMH20PerformativePrediction.Theorem310CorrectedRERMAllRoundTrajectoryData Parameter → Set Parameter

field hconvex:
∀ {Parameter : Type u_1} [inst : MeasurableSpace Parameter] [inst_1 : NormedAddCommGroup Parameter]
  [inst_2 : InnerProductSpace ℝ Parameter] [inst_3 : CompleteSpace Parameter]
  (self : PZMH20PerformativePrediction.Theorem310CorrectedRERMAllRoundTrajectoryData Parameter), Convex ℝ self.domain

field hcountPositive:
∀ {Parameter : Type u_1} [inst : MeasurableSpace Parameter] [inst_1 : NormedAddCommGroup Parameter]
  [inst_2 : InnerProductSpace ℝ Parameter] [inst_3 : CompleteSpace Parameter]
  (self : PZMH20PerformativePrediction.Theorem310CorrectedRERMAllRoundTrajectoryData Parameter) (round : ℕ),
  0 <
    PZMH20PerformativePrediction.theorem310CorrectedRERMCountSchedule self.dimension self.cutoff self.eta self.alpha
      self.gamma self.momentBound self.tailBound self.scaledDeviation self.p self.headTolerance round

field deployedOfHistory:
{Parameter : Type u_1} →
  [inst : MeasurableSpace Parameter] →
    [inst_1 : NormedAddCommGroup Parameter] →
      [inst_2 : InnerProductSpace ℝ Parameter] →
        [inst_3 : CompleteSpace Parameter] →
          (self : PZMH20PerformativePrediction.Theorem310CorrectedRERMAllRoundTrajectoryData Parameter) →
            (iteration : ℕ) →
              AppliedModelingLib.HeterogeneousBatchTrace
                  (fun (round : ℕ) =>
                    Fin
                      (PZMH20PerformativePrediction.theorem310CorrectedRERMCountSchedule self.dimension self.cutoff
                        self.eta self.alpha self.gamma self.momentBound self.tailBound self.scaledDeviation self.p
                        self.headTolerance round))
                  (EuclideanSpace ℝ (Fin self.dimension)) iteration →
                Parameter

field hmeasurableDeployed:
∀ {Parameter : Type u_1} [inst : MeasurableSpace Parameter] [inst_1 : NormedAddCommGroup Parameter]
  [inst_2 : InnerProductSpace ℝ Parameter] [inst_3 : CompleteSpace Parameter]
  (self : PZMH20PerformativePrediction.Theorem310CorrectedRERMAllRoundTrajectoryData Parameter) (iteration : ℕ),
  Measurable (self.deployedOfHistory iteration)

field hempirical:
∀ {Parameter : Type u_1} [inst : MeasurableSpace Parameter] [inst_1 : NormedAddCommGroup Parameter]
  [inst_2 : InnerProductSpace ℝ Parameter] [inst_3 : CompleteSpace Parameter]
  (self : PZMH20PerformativePrediction.Theorem310CorrectedRERMAllRoundTrajectoryData Parameter),
  AppliedModelingLib.IsHeterogeneousMeasureRepeatedEmpiricalRiskMinimizationTrajectoryOn self.model self.domain
    (PZMH20PerformativePrediction.theorem310CorrectedRERMCountSchedule self.dimension self.cutoff self.eta self.alpha
      self.gamma self.momentBound self.tailBound self.scaledDeviation self.p self.headTolerance)
    self.hcountPositive self.deployedOfHistory

field populationUpdate:
{Parameter : Type u_1} →
  [inst : MeasurableSpace Parameter] →
    [inst_1 : NormedAddCommGroup Parameter] →
      [inst_2 : InnerProductSpace ℝ Parameter] →
        [inst_3 : CompleteSpace Parameter] →
          PZMH20PerformativePrediction.Theorem310CorrectedRERMAllRoundTrajectoryData Parameter → Parameter → Parameter

field hpopulation:
∀ {Parameter : Type u_1} [inst : MeasurableSpace Parameter] [inst_1 : NormedAddCommGroup Parameter]
  [inst_2 : InnerProductSpace ℝ Parameter] [inst_3 : CompleteSpace Parameter]
  (self : PZMH20PerformativePrediction.Theorem310CorrectedRERMAllRoundTrajectoryData Parameter),
  AppliedModelingLib.IsMeasureRepeatedRiskMinimizationOn self.model self.domain self.populationUpdate

field hloss_meas:
∀ {Parameter : Type u_1} [inst : MeasurableSpace Parameter] [inst_1 : NormedAddCommGroup Parameter]
  [inst_2 : InnerProductSpace ℝ Parameter] [inst_3 : CompleteSpace Parameter]
  (self : PZMH20PerformativePrediction.Theorem310CorrectedRERMAllRoundTrajectoryData Parameter)
  (distributionParameter evaluatedParameter : Parameter),
  ∀ᶠ (parameter : Parameter) in nhds evaluatedParameter,
    MeasureTheory.AEStronglyMeasurable
      (fun (datum : EuclideanSpace ℝ (Fin self.dimension)) => self.model.loss datum parameter)
      ↑(self.model.dataLaw distributionParameter)

field hpointwise:
∀ {Parameter : Type u_1} [inst : MeasurableSpace Parameter] [inst_1 : NormedAddCommGroup Parameter]
  [inst_2 : InnerProductSpace ℝ Parameter] [inst_3 : CompleteSpace Parameter]
  (self : PZMH20PerformativePrediction.Theorem310CorrectedRERMAllRoundTrajectoryData Parameter)
  (datum : EuclideanSpace ℝ (Fin self.dimension)) (parameter : Parameter),
  HasGradientAt (fun (theta : Parameter) => self.model.loss datum theta) (self.gradient datum parameter) parameter

field stable:
{Parameter : Type u_1} →
  [inst : MeasurableSpace Parameter] →
    [inst_1 : NormedAddCommGroup Parameter] →
      [inst_2 : InnerProductSpace ℝ Parameter] →
        [inst_3 : CompleteSpace Parameter] →
          PZMH20PerformativePrediction.Theorem310CorrectedRERMAllRoundTrajectoryData Parameter → Parameter

field hstablePerformative:
∀ {Parameter : Type u_1} [inst : MeasurableSpace Parameter] [inst_1 : NormedAddCommGroup Parameter]
  [inst_2 : InnerProductSpace ℝ Parameter] [inst_3 : CompleteSpace Parameter]
  (self : PZMH20PerformativePrediction.Theorem310CorrectedRERMAllRoundTrajectoryData Parameter),
  AppliedModelingLib.IsMeasurePerformativelyStableOn self.model self.domain self.stable

field contraction:
{Parameter : Type u_1} →
  [inst : MeasurableSpace Parameter] →
    [inst_1 : NormedAddCommGroup Parameter] →
      [inst_2 : InnerProductSpace ℝ Parameter] →
        [inst_3 : CompleteSpace Parameter] →
          PZMH20PerformativePrediction.Theorem310CorrectedRERMAllRoundTrajectoryData Parameter → ℝ

field outerContraction:
{Parameter : Type u_1} →
  [inst : MeasurableSpace Parameter] →
    [inst_1 : NormedAddCommGroup Parameter] →
      [inst_2 : InnerProductSpace ℝ Parameter] →
        [inst_3 : CompleteSpace Parameter] →
          PZMH20PerformativePrediction.Theorem310CorrectedRERMAllRoundTrajectoryData Parameter → ℝ

field wassersteinBound:
{Parameter : Type u_1} →
  [inst : MeasurableSpace Parameter] →
    [inst_1 : NormedAddCommGroup Parameter] →
      [inst_2 : InnerProductSpace ℝ Parameter] →
        [inst_3 : CompleteSpace Parameter] →
          PZMH20PerformativePrediction.Theorem310CorrectedRERMAllRoundTrajectoryData Parameter → ℝ

field tolerance:
{Parameter : Type u_1} →
  [inst : MeasurableSpace Parameter] →
    [inst_1 : NormedAddCommGroup Parameter] →
      [inst_2 : InnerProductSpace ℝ Parameter] →
        [inst_3 : CompleteSpace Parameter] →
          PZMH20PerformativePrediction.Theorem310CorrectedRERMAllRoundTrajectoryData Parameter → ℝ

field radius:
{Parameter : Type u_1} →
  [inst : MeasurableSpace Parameter] →
    [inst_1 : NormedAddCommGroup Parameter] →
      [inst_2 : InnerProductSpace ℝ Parameter] →
        [inst_3 : CompleteSpace Parameter] →
          PZMH20PerformativePrediction.Theorem310CorrectedRERMAllRoundTrajectoryData Parameter → ℝ

field hcontraction_eq:
∀ {Parameter : Type u_1} [inst : MeasurableSpace Parameter] [inst_1 : NormedAddCommGroup Parameter]
  [inst_2 : InnerProductSpace ℝ Parameter] [inst_3 : CompleteSpace Parameter]
  (self : PZMH20PerformativePrediction.Theorem310CorrectedRERMAllRoundTrajectoryData Parameter),
  self.contraction = self.sensitivity * ↑self.smoothness / self.modulus

field hcontraction_le_outer:
∀ {Parameter : Type u_1} [inst : MeasurableSpace Parameter] [inst_1 : NormedAddCommGroup Parameter]
  [inst_2 : InnerProductSpace ℝ Parameter] [inst_3 : CompleteSpace Parameter]
  (self : PZMH20PerformativePrediction.Theorem310CorrectedRERMAllRoundTrajectoryData Parameter),
  self.contraction ≤ self.outerContraction

field htolerance_nonneg:
∀ {Parameter : Type u_1} [inst : MeasurableSpace Parameter] [inst_1 : NormedAddCommGroup Parameter]
  [inst_2 : InnerProductSpace ℝ Parameter] [inst_3 : CompleteSpace Parameter]
  (self : PZMH20PerformativePrediction.Theorem310CorrectedRERMAllRoundTrajectoryData Parameter), 0 ≤ self.tolerance

field hinitial:
∀ {Parameter : Type u_1} [inst : MeasurableSpace Parameter] [inst_1 : NormedAddCommGroup Parameter]
  [inst_2 : InnerProductSpace ℝ Parameter] [inst_3 : CompleteSpace Parameter]
  (self : PZMH20PerformativePrediction.Theorem310CorrectedRERMAllRoundTrajectoryData Parameter)
  (history :
    AppliedModelingLib.HeterogeneousBatchTrace
      (fun (round : ℕ) =>
        Fin
          (PZMH20PerformativePrediction.theorem310CorrectedRERMCountSchedule self.dimension self.cutoff self.eta
            self.alpha self.gamma self.momentBound self.tailBound self.scaledDeviation self.p self.headTolerance round))
      (EuclideanSpace ℝ (Fin self.dimension)) 0),
  self.deployedOfHistory 0 history ∈ self.domain

field hbudget:
∀ {Parameter : Type u_1} [inst : MeasurableSpace Parameter] [inst_1 : NormedAddCommGroup Parameter]
  [inst_2 : InnerProductSpace ℝ Parameter] [inst_3 : CompleteSpace Parameter]
  (self : PZMH20PerformativePrediction.Theorem310CorrectedRERMAllRoundTrajectoryData Parameter),
  √↑self.dimension * self.deviation + 2 * √↑self.dimension * (32 / 15 * self.tailBound / 16 ^ self.cutoff) +
      self.headTolerance ≤
    self.wassersteinBound

field herror_outer:
∀ {Parameter : Type u_1} [inst : MeasurableSpace Parameter] [inst_1 : NormedAddCommGroup Parameter]
  [inst_2 : InnerProductSpace ℝ Parameter] [inst_3 : CompleteSpace Parameter]
  (self : PZMH20PerformativePrediction.Theorem310CorrectedRERMAllRoundTrajectoryData Parameter),
  √(4 * self.tolerance / self.modulus) ≤ (self.outerContraction - self.contraction) * self.radius

field herror_absorbed:
∀ {Parameter : Type u_1} [inst : MeasurableSpace Parameter] [inst_1 : NormedAddCommGroup Parameter]
  [inst_2 : InnerProductSpace ℝ Parameter] [inst_3 : CompleteSpace Parameter]
  (self : PZMH20PerformativePrediction.Theorem310CorrectedRERMAllRoundTrajectoryData Parameter),
  √(4 * self.tolerance / self.modulus) ≤ (1 - self.contraction) * self.radius

field initialDistance:
{Parameter : Type u_1} →
  [inst : MeasurableSpace Parameter] →
    [inst_1 : NormedAddCommGroup Parameter] →
      [inst_2 : InnerProductSpace ℝ Parameter] →
        [inst_3 : CompleteSpace Parameter] →
          PZMH20PerformativePrediction.Theorem310CorrectedRERMAllRoundTrajectoryData Parameter → ℝ

field houter_lt_one:
∀ {Parameter : Type u_1} [inst : MeasurableSpace Parameter] [inst_1 : NormedAddCommGroup Parameter]
  [inst_2 : InnerProductSpace ℝ Parameter] [inst_3 : CompleteSpace Parameter]
  (self : PZMH20PerformativePrediction.Theorem310CorrectedRERMAllRoundTrajectoryData Parameter),
  self.outerContraction < 1

field hradius_pos:
∀ {Parameter : Type u_1} [inst : MeasurableSpace Parameter] [inst_1 : NormedAddCommGroup Parameter]
  [inst_2 : InnerProductSpace ℝ Parameter] [inst_3 : CompleteSpace Parameter]
  (self : PZMH20PerformativePrediction.Theorem310CorrectedRERMAllRoundTrajectoryData Parameter), 0 < self.radius

field hinitialDistance:
∀ {Parameter : Type u_1} [inst : MeasurableSpace Parameter] [inst_1 : NormedAddCommGroup Parameter]
  [inst_2 : InnerProductSpace ℝ Parameter] [inst_3 : CompleteSpace Parameter]
  (self : PZMH20PerformativePrediction.Theorem310CorrectedRERMAllRoundTrajectoryData Parameter)
  (history :
    AppliedModelingLib.HeterogeneousBatchTrace
      (fun (round : ℕ) =>
        Fin
          (PZMH20PerformativePrediction.theorem310CorrectedRERMCountSchedule self.dimension self.cutoff self.eta
            self.alpha self.gamma self.momentBound self.tailBound self.scaledDeviation self.p self.headTolerance round))
      (EuclideanSpace ℝ (Fin self.dimension)) 0),
  dist (self.deployedOfHistory 0 history) self.stable ≤ self.initialDistance

field hmoment:
∀ {Parameter : Type u_1} [inst : MeasurableSpace Parameter] [inst_1 : NormedAddCommGroup Parameter]
  [inst_2 : InnerProductSpace ℝ Parameter] [inst_3 : CompleteSpace Parameter]
  (self : PZMH20PerformativePrediction.Theorem310CorrectedRERMAllRoundTrajectoryData Parameter) (iteration : ℕ)
  (history :
    AppliedModelingLib.HeterogeneousBatchTrace
      (fun (round : ℕ) =>
        Fin
          (PZMH20PerformativePrediction.theorem310CorrectedRERMCountSchedule self.dimension self.cutoff self.eta
            self.alpha self.gamma self.momentBound self.tailBound self.scaledDeviation self.p self.headTolerance round))
      (EuclideanSpace ℝ (Fin self.dimension)) iteration),
  AppliedModelingLib.Probability.HasBoundedExponentialRadialMoment
    (self.model.dataLaw (self.deployedOfHistory iteration history)) self.alpha self.gamma self.momentBound

field hintegrableEmpiricalLoss:
∀ {Parameter : Type u_1} [inst : MeasurableSpace Parameter] [inst_1 : NormedAddCommGroup Parameter]
  [inst_2 : InnerProductSpace ℝ Parameter] [inst_3 : CompleteSpace Parameter]
  (self : PZMH20PerformativePrediction.Theorem310CorrectedRERMAllRoundTrajectoryData Parameter) (iteration : ℕ)
  (history :
    AppliedModelingLib.HeterogeneousBatchTrace
      (fun (round : ℕ) =>
        Fin
          (PZMH20PerformativePrediction.theorem310CorrectedRERMCountSchedule self.dimension self.cutoff self.eta
            self.alpha self.gamma self.momentBound self.tailBound self.scaledDeviation self.p self.headTolerance round))
      (EuclideanSpace ℝ (Fin self.dimension)) iteration)
  (batch :
    Fin
        (PZMH20PerformativePrediction.theorem310CorrectedRERMCountSchedule self.dimension self.cutoff self.eta
          self.alpha self.gamma self.momentBound self.tailBound self.scaledDeviation self.p self.headTolerance
          iteration) →
      EuclideanSpace ℝ (Fin self.dimension)),
  ∀ candidate ∈ self.domain,
    MeasureTheory.Integrable (fun (datum : EuclideanSpace ℝ (Fin self.dimension)) => self.model.loss datum candidate)
      ↑(AppliedModelingLib.empiricalSampleProbabilityMeasureOfPos (self.hcountPositive iteration) batch)

field hbound_nonneg:
∀ {Parameter : Type u_1} [inst : MeasurableSpace Parameter] [inst_1 : NormedAddCommGroup Parameter]
  [inst_2 : InnerProductSpace ℝ Parameter] [inst_3 : CompleteSpace Parameter]
  (self : PZMH20PerformativePrediction.Theorem310CorrectedRERMAllRoundTrajectoryData Parameter),
  0 ≤ self.wassersteinBound

field htolerance:
∀ {Parameter : Type u_1} [inst : MeasurableSpace Parameter] [inst_1 : NormedAddCommGroup Parameter]
  [inst_2 : InnerProductSpace ℝ Parameter] [inst_3 : CompleteSpace Parameter]
  (self : PZMH20PerformativePrediction.Theorem310CorrectedRERMAllRoundTrajectoryData Parameter),
  ↑self.constant * self.wassersteinBound ≤ self.tolerance

field hevent:
∀ {Parameter : Type u_1} [inst : MeasurableSpace Parameter] [inst_1 : NormedAddCommGroup Parameter]
  [inst_2 : InnerProductSpace ℝ Parameter] [inst_3 : CompleteSpace Parameter]
  (self : PZMH20PerformativePrediction.Theorem310CorrectedRERMAllRoundTrajectoryData Parameter) (iteration : ℕ),
  MeasurableSet
    (AppliedModelingLib.heterogeneousBatchTraceBadEventPair
      (fun (round : ℕ) =>
        Fin
          (PZMH20PerformativePrediction.theorem310CorrectedRERMCountSchedule self.dimension self.cutoff self.eta
            self.alpha self.gamma self.momentBound self.tailBound self.scaledDeviation self.p self.headTolerance round))
      (AppliedModelingLib.pOneFournierGuillinAdaptiveEuclideanSelectedShellPartitionCertificateBadEvent self.dimension
        self.cutoff self.eta
        (AppliedModelingLib.pOneFournierGuillinTheorem310CappedDeviation self.dimension self.eta self.alpha self.gamma
          self.momentBound self.scaledDeviation)
        (AppliedModelingLib.pOneFournierGuillinTheorem310SupercriticalConfidenceSchedule self.p)
        (PZMH20PerformativePrediction.theorem310CorrectedRERMCountSchedule self.dimension self.cutoff self.eta
          self.alpha self.gamma self.momentBound self.tailBound self.scaledDeviation self.p self.headTolerance)
        (AppliedModelingLib.measurePerformativeDeployedLaw self.model
          (fun (round : ℕ) =>
            Fin
              (PZMH20PerformativePrediction.theorem310CorrectedRERMCountSchedule self.dimension self.cutoff self.eta
                self.alpha self.gamma self.momentBound self.tailBound self.scaledDeviation self.p self.headTolerance
                round))
          self.deployedOfHistory))
      iteration)
\end{ReviewVerbatim}

\paragraph{Direct retained dependencies}
\begin{itemize}\raggedright
\item \hyperlink{governing-declaration-307f711c4bf618c0}{Applied\allowbreak{}Modeling\allowbreak{}Lib.\allowbreak{}Heterogeneous\allowbreak{}Batch\allowbreak{}Trace}
\item \hyperlink{governing-declaration-37d7ae049de215a7}{Applied\allowbreak{}Modeling\allowbreak{}Lib.\allowbreak{}Heterogeneous\allowbreak{}Batch\allowbreak{}Trace\allowbreak{}Coordinate}
\item \hyperlink{governing-declaration-50fa794af5bef559}{Applied\allowbreak{}Modeling\allowbreak{}Lib.\allowbreak{}Is\allowbreak{}Heterogeneous\allowbreak{}Measure\allowbreak{}Repeated\allowbreak{}Empirical\allowbreak{}Risk\allowbreak{}Minimization\allowbreak{}Trajectory\allowbreak{}On}
\item \hyperlink{library-prerequisite-2e05b739f5b71f23}{Applied\allowbreak{}Modeling\allowbreak{}Lib.\allowbreak{}Is\allowbreak{}Jointly\allowbreak{}Smooth\allowbreak{}Gradient}
\item \hyperlink{library-prerequisite-49496bb3a06c52e5}{Applied\allowbreak{}Modeling\allowbreak{}Lib.\allowbreak{}Is\allowbreak{}Measure\allowbreak{}Data\allowbreak{}Loss\allowbreak{}Lipschitz}
\item \hyperlink{governing-declaration-522b06e211b1e40a}{Applied\allowbreak{}Modeling\allowbreak{}Lib.\allowbreak{}Is\allowbreak{}Measure\allowbreak{}Performatively\allowbreak{}Stable\allowbreak{}On}
\item \hyperlink{governing-declaration-6ae86e5703739171}{Applied\allowbreak{}Modeling\allowbreak{}Lib.\allowbreak{}Is\allowbreak{}Measure\allowbreak{}Repeated\allowbreak{}Risk\allowbreak{}Minimization\allowbreak{}On}
\item \hyperlink{library-prerequisite-96553d4fbb82f1ac}{Applied\allowbreak{}Modeling\allowbreak{}Lib.\allowbreak{}Measure\allowbreak{}Performative\allowbreak{}Model}
\item \hyperlink{governing-declaration-527443e0fb077287}{Applied\allowbreak{}Modeling\allowbreak{}Lib.\allowbreak{}Measure\allowbreak{}Performative\allowbreak{}Sampling\allowbreak{}Kernel}
\item \hyperlink{governing-declaration-9604ed56aea35d91}{Applied\allowbreak{}Modeling\allowbreak{}Lib.\allowbreak{}Probability.\allowbreak{}Has\allowbreak{}Bounded\allowbreak{}Exponential\allowbreak{}Radial\allowbreak{}Moment}
\item \hyperlink{governing-declaration-93c495b865c4b878}{Applied\allowbreak{}Modeling\allowbreak{}Lib.\allowbreak{}Probability.\allowbreak{}fifth\allowbreak{}Order\allowbreak{}Exponential\allowbreak{}Radial\allowbreak{}Tail\allowbreak{}Constant}
\item \hyperlink{governing-declaration-cbd1f5f969456e5f}{Applied\allowbreak{}Modeling\allowbreak{}Lib.\allowbreak{}empirical\allowbreak{}Sample\allowbreak{}Probability\allowbreak{}Measure\allowbreak{}Of\allowbreak{}Pos}
\item \hyperlink{governing-declaration-68ef005103543934}{Applied\allowbreak{}Modeling\allowbreak{}Lib.\allowbreak{}heterogeneous\allowbreak{}Batch\allowbreak{}Trace\allowbreak{}Bad\allowbreak{}Event\allowbreak{}Pair}
\item \hyperlink{governing-declaration-92d1582063c434f3}{Applied\allowbreak{}Modeling\allowbreak{}Lib.\allowbreak{}heterogeneous\allowbreak{}Batch\allowbreak{}Trace\allowbreak{}Coordinate\allowbreak{}Measurable\allowbreak{}Space}
\item \hyperlink{governing-declaration-caf4e7d2e1526646}{Applied\allowbreak{}Modeling\allowbreak{}Lib.\allowbreak{}heterogeneous\allowbreak{}Batch\allowbreak{}Trace\allowbreak{}Measurable\allowbreak{}Space}
\item \hyperlink{governing-declaration-7bc25db7b9b1aaf4}{Applied\allowbreak{}Modeling\allowbreak{}Lib.\allowbreak{}measure\allowbreak{}Performative\allowbreak{}Deployed\allowbreak{}Law}
\item \hyperlink{governing-declaration-30e7d69ebe63768c}{Applied\allowbreak{}Modeling\allowbreak{}Lib.\allowbreak{}p\allowbreak{}One\allowbreak{}Fournier\allowbreak{}Guillin\allowbreak{}Adaptive\allowbreak{}Euclidean\allowbreak{}Selected\allowbreak{}Shell\allowbreak{}Partition\allowbreak{}Certificate\allowbreak{}Bad\allowbreak{}Event}
\item \hyperlink{governing-declaration-25a90c9551dbd2b1}{Applied\allowbreak{}Modeling\allowbreak{}Lib.\allowbreak{}p\allowbreak{}One\allowbreak{}Fournier\allowbreak{}Guillin\allowbreak{}Theorem310\allowbreak{}Capped\allowbreak{}Deviation}
\item \hyperlink{governing-declaration-6a2f3c5b322d78da}{Applied\allowbreak{}Modeling\allowbreak{}Lib.\allowbreak{}p\allowbreak{}One\allowbreak{}Fournier\allowbreak{}Guillin\allowbreak{}Theorem310\allowbreak{}Supercritical\allowbreak{}Confidence\allowbreak{}Schedule}
\item \hyperlink{governing-declaration-11c9ab6e3c6f528b}{PZMH20\allowbreak{}Performative\allowbreak{}Prediction.\allowbreak{}Is\allowbreak{}Measure\allowbreak{}Pointwise\allowbreak{}Gradient\allowbreak{}Strongly\allowbreak{}Convex}
\item \hyperlink{paper-prerequisite-2ea2f5769822c5c4}{PZMH20\allowbreak{}Performative\allowbreak{}Prediction.\allowbreak{}Is\allowbreak{}Measure\allowbreak{}Wasserstein\allowbreak{}Sensitive}
\item \hyperlink{governing-declaration-f4e329b046167e0c}{PZMH20\allowbreak{}Performative\allowbreak{}Prediction.\allowbreak{}theorem310\allowbreak{}Corrected\allowbreak{}RERM\allowbreak{}Count\allowbreak{}Schedule}
\end{itemize}
\paragraph{Recorded source-to-declaration screening}
\textbf{Verdict:} \texttt{matches\_approved\_corrected\_target}
\\
{\raggedright
\textbf{Reason:} The expanded carrier matches the disclosed corrected RERM premise bundle:\allowbreak{} dimension greater than two, the capped count and inverse-\allowbreak{}square failure schedules, a uniform adaptive exponential-\allowbreak{}moment/\allowbreak{}tail envelope, measurable sampling and bad events, globally scoped loss/\allowbreak{}gradient regularity and integrability, empirical and population RRM data, a supplied stable point, and the stated contraction, W1-\allowbreak{}budget, absorption, radius, and initial-\allowbreak{}distance conditions.\allowbreak{} The graph-\allowbreak{}bound use consumes this carrier to conclude a deterministic finite entry iteration followed by all-\allowbreak{}round containment with probability at least 1 -\allowbreak{} p;\allowbreak{} the carrier supplies neither that entry nor the removed small-\allowbreak{}deviation gate.\allowbreak{}
\par}
\reviewmatch{Matches formalized target}
\noindent\textbf{Reviewer annotation}\par
\reviewerbox
\clearpage
\hypertarget{paper-prerequisite-5a443fd08b184a8b}{}
\subsection*{\small\ttfamily\raggedright PZMH20\allowbreak{}Performative\allowbreak{}Prediction.\allowbreak{}Theorem310\allowbreak{}Corrected\allowbreak{}RGD\allowbreak{}All\allowbreak{}Round\allowbreak{}Trajectory\allowbreak{}Data}
\noindent\textbf{Paper-local declaration:}\par
{\footnotesize\ttfamily\raggedright PZMH20\allowbreak{}Performative\allowbreak{}Prediction.\allowbreak{}Theorem310\allowbreak{}Corrected\allowbreak{}RGD\allowbreak{}All\allowbreak{}Round\allowbreak{}Trajectory\allowbreak{}Data\par}
\textbf{Source locator:} {\footnotesize\raggedright cited publication, lines 373-\allowbreak{}415\par}
\paragraph{Verbatim paper-source connection}
\begin{ReviewVerbatim}
Theorem 3.10. Suppose that the loss `(z; θ) is β-jointly
Theorem 3.8. Suppose that the loss `(z; θ) is β-jointly                       smooth (A1) and γ-strongly convex (A2), and that there exist
smooth (A1) and γ-strongly convex (A2). If the distribution                                                def R        α
                                        γ                                     α > 1, µ > 0 such that ξα,µ = Rm eµ|x| dD(θ) is finite
map D(·) is ε-sensitive with ε < (β+γ)(1+1.5ηβ) , then RGD
                                                                              ∀θ ∈ Θ. Let δ ∈ (0, 1) be a radius of convergence.     Con-
                     2
with step size η 6 β+γ satisfies the following:
                                                                                                                                    � t
                                                                                                                            1
                                                                              sider running RERM or RGD with nt = O (εδ)m log p
(a) kGgd (θ) − Ggd (θ0 )k2                                                    samples at time t.
                                                                                                                            γ
                                                                              (a) If the map D(·) is ε-sensitive with ε < 2β  , then with
                         βγ                                          0
      6       1−η           − ε(1.5ηβ 2 + β)             kθ − θ k2 .              probability 1 − p, RERM satisfies,
                        β+γ
                                                                                                                  log 1δ kθ0 − θPS k2
                                                                                                                     �
(b) The iterates θt of RGD converge to a unique performa-                           kθt − θPS k2 6 δ, for all t >                       .
    tively stable point θPS at a linear rate,                                                                           1 − 2εβ
                                                                                                                              γ

                                                                                                                                 γ
                                               �1                             (b) If the map D(·) is ε-sensitive with ε < (β+γ)(1+1.5ηβ)  ,
                                         log    δ kθ0 − θPS k2
     kθt − θPS k2 6 δ for t >                                             .       then with probability 1 − p, REGD with satisfies,
                                          βγ          2
                                 η       β+γ − ε(1.5ηβ + β)
                                                                                                                    log 1δ kθ0 − θPS k2
                                                                                                                        �
                                                                                  kθt −θPS k2 6 δ, for all t >                              ,
The conclusion of Theorem 3.8 is a strict generalization of a                                                        βγ
                                                                                                                 η β+γ    − ε(3ηβ 2 + 2β)
classical optimization result which considers a static objec-

                                                        βγ
tive, in which case the rate of contraction is 1 − η β+γ                                                                   2
                                                                                   for a constant choice of step size η 6 β+γ .
(see for example Theorem 2.1.15 in Nesterov (2013) or
Lemma 3.7 in Hardt et al. (2016b)). Our rate exactly                          Proof sketch. The basic idea behind these results is the fol-
matches this standard result in the case that ε = 0. The                      lowing. While kθt − θPS k2 > δ, the sample size nt is
proof of Theorem 3.8 can be found in Appendix E.3.                            sufficiently large to ensure a behavior similar to that on a
\end{ReviewVerbatim}


\paragraph{Formalized review target}
The archival source above is not asserted equivalent to this formalized target.
\begin{ReviewVerbatim}
The REGD premise bundle for Theorem 3.10(b) specifies data dimension greater than two, the concrete capped count and inverse-square confidence schedules, a shared bounded exponential moment and tail envelope, measurable adaptive sampling and bad events, joint-gradient smoothness and pointwise strong convexity with Wasserstein sensitivity, a convex domain with a variational nonexpansive projection, gradient integrability and pointwise loss differentiation, a checked empirical-gradient trajectory and a stable point, and the step-size, numerical W1 budget, contraction slack, absorbing-radius, and initial-distance conditions. The trajectory theorem derives population contraction and a finite entry iteration; neither a caller-supplied small-deviation gate nor a supplied burn-in bound is a premise. The model and expected-loss integrability are defined for every ambient parameter. Joint-gradient smoothness and pointwise strong convexity quantify over all ambient parameters and data points; Wasserstein sensitivity, gradient integrability, and loss differentiation quantify over all ambient parameters. These global hypotheses are stronger in scope than the source assumptions on the parameter domain and the union of its induced supports; their necessity for the source conclusion is unresolved.
\end{ReviewVerbatim}

\textbf{Archival source anchor:} {\footnotesize\raggedright cited publication, lines 373-\allowbreak{}415\par}
\paragraph{Lean-elaborated paper signature}
\begin{ReviewVerbatim}
field dimension:
{Parameter : Type u_1} →
  [inst : MeasurableSpace Parameter] →
    [inst_1 : NormedAddCommGroup Parameter] →
      [inst_2 : InnerProductSpace ℝ Parameter] →
        [inst_3 : CompleteSpace Parameter] →
          PZMH20PerformativePrediction.Theorem310CorrectedRGDAllRoundTrajectoryData Parameter → ℕ

field cutoff:
{Parameter : Type u_1} →
  [inst : MeasurableSpace Parameter] →
    [inst_1 : NormedAddCommGroup Parameter] →
      [inst_2 : InnerProductSpace ℝ Parameter] →
        [inst_3 : CompleteSpace Parameter] →
          PZMH20PerformativePrediction.Theorem310CorrectedRGDAllRoundTrajectoryData Parameter → ℕ

field hdimension:
∀ {Parameter : Type u_1} [inst : MeasurableSpace Parameter] [inst_1 : NormedAddCommGroup Parameter]
  [inst_2 : InnerProductSpace ℝ Parameter] [inst_3 : CompleteSpace Parameter]
  (self : PZMH20PerformativePrediction.Theorem310CorrectedRGDAllRoundTrajectoryData Parameter), 2 < self.dimension

field eta:
{Parameter : Type u_1} →
  [inst : MeasurableSpace Parameter] →
    [inst_1 : NormedAddCommGroup Parameter] →
      [inst_2 : InnerProductSpace ℝ Parameter] →
        [inst_3 : CompleteSpace Parameter] →
          PZMH20PerformativePrediction.Theorem310CorrectedRGDAllRoundTrajectoryData Parameter → ℝ

field alpha:
{Parameter : Type u_1} →
  [inst : MeasurableSpace Parameter] →
    [inst_1 : NormedAddCommGroup Parameter] →
      [inst_2 : InnerProductSpace ℝ Parameter] →
        [inst_3 : CompleteSpace Parameter] →
          PZMH20PerformativePrediction.Theorem310CorrectedRGDAllRoundTrajectoryData Parameter → ℝ

field gamma:
{Parameter : Type u_1} →
  [inst : MeasurableSpace Parameter] →
    [inst_1 : NormedAddCommGroup Parameter] →
      [inst_2 : InnerProductSpace ℝ Parameter] →
        [inst_3 : CompleteSpace Parameter] →
          PZMH20PerformativePrediction.Theorem310CorrectedRGDAllRoundTrajectoryData Parameter → ℝ

field momentBound:
{Parameter : Type u_1} →
  [inst : MeasurableSpace Parameter] →
    [inst_1 : NormedAddCommGroup Parameter] →
      [inst_2 : InnerProductSpace ℝ Parameter] →
        [inst_3 : CompleteSpace Parameter] →
          PZMH20PerformativePrediction.Theorem310CorrectedRGDAllRoundTrajectoryData Parameter → ℝ

field deviation:
{Parameter : Type u_1} →
  [inst : MeasurableSpace Parameter] →
    [inst_1 : NormedAddCommGroup Parameter] →
      [inst_2 : InnerProductSpace ℝ Parameter] →
        [inst_3 : CompleteSpace Parameter] →
          PZMH20PerformativePrediction.Theorem310CorrectedRGDAllRoundTrajectoryData Parameter → ℝ

field scaledDeviation:
{Parameter : Type u_1} →
  [inst : MeasurableSpace Parameter] →
    [inst_1 : NormedAddCommGroup Parameter] →
      [inst_2 : InnerProductSpace ℝ Parameter] →
        [inst_3 : CompleteSpace Parameter] →
          PZMH20PerformativePrediction.Theorem310CorrectedRGDAllRoundTrajectoryData Parameter → ℝ

field tailBound:
{Parameter : Type u_1} →
  [inst : MeasurableSpace Parameter] →
    [inst_1 : NormedAddCommGroup Parameter] →
      [inst_2 : InnerProductSpace ℝ Parameter] →
        [inst_3 : CompleteSpace Parameter] →
          PZMH20PerformativePrediction.Theorem310CorrectedRGDAllRoundTrajectoryData Parameter → ℝ

field heta_pos:
∀ {Parameter : Type u_1} [inst : MeasurableSpace Parameter] [inst_1 : NormedAddCommGroup Parameter]
  [inst_2 : InnerProductSpace ℝ Parameter] [inst_3 : CompleteSpace Parameter]
  (self : PZMH20PerformativePrediction.Theorem310CorrectedRGDAllRoundTrajectoryData Parameter), 0 < self.eta

field heta_le_one:
∀ {Parameter : Type u_1} [inst : MeasurableSpace Parameter] [inst_1 : NormedAddCommGroup Parameter]
  [inst_2 : InnerProductSpace ℝ Parameter] [inst_3 : CompleteSpace Parameter]
  (self : PZMH20PerformativePrediction.Theorem310CorrectedRGDAllRoundTrajectoryData Parameter), self.eta ≤ 1

field halpha_pos:
∀ {Parameter : Type u_1} [inst : MeasurableSpace Parameter] [inst_1 : NormedAddCommGroup Parameter]
  [inst_2 : InnerProductSpace ℝ Parameter] [inst_3 : CompleteSpace Parameter]
  (self : PZMH20PerformativePrediction.Theorem310CorrectedRGDAllRoundTrajectoryData Parameter), 0 < self.alpha

field hgamma_pos:
∀ {Parameter : Type u_1} [inst : MeasurableSpace Parameter] [inst_1 : NormedAddCommGroup Parameter]
  [inst_2 : InnerProductSpace ℝ Parameter] [inst_3 : CompleteSpace Parameter]
  (self : PZMH20PerformativePrediction.Theorem310CorrectedRGDAllRoundTrajectoryData Parameter), 0 < self.gamma

field halpha_gap:
∀ {Parameter : Type u_1} [inst : MeasurableSpace Parameter] [inst_1 : NormedAddCommGroup Parameter]
  [inst_2 : InnerProductSpace ℝ Parameter] [inst_3 : CompleteSpace Parameter]
  (self : PZMH20PerformativePrediction.Theorem310CorrectedRGDAllRoundTrajectoryData Parameter),
  1 + self.eta < self.alpha

field hscaled:
∀ {Parameter : Type u_1} [inst : MeasurableSpace Parameter] [inst_1 : NormedAddCommGroup Parameter]
  [inst_2 : InnerProductSpace ℝ Parameter] [inst_3 : CompleteSpace Parameter]
  (self : PZMH20PerformativePrediction.Theorem310CorrectedRGDAllRoundTrajectoryData Parameter),
  self.scaledDeviation = self.deviation * Real.rpow 2 (-(1 + self.eta))

field hscaled_pos:
∀ {Parameter : Type u_1} [inst : MeasurableSpace Parameter] [inst_1 : NormedAddCommGroup Parameter]
  [inst_2 : InnerProductSpace ℝ Parameter] [inst_3 : CompleteSpace Parameter]
  (self : PZMH20PerformativePrediction.Theorem310CorrectedRGDAllRoundTrajectoryData Parameter), 0 < self.scaledDeviation

field hscaled_le_one:
∀ {Parameter : Type u_1} [inst : MeasurableSpace Parameter] [inst_1 : NormedAddCommGroup Parameter]
  [inst_2 : InnerProductSpace ℝ Parameter] [inst_3 : CompleteSpace Parameter]
  (self : PZMH20PerformativePrediction.Theorem310CorrectedRGDAllRoundTrajectoryData Parameter), self.scaledDeviation ≤ 1

field htailBound_pos:
∀ {Parameter : Type u_1} [inst : MeasurableSpace Parameter] [inst_1 : NormedAddCommGroup Parameter]
  [inst_2 : InnerProductSpace ℝ Parameter] [inst_3 : CompleteSpace Parameter]
  (self : PZMH20PerformativePrediction.Theorem310CorrectedRGDAllRoundTrajectoryData Parameter), 0 < self.tailBound

field htailEnvelope:
∀ {Parameter : Type u_1} [inst : MeasurableSpace Parameter] [inst_1 : NormedAddCommGroup Parameter]
  [inst_2 : InnerProductSpace ℝ Parameter] [inst_3 : CompleteSpace Parameter]
  (self : PZMH20PerformativePrediction.Theorem310CorrectedRGDAllRoundTrajectoryData Parameter)
  (law : MeasureTheory.ProbabilityMeasure (EuclideanSpace ℝ (Fin self.dimension))),
  AppliedModelingLib.Probability.HasBoundedExponentialRadialMoment law self.alpha self.gamma self.momentBound →
    AppliedModelingLib.Probability.fifthOrderExponentialRadialTailConstant law self.alpha self.gamma ≤ self.tailBound

field p:
{Parameter : Type u_1} →
  [inst : MeasurableSpace Parameter] →
    [inst_1 : NormedAddCommGroup Parameter] →
      [inst_2 : InnerProductSpace ℝ Parameter] →
        [inst_3 : CompleteSpace Parameter] →
          PZMH20PerformativePrediction.Theorem310CorrectedRGDAllRoundTrajectoryData Parameter → ℝ

field headTolerance:
{Parameter : Type u_1} →
  [inst : MeasurableSpace Parameter] →
    [inst_1 : NormedAddCommGroup Parameter] →
      [inst_2 : InnerProductSpace ℝ Parameter] →
        [inst_3 : CompleteSpace Parameter] →
          PZMH20PerformativePrediction.Theorem310CorrectedRGDAllRoundTrajectoryData Parameter → ℝ

field hp:
∀ {Parameter : Type u_1} [inst : MeasurableSpace Parameter] [inst_1 : NormedAddCommGroup Parameter]
  [inst_2 : InnerProductSpace ℝ Parameter] [inst_3 : CompleteSpace Parameter]
  (self : PZMH20PerformativePrediction.Theorem310CorrectedRGDAllRoundTrajectoryData Parameter), 0 < self.p

field hheadTolerance:
∀ {Parameter : Type u_1} [inst : MeasurableSpace Parameter] [inst_1 : NormedAddCommGroup Parameter]
  [inst_2 : InnerProductSpace ℝ Parameter] [inst_3 : CompleteSpace Parameter]
  (self : PZMH20PerformativePrediction.Theorem310CorrectedRGDAllRoundTrajectoryData Parameter), 0 < self.headTolerance

field htailBound_moment:
∀ {Parameter : Type u_1} [inst : MeasurableSpace Parameter] [inst_1 : NormedAddCommGroup Parameter]
  [inst_2 : InnerProductSpace ℝ Parameter] [inst_3 : CompleteSpace Parameter]
  (self : PZMH20PerformativePrediction.Theorem310CorrectedRGDAllRoundTrajectoryData Parameter),
  self.tailBound = ↑(Nat.factorial 5) * self.gamma⁻¹ ^ 5 * self.momentBound

field model:
{Parameter : Type u_1} →
  [inst : MeasurableSpace Parameter] →
    [inst_1 : NormedAddCommGroup Parameter] →
      [inst_2 : InnerProductSpace ℝ Parameter] →
        [inst_3 : CompleteSpace Parameter] →
          (self : PZMH20PerformativePrediction.Theorem310CorrectedRGDAllRoundTrajectoryData Parameter) →
            AppliedModelingLib.MeasurePerformativeModel Parameter (EuclideanSpace ℝ (Fin self.dimension))

field sampling:
{Parameter : Type u_1} →
  [inst : MeasurableSpace Parameter] →
    [inst_1 : NormedAddCommGroup Parameter] →
      [inst_2 : InnerProductSpace ℝ Parameter] →
        [inst_3 : CompleteSpace Parameter] →
          (self : PZMH20PerformativePrediction.Theorem310CorrectedRGDAllRoundTrajectoryData Parameter) →
            AppliedModelingLib.MeasurePerformativeSamplingKernel self.model

field gradient:
{Parameter : Type u_1} →
  [inst : MeasurableSpace Parameter] →
    [inst_1 : NormedAddCommGroup Parameter] →
      [inst_2 : InnerProductSpace ℝ Parameter] →
        [inst_3 : CompleteSpace Parameter] →
          (self : PZMH20PerformativePrediction.Theorem310CorrectedRGDAllRoundTrajectoryData Parameter) →
            EuclideanSpace ℝ (Fin self.dimension) → Parameter → Parameter

field smoothness:
{Parameter : Type u_1} →
  [inst : MeasurableSpace Parameter] →
    [inst_1 : NormedAddCommGroup Parameter] →
      [inst_2 : InnerProductSpace ℝ Parameter] →
        [inst_3 : CompleteSpace Parameter] →
          PZMH20PerformativePrediction.Theorem310CorrectedRGDAllRoundTrajectoryData Parameter → NNReal

field hjoint:
∀ {Parameter : Type u_1} [inst : MeasurableSpace Parameter] [inst_1 : NormedAddCommGroup Parameter]
  [inst_2 : InnerProductSpace ℝ Parameter] [inst_3 : CompleteSpace Parameter]
  (self : PZMH20PerformativePrediction.Theorem310CorrectedRGDAllRoundTrajectoryData Parameter),
  AppliedModelingLib.IsJointlySmoothGradient self.gradient self.smoothness

field hsmoothness:
∀ {Parameter : Type u_1} [inst : MeasurableSpace Parameter] [inst_1 : NormedAddCommGroup Parameter]
  [inst_2 : InnerProductSpace ℝ Parameter] [inst_3 : CompleteSpace Parameter]
  (self : PZMH20PerformativePrediction.Theorem310CorrectedRGDAllRoundTrajectoryData Parameter), self.smoothness ≠ 0

field sensitivity:
{Parameter : Type u_1} →
  [inst : MeasurableSpace Parameter] →
    [inst_1 : NormedAddCommGroup Parameter] →
      [inst_2 : InnerProductSpace ℝ Parameter] →
        [inst_3 : CompleteSpace Parameter] →
          PZMH20PerformativePrediction.Theorem310CorrectedRGDAllRoundTrajectoryData Parameter → ℝ

field modulus:
{Parameter : Type u_1} →
  [inst : MeasurableSpace Parameter] →
    [inst_1 : NormedAddCommGroup Parameter] →
      [inst_2 : InnerProductSpace ℝ Parameter] →
        [inst_3 : CompleteSpace Parameter] →
          PZMH20PerformativePrediction.Theorem310CorrectedRGDAllRoundTrajectoryData Parameter → ℝ

field hsensitivity_nonneg:
∀ {Parameter : Type u_1} [inst : MeasurableSpace Parameter] [inst_1 : NormedAddCommGroup Parameter]
  [inst_2 : InnerProductSpace ℝ Parameter] [inst_3 : CompleteSpace Parameter]
  (self : PZMH20PerformativePrediction.Theorem310CorrectedRGDAllRoundTrajectoryData Parameter), 0 ≤ self.sensitivity

field hsensitive:
∀ {Parameter : Type u_1} [inst : MeasurableSpace Parameter] [inst_1 : NormedAddCommGroup Parameter]
  [inst_2 : InnerProductSpace ℝ Parameter] [inst_3 : CompleteSpace Parameter]
  (self : PZMH20PerformativePrediction.Theorem310CorrectedRGDAllRoundTrajectoryData Parameter),
  PZMH20PerformativePrediction.IsMeasureWassersteinSensitive self.model self.sensitivity

field hmodulus_pos:
∀ {Parameter : Type u_1} [inst : MeasurableSpace Parameter] [inst_1 : NormedAddCommGroup Parameter]
  [inst_2 : InnerProductSpace ℝ Parameter] [inst_3 : CompleteSpace Parameter]
  (self : PZMH20PerformativePrediction.Theorem310CorrectedRGDAllRoundTrajectoryData Parameter), 0 < self.modulus

field hmodulus_le:
∀ {Parameter : Type u_1} [inst : MeasurableSpace Parameter] [inst_1 : NormedAddCommGroup Parameter]
  [inst_2 : InnerProductSpace ℝ Parameter] [inst_3 : CompleteSpace Parameter]
  (self : PZMH20PerformativePrediction.Theorem310CorrectedRGDAllRoundTrajectoryData Parameter),
  self.modulus ≤ ↑self.smoothness

field hstrong:
∀ {Parameter : Type u_1} [inst : MeasurableSpace Parameter] [inst_1 : NormedAddCommGroup Parameter]
  [inst_2 : InnerProductSpace ℝ Parameter] [inst_3 : CompleteSpace Parameter]
  (self : PZMH20PerformativePrediction.Theorem310CorrectedRGDAllRoundTrajectoryData Parameter),
  PZMH20PerformativePrediction.IsMeasurePointwiseGradientStronglyConvex self.model self.gradient self.modulus

field domain:
{Parameter : Type u_1} →
  [inst : MeasurableSpace Parameter] →
    [inst_1 : NormedAddCommGroup Parameter] →
      [inst_2 : InnerProductSpace ℝ Parameter] →
        [inst_3 : CompleteSpace Parameter] →
          PZMH20PerformativePrediction.Theorem310CorrectedRGDAllRoundTrajectoryData Parameter → Set Parameter

field hconvex:
∀ {Parameter : Type u_1} [inst : MeasurableSpace Parameter] [inst_1 : NormedAddCommGroup Parameter]
  [inst_2 : InnerProductSpace ℝ Parameter] [inst_3 : CompleteSpace Parameter]
  (self : PZMH20PerformativePrediction.Theorem310CorrectedRGDAllRoundTrajectoryData Parameter), Convex ℝ self.domain

field project:
{Parameter : Type u_1} →
  [inst : MeasurableSpace Parameter] →
    [inst_1 : NormedAddCommGroup Parameter] →
      [inst_2 : InnerProductSpace ℝ Parameter] →
        [inst_3 : CompleteSpace Parameter] →
          PZMH20PerformativePrediction.Theorem310CorrectedRGDAllRoundTrajectoryData Parameter → Parameter → Parameter

field hproject:
∀ {Parameter : Type u_1} [inst : MeasurableSpace Parameter] [inst_1 : NormedAddCommGroup Parameter]
  [inst_2 : InnerProductSpace ℝ Parameter] [inst_3 : CompleteSpace Parameter]
  (self : PZMH20PerformativePrediction.Theorem310CorrectedRGDAllRoundTrajectoryData Parameter),
  LipschitzWith 1 self.project

field hprojectVariational:
∀ {Parameter : Type u_1} [inst : MeasurableSpace Parameter] [inst_1 : NormedAddCommGroup Parameter]
  [inst_2 : InnerProductSpace ℝ Parameter] [inst_3 : CompleteSpace Parameter]
  (self : PZMH20PerformativePrediction.Theorem310CorrectedRGDAllRoundTrajectoryData Parameter),
  AppliedModelingLib.IsVariationalEuclideanProjectionOn self.domain self.project

field stepSize:
{Parameter : Type u_1} →
  [inst : MeasurableSpace Parameter] →
    [inst_1 : NormedAddCommGroup Parameter] →
      [inst_2 : InnerProductSpace ℝ Parameter] →
        [inst_3 : CompleteSpace Parameter] →
          PZMH20PerformativePrediction.Theorem310CorrectedRGDAllRoundTrajectoryData Parameter → ℝ

field hstepSize:
∀ {Parameter : Type u_1} [inst : MeasurableSpace Parameter] [inst_1 : NormedAddCommGroup Parameter]
  [inst_2 : InnerProductSpace ℝ Parameter] [inst_3 : CompleteSpace Parameter]
  (self : PZMH20PerformativePrediction.Theorem310CorrectedRGDAllRoundTrajectoryData Parameter), 0 ≤ self.stepSize

field hstepSize_pos:
∀ {Parameter : Type u_1} [inst : MeasurableSpace Parameter] [inst_1 : NormedAddCommGroup Parameter]
  [inst_2 : InnerProductSpace ℝ Parameter] [inst_3 : CompleteSpace Parameter]
  (self : PZMH20PerformativePrediction.Theorem310CorrectedRGDAllRoundTrajectoryData Parameter), 0 < self.stepSize

field hstepSize_le:
∀ {Parameter : Type u_1} [inst : MeasurableSpace Parameter] [inst_1 : NormedAddCommGroup Parameter]
  [inst_2 : InnerProductSpace ℝ Parameter] [inst_3 : CompleteSpace Parameter]
  (self : PZMH20PerformativePrediction.Theorem310CorrectedRGDAllRoundTrajectoryData Parameter),
  self.stepSize ≤ 2 / (self.modulus + ↑self.smoothness)

field hsensitivity_small:
∀ {Parameter : Type u_1} [inst : MeasurableSpace Parameter] [inst_1 : NormedAddCommGroup Parameter]
  [inst_2 : InnerProductSpace ℝ Parameter] [inst_3 : CompleteSpace Parameter]
  (self : PZMH20PerformativePrediction.Theorem310CorrectedRGDAllRoundTrajectoryData Parameter),
  self.sensitivity < self.modulus / ((self.modulus + ↑self.smoothness) * (1 + 3 / 2 * self.stepSize * ↑self.smoothness))

field hintegrable:
∀ {Parameter : Type u_1} [inst : MeasurableSpace Parameter] [inst_1 : NormedAddCommGroup Parameter]
  [inst_2 : InnerProductSpace ℝ Parameter] [inst_3 : CompleteSpace Parameter]
  (self : PZMH20PerformativePrediction.Theorem310CorrectedRGDAllRoundTrajectoryData Parameter)
  (distributionParameter evaluatedParameter : Parameter),
  MeasureTheory.Integrable
    (fun (datum : EuclideanSpace ℝ (Fin self.dimension)) => self.gradient datum evaluatedParameter)
    ↑(self.model.dataLaw distributionParameter)

field hloss_meas:
∀ {Parameter : Type u_1} [inst : MeasurableSpace Parameter] [inst_1 : NormedAddCommGroup Parameter]
  [inst_2 : InnerProductSpace ℝ Parameter] [inst_3 : CompleteSpace Parameter]
  (self : PZMH20PerformativePrediction.Theorem310CorrectedRGDAllRoundTrajectoryData Parameter)
  (distributionParameter evaluatedParameter : Parameter),
  ∀ᶠ (parameter : Parameter) in nhds evaluatedParameter,
    MeasureTheory.AEStronglyMeasurable
      (fun (datum : EuclideanSpace ℝ (Fin self.dimension)) => self.model.loss datum parameter)
      ↑(self.model.dataLaw distributionParameter)

field hpointwise:
∀ {Parameter : Type u_1} [inst : MeasurableSpace Parameter] [inst_1 : NormedAddCommGroup Parameter]
  [inst_2 : InnerProductSpace ℝ Parameter] [inst_3 : CompleteSpace Parameter]
  (self : PZMH20PerformativePrediction.Theorem310CorrectedRGDAllRoundTrajectoryData Parameter)
  (datum : EuclideanSpace ℝ (Fin self.dimension)) (parameter : Parameter),
  HasGradientAt (fun (theta : Parameter) => self.model.loss datum theta) (self.gradient datum parameter) parameter

field hcountPositive:
∀ {Parameter : Type u_1} [inst : MeasurableSpace Parameter] [inst_1 : NormedAddCommGroup Parameter]
  [inst_2 : InnerProductSpace ℝ Parameter] [inst_3 : CompleteSpace Parameter]
  (self : PZMH20PerformativePrediction.Theorem310CorrectedRGDAllRoundTrajectoryData Parameter) (round : ℕ),
  0 <
    AppliedModelingLib.pOneFournierGuillinTheorem310SupercriticalCappedConcreteCountSchedule self.dimension self.cutoff
      self.eta self.alpha self.gamma self.momentBound self.tailBound self.scaledDeviation self.p self.headTolerance
      round

field deployedOfHistory:
{Parameter : Type u_1} →
  [inst : MeasurableSpace Parameter] →
    [inst_1 : NormedAddCommGroup Parameter] →
      [inst_2 : InnerProductSpace ℝ Parameter] →
        [inst_3 : CompleteSpace Parameter] →
          (self : PZMH20PerformativePrediction.Theorem310CorrectedRGDAllRoundTrajectoryData Parameter) →
            (iteration : ℕ) →
              AppliedModelingLib.HeterogeneousBatchTrace
                  (fun (round : ℕ) =>
                    Fin
                      (AppliedModelingLib.pOneFournierGuillinTheorem310SupercriticalCappedConcreteCountSchedule
                        self.dimension self.cutoff self.eta self.alpha self.gamma self.momentBound self.tailBound
                        self.scaledDeviation self.p self.headTolerance round))
                  (EuclideanSpace ℝ (Fin self.dimension)) iteration →
                Parameter

field hmeasurableDeployed:
∀ {Parameter : Type u_1} [inst : MeasurableSpace Parameter] [inst_1 : NormedAddCommGroup Parameter]
  [inst_2 : InnerProductSpace ℝ Parameter] [inst_3 : CompleteSpace Parameter]
  (self : PZMH20PerformativePrediction.Theorem310CorrectedRGDAllRoundTrajectoryData Parameter) (iteration : ℕ),
  Measurable (self.deployedOfHistory iteration)

field hempirical:
∀ {Parameter : Type u_1} [inst : MeasurableSpace Parameter] [inst_1 : NormedAddCommGroup Parameter]
  [inst_2 : InnerProductSpace ℝ Parameter] [inst_3 : CompleteSpace Parameter]
  (self : PZMH20PerformativePrediction.Theorem310CorrectedRGDAllRoundTrajectoryData Parameter) (iteration : ℕ)
  (history :
    AppliedModelingLib.HeterogeneousBatchTrace
      (fun (round : ℕ) =>
        Fin
          (AppliedModelingLib.pOneFournierGuillinTheorem310SupercriticalCappedConcreteCountSchedule self.dimension
            self.cutoff self.eta self.alpha self.gamma self.momentBound self.tailBound self.scaledDeviation self.p
            self.headTolerance round))
      (EuclideanSpace ℝ (Fin self.dimension)) iteration)
  (batch :
    Fin
        (AppliedModelingLib.pOneFournierGuillinTheorem310SupercriticalCappedConcreteCountSchedule self.dimension
          self.cutoff self.eta self.alpha self.gamma self.momentBound self.tailBound self.scaledDeviation self.p
          self.headTolerance iteration) →
      EuclideanSpace ℝ (Fin self.dimension)),
  self.deployedOfHistory (iteration + 1) (AppliedModelingLib.heterogeneousBatchTraceSnoc history batch) =
    AppliedModelingLib.measureSampledGradientDescentUpdate self.gradient self.project self.stepSize
      (self.hcountPositive iteration) batch (self.deployedOfHistory iteration history)

field hintegrableEmpirical:
∀ {Parameter : Type u_1} [inst : MeasurableSpace Parameter] [inst_1 : NormedAddCommGroup Parameter]
  [inst_2 : InnerProductSpace ℝ Parameter] [inst_3 : CompleteSpace Parameter]
  (self : PZMH20PerformativePrediction.Theorem310CorrectedRGDAllRoundTrajectoryData Parameter) (iteration : ℕ)
  (history :
    AppliedModelingLib.HeterogeneousBatchTrace
      (fun (round : ℕ) =>
        Fin
          (AppliedModelingLib.pOneFournierGuillinTheorem310SupercriticalCappedConcreteCountSchedule self.dimension
            self.cutoff self.eta self.alpha self.gamma self.momentBound self.tailBound self.scaledDeviation self.p
            self.headTolerance round))
      (EuclideanSpace ℝ (Fin self.dimension)) iteration)
  (batch :
    Fin
        (AppliedModelingLib.pOneFournierGuillinTheorem310SupercriticalCappedConcreteCountSchedule self.dimension
          self.cutoff self.eta self.alpha self.gamma self.momentBound self.tailBound self.scaledDeviation self.p
          self.headTolerance iteration) →
      EuclideanSpace ℝ (Fin self.dimension)),
  MeasureTheory.Integrable
    (fun (datum : EuclideanSpace ℝ (Fin self.dimension)) =>
      self.gradient datum (self.deployedOfHistory iteration history))
    ↑(AppliedModelingLib.empiricalSampleProbabilityMeasureOfPos (self.hcountPositive iteration) batch)

field stable:
{Parameter : Type u_1} →
  [inst : MeasurableSpace Parameter] →
    [inst_1 : NormedAddCommGroup Parameter] →
      [inst_2 : InnerProductSpace ℝ Parameter] →
        [inst_3 : CompleteSpace Parameter] →
          PZMH20PerformativePrediction.Theorem310CorrectedRGDAllRoundTrajectoryData Parameter → Parameter

field hstablePerformative:
∀ {Parameter : Type u_1} [inst : MeasurableSpace Parameter] [inst_1 : NormedAddCommGroup Parameter]
  [inst_2 : InnerProductSpace ℝ Parameter] [inst_3 : CompleteSpace Parameter]
  (self : PZMH20PerformativePrediction.Theorem310CorrectedRGDAllRoundTrajectoryData Parameter),
  AppliedModelingLib.IsMeasurePerformativelyStableOn self.model self.domain self.stable

field contraction:
{Parameter : Type u_1} →
  [inst : MeasurableSpace Parameter] →
    [inst_1 : NormedAddCommGroup Parameter] →
      [inst_2 : InnerProductSpace ℝ Parameter] →
        [inst_3 : CompleteSpace Parameter] →
          PZMH20PerformativePrediction.Theorem310CorrectedRGDAllRoundTrajectoryData Parameter → ℝ

field outerContraction:
{Parameter : Type u_1} →
  [inst : MeasurableSpace Parameter] →
    [inst_1 : NormedAddCommGroup Parameter] →
      [inst_2 : InnerProductSpace ℝ Parameter] →
        [inst_3 : CompleteSpace Parameter] →
          PZMH20PerformativePrediction.Theorem310CorrectedRGDAllRoundTrajectoryData Parameter → ℝ

field wassersteinBound:
{Parameter : Type u_1} →
  [inst : MeasurableSpace Parameter] →
    [inst_1 : NormedAddCommGroup Parameter] →
      [inst_2 : InnerProductSpace ℝ Parameter] →
        [inst_3 : CompleteSpace Parameter] →
          PZMH20PerformativePrediction.Theorem310CorrectedRGDAllRoundTrajectoryData Parameter → ℝ

field radius:
{Parameter : Type u_1} →
  [inst : MeasurableSpace Parameter] →
    [inst_1 : NormedAddCommGroup Parameter] →
      [inst_2 : InnerProductSpace ℝ Parameter] →
        [inst_3 : CompleteSpace Parameter] →
          PZMH20PerformativePrediction.Theorem310CorrectedRGDAllRoundTrajectoryData Parameter → ℝ

field hcontraction_eq:
∀ {Parameter : Type u_1} [inst : MeasurableSpace Parameter] [inst_1 : NormedAddCommGroup Parameter]
  [inst_2 : InnerProductSpace ℝ Parameter] [inst_3 : CompleteSpace Parameter]
  (self : PZMH20PerformativePrediction.Theorem310CorrectedRGDAllRoundTrajectoryData Parameter),
  self.contraction =
    1 -
      self.stepSize *
        (self.modulus * ↑self.smoothness / (self.modulus + ↑self.smoothness) -
          self.sensitivity * (3 / 2 * self.stepSize * ↑self.smoothness ^ 2 + ↑self.smoothness))

field hcontraction_le_outer:
∀ {Parameter : Type u_1} [inst : MeasurableSpace Parameter] [inst_1 : NormedAddCommGroup Parameter]
  [inst_2 : InnerProductSpace ℝ Parameter] [inst_3 : CompleteSpace Parameter]
  (self : PZMH20PerformativePrediction.Theorem310CorrectedRGDAllRoundTrajectoryData Parameter),
  self.contraction ≤ self.outerContraction

field hbound_nonneg:
∀ {Parameter : Type u_1} [inst : MeasurableSpace Parameter] [inst_1 : NormedAddCommGroup Parameter]
  [inst_2 : InnerProductSpace ℝ Parameter] [inst_3 : CompleteSpace Parameter]
  (self : PZMH20PerformativePrediction.Theorem310CorrectedRGDAllRoundTrajectoryData Parameter),
  0 ≤ self.wassersteinBound

field hbudget:
∀ {Parameter : Type u_1} [inst : MeasurableSpace Parameter] [inst_1 : NormedAddCommGroup Parameter]
  [inst_2 : InnerProductSpace ℝ Parameter] [inst_3 : CompleteSpace Parameter]
  (self : PZMH20PerformativePrediction.Theorem310CorrectedRGDAllRoundTrajectoryData Parameter),
  √↑self.dimension * self.deviation + 2 * √↑self.dimension * (32 / 15 * self.tailBound / 16 ^ self.cutoff) +
      self.headTolerance ≤
    self.wassersteinBound

field herror_outer:
∀ {Parameter : Type u_1} [inst : MeasurableSpace Parameter] [inst_1 : NormedAddCommGroup Parameter]
  [inst_2 : InnerProductSpace ℝ Parameter] [inst_3 : CompleteSpace Parameter]
  (self : PZMH20PerformativePrediction.Theorem310CorrectedRGDAllRoundTrajectoryData Parameter),
  self.stepSize * ↑self.smoothness * self.wassersteinBound ≤ (self.outerContraction - self.contraction) * self.radius

field herror_absorbed:
∀ {Parameter : Type u_1} [inst : MeasurableSpace Parameter] [inst_1 : NormedAddCommGroup Parameter]
  [inst_2 : InnerProductSpace ℝ Parameter] [inst_3 : CompleteSpace Parameter]
  (self : PZMH20PerformativePrediction.Theorem310CorrectedRGDAllRoundTrajectoryData Parameter),
  self.stepSize * ↑self.smoothness * self.wassersteinBound ≤ (1 - self.contraction) * self.radius

field initialDistance:
{Parameter : Type u_1} →
  [inst : MeasurableSpace Parameter] →
    [inst_1 : NormedAddCommGroup Parameter] →
      [inst_2 : InnerProductSpace ℝ Parameter] →
        [inst_3 : CompleteSpace Parameter] →
          PZMH20PerformativePrediction.Theorem310CorrectedRGDAllRoundTrajectoryData Parameter → ℝ

field houter_lt_one:
∀ {Parameter : Type u_1} [inst : MeasurableSpace Parameter] [inst_1 : NormedAddCommGroup Parameter]
  [inst_2 : InnerProductSpace ℝ Parameter] [inst_3 : CompleteSpace Parameter]
  (self : PZMH20PerformativePrediction.Theorem310CorrectedRGDAllRoundTrajectoryData Parameter),
  self.outerContraction < 1

field hradius_pos:
∀ {Parameter : Type u_1} [inst : MeasurableSpace Parameter] [inst_1 : NormedAddCommGroup Parameter]
  [inst_2 : InnerProductSpace ℝ Parameter] [inst_3 : CompleteSpace Parameter]
  (self : PZMH20PerformativePrediction.Theorem310CorrectedRGDAllRoundTrajectoryData Parameter), 0 < self.radius

field hinitialDistance:
∀ {Parameter : Type u_1} [inst : MeasurableSpace Parameter] [inst_1 : NormedAddCommGroup Parameter]
  [inst_2 : InnerProductSpace ℝ Parameter] [inst_3 : CompleteSpace Parameter]
  (self : PZMH20PerformativePrediction.Theorem310CorrectedRGDAllRoundTrajectoryData Parameter)
  (history :
    AppliedModelingLib.HeterogeneousBatchTrace
      (fun (round : ℕ) =>
        Fin
          (AppliedModelingLib.pOneFournierGuillinTheorem310SupercriticalCappedConcreteCountSchedule self.dimension
            self.cutoff self.eta self.alpha self.gamma self.momentBound self.tailBound self.scaledDeviation self.p
            self.headTolerance round))
      (EuclideanSpace ℝ (Fin self.dimension)) 0),
  dist (self.deployedOfHistory 0 history) self.stable ≤ self.initialDistance

field hmoment:
∀ {Parameter : Type u_1} [inst : MeasurableSpace Parameter] [inst_1 : NormedAddCommGroup Parameter]
  [inst_2 : InnerProductSpace ℝ Parameter] [inst_3 : CompleteSpace Parameter]
  (self : PZMH20PerformativePrediction.Theorem310CorrectedRGDAllRoundTrajectoryData Parameter) (iteration : ℕ)
  (history :
    AppliedModelingLib.HeterogeneousBatchTrace
      (fun (round : ℕ) =>
        Fin
          (AppliedModelingLib.pOneFournierGuillinTheorem310SupercriticalCappedConcreteCountSchedule self.dimension
            self.cutoff self.eta self.alpha self.gamma self.momentBound self.tailBound self.scaledDeviation self.p
            self.headTolerance round))
      (EuclideanSpace ℝ (Fin self.dimension)) iteration),
  AppliedModelingLib.Probability.HasBoundedExponentialRadialMoment
    (self.model.dataLaw (self.deployedOfHistory iteration history)) self.alpha self.gamma self.momentBound

field hevent:
∀ {Parameter : Type u_1} [inst : MeasurableSpace Parameter] [inst_1 : NormedAddCommGroup Parameter]
  [inst_2 : InnerProductSpace ℝ Parameter] [inst_3 : CompleteSpace Parameter]
  (self : PZMH20PerformativePrediction.Theorem310CorrectedRGDAllRoundTrajectoryData Parameter) (iteration : ℕ),
  MeasurableSet
    (AppliedModelingLib.heterogeneousBatchTraceBadEventPair
      (fun (round : ℕ) =>
        Fin
          (AppliedModelingLib.pOneFournierGuillinTheorem310SupercriticalCappedConcreteCountSchedule self.dimension
            self.cutoff self.eta self.alpha self.gamma self.momentBound self.tailBound self.scaledDeviation self.p
            self.headTolerance round))
      (AppliedModelingLib.pOneFournierGuillinAdaptiveEuclideanSelectedShellPartitionCertificateBadEvent self.dimension
        self.cutoff self.eta
        (AppliedModelingLib.pOneFournierGuillinTheorem310CappedDeviation self.dimension self.eta self.alpha self.gamma
          self.momentBound self.scaledDeviation)
        (AppliedModelingLib.pOneFournierGuillinTheorem310SupercriticalConfidenceSchedule self.p)
        (AppliedModelingLib.pOneFournierGuillinTheorem310SupercriticalCappedConcreteCountSchedule self.dimension
          self.cutoff self.eta self.alpha self.gamma self.momentBound self.tailBound self.scaledDeviation self.p
          self.headTolerance)
        (AppliedModelingLib.measurePerformativeDeployedLaw self.model
          (fun (round : ℕ) =>
            Fin
              (AppliedModelingLib.pOneFournierGuillinTheorem310SupercriticalCappedConcreteCountSchedule self.dimension
                self.cutoff self.eta self.alpha self.gamma self.momentBound self.tailBound self.scaledDeviation self.p
                self.headTolerance round))
          self.deployedOfHistory))
      iteration)
\end{ReviewVerbatim}

\paragraph{Direct retained dependencies}
\begin{itemize}\raggedright
\item \hyperlink{governing-declaration-307f711c4bf618c0}{Applied\allowbreak{}Modeling\allowbreak{}Lib.\allowbreak{}Heterogeneous\allowbreak{}Batch\allowbreak{}Trace}
\item \hyperlink{governing-declaration-37d7ae049de215a7}{Applied\allowbreak{}Modeling\allowbreak{}Lib.\allowbreak{}Heterogeneous\allowbreak{}Batch\allowbreak{}Trace\allowbreak{}Coordinate}
\item \hyperlink{library-prerequisite-2e05b739f5b71f23}{Applied\allowbreak{}Modeling\allowbreak{}Lib.\allowbreak{}Is\allowbreak{}Jointly\allowbreak{}Smooth\allowbreak{}Gradient}
\item \hyperlink{governing-declaration-522b06e211b1e40a}{Applied\allowbreak{}Modeling\allowbreak{}Lib.\allowbreak{}Is\allowbreak{}Measure\allowbreak{}Performatively\allowbreak{}Stable\allowbreak{}On}
\item \hyperlink{governing-declaration-41e49a877b22dbe7}{Applied\allowbreak{}Modeling\allowbreak{}Lib.\allowbreak{}Is\allowbreak{}Variational\allowbreak{}Euclidean\allowbreak{}Projection\allowbreak{}On}
\item \hyperlink{library-prerequisite-96553d4fbb82f1ac}{Applied\allowbreak{}Modeling\allowbreak{}Lib.\allowbreak{}Measure\allowbreak{}Performative\allowbreak{}Model}
\item \hyperlink{governing-declaration-527443e0fb077287}{Applied\allowbreak{}Modeling\allowbreak{}Lib.\allowbreak{}Measure\allowbreak{}Performative\allowbreak{}Sampling\allowbreak{}Kernel}
\item \hyperlink{governing-declaration-9604ed56aea35d91}{Applied\allowbreak{}Modeling\allowbreak{}Lib.\allowbreak{}Probability.\allowbreak{}Has\allowbreak{}Bounded\allowbreak{}Exponential\allowbreak{}Radial\allowbreak{}Moment}
\item \hyperlink{governing-declaration-93c495b865c4b878}{Applied\allowbreak{}Modeling\allowbreak{}Lib.\allowbreak{}Probability.\allowbreak{}fifth\allowbreak{}Order\allowbreak{}Exponential\allowbreak{}Radial\allowbreak{}Tail\allowbreak{}Constant}
\item \hyperlink{governing-declaration-cbd1f5f969456e5f}{Applied\allowbreak{}Modeling\allowbreak{}Lib.\allowbreak{}empirical\allowbreak{}Sample\allowbreak{}Probability\allowbreak{}Measure\allowbreak{}Of\allowbreak{}Pos}
\item \hyperlink{governing-declaration-68ef005103543934}{Applied\allowbreak{}Modeling\allowbreak{}Lib.\allowbreak{}heterogeneous\allowbreak{}Batch\allowbreak{}Trace\allowbreak{}Bad\allowbreak{}Event\allowbreak{}Pair}
\item \hyperlink{governing-declaration-92d1582063c434f3}{Applied\allowbreak{}Modeling\allowbreak{}Lib.\allowbreak{}heterogeneous\allowbreak{}Batch\allowbreak{}Trace\allowbreak{}Coordinate\allowbreak{}Measurable\allowbreak{}Space}
\item \hyperlink{governing-declaration-caf4e7d2e1526646}{Applied\allowbreak{}Modeling\allowbreak{}Lib.\allowbreak{}heterogeneous\allowbreak{}Batch\allowbreak{}Trace\allowbreak{}Measurable\allowbreak{}Space}
\item \hyperlink{governing-declaration-73f8e033eb247305}{Applied\allowbreak{}Modeling\allowbreak{}Lib.\allowbreak{}heterogeneous\allowbreak{}Batch\allowbreak{}Trace\allowbreak{}Snoc}
\item \hyperlink{governing-declaration-7bc25db7b9b1aaf4}{Applied\allowbreak{}Modeling\allowbreak{}Lib.\allowbreak{}measure\allowbreak{}Performative\allowbreak{}Deployed\allowbreak{}Law}
\item \hyperlink{governing-declaration-46a8b190a3977782}{Applied\allowbreak{}Modeling\allowbreak{}Lib.\allowbreak{}measure\allowbreak{}Sampled\allowbreak{}Gradient\allowbreak{}Descent\allowbreak{}Update}
\item \hyperlink{governing-declaration-30e7d69ebe63768c}{Applied\allowbreak{}Modeling\allowbreak{}Lib.\allowbreak{}p\allowbreak{}One\allowbreak{}Fournier\allowbreak{}Guillin\allowbreak{}Adaptive\allowbreak{}Euclidean\allowbreak{}Selected\allowbreak{}Shell\allowbreak{}Partition\allowbreak{}Certificate\allowbreak{}Bad\allowbreak{}Event}
\item \hyperlink{governing-declaration-25a90c9551dbd2b1}{Applied\allowbreak{}Modeling\allowbreak{}Lib.\allowbreak{}p\allowbreak{}One\allowbreak{}Fournier\allowbreak{}Guillin\allowbreak{}Theorem310\allowbreak{}Capped\allowbreak{}Deviation}
\item \hyperlink{governing-declaration-3cf6b14c7bcca80c}{Applied\allowbreak{}Modeling\allowbreak{}Lib.\allowbreak{}p\allowbreak{}One\allowbreak{}Fournier\allowbreak{}Guillin\allowbreak{}Theorem310\allowbreak{}Supercritical\allowbreak{}Capped\allowbreak{}Concrete\allowbreak{}Count\allowbreak{}Schedule}
\item \hyperlink{governing-declaration-6a2f3c5b322d78da}{Applied\allowbreak{}Modeling\allowbreak{}Lib.\allowbreak{}p\allowbreak{}One\allowbreak{}Fournier\allowbreak{}Guillin\allowbreak{}Theorem310\allowbreak{}Supercritical\allowbreak{}Confidence\allowbreak{}Schedule}
\item \hyperlink{governing-declaration-11c9ab6e3c6f528b}{PZMH20\allowbreak{}Performative\allowbreak{}Prediction.\allowbreak{}Is\allowbreak{}Measure\allowbreak{}Pointwise\allowbreak{}Gradient\allowbreak{}Strongly\allowbreak{}Convex}
\item \hyperlink{paper-prerequisite-2ea2f5769822c5c4}{PZMH20\allowbreak{}Performative\allowbreak{}Prediction.\allowbreak{}Is\allowbreak{}Measure\allowbreak{}Wasserstein\allowbreak{}Sensitive}
\end{itemize}
\paragraph{Recorded source-to-declaration screening}
\textbf{Verdict:} \texttt{matches\_approved\_corrected\_target}
\\
{\raggedright
\textbf{Reason:} The expanded carrier matches the disclosed corrected REGD premise bundle:\allowbreak{} dimension greater than two, the capped count and inverse-\allowbreak{}square failure schedules, a uniform adaptive exponential-\allowbreak{}moment/\allowbreak{}tail envelope, measurable sampling and bad events, globally scoped gradient regularity and integrability, the convex-\allowbreak{}domain variational nonexpansive projection, the source step-\allowbreak{}size and sensitivity inequalities, the empirical REGD recurrence, a supplied stable point, and the stated contraction, W1-\allowbreak{}budget, absorption, radius, and initial-\allowbreak{}distance conditions.\allowbreak{} Its graph-\allowbreak{}bound use derives a deterministic finite entry iteration and all-\allowbreak{}subsequent-\allowbreak{}round containment with probability at least 1 -\allowbreak{} p;\allowbreak{} no entry bound or small-\allowbreak{}deviation gate is supplied by the caller.\allowbreak{}
\par}
\reviewmatch{Matches formalized target}
\noindent\textbf{Reviewer annotation}\par
\reviewerbox

\clearpage
\hypertarget{source-claim-5b4f70cb8ea743cb}{}
\hypertarget{source-section-0c7af8265958a9ea}{}
\section*{Main-text source claims}
\subsection*{Review row 1}
\textbf{Source locator:} {\footnotesize\raggedright cited publication, lines 264-\allowbreak{}269\par}
\paragraph{Verbatim source input}
\begin{ReviewVerbatim}
In addition to this assumption on the distribution map, we          (b) If ε < βγ , the iterates θt of RRM converge to a unique
will often make standard assumptions on the loss function               performatively stable point θPS at a linear rate,
`(z; θ) which hold for broad classes of losses. Each tech-
nical result to follow will invoke some subset of them. To                                                log( 1δ kθ0 − θPS k2 )
                                 def                                         kθt − θPS k2 6 δ for t >                              .
simplify our presentation, let Z = ∪θ∈Θ supp(D(θ)).                                                             (1 − ε βγ )

[Next verbatim source excerpt]

Theorem 3.5. Suppose that the loss `(z; θ) is β-jointly
                                                                   Proposition 3.6 suggests a fundamental difference between
smooth (A1) and γ-strongly convex (A2). If the distribution
                                                                   strong and weak convexity in our framing of performative
map D(·) is ε-sensitive, then the following is true:
                                                                   prediction (weak meaning γ = 0). In supervised learn-
(a) kG(θ) − G(θ0 )k2 6 ε βγ kθ − θ0 k2 , ∀θ, θ0 ∈ Θ.               ing, using strongly convex losses generally guarantees a

[form-feed in source extraction]
                                                           Performative Prediction

faster rate of optimization, yet asymptotically, the solution                 3.4. Finite-Sample Analysis
achieved with either strongly or weakly convex losses is
                                                                              We now extend our main results regarding the convergence
globally optimal. However, in our framework, strong con-
                                                                              of RRM and RGD to the finite-sample regime. To do so,
vexity is in fact necessary to guarantee convergence of re-
                                                                              we leverage the fact that, under mild regularity conditions,
peated risk minimization, even for arbitrarily smooth losses
                                                                              the empirical distribution Dn given by n samples drawn
and an arbitrarily small sensitivity parameter.
                                                                              i.i.d. from a true distribution D is with high probability
                                                                              close to D in earth mover’s distance (Fournier & Guillin,
\end{ReviewVerbatim}


\paragraph{Formalized review target}
The archival source above is not asserted equivalent to this formalized target.
\begin{ReviewVerbatim}
For an arbitrary-law measure-performative model on a complete real inner-product ambient parameter space, let the source parameter domain be closed and convex and let the repeated-risk-minimization update minimize frozen risk on that domain. Assume expected-loss and gradient integrability for every ambient deployed/evaluated parameter pair; A1 joint gradient smoothness and A2 pointwise strong convexity for every ambient parameter and every datum in the chosen data carrier; nonnegative Wasserstein sensitivity for every ambient parameter pair; local loss measurability and pointwise parameter differentiability for every ambient parameter and datum; and strict epsilon beta over gamma below one. Then, for each initial point in the domain, the update has the displayed epsilon beta over gamma contraction on the domain and there exists a domain-stable point, unique among domain-stable points, such that the iterates from that initial point converge to it with the displayed geometric rate.
\end{ReviewVerbatim}

\textbf{Archival source anchor:} {\footnotesize\raggedright cited publication, lines 264-\allowbreak{}330\par}
\paragraph{Expanded PaperInterface specification}
\begin{ReviewVerbatim}
∀ {Parameter : Type u_1} {Data : Type u_2} [inst : NormedAddCommGroup Parameter]
  [inst_1 : InnerProductSpace ℝ Parameter] [inst_2 : CompleteSpace Parameter] [inst_3 : MeasurableSpace Data]
  [inst_4 : MetricSpace Data] (model : AppliedModelingLib.MeasurePerformativeModel Parameter Data)
  (gradient : Data → Parameter → Parameter) {smoothness : NNReal},
  AppliedModelingLib.IsJointlySmoothGradient gradient smoothness →
    (∀ (distributionParameter evaluatedParameter : Parameter),
        MeasureTheory.Integrable (fun (datum : Data) => gradient datum evaluatedParameter)
          ↑(model.dataLaw distributionParameter)) →
      ∀ {sensitivity modulus : ℝ},
        0 ≤ sensitivity →
          0 < modulus →
            PZMH20PerformativePrediction.IsMeasureWassersteinSensitive model sensitivity →
              ∀ (domain : Set Parameter),
                IsClosed domain →
                  Convex ℝ domain →
                    ∀ (update : Parameter → Parameter),
                      PZMH20PerformativePrediction.IsMeasureRRMOn model domain update →
                        (∀ (distributionParameter evaluatedParameter : Parameter),
                            ∀ᶠ (parameter : Parameter) in nhds evaluatedParameter,
                              MeasureTheory.AEStronglyMeasurable (fun (datum : Data) => model.loss datum parameter)
                                ↑(model.dataLaw distributionParameter)) →
                          (∀ (datum : Data) (parameter : Parameter),
                              HasGradientAt (fun (theta : Parameter) => model.loss datum theta)
                                (gradient datum parameter) parameter) →
                            PZMH20PerformativePrediction.IsMeasurePointwiseGradientStronglyConvex model gradient
                                modulus →
                              sensitivity * ↑smoothness / modulus < 1 →
                                ∀ initial ∈ domain,
                                  (∀ first ∈ domain,
                                      ∀ second ∈ domain,
                                        ‖update first - update second‖ ≤
                                          sensitivity * ↑smoothness / modulus * ‖first - second‖) ∧
                                    ∃ stable ∈ domain,
                                      PZMH20PerformativePrediction.IsMeasureStableOn model domain stable ∧
                                        (∀ other ∈ domain,
                                            PZMH20PerformativePrediction.IsMeasureStableOn model domain other →
                                              other = stable) ∧
                                          Filter.Tendsto (fun (iteration : ℕ) => update^[iteration] initial)
                                              Filter.atTop (nhds stable) ∧
                                            ∀ (iteration : ℕ),
                                              ‖update^[iteration] initial - stable‖ ≤
                                                (sensitivity * ↑smoothness / modulus) ^ iteration * ‖initial - stable‖
\end{ReviewVerbatim}

\paragraph{Retained governing declarations}
\begin{itemize}\raggedright
\item \hyperlink{library-prerequisite-2e05b739f5b71f23}{Applied\allowbreak{}Modeling\allowbreak{}Lib.\allowbreak{}Is\allowbreak{}Jointly\allowbreak{}Smooth\allowbreak{}Gradient}
\item \hyperlink{library-prerequisite-96553d4fbb82f1ac}{Applied\allowbreak{}Modeling\allowbreak{}Lib.\allowbreak{}Measure\allowbreak{}Performative\allowbreak{}Model}
\item \hyperlink{governing-declaration-11c9ab6e3c6f528b}{PZMH20\allowbreak{}Performative\allowbreak{}Prediction.\allowbreak{}Is\allowbreak{}Measure\allowbreak{}Pointwise\allowbreak{}Gradient\allowbreak{}Strongly\allowbreak{}Convex}
\item \hyperlink{paper-prerequisite-28cbc41dc9d12634}{PZMH20\allowbreak{}Performative\allowbreak{}Prediction.\allowbreak{}Is\allowbreak{}Measure\allowbreak{}RRM\allowbreak{}On}
\item \hyperlink{governing-declaration-d4ecd5e34f1d7939}{PZMH20\allowbreak{}Performative\allowbreak{}Prediction.\allowbreak{}Is\allowbreak{}Measure\allowbreak{}Stable\allowbreak{}On}
\item \hyperlink{paper-prerequisite-2ea2f5769822c5c4}{PZMH20\allowbreak{}Performative\allowbreak{}Prediction.\allowbreak{}Is\allowbreak{}Measure\allowbreak{}Wasserstein\allowbreak{}Sensitive}
\end{itemize}
\paragraph{Lean-elaborated claim roles}\begin{ReviewVerbatim}
1. parameter: Type u_1
2. parameter: Type u_2
3. parameter: NormedAddCommGroup Parameter
4. parameter: InnerProductSpace ℝ Parameter
5. assumption: CompleteSpace Parameter
6. parameter: MeasurableSpace Data
7. parameter: MetricSpace Data
8. parameter: AppliedModelingLib.MeasurePerformativeModel Parameter Data
9. parameter: Data → Parameter → Parameter
10. parameter: NNReal
11. assumption: AppliedModelingLib.IsJointlySmoothGradient gradient smoothness
12. assumption: ∀ (distributionParameter evaluatedParameter : Parameter),
  MeasureTheory.Integrable (fun (datum : Data) => gradient datum evaluatedParameter)
    ↑(model.dataLaw distributionParameter)
13. parameter: ℝ
14. parameter: ℝ
15. assumption: 0 ≤ sensitivity
16. assumption: 0 < modulus
17. assumption: PZMH20PerformativePrediction.IsMeasureWassersteinSensitive model sensitivity
18. parameter: Set Parameter
19. assumption: IsClosed domain
20. assumption: Convex ℝ domain
21. parameter: Parameter → Parameter
22. assumption: PZMH20PerformativePrediction.IsMeasureRRMOn model domain update
23. assumption: ∀ (distributionParameter evaluatedParameter : Parameter),
  ∀ᶠ (parameter : Parameter) in nhds evaluatedParameter,
    MeasureTheory.AEStronglyMeasurable (fun (datum : Data) => model.loss datum parameter)
      ↑(model.dataLaw distributionParameter)
24. assumption: ∀ (datum : Data) (parameter : Parameter),
  HasGradientAt (fun (theta : Parameter) => model.loss datum theta) (gradient datum parameter) parameter
25. assumption: PZMH20PerformativePrediction.IsMeasurePointwiseGradientStronglyConvex model gradient modulus
26. assumption: sensitivity * ↑smoothness / modulus < 1
27. parameter: Parameter
28. assumption: initial ∈ domain
29. conclusion: (∀ first ∈ domain,
    ∀ second ∈ domain, ‖update first - update second‖ ≤ sensitivity * ↑smoothness / modulus * ‖first - second‖) ∧
  ∃ stable ∈ domain,
    PZMH20PerformativePrediction.IsMeasureStableOn model domain stable ∧
      (∀ other ∈ domain, PZMH20PerformativePrediction.IsMeasureStableOn model domain other → other = stable) ∧
        Filter.Tendsto (fun (iteration : ℕ) => update^[iteration] initial) Filter.atTop (nhds stable) ∧
          ∀ (iteration : ℕ),
            ‖update^[iteration] initial - stable‖ ≤
              (sensitivity * ↑smoothness / modulus) ^ iteration * ‖initial - stable‖
\end{ReviewVerbatim}

\paragraph{Lean proof endpoint (built separately)}\begin{ReviewVerbatim}
PZMH20PerformativePrediction.measureWassersteinConvexRRMConvergesLinearlyProof
\end{ReviewVerbatim}

\paragraph{Recorded source-to-Spec screening}
\textbf{Verdict:} \texttt{matches\_approved\_corrected\_target}
\\
{\raggedright
\textbf{Reason:} Compared Theorem 3.\allowbreak{}5 at cited publication:\allowbreak{}264-\allowbreak{}269 and 324-\allowbreak{}344 under the bound corrected target for theorem-\allowbreak{}3.\allowbreak{}5-\allowbreak{}general-\allowbreak{}measure-\allowbreak{}analytic-\allowbreak{}bridge with the complete expanded target, roles 1-\allowbreak{}29, and every listed support display.\allowbreak{} The probability-\allowbreak{}law model exposes expected-\allowbreak{}loss integrability for every ambient deployed/\allowbreak{}evaluated pair;\allowbreak{} the remaining premises expose global gradient integrability, local loss measurability, pointwise differentiation, both beta-\allowbreak{}gradient Lipschitz inequalities, gamma-\allowbreak{}strong convexity, and nonnegative epsilon sensitivity.\allowbreak{} The Wasserstein predicate expands to a finite expected-\allowbreak{}distance coupling witness with prescribed marginals and infimum cost bounded by epsilon times parameter distance.\allowbreak{} These universal ambient/\allowbreak{}all-\allowbreak{}data conditions agree with the corrected scope, rather than being inferred from conditions merely on Theta or Z.\allowbreak{} The RRM predicate maps each deployed point in the closed convex domain back into that domain and minimizes its frozen expected risk over all domain candidates.\allowbreak{} The selected initial point supplies nonemptiness.\allowbreak{} With gamma > 0 and q = epsilon beta/\allowbreak{}gamma < 1, the conclusion states the q contraction on the domain, a domain-\allowbreak{}stable minimizer, uniqueness among all domain-\allowbreak{}stable points, convergence of the selected orbit, and the exact q\textasciicircum{}t geometric bound, which yields the source logarithmic time bound for 0 <= q < 1, including immediate convergence when q = 0.\allowbreak{} Requiring the strict factor for this combined endpoint is explicit in the corrected target;\allowbreak{} archival equivalence is not asserted.\allowbreak{}
\par}
\reviewmatch{Matches formalized target}
\noindent\textbf{Reviewer annotation}\par
\reviewerbox
\clearpage
\hypertarget{source-claim-d10117a376c0a20b}{}
\section*{Review row 2}
\textbf{Source locator:} {\footnotesize\raggedright cited publication, lines 373-\allowbreak{}415\par}
\paragraph{Verbatim source input}
\begin{ReviewVerbatim}
Theorem 3.10. Suppose that the loss `(z; θ) is β-jointly
Theorem 3.8. Suppose that the loss `(z; θ) is β-jointly                       smooth (A1) and γ-strongly convex (A2), and that there exist
smooth (A1) and γ-strongly convex (A2). If the distribution                                                def R        α
                                        γ                                     α > 1, µ > 0 such that ξα,µ = Rm eµ|x| dD(θ) is finite
map D(·) is ε-sensitive with ε < (β+γ)(1+1.5ηβ) , then RGD
                                                                              ∀θ ∈ Θ. Let δ ∈ (0, 1) be a radius of convergence.     Con-
                     2
with step size η 6 β+γ satisfies the following:
                                                                                                                                    � t
                                                                                                                            1
                                                                              sider running RERM or RGD with nt = O (εδ)m log p
(a) kGgd (θ) − Ggd (θ0 )k2                                                    samples at time t.
                                                                                                                            γ
                                                                              (a) If the map D(·) is ε-sensitive with ε < 2β  , then with
                         βγ                                          0
      6       1−η           − ε(1.5ηβ 2 + β)             kθ − θ k2 .              probability 1 − p, RERM satisfies,
                        β+γ
                                                                                                                  log 1δ kθ0 − θPS k2
                                                                                                                     �
(b) The iterates θt of RGD converge to a unique performa-                           kθt − θPS k2 6 δ, for all t >                       .
    tively stable point θPS at a linear rate,                                                                           1 − 2εβ
                                                                                                                              γ

                                                                                                                                 γ
                                               �1                             (b) If the map D(·) is ε-sensitive with ε < (β+γ)(1+1.5ηβ)  ,
                                         log    δ kθ0 − θPS k2
     kθt − θPS k2 6 δ for t >                                             .       then with probability 1 − p, REGD with satisfies,
                                          βγ          2
                                 η       β+γ − ε(1.5ηβ + β)
                                                                                                                    log 1δ kθ0 − θPS k2
                                                                                                                        �
                                                                                  kθt −θPS k2 6 δ, for all t >                              ,
The conclusion of Theorem 3.8 is a strict generalization of a                                                        βγ
                                                                                                                 η β+γ    − ε(3ηβ 2 + 2β)
classical optimization result which considers a static objec-

                                                        βγ
tive, in which case the rate of contraction is 1 − η β+γ                                                                   2
                                                                                   for a constant choice of step size η 6 β+γ .
(see for example Theorem 2.1.15 in Nesterov (2013) or
Lemma 3.7 in Hardt et al. (2016b)). Our rate exactly                          Proof sketch. The basic idea behind these results is the fol-
matches this standard result in the case that ε = 0. The                      lowing. While kθt − θPS k2 > δ, the sample size nt is
proof of Theorem 3.8 can be found in Appendix E.3.                            sufficiently large to ensure a behavior similar to that on a
\end{ReviewVerbatim}


\paragraph{Formalized review target}
The archival source above is not asserted equivalent to this formalized target.
\begin{ReviewVerbatim}
For an arbitrary-law measure-performative model on a complete real inner-product ambient parameter space, let the source parameter domain be nonempty, closed, and convex with a variational nonexpansive projection. Assume expected-loss and gradient integrability for every ambient deployed/evaluated parameter pair; A1 joint gradient smoothness and A2 pointwise strong convexity for every ambient parameter and every datum in the chosen data carrier; nonnegative Wasserstein sensitivity for every ambient parameter pair; local loss measurability and pointwise parameter differentiability for every ambient parameter and datum; and the exact inequalities 0 < gamma <= beta, 0 < eta <= 2 / (gamma + beta), and epsilon < gamma / ((gamma + beta) * (1 + 3 * eta * beta / 2)). Then repeated projected gradient descent has the displayed contraction on the domain, a unique domain-stable point, convergence from every domain initial point, and the displayed geometric rate.
\end{ReviewVerbatim}

\textbf{Archival source anchor:} {\footnotesize\raggedright cited publication, lines 373-\allowbreak{}415\par}
\paragraph{Expanded PaperInterface specification}
\begin{ReviewVerbatim}
∀ {Parameter : Type u_1} {Data : Type u_2} [inst : NormedAddCommGroup Parameter]
  [inst_1 : InnerProductSpace ℝ Parameter] [inst_2 : CompleteSpace Parameter] [inst_3 : MeasurableSpace Data]
  [inst_4 : MetricSpace Data] (model : AppliedModelingLib.MeasurePerformativeModel Parameter Data)
  (gradient : Data → Parameter → Parameter) {smoothness : NNReal},
  AppliedModelingLib.IsJointlySmoothGradient gradient smoothness →
    (∀ (distributionParameter evaluatedParameter : Parameter),
        MeasureTheory.Integrable (fun (datum : Data) => gradient datum evaluatedParameter)
          ↑(model.dataLaw distributionParameter)) →
      ∀ {sensitivity modulus : ℝ},
        0 ≤ sensitivity →
          PZMH20PerformativePrediction.IsMeasureWassersteinSensitive model sensitivity →
            0 < modulus →
              modulus ≤ ↑smoothness →
                PZMH20PerformativePrediction.IsMeasurePointwiseGradientStronglyConvex model gradient modulus →
                  ∀ (domain : Set Parameter),
                    IsClosed domain →
                      domain.Nonempty →
                        Convex ℝ domain →
                          ∀ (project : Parameter → Parameter),
                            AppliedModelingLib.IsVariationalEuclideanProjectionOn domain project →
                              LipschitzWith 1 project →
                                ∀ (stepSize : ℝ),
                                  0 < stepSize →
                                    stepSize ≤ 2 / (modulus + ↑smoothness) →
                                      sensitivity <
                                          modulus / ((modulus + ↑smoothness) * (1 + 3 / 2 * stepSize * ↑smoothness)) →
                                        (∀ (distributionParameter evaluatedParameter : Parameter),
                                            ∀ᶠ (parameter : Parameter) in nhds evaluatedParameter,
                                              MeasureTheory.AEStronglyMeasurable
                                                (fun (datum : Data) => model.loss datum parameter)
                                                ↑(model.dataLaw distributionParameter)) →
                                          (∀ (datum : Data) (parameter : Parameter),
                                              HasGradientAt (fun (theta : Parameter) => model.loss datum theta)
                                                (gradient datum parameter) parameter) →
                                            ∀ initial ∈ domain,
                                              (∀ first ∈ domain,
                                                  ∀ second ∈ domain,
                                                    ‖AppliedModelingLib.measureRepeatedGradientDescentUpdate model
                                                            gradient project stepSize first -
                                                          AppliedModelingLib.measureRepeatedGradientDescentUpdate model
                                                            gradient project stepSize second‖ ≤
                                                      (1 -
                                                          stepSize *
                                                            (modulus * ↑smoothness / (modulus + ↑smoothness) -
                                                              sensitivity *
                                                                (3 / 2 * stepSize * ↑smoothness ^ 2 + ↑smoothness))) *
                                                        ‖first - second‖) ∧
                                                ∃ stable ∈ domain,
                                                  PZMH20PerformativePrediction.IsMeasureStableOn model domain stable ∧
                                                    (∀ other ∈ domain,
                                                        PZMH20PerformativePrediction.IsMeasureStableOn model domain
                                                            other →
                                                          other = stable) ∧
                                                      Filter.Tendsto
                                                          (fun (iteration : ℕ) =>
                                                            (AppliedModelingLib.measureRepeatedGradientDescentUpdate
                                                                  model gradient project stepSize)^[iteration]
                                                              initial)
                                                          Filter.atTop (nhds stable) ∧
                                                        ∀ (iteration : ℕ),
                                                          ‖(AppliedModelingLib.measureRepeatedGradientDescentUpdate
                                                                      model gradient project stepSize)^[iteration]
                                                                  initial -
                                                                stable‖ ≤
                                                            (1 -
                                                                  stepSize *
                                                                    (modulus * ↑smoothness / (modulus + ↑smoothness) -
                                                                      sensitivity *
                                                                        (3 / 2 * stepSize * ↑smoothness ^ 2 +
                                                                          ↑smoothness))) ^
                                                                iteration *
                                                              ‖initial - stable‖
\end{ReviewVerbatim}

\paragraph{Retained governing declarations}
\begin{itemize}\raggedright
\item \hyperlink{library-prerequisite-2e05b739f5b71f23}{Applied\allowbreak{}Modeling\allowbreak{}Lib.\allowbreak{}Is\allowbreak{}Jointly\allowbreak{}Smooth\allowbreak{}Gradient}
\item \hyperlink{governing-declaration-41e49a877b22dbe7}{Applied\allowbreak{}Modeling\allowbreak{}Lib.\allowbreak{}Is\allowbreak{}Variational\allowbreak{}Euclidean\allowbreak{}Projection\allowbreak{}On}
\item \hyperlink{library-prerequisite-96553d4fbb82f1ac}{Applied\allowbreak{}Modeling\allowbreak{}Lib.\allowbreak{}Measure\allowbreak{}Performative\allowbreak{}Model}
\item \hyperlink{governing-declaration-836bd13b5596fdaf}{Applied\allowbreak{}Modeling\allowbreak{}Lib.\allowbreak{}measure\allowbreak{}Repeated\allowbreak{}Gradient\allowbreak{}Descent\allowbreak{}Update}
\item \hyperlink{governing-declaration-11c9ab6e3c6f528b}{PZMH20\allowbreak{}Performative\allowbreak{}Prediction.\allowbreak{}Is\allowbreak{}Measure\allowbreak{}Pointwise\allowbreak{}Gradient\allowbreak{}Strongly\allowbreak{}Convex}
\item \hyperlink{governing-declaration-d4ecd5e34f1d7939}{PZMH20\allowbreak{}Performative\allowbreak{}Prediction.\allowbreak{}Is\allowbreak{}Measure\allowbreak{}Stable\allowbreak{}On}
\item \hyperlink{paper-prerequisite-2ea2f5769822c5c4}{PZMH20\allowbreak{}Performative\allowbreak{}Prediction.\allowbreak{}Is\allowbreak{}Measure\allowbreak{}Wasserstein\allowbreak{}Sensitive}
\end{itemize}
\paragraph{Lean-elaborated claim roles}\begin{ReviewVerbatim}
1. parameter: Type u_1
2. parameter: Type u_2
3. parameter: NormedAddCommGroup Parameter
4. parameter: InnerProductSpace ℝ Parameter
5. assumption: CompleteSpace Parameter
6. parameter: MeasurableSpace Data
7. parameter: MetricSpace Data
8. parameter: AppliedModelingLib.MeasurePerformativeModel Parameter Data
9. parameter: Data → Parameter → Parameter
10. parameter: NNReal
11. assumption: AppliedModelingLib.IsJointlySmoothGradient gradient smoothness
12. assumption: ∀ (distributionParameter evaluatedParameter : Parameter),
  MeasureTheory.Integrable (fun (datum : Data) => gradient datum evaluatedParameter)
    ↑(model.dataLaw distributionParameter)
13. parameter: ℝ
14. parameter: ℝ
15. assumption: 0 ≤ sensitivity
16. assumption: PZMH20PerformativePrediction.IsMeasureWassersteinSensitive model sensitivity
17. assumption: 0 < modulus
18. assumption: modulus ≤ ↑smoothness
19. assumption: PZMH20PerformativePrediction.IsMeasurePointwiseGradientStronglyConvex model gradient modulus
20. parameter: Set Parameter
21. assumption: IsClosed domain
22. assumption: domain.Nonempty
23. assumption: Convex ℝ domain
24. parameter: Parameter → Parameter
25. assumption: AppliedModelingLib.IsVariationalEuclideanProjectionOn domain project
26. assumption: LipschitzWith 1 project
27. parameter: ℝ
28. assumption: 0 < stepSize
29. assumption: stepSize ≤ 2 / (modulus + ↑smoothness)
30. assumption: sensitivity < modulus / ((modulus + ↑smoothness) * (1 + 3 / 2 * stepSize * ↑smoothness))
31. assumption: ∀ (distributionParameter evaluatedParameter : Parameter),
  ∀ᶠ (parameter : Parameter) in nhds evaluatedParameter,
    MeasureTheory.AEStronglyMeasurable (fun (datum : Data) => model.loss datum parameter)
      ↑(model.dataLaw distributionParameter)
32. assumption: ∀ (datum : Data) (parameter : Parameter),
  HasGradientAt (fun (theta : Parameter) => model.loss datum theta) (gradient datum parameter) parameter
33. parameter: Parameter
34. assumption: initial ∈ domain
35. conclusion: (∀ first ∈ domain,
    ∀ second ∈ domain,
      ‖AppliedModelingLib.measureRepeatedGradientDescentUpdate model gradient project stepSize first -
            AppliedModelingLib.measureRepeatedGradientDescentUpdate model gradient project stepSize second‖ ≤
        (1 -
            stepSize *
              (modulus * ↑smoothness / (modulus + ↑smoothness) -
                sensitivity * (3 / 2 * stepSize * ↑smoothness ^ 2 + ↑smoothness))) *
          ‖first - second‖) ∧
  ∃ stable ∈ domain,
    PZMH20PerformativePrediction.IsMeasureStableOn model domain stable ∧
      (∀ other ∈ domain, PZMH20PerformativePrediction.IsMeasureStableOn model domain other → other = stable) ∧
        Filter.Tendsto
            (fun (iteration : ℕ) =>
              (AppliedModelingLib.measureRepeatedGradientDescentUpdate model gradient project stepSize)^[iteration]
                initial)
            Filter.atTop (nhds stable) ∧
          ∀ (iteration : ℕ),
            ‖(AppliedModelingLib.measureRepeatedGradientDescentUpdate model gradient project stepSize)^[iteration]
                    initial -
                  stable‖ ≤
              (1 -
                    stepSize *
                      (modulus * ↑smoothness / (modulus + ↑smoothness) -
                        sensitivity * (3 / 2 * stepSize * ↑smoothness ^ 2 + ↑smoothness))) ^
                  iteration *
                ‖initial - stable‖
\end{ReviewVerbatim}

\paragraph{Lean proof endpoint (built separately)}\begin{ReviewVerbatim}
PZMH20PerformativePrediction.measureWassersteinConvexRGDConvergesLinearlyProof
\end{ReviewVerbatim}

\paragraph{Recorded source-to-Spec screening}
\textbf{Verdict:} \texttt{matches\_approved\_corrected\_target}
\\
{\raggedright
\textbf{Reason:} Compared Theorem 3.\allowbreak{}8 at cited publication:\allowbreak{}373-\allowbreak{}415 under the bound corrected target for theorem-\allowbreak{}3.\allowbreak{}8-\allowbreak{}general-\allowbreak{}measure-\allowbreak{}analytic-\allowbreak{}bridge with the full expanded target, roles 1-\allowbreak{}35, and every listed support display.\allowbreak{} The model and analytic premises quantify expected-\allowbreak{}loss/\allowbreak{}gradient integrability, local loss measurability, actual pointwise gradients, two-\allowbreak{}coordinate beta-\allowbreak{}gradient Lipschitz bounds, gamma-\allowbreak{}strong convexity, and epsilon-\allowbreak{}Wasserstein sensitivity over all ambient parameters and the chosen data carrier, as the corrected scope expressly requires.\allowbreak{} The coupling expansion uses prescribed probability marginals and the infimum of finite integrable expected-\allowbreak{}distance costs.\allowbreak{} The domain is nonempty, closed and convex;\allowbreak{} the projection lands in it, satisfies the stated variational inequality, and is nonexpansive.\allowbreak{} The update expands to projection of theta minus eta times the loss gradient averaged under D(theta).\allowbreak{} The numerical assumptions are exactly 0 < gamma <= beta, 0 < eta <= 2/\allowbreak{}(gamma+beta), and epsilon < gamma/\allowbreak{}((gamma+beta)(1+3 eta beta/\allowbreak{}2)).\allowbreak{} Both the pairwise bound and iterate bound use q = 1 -\allowbreak{} eta*(gamma beta/\allowbreak{}(gamma+beta) -\allowbreak{} epsilon*(3 eta beta\textasciicircum{}2/\allowbreak{}2+beta)).\allowbreak{} These premises give 0 <= q < 1, and the q\textasciicircum{}t bound implies the displayed logarithmic convergence-\allowbreak{}time guarantee.\allowbreak{} The conclusion includes existence and uniqueness among all domain-\allowbreak{}stable points and convergence from each domain initial point.\allowbreak{} The ambient analytic scope is accepted only as the disclosed non-\allowbreak{}equivalent corrected target.\allowbreak{}
\par}
\reviewmatch{Matches formalized target}
\noindent\textbf{Reviewer annotation}\par
\reviewerbox
\clearpage
\hypertarget{source-claim-bb0077763720ebbf}{}
\section*{Review row 3}
\textbf{Source locator:} {\footnotesize\raggedright cited publication, lines 442-\allowbreak{}445\par}
\paragraph{Verbatim source input}
\begin{ReviewVerbatim}
Proposition 4.1. Let the distribution map D(·) be ε-                  within the context of strategic classification and discuss
sensitive and Θ ⊂ Rd be compact. If the loss `(z; θ) is               some of the implications of our theorems for this setting.
convex and jointly continuous in (z, θ), then there exists a          We begin by formally establishing how strategic classifica-
performatively stable classifier.                                      tion can be cast as a performative prediction problem and
\end{ReviewVerbatim}


\paragraph{Formalized review target}
The archival source above is not asserted equivalent to this formalized target.
\begin{ReviewVerbatim}
For a globally defined arbitrary-law measure-performative model on a finite-dimensional real ambient parameter space, assume expected-loss integrability for every ambient deployed/evaluated parameter pair. Let the source parameter domain be nonempty, compact, and convex. If the decoupled expected risk is jointly continuous on the entire ambient parameter product and every pointwise loss is convex on the domain for every datum in the chosen data carrier, then a performatively stable parameter exists in the domain.
\end{ReviewVerbatim}

\textbf{Archival source anchor:} {\footnotesize\raggedright cited publication, lines 442-\allowbreak{}445\par}
\paragraph{Expanded PaperInterface specification}
\begin{ReviewVerbatim}
∀ {Parameter : Type u_1} {Data : Type u_2} [inst : NormedAddCommGroup Parameter] [inst_1 : NormedSpace ℝ Parameter]
  [FiniteDimensional ℝ Parameter] [inst_3 : MeasurableSpace Data]
  (model : AppliedModelingLib.MeasurePerformativeModel Parameter Data) (domain : Set Parameter),
  IsCompact domain →
    domain.Nonempty →
      Convex ℝ domain →
        (Continuous fun (point : Parameter × Parameter) =>
            AppliedModelingLib.measureDecoupledPerformativeRisk model point.1 point.2) →
          (∀ (datum : Data), ConvexOn ℝ domain (model.loss datum)) →
            ∃ parameter ∈ domain, AppliedModelingLib.IsMeasurePerformativelyStableOn model domain parameter
\end{ReviewVerbatim}

\paragraph{Retained governing declarations}
\begin{itemize}\raggedright
\item \hyperlink{governing-declaration-522b06e211b1e40a}{Applied\allowbreak{}Modeling\allowbreak{}Lib.\allowbreak{}Is\allowbreak{}Measure\allowbreak{}Performatively\allowbreak{}Stable\allowbreak{}On}
\item \hyperlink{library-prerequisite-96553d4fbb82f1ac}{Applied\allowbreak{}Modeling\allowbreak{}Lib.\allowbreak{}Measure\allowbreak{}Performative\allowbreak{}Model}
\item \hyperlink{library-prerequisite-d4de05e57b5947f0}{Applied\allowbreak{}Modeling\allowbreak{}Lib.\allowbreak{}measure\allowbreak{}Decoupled\allowbreak{}Performative\allowbreak{}Risk}
\end{itemize}
\paragraph{Lean-elaborated claim roles}\begin{ReviewVerbatim}
1. parameter: Type u_1
2. parameter: Type u_2
3. parameter: NormedAddCommGroup Parameter
4. parameter: NormedSpace ℝ Parameter
5. assumption: FiniteDimensional ℝ Parameter
6. parameter: MeasurableSpace Data
7. parameter: AppliedModelingLib.MeasurePerformativeModel Parameter Data
8. parameter: Set Parameter
9. assumption: IsCompact domain
10. assumption: domain.Nonempty
11. assumption: Convex ℝ domain
12. assumption: Continuous fun (point : Parameter × Parameter) =>
  AppliedModelingLib.measureDecoupledPerformativeRisk model point.1 point.2
13. assumption: ∀ (datum : Data), ConvexOn ℝ domain (model.loss datum)
14. conclusion: ∃ parameter ∈ domain, AppliedModelingLib.IsMeasurePerformativelyStableOn model domain parameter
\end{ReviewVerbatim}

\paragraph{Lean proof endpoint (built separately)}\begin{ReviewVerbatim}
PZMH20PerformativePrediction.measureCompactConvexStablePointExistsProof
\end{ReviewVerbatim}

\paragraph{Recorded source-to-Spec screening}
\textbf{Verdict:} \texttt{matches\_approved\_corrected\_target}
\\
{\raggedright
\textbf{Reason:} Compared Proposition 4.\allowbreak{}1 at cited publication:\allowbreak{}442-\allowbreak{}445 with its identity-\allowbreak{}bound corrected target (defects proposition-\allowbreak{}4.\allowbreak{}1-\allowbreak{}expected-\allowbreak{}risk-\allowbreak{}continuity-\allowbreak{}bridge and proposition-\allowbreak{}4.\allowbreak{}1-\allowbreak{}convex-\allowbreak{}domain-\allowbreak{}omission), the full expanded target and roles 1-\allowbreak{}14, and all three displayed prerequisites.\allowbreak{} The Lean model supplies a probability law and expected-\allowbreak{}loss integrability for every ambient deployed/\allowbreak{}evaluated pair;\allowbreak{} the ambient parameter space is finite-\allowbreak{}dimensional over the reals.\allowbreak{} The hypotheses require a nonempty compact convex domain, joint continuity of the decoupled expected risk on the entire ambient product, and convexity on the domain for every datum in the chosen carrier, exactly as the corrected statement and scope note specify.\allowbreak{} Expanding the risk and stability predicates gives an in-\allowbreak{}domain parameter minimizing its own induced-\allowbreak{}law expected loss over every in-\allowbreak{}domain candidate.\allowbreak{} Thus the conclusion is the corrected domain-\allowbreak{}stable existence assertion.\allowbreak{} The archival distribution-\allowbreak{}sensitivity and pointwise-\allowbreak{}continuity route is replaced by the explicitly approved expected-\allowbreak{}risk continuity bridge;\allowbreak{} the global integrability/\allowbreak{}continuity scope and all-\allowbreak{}data convexity are retained as disclosed corrected-\allowbreak{}target conditions, without asserting their necessity or archival equivalence.\allowbreak{}
\par}
\reviewmatch{Matches formalized target}
\noindent\textbf{Reviewer annotation}\par
\reviewerbox
\clearpage
\hypertarget{source-claim-a3115b33fe478ce6}{}
\section*{Review row 4}
\textbf{Source locator:} {\footnotesize\raggedright cited publication, lines 460-\allowbreak{}475\par}
\paragraph{Verbatim source input}
\begin{ReviewVerbatim}
Proposition 4.2. The performative risk PR(θ) can be con-
                                                                      their features in order to improve their outcomes.
cave in θ, even if the loss `(z; θ) is β-jointly smooth (A1),
γ-strongly convex (A2), and the distribution map D(·) is              A classic example of this setting is that of a bank which
ε-sensitive with ε < βγ .                                             uses a machine learning classifier to predict whether or not a
                                                                      loan applicant is creditworthy. Individual applicants react to
However, what we can show is that there are cases where               the bank’s classifier by manipulating their features with the
finding performatively stable points is sufficient to guaran-           hopes of inducing a favorable classification. This game is
tee that the resulting classifier has low performative risk. In        said to have a Stackelberg structure since agents adapt their
particular, our next result demonstrates that if the loss func-       features only after the bank has deployed their classifier.
tion `(z; θ) is Lipschitz in z and γ-strongly convex, then
all performatively stable points and performative optima              The optimal strategy for the institution in a strategic classifi-
lie in a small neighborhood around each other. Moreover,              cation setting is to deploy the solution corresponding to the
the theorem holds for cases where performative optima and             Stackelberg equilibrium, defined as the classifier fθ which
performatively stable points are not necessarily unique.              achieves minimal loss over the induced distribution D(θ)
                                                                      in which agents have strategically adapted their features in
\end{ReviewVerbatim}



\paragraph{Expanded PaperInterface specification}
\begin{ReviewVerbatim}
∃ (mu : ℝ) (epsilon : ℝ) (hmu : |mu| ≤ 1 / 2) (hmuEpsilon : |mu + epsilon| ≤ 1 / 2),
  1 / 2 < epsilon ∧
    epsilon < 2 / 2 ∧
      PZMH20PerformativePrediction.IsFiniteWassersteinSensitive
          (AppliedModelingLib.correctedProposition42Model mu epsilon hmu hmuEpsilon) epsilon ∧
        (∀ (datum : Bool × Bool) (candidate : ℝ),
            HasGradientAt (fun (parameter : ℝ) => AppliedModelingLib.proposition42Loss parameter datum)
              (AppliedModelingLib.proposition42Gradient datum candidate) candidate) ∧
          (∀ (datum : Bool × Bool) (center candidate : ℝ),
              AppliedModelingLib.proposition42Loss candidate datum ≥
                AppliedModelingLib.proposition42Loss center datum +
                    AppliedModelingLib.proposition42Gradient datum center * (candidate - center) +
                  2 / 2 * (candidate - center) ^ 2) ∧
            AppliedModelingLib.IsJointlySmoothGradient AppliedModelingLib.proposition42Gradient 2 ∧
              (∀ (theta : { theta : ℝ // theta ∈ Set.Icc 0 1 }),
                  AppliedModelingLib.performativeRisk
                      (AppliedModelingLib.correctedProposition42Model mu epsilon hmu hmuEpsilon) theta =
                    1 / 4 - 2 * ↑theta * mu + (1 - 2 * epsilon) * ↑theta ^ 2) ∧
                ConcaveOn ℝ Set.univ fun (theta : ℝ) => 1 / 4 - 2 * theta * mu + (1 - 2 * epsilon) * theta ^ 2
\end{ReviewVerbatim}

\paragraph{Retained governing declarations}
\begin{itemize}\raggedright
\item \hyperlink{library-prerequisite-2e05b739f5b71f23}{Applied\allowbreak{}Modeling\allowbreak{}Lib.\allowbreak{}Is\allowbreak{}Jointly\allowbreak{}Smooth\allowbreak{}Gradient}
\item \hyperlink{governing-declaration-a540759b18439129}{Applied\allowbreak{}Modeling\allowbreak{}Lib.\allowbreak{}corrected\allowbreak{}Proposition42\allowbreak{}Model}
\item \hyperlink{governing-declaration-4a2c8294f0669302}{Applied\allowbreak{}Modeling\allowbreak{}Lib.\allowbreak{}performative\allowbreak{}Risk}
\item \hyperlink{governing-declaration-d4d35ede9cba1f94}{Applied\allowbreak{}Modeling\allowbreak{}Lib.\allowbreak{}proposition42\allowbreak{}Bool\allowbreak{}Metric\allowbreak{}Space}
\item \hyperlink{governing-declaration-8c485a232d3e9b18}{Applied\allowbreak{}Modeling\allowbreak{}Lib.\allowbreak{}proposition42\allowbreak{}Gradient}
\item \hyperlink{governing-declaration-ded6eb207c916fe6}{Applied\allowbreak{}Modeling\allowbreak{}Lib.\allowbreak{}proposition42\allowbreak{}Loss}
\item \hyperlink{governing-declaration-12799a7660a5d8a2}{PZMH20\allowbreak{}Performative\allowbreak{}Prediction.\allowbreak{}Is\allowbreak{}Finite\allowbreak{}Wasserstein\allowbreak{}Sensitive}
\end{itemize}
\paragraph{Lean-elaborated claim roles}\begin{ReviewVerbatim}
1. conclusion: ∃ (mu : ℝ) (epsilon : ℝ) (hmu : |mu| ≤ 1 / 2) (hmuEpsilon : |mu + epsilon| ≤ 1 / 2),
  1 / 2 < epsilon ∧
    epsilon < 2 / 2 ∧
      PZMH20PerformativePrediction.IsFiniteWassersteinSensitive
          (AppliedModelingLib.correctedProposition42Model mu epsilon hmu hmuEpsilon) epsilon ∧
        (∀ (datum : Bool × Bool) (candidate : ℝ),
            HasGradientAt (fun (parameter : ℝ) => AppliedModelingLib.proposition42Loss parameter datum)
              (AppliedModelingLib.proposition42Gradient datum candidate) candidate) ∧
          (∀ (datum : Bool × Bool) (center candidate : ℝ),
              AppliedModelingLib.proposition42Loss candidate datum ≥
                AppliedModelingLib.proposition42Loss center datum +
                    AppliedModelingLib.proposition42Gradient datum center * (candidate - center) +
                  2 / 2 * (candidate - center) ^ 2) ∧
            AppliedModelingLib.IsJointlySmoothGradient AppliedModelingLib.proposition42Gradient 2 ∧
              (∀ (theta : { theta : ℝ // theta ∈ Set.Icc 0 1 }),
                  AppliedModelingLib.performativeRisk
                      (AppliedModelingLib.correctedProposition42Model mu epsilon hmu hmuEpsilon) theta =
                    1 / 4 - 2 * ↑theta * mu + (1 - 2 * epsilon) * ↑theta ^ 2) ∧
                ConcaveOn ℝ Set.univ fun (theta : ℝ) => 1 / 4 - 2 * theta * mu + (1 - 2 * epsilon) * theta ^ 2
\end{ReviewVerbatim}

\paragraph{Lean proof endpoint (built separately)}\begin{ReviewVerbatim}
PZMH20PerformativePrediction.proposition42CorrectedFiniteWitness
\end{ReviewVerbatim}

\paragraph{Recorded source-to-Spec screening}
\textbf{Verdict:} \texttt{matches}
\\
{\raggedright
\textbf{Reason:} The source makes an existential claim that performative risk can be concave while the loss is jointly smooth and strongly convex and the distribution map is epsilon-\allowbreak{}sensitive below gamma/\allowbreak{}beta.\allowbreak{} The expanded target supplies a valid finite model with beta = gamma = 2 and 1/\allowbreak{}2 < epsilon < 1, ties proposition42Gradient to proposition42Loss by HasGradientAt, states the 2-\allowbreak{}strong-\allowbreak{}convexity inequality and 2-\allowbreak{}joint-\allowbreak{}smoothness condition, establishes epsilon-\allowbreak{}Wasserstein sensitivity, identifies the exact performative-\allowbreak{}risk polynomial on the parameter interval, and proves that polynomial concave.\allowbreak{} The displayed finite-\allowbreak{}law, expectation, metric, coupling, and sensitivity definitions give these clauses their stated meanings, so the witness establishes the source conclusion.\allowbreak{}
\par}
\reviewmatch{Matches source input}
\noindent\textbf{Reviewer annotation}\par
\reviewerbox
\clearpage
\hypertarget{source-claim-cef541336c7da1ef}{}
\section*{Review row 5}
\textbf{Source locator:} {\footnotesize\raggedright cited publication, lines 418-\allowbreak{}423\par}
\paragraph{Verbatim source input}
\begin{ReviewVerbatim}
population level: as in Theorems 3.5 and 3.8, the iterates θt         D(·) is ε-sensitive. Then, for every performatively stable
contract toward θPS . This implies that θt eventually enters a        point θPS and every performative optimum θPO :
δ-ball around θPS , for some large enough t. Once this hap-
                                                                                                            2Lz ε
pens, contrary to population-level results, a contraction is                            kθPO − θPS k2 6           .
no longer guaranteed due to the noise inherent in observing                                                  γ

[Next verbatim source excerpt]

Theorem 4.3. Suppose that the loss `(z; θ) is Lz -Lipschitz           response to fθ . In fact, we see that this equilibrium notion
in z, γ-strongly convex (A2), and that the distribution map           exactly matches our definition of performative optimality:
    1
      In particular, they are guaranteed to exist over the extended
real line, i.e. we allow θ ∈ (R ∪ {±∞})d .                            fθSE is a Stackelberg equilibrium ⇔ θSE ∈ arg min PR(θ).
                                                                                                                         θ

[form-feed in source extraction]
                                                                   Performative Prediction



                                                                                                         -2
                                -2                                                                  10
                           10
        c ¢ kµt + 1 ¡ µt k 2




                                                                                 c ¢ kµt + 1 ¡ µt k 2
                                                                  " = 0.01
                                                                  "=1                                    -7                                 " = 0.01
                           10
                                -7                                                                  10
                                                                  " = 100                                                                   "=1
                               -12                                " = 1000                      10
                                                                                                        -12                                 " = 100
                       10
                                                                                                                                            " = 1000

                                     0       20          40       60                                          0        1000          2000        3000
                                                   Iteration t                                                             Iteration t

                                 (a) Repeated Risk Minimization (RRM)                                     (b) Repeated Gradient Descent (RGD)
\end{ReviewVerbatim}



\paragraph{Expanded PaperInterface specification}
\begin{ReviewVerbatim}
∀ {Parameter : Type u_1} {Data : Type u_2} [inst : NormedAddCommGroup Parameter]
  [inst_1 : InnerProductSpace ℝ Parameter] [inst_2 : MeasurableSpace Data] [inst_3 : MetricSpace Data]
  {domain : Set Parameter} (model : AppliedModelingLib.MeasurePerformativeModelOn Parameter Data domain)
  (gradient : Data → ↑domain → Parameter),
  IsClosed domain →
    Convex ℝ domain →
      ∀ {modulus : ℝ},
        0 < modulus →
          model.IsPointwiseGradientStronglyConvex gradient modulus →
            ∀ {constant : NNReal},
              model.IsDataLossLipschitz constant →
                ∀ {sensitivity : ℝ},
                  0 ≤ sensitivity →
                    model.IsWassersteinSensitive sensitivity →
                      ∀ (optimal stable : ↑domain),
                        model.IsPerformativelyOptimal optimal →
                          model.IsPerformativelyStable stable →
                            ‖↑optimal - ↑stable‖ ≤ 2 * ↑constant * sensitivity / modulus
\end{ReviewVerbatim}

\paragraph{Retained governing declarations}
\begin{itemize}\raggedright
\item \hyperlink{library-prerequisite-e14cc5e8b1e1460e}{Applied\allowbreak{}Modeling\allowbreak{}Lib.\allowbreak{}Measure\allowbreak{}Performative\allowbreak{}Model\allowbreak{}On}
\item \hyperlink{library-prerequisite-0de79b655cc05174}{Applied\allowbreak{}Modeling\allowbreak{}Lib.\allowbreak{}Measure\allowbreak{}Performative\allowbreak{}Model\allowbreak{}On.\allowbreak{}Is\allowbreak{}Data\allowbreak{}Loss\allowbreak{}Lipschitz}
\item \hyperlink{library-prerequisite-555f831d5bfa2a27}{Applied\allowbreak{}Modeling\allowbreak{}Lib.\allowbreak{}Measure\allowbreak{}Performative\allowbreak{}Model\allowbreak{}On.\allowbreak{}Is\allowbreak{}Performatively\allowbreak{}Optimal}
\item \hyperlink{library-prerequisite-4af6c70f954a80a0}{Applied\allowbreak{}Modeling\allowbreak{}Lib.\allowbreak{}Measure\allowbreak{}Performative\allowbreak{}Model\allowbreak{}On.\allowbreak{}Is\allowbreak{}Performatively\allowbreak{}Stable}
\item \hyperlink{library-prerequisite-eb9ab8e4235e1d28}{Applied\allowbreak{}Modeling\allowbreak{}Lib.\allowbreak{}Measure\allowbreak{}Performative\allowbreak{}Model\allowbreak{}On.\allowbreak{}Is\allowbreak{}Pointwise\allowbreak{}Gradient\allowbreak{}Strongly\allowbreak{}Convex}
\item \hyperlink{library-prerequisite-0e1e192eee739a4b}{Applied\allowbreak{}Modeling\allowbreak{}Lib.\allowbreak{}Measure\allowbreak{}Performative\allowbreak{}Model\allowbreak{}On.\allowbreak{}Is\allowbreak{}Wasserstein\allowbreak{}Sensitive}
\end{itemize}
\paragraph{Lean-elaborated claim roles}\begin{ReviewVerbatim}
1. parameter: Type u_1
2. parameter: Type u_2
3. parameter: NormedAddCommGroup Parameter
4. parameter: InnerProductSpace ℝ Parameter
5. parameter: MeasurableSpace Data
6. parameter: MetricSpace Data
7. parameter: Set Parameter
8. parameter: AppliedModelingLib.MeasurePerformativeModelOn Parameter Data domain
9. parameter: Data → ↑domain → Parameter
10. assumption: IsClosed domain
11. assumption: Convex ℝ domain
12. parameter: ℝ
13. assumption: 0 < modulus
14. assumption: model.IsPointwiseGradientStronglyConvex gradient modulus
15. parameter: NNReal
16. assumption: model.IsDataLossLipschitz constant
17. parameter: ℝ
18. assumption: 0 ≤ sensitivity
19. assumption: model.IsWassersteinSensitive sensitivity
20. parameter: ↑domain
21. parameter: ↑domain
22. assumption: model.IsPerformativelyOptimal optimal
23. assumption: model.IsPerformativelyStable stable
24. conclusion: Real.le✝ ‖↑optimal - ↑stable‖ (2 * ↑constant * sensitivity / modulus)
\end{ReviewVerbatim}

\paragraph{Lean proof endpoint (built separately)}\begin{ReviewVerbatim}
PZMH20PerformativePrediction.measureWassersteinOptimumStableDistanceProof
\end{ReviewVerbatim}

\paragraph{Recorded source-to-Spec screening}
\textbf{Verdict:} \texttt{matches}
\\
{\raggedright
\textbf{Reason:} Compared theorem\_\allowbreak{}4\_\allowbreak{}3\_\allowbreak{}optimum\_\allowbreak{}stable\_\allowbreak{}distance, cited publication:\allowbreak{}418-\allowbreak{}423,476-\allowbreak{}508, directly with the expanded Lean proposition, all 24 elaborated claim roles, and all 11 displayed material prerequisites.\allowbreak{} The source quantifies over every performatively stable point and every performative optimum and concludes norm(thetaPO-\allowbreak{}thetaPS) <= 2*Lz*epsilon/\allowbreak{}gamma.\allowbreak{} The target has precisely that universal pair quantification and constant;\allowbreak{} reversing the order of the independent pair binders changes no assertion.\allowbreak{} The supporting risk definitions use the integral under the deployed law with a separately evaluated parameter;\allowbreak{} optimality compares the diagonals over all of Theta, while stability fixes its own induced law and compares every feasible evaluated parameter.\allowbreak{} They do not impose uniqueness or assert existence.\allowbreak{} The model supplies real integrable expectations for all feasible pairs and no values outside Theta.\allowbreak{} The positive modulus and the complete A2 prerequisite require, for every datum and feasible center, the loss derivative within Theta together with the quadratic lower bound using that same derivative and gamma/\allowbreak{}2.\allowbreak{} Its within-\allowbreak{}domain zero extension adds no off-\allowbreak{}domain differentiability requirement.\allowbreak{} The uniform NNReal data-\allowbreak{}Lipschitz constant and nonnegative sensitivity include their zero cases.\allowbreak{} The coupling support displays define a probability law with the prescribed marginals, integrable expected distance costs, and their real infimum;\allowbreak{} the finite-\allowbreak{}cost witness therefore supplies the source W1 inequality with the same epsilon times parameter distance.\allowbreak{} The target retains the closed convex feasible domain and adds neither A1 joint smoothness, a gamma/\allowbreak{}beta sensitivity threshold, finite-\allowbreak{}support restrictions, gradient integrability, nor an attained transport infimum.\allowbreak{} The concluding claim role is exactly the source distance bound, with no dropped conclusion or additional optimal/\allowbreak{}stable uniqueness premise.\allowbreak{} No corrected-\allowbreak{}target record is invoked.\allowbreak{}
\par}
\reviewmatch{Matches source input}
\noindent\textbf{Reviewer annotation}\par
\reviewerbox
\clearpage
\hypertarget{source-claim-a22662c7824bd286}{}
\section*{Review row 6}
\textbf{Source locator:} {\footnotesize\raggedright cited publication, lines 543-\allowbreak{}561\par}
\paragraph{Verbatim source input}
\begin{ReviewVerbatim}
Corollary 5.1. Let the institution’s loss be Lz - and Lθ -                      sion classifier. We assume that individuals have linear
Lipschitz in z and θ respectively, β-jointly smooth (A1), and                   utilities u(θ, x) = −hθ, xi and quadratic costs c(x0 , x) =
                                                                                 1     0     2
γ-strongly convex (A2). If the induced distribution map is                      2ε kx − xk2 , where ε is a positive constant that regulates
ε-sensitive, with ε < βγ , then RRM converges to a stable                       the cost incurred by changing features, and hence the sensi-
classifier θPS that is 2Lz ε(Lθ + Lz ε)γ −1 close in objective                   tivity of the distribution map. Linear utilities indicate that
value to the Stackelberg equilibrium.                                           agents wish to minimize their assigned probability of de-
                                                                                fault. We divide the set of features into strategic features
                                                                                S ⊆ [m − 1], such as the number of open credit lines, and
                                                                                non-strategic features (e.g., age). Solving the optimization
                                                                                problem described in Figure 2, the best response for an indi-
 Input: base distribution D, classifier fθ , cost function c                     vidual corresponds to the following update, x0S = xS − εθS ,
 and utility function u                                                         where xS , x0S , θS ∈ R|S| . As per convention in the literature
 Sampling procedure for D(θ):                                                   (Brückner et al., 2012; Hardt et al., 2016a; Milli et al., 2019),
   1. Sample (x, y) ∼ D                                                         individual outcomes y are unaffected by manipulation.
   2. Compute xBR ← arg maxx0 u(x0 , θ) − c(x0 , x)
   3. Output sample (xBR , y)                                                   Intuitively, this data-generating process is ε-sensitive since
                                                                                for a given choice of classifiers, fθ and fθ0 , an individual
                                                                                feature vector is shifted to xS − εθS and to xS − εθS0 , re-
\end{ReviewVerbatim}


\paragraph{Formalized review target}
The archival source above is not asserted equivalent to this formalized target.
\begin{ReviewVerbatim}
For an arbitrary-law measure-performative model on a complete real inner-product ambient parameter space, take the admissible parameter domain to be the entire ambient carrier. Assume expected-loss and gradient integrability for every ambient deployed/evaluated parameter pair; A1 joint gradient smoothness, A2 pointwise strong convexity, and data- and parameter-coordinate Lipschitz loss for every ambient parameter and every datum in the chosen data carrier; nonnegative Wasserstein sensitivity for every ambient parameter pair; local loss measurability and pointwise parameter differentiability for every ambient parameter and datum; a whole-space repeated-risk-minimization update with strict epsilon beta over gamma below one; and a whole-space performative optimum. Then, from every ambient initial point, the update converges to a unique whole-space stable point whose performative objective gap to that optimum is at most two Lz epsilon times (Ltheta plus Lz epsilon) over gamma.
\end{ReviewVerbatim}

\textbf{Archival source anchor:} {\footnotesize\raggedright cited publication, lines 543-\allowbreak{}561\par}
\paragraph{Expanded PaperInterface specification}
\begin{ReviewVerbatim}
∀ {Parameter : Type u_1} {Data : Type u_2} [inst : NormedAddCommGroup Parameter]
  [inst_1 : InnerProductSpace ℝ Parameter] [inst_2 : CompleteSpace Parameter] [inst_3 : MeasurableSpace Data]
  [inst_4 : MetricSpace Data] (model : AppliedModelingLib.MeasurePerformativeModel Parameter Data)
  (gradient : Data → Parameter → Parameter) {smoothness : NNReal},
  AppliedModelingLib.IsJointlySmoothGradient gradient smoothness →
    (∀ (distributionParameter evaluatedParameter : Parameter),
        MeasureTheory.Integrable (fun (datum : Data) => gradient datum evaluatedParameter)
          ↑(model.dataLaw distributionParameter)) →
      (∀ (distributionParameter evaluatedParameter : Parameter),
          ∀ᶠ (parameter : Parameter) in nhds evaluatedParameter,
            MeasureTheory.AEStronglyMeasurable (fun (datum : Data) => model.loss datum parameter)
              ↑(model.dataLaw distributionParameter)) →
        (∀ (datum : Data) (parameter : Parameter),
            HasGradientAt (fun (theta : Parameter) => model.loss datum theta) (gradient datum parameter) parameter) →
          ∀ {modulus : ℝ},
            0 < modulus →
              PZMH20PerformativePrediction.IsMeasurePointwiseGradientStronglyConvex model gradient modulus →
                ∀ {dataConstant parameterConstant : NNReal},
                  AppliedModelingLib.IsMeasureDataLossLipschitz model dataConstant →
                    AppliedModelingLib.IsMeasureParameterLossLipschitz model parameterConstant →
                      ∀ {sensitivity : ℝ},
                        0 ≤ sensitivity →
                          PZMH20PerformativePrediction.IsMeasureWassersteinSensitive model sensitivity →
                            ∀ (update : Parameter → Parameter),
                              PZMH20PerformativePrediction.IsMeasureRRM model update →
                                sensitivity * ↑smoothness / modulus < 1 →
                                  ∀ (initial optimal : Parameter),
                                    PZMH20PerformativePrediction.IsMeasurePerformativelyOptimal model optimal →
                                      ∃ (stable : Parameter),
                                        PZMH20PerformativePrediction.IsMeasureStable model stable ∧
                                          (∀ (other : Parameter),
                                              PZMH20PerformativePrediction.IsMeasureStable model other →
                                                other = stable) ∧
                                            Filter.Tendsto (fun (iteration : ℕ) => update^[iteration] initial)
                                                Filter.atTop (nhds stable) ∧
                                              (∀ (iteration : ℕ),
                                                  ‖update^[iteration] initial - stable‖ ≤
                                                    (sensitivity * ↑smoothness / modulus) ^ iteration *
                                                      ‖initial - stable‖) ∧
                                                AppliedModelingLib.measurePerformativeRisk model stable -
                                                    AppliedModelingLib.measurePerformativeRisk model optimal ≤
                                                  2 * ↑dataConstant * sensitivity *
                                                      (↑parameterConstant + ↑dataConstant * sensitivity) /
                                                    modulus
\end{ReviewVerbatim}

\paragraph{Retained governing declarations}
\begin{itemize}\raggedright
\item \hyperlink{library-prerequisite-2e05b739f5b71f23}{Applied\allowbreak{}Modeling\allowbreak{}Lib.\allowbreak{}Is\allowbreak{}Jointly\allowbreak{}Smooth\allowbreak{}Gradient}
\item \hyperlink{library-prerequisite-49496bb3a06c52e5}{Applied\allowbreak{}Modeling\allowbreak{}Lib.\allowbreak{}Is\allowbreak{}Measure\allowbreak{}Data\allowbreak{}Loss\allowbreak{}Lipschitz}
\item \hyperlink{governing-declaration-ef89f6befb052e48}{Applied\allowbreak{}Modeling\allowbreak{}Lib.\allowbreak{}Is\allowbreak{}Measure\allowbreak{}Parameter\allowbreak{}Loss\allowbreak{}Lipschitz}
\item \hyperlink{library-prerequisite-96553d4fbb82f1ac}{Applied\allowbreak{}Modeling\allowbreak{}Lib.\allowbreak{}Measure\allowbreak{}Performative\allowbreak{}Model}
\item \hyperlink{governing-declaration-13fb93228218a38e}{Applied\allowbreak{}Modeling\allowbreak{}Lib.\allowbreak{}measure\allowbreak{}Performative\allowbreak{}Risk}
\item \hyperlink{paper-prerequisite-48fa8b1c1ddbaa9c}{PZMH20\allowbreak{}Performative\allowbreak{}Prediction.\allowbreak{}Is\allowbreak{}Measure\allowbreak{}Performatively\allowbreak{}Optimal}
\item \hyperlink{governing-declaration-11c9ab6e3c6f528b}{PZMH20\allowbreak{}Performative\allowbreak{}Prediction.\allowbreak{}Is\allowbreak{}Measure\allowbreak{}Pointwise\allowbreak{}Gradient\allowbreak{}Strongly\allowbreak{}Convex}
\item \hyperlink{governing-declaration-32f609ab3d5aaf05}{PZMH20\allowbreak{}Performative\allowbreak{}Prediction.\allowbreak{}Is\allowbreak{}Measure\allowbreak{}RRM}
\item \hyperlink{paper-prerequisite-f6b3bb0ef1cf7d32}{PZMH20\allowbreak{}Performative\allowbreak{}Prediction.\allowbreak{}Is\allowbreak{}Measure\allowbreak{}Stable}
\item \hyperlink{paper-prerequisite-2ea2f5769822c5c4}{PZMH20\allowbreak{}Performative\allowbreak{}Prediction.\allowbreak{}Is\allowbreak{}Measure\allowbreak{}Wasserstein\allowbreak{}Sensitive}
\end{itemize}
\paragraph{Lean-elaborated claim roles}\begin{ReviewVerbatim}
1. parameter: Type u_1
2. parameter: Type u_2
3. parameter: NormedAddCommGroup Parameter
4. parameter: InnerProductSpace ℝ Parameter
5. assumption: CompleteSpace Parameter
6. parameter: MeasurableSpace Data
7. parameter: MetricSpace Data
8. parameter: AppliedModelingLib.MeasurePerformativeModel Parameter Data
9. parameter: Data → Parameter → Parameter
10. parameter: NNReal
11. assumption: AppliedModelingLib.IsJointlySmoothGradient gradient smoothness
12. assumption: ∀ (distributionParameter evaluatedParameter : Parameter),
  MeasureTheory.Integrable (fun (datum : Data) => gradient datum evaluatedParameter)
    ↑(model.dataLaw distributionParameter)
13. assumption: ∀ (distributionParameter evaluatedParameter : Parameter),
  ∀ᶠ (parameter : Parameter) in nhds evaluatedParameter,
    MeasureTheory.AEStronglyMeasurable (fun (datum : Data) => model.loss datum parameter)
      ↑(model.dataLaw distributionParameter)
14. assumption: ∀ (datum : Data) (parameter : Parameter),
  HasGradientAt (fun (theta : Parameter) => model.loss datum theta) (gradient datum parameter) parameter
15. parameter: ℝ
16. assumption: 0 < modulus
17. assumption: PZMH20PerformativePrediction.IsMeasurePointwiseGradientStronglyConvex model gradient modulus
18. parameter: NNReal
19. parameter: NNReal
20. assumption: AppliedModelingLib.IsMeasureDataLossLipschitz model dataConstant
21. assumption: AppliedModelingLib.IsMeasureParameterLossLipschitz model parameterConstant
22. parameter: ℝ
23. assumption: 0 ≤ sensitivity
24. assumption: PZMH20PerformativePrediction.IsMeasureWassersteinSensitive model sensitivity
25. parameter: Parameter → Parameter
26. assumption: PZMH20PerformativePrediction.IsMeasureRRM model update
27. assumption: sensitivity * ↑smoothness / modulus < 1
28. parameter: Parameter
29. parameter: Parameter
30. assumption: PZMH20PerformativePrediction.IsMeasurePerformativelyOptimal model optimal
31. conclusion: ∃ (stable : Parameter),
  PZMH20PerformativePrediction.IsMeasureStable model stable ∧
    (∀ (other : Parameter), PZMH20PerformativePrediction.IsMeasureStable model other → other = stable) ∧
      Filter.Tendsto (fun (iteration : ℕ) => update^[iteration] initial) Filter.atTop (nhds stable) ∧
        (∀ (iteration : ℕ),
            ‖update^[iteration] initial - stable‖ ≤
              (sensitivity * ↑smoothness / modulus) ^ iteration * ‖initial - stable‖) ∧
          AppliedModelingLib.measurePerformativeRisk model stable -
              AppliedModelingLib.measurePerformativeRisk model optimal ≤
            2 * ↑dataConstant * sensitivity * (↑parameterConstant + ↑dataConstant * sensitivity) / modulus
\end{ReviewVerbatim}

\paragraph{Lean proof endpoint (built separately)}\begin{ReviewVerbatim}
PZMH20PerformativePrediction.measureWassersteinStableObjectiveGapProof
\end{ReviewVerbatim}

\paragraph{Recorded source-to-Spec screening}
\textbf{Verdict:} \texttt{matches\_approved\_corrected\_target}
\\
{\raggedright
\textbf{Reason:} Compared Corollary 5.\allowbreak{}1 at cited publication:\allowbreak{}543-\allowbreak{}561 under the bound corrected target for corollary-\allowbreak{}5.\allowbreak{}1-\allowbreak{}general-\allowbreak{}measure-\allowbreak{}analytic-\allowbreak{}bridge with the full expanded target, roles 1-\allowbreak{}31, and all listed prerequisites.\allowbreak{} The corrected scope expressly chooses the entire ambient parameter carrier.\allowbreak{} Accordingly, the expanded RRM, stability and performative-\allowbreak{}optimality predicates minimize over every ambient candidate, and the initial point is unrestricted;\allowbreak{} no source-\allowbreak{}domain reduction is claimed.\allowbreak{} The loss model, analytic premises, A1/\allowbreak{}A2 predicates, data-\allowbreak{} and parameter-\allowbreak{}coordinate Lipschitz predicates, and finite expected-\allowbreak{}distance Wasserstein coupling condition carry the stated global quantifiers and constants.\allowbreak{} The RRM contraction factor is q = epsilon beta/\allowbreak{}gamma with gamma > 0, nonnegative epsilon and q < 1.\allowbreak{} The conclusion supplies a unique whole-\allowbreak{}space stable point, convergence from the given arbitrary initial point, the q\textasciicircum{}t bound, and performative risk(stable) minus performative risk(optimal) <= 2 Lz epsilon*(Ltheta + Lz epsilon)/\allowbreak{}gamma.\allowbreak{} Expanding performative risk evaluates each loss under its own induced law;\allowbreak{} whole-\allowbreak{}space optimality gives the reverse nonnegative-\allowbreak{}gap inequality, so this is the claimed closeness in objective value.\allowbreak{} The whole-\allowbreak{}space and ambient analytic restrictions are the disclosed non-\allowbreak{}equivalent corrected target, with necessity left unresolved by its scope note.\allowbreak{}
\par}
\reviewmatch{Matches formalized target}
\noindent\textbf{Reviewer annotation}\par
\reviewerbox
\clearpage
\hypertarget{source-claim-1127efac5a38e232}{}
\section*{Review row 7}
\textbf{Source locator:} {\footnotesize\raggedright cited publication, lines 315-\allowbreak{}323\par}
\paragraph{Verbatim source input}
\begin{ReviewVerbatim}
if θ > 12 , and a point mass at 1 if θ < 21 . Clearly there is     Proposition 3.6. Suppose that the distribution map D(·) is
no performatively stable point, and RRM will simply result         ε-sensitive with ε > 0. RRM can fail to converge at all in
in the alternating sequence 1, 0, 1, 0, . . . . The performative   any of the following cases, for any choice of β, γ > 0:
optimum in this case is θPO = 12 .
                                                                    (a) The loss is β-jointly smooth and convex, but not
To show convergence of retraining schemes, it is hence                  strongly convex.
necessary to make a regularity assumption on D(·), such             (b) The loss is γ-strongly convex, but not jointly smooth.
as ε-sensitivity. We are now ready to state our main result         (c) The loss is β-jointly smooth and γ-strongly convex, but
regarding the convergence of repeated risk minimization.                ε > βγ .
\end{ReviewVerbatim}



\paragraph{Expanded PaperInterface specification}
\begin{ReviewVerbatim}
∀ (epsilon beta gamma : ℝ),
  0 < epsilon →
    0 < beta →
      0 < gamma →
        AppliedModelingLib.IsMeasureRepeatedRiskMinimizationOn
            (AppliedModelingLib.proposition36aMeasureModel epsilon beta) AppliedModelingLib.proposition36aDomain
            AppliedModelingLib.proposition36aRRMUpdate ∧
          (∀ (datum candidate : ℝ),
              HasGradientAt (fun (parameter : ℝ) => AppliedModelingLib.proposition36aMeasureLoss beta datum parameter)
                (AppliedModelingLib.proposition36aMeasureGradient beta datum candidate) candidate) ∧
            (∀ (datum : ℝ),
                ConvexOn ℝ Set.univ fun (candidate : ℝ) =>
                  AppliedModelingLib.proposition36aMeasureLoss beta datum candidate) ∧
              AppliedModelingLib.IsJointlySmoothGradient (AppliedModelingLib.proposition36aMeasureGradient beta)
                  beta.toNNReal ∧
                PZMH20PerformativePrediction.IsMeasureWassersteinSensitive
                    (AppliedModelingLib.proposition36aMeasureModel epsilon beta) epsilon ∧
                  (∀ (modulus : ℝ),
                      0 < modulus →
                        ¬PZMH20PerformativePrediction.IsMeasurePointwiseGradientStronglyConvex
                            (AppliedModelingLib.proposition36aMeasureModel epsilon beta)
                            (AppliedModelingLib.proposition36aMeasureGradient beta) modulus) ∧
                    (∀ (theta : ℝ),
                        AppliedModelingLib.IsMeasurePerformativelyStableOn
                            (AppliedModelingLib.proposition36aMeasureModel epsilon beta)
                            AppliedModelingLib.proposition36aDomain theta ↔
                          theta = 0) ∧
                      ¬∃ (limit : ℝ),
                          Filter.Tendsto
                            (fun (iteration : ℕ) => AppliedModelingLib.proposition36aRRMUpdate^[iteration] 1)
                            Filter.atTop (nhds limit)
\end{ReviewVerbatim}

\paragraph{Retained governing declarations}
\begin{itemize}\raggedright
\item \hyperlink{library-prerequisite-2e05b739f5b71f23}{Applied\allowbreak{}Modeling\allowbreak{}Lib.\allowbreak{}Is\allowbreak{}Jointly\allowbreak{}Smooth\allowbreak{}Gradient}
\item \hyperlink{governing-declaration-522b06e211b1e40a}{Applied\allowbreak{}Modeling\allowbreak{}Lib.\allowbreak{}Is\allowbreak{}Measure\allowbreak{}Performatively\allowbreak{}Stable\allowbreak{}On}
\item \hyperlink{governing-declaration-6ae86e5703739171}{Applied\allowbreak{}Modeling\allowbreak{}Lib.\allowbreak{}Is\allowbreak{}Measure\allowbreak{}Repeated\allowbreak{}Risk\allowbreak{}Minimization\allowbreak{}On}
\item \hyperlink{governing-declaration-965db1948b41d1e4}{Applied\allowbreak{}Modeling\allowbreak{}Lib.\allowbreak{}proposition36a\allowbreak{}Domain}
\item \hyperlink{governing-declaration-6b5ac0170880188f}{Applied\allowbreak{}Modeling\allowbreak{}Lib.\allowbreak{}proposition36a\allowbreak{}Measure\allowbreak{}Gradient}
\item \hyperlink{governing-declaration-10c605398ed70a9b}{Applied\allowbreak{}Modeling\allowbreak{}Lib.\allowbreak{}proposition36a\allowbreak{}Measure\allowbreak{}Loss}
\item \hyperlink{governing-declaration-3d85baf63239a428}{Applied\allowbreak{}Modeling\allowbreak{}Lib.\allowbreak{}proposition36a\allowbreak{}Measure\allowbreak{}Model}
\item \hyperlink{governing-declaration-03f87e42b92f1bad}{Applied\allowbreak{}Modeling\allowbreak{}Lib.\allowbreak{}proposition36a\allowbreak{}RRM\allowbreak{}Update}
\item \hyperlink{governing-declaration-11c9ab6e3c6f528b}{PZMH20\allowbreak{}Performative\allowbreak{}Prediction.\allowbreak{}Is\allowbreak{}Measure\allowbreak{}Pointwise\allowbreak{}Gradient\allowbreak{}Strongly\allowbreak{}Convex}
\item \hyperlink{paper-prerequisite-2ea2f5769822c5c4}{PZMH20\allowbreak{}Performative\allowbreak{}Prediction.\allowbreak{}Is\allowbreak{}Measure\allowbreak{}Wasserstein\allowbreak{}Sensitive}
\end{itemize}
\paragraph{Lean-elaborated claim roles}\begin{ReviewVerbatim}
1. parameter: ℝ
2. parameter: ℝ
3. parameter: ℝ
4. assumption: 0 < epsilon
5. assumption: 0 < beta
6. assumption: 0 < gamma
7. conclusion: AppliedModelingLib.IsMeasureRepeatedRiskMinimizationOn (AppliedModelingLib.proposition36aMeasureModel epsilon beta)
    AppliedModelingLib.proposition36aDomain AppliedModelingLib.proposition36aRRMUpdate ∧
  (∀ (datum candidate : ℝ),
      HasGradientAt (fun (parameter : ℝ) => AppliedModelingLib.proposition36aMeasureLoss beta datum parameter)
        (AppliedModelingLib.proposition36aMeasureGradient beta datum candidate) candidate) ∧
    (∀ (datum : ℝ),
        ConvexOn ℝ Set.univ fun (candidate : ℝ) => AppliedModelingLib.proposition36aMeasureLoss beta datum candidate) ∧
      AppliedModelingLib.IsJointlySmoothGradient (AppliedModelingLib.proposition36aMeasureGradient beta) beta.toNNReal ∧
        PZMH20PerformativePrediction.IsMeasureWassersteinSensitive
            (AppliedModelingLib.proposition36aMeasureModel epsilon beta) epsilon ∧
          (∀ (modulus : ℝ),
              0 < modulus →
                ¬PZMH20PerformativePrediction.IsMeasurePointwiseGradientStronglyConvex
                    (AppliedModelingLib.proposition36aMeasureModel epsilon beta)
                    (AppliedModelingLib.proposition36aMeasureGradient beta) modulus) ∧
            (∀ (theta : ℝ),
                AppliedModelingLib.IsMeasurePerformativelyStableOn
                    (AppliedModelingLib.proposition36aMeasureModel epsilon beta) AppliedModelingLib.proposition36aDomain
                    theta ↔
                  theta = 0) ∧
              ¬∃ (limit : ℝ),
                  Filter.Tendsto (fun (iteration : ℕ) => AppliedModelingLib.proposition36aRRMUpdate^[iteration] 1)
                    Filter.atTop (nhds limit)
\end{ReviewVerbatim}

\paragraph{Lean proof endpoint (built separately)}\begin{ReviewVerbatim}
PZMH20PerformativePrediction.proposition36aExactDiracCounterexample
\end{ReviewVerbatim}

\paragraph{Recorded source-to-Spec screening}
\textbf{Verdict:} \texttt{matches}
\\
{\raggedright
\textbf{Reason:} Independent review:\allowbreak{} the exact Dirac linear-\allowbreak{}loss witness gives beta-\allowbreak{}smooth convex, non-\allowbreak{}strongly-\allowbreak{}convex RRM nonconvergence for the positive source constants.\allowbreak{}
\par}
\reviewmatch{Matches source input}
\noindent\textbf{Reviewer annotation}\par
\reviewerbox
\clearpage
\hypertarget{source-claim-c4686335fcfce867}{}
\section*{Review row 8}
\textbf{Source locator:} {\footnotesize\raggedright cited publication, lines 315-\allowbreak{}323\par}
\paragraph{Verbatim source input}
\begin{ReviewVerbatim}
if θ > 12 , and a point mass at 1 if θ < 21 . Clearly there is     Proposition 3.6. Suppose that the distribution map D(·) is
no performatively stable point, and RRM will simply result         ε-sensitive with ε > 0. RRM can fail to converge at all in
in the alternating sequence 1, 0, 1, 0, . . . . The performative   any of the following cases, for any choice of β, γ > 0:
optimum in this case is θPO = 12 .
                                                                    (a) The loss is β-jointly smooth and convex, but not
To show convergence of retraining schemes, it is hence                  strongly convex.
necessary to make a regularity assumption on D(·), such             (b) The loss is γ-strongly convex, but not jointly smooth.
as ε-sensitivity. We are now ready to state our main result         (c) The loss is β-jointly smooth and γ-strongly convex, but
regarding the convergence of repeated risk minimization.                ε > βγ .
\end{ReviewVerbatim}



\paragraph{Expanded PaperInterface specification}
\begin{ReviewVerbatim}
∀ (epsilon beta gamma : ℝ),
  0 < epsilon →
    0 < beta →
      0 < gamma →
        ∃ (penalty : ℝ),
          AppliedModelingLib.IsProposition36bLargePenalty epsilon gamma penalty ∧
            PZMH20PerformativePrediction.proposition36bExactDiracEndpointCounterexampleSpec epsilon gamma penalty
\end{ReviewVerbatim}

\paragraph{Retained governing declarations}
\begin{itemize}\raggedright
\item \hyperlink{governing-declaration-ee9adf0c7d04f002}{Applied\allowbreak{}Modeling\allowbreak{}Lib.\allowbreak{}Is\allowbreak{}Proposition36b\allowbreak{}Large\allowbreak{}Penalty}
\item \hyperlink{governing-declaration-e0c24b576a9d5f99}{PZMH20\allowbreak{}Performative\allowbreak{}Prediction.\allowbreak{}proposition36b\allowbreak{}Exact\allowbreak{}Dirac\allowbreak{}Endpoint\allowbreak{}Counterexample\allowbreak{}Spec}
\end{itemize}
\paragraph{Lean-elaborated claim roles}\begin{ReviewVerbatim}
1. parameter: ℝ
2. parameter: ℝ
3. parameter: ℝ
4. assumption: 0 < epsilon
5. assumption: 0 < beta
6. assumption: 0 < gamma
7. conclusion: ∃ (penalty : ℝ),
  AppliedModelingLib.IsProposition36bLargePenalty epsilon gamma penalty ∧
    PZMH20PerformativePrediction.proposition36bExactDiracEndpointCounterexampleSpec epsilon gamma penalty
\end{ReviewVerbatim}

\paragraph{Lean proof endpoint (built separately)}\begin{ReviewVerbatim}
PZMH20PerformativePrediction.proposition36bSourceUniformCounterexample
\end{ReviewVerbatim}

\paragraph{Recorded source-to-Spec screening}
\textbf{Verdict:} \texttt{matches}
\\
{\raggedright
\textbf{Reason:} Independent review:\allowbreak{} the uniform large-\allowbreak{}penalty witness supplies the positive-\allowbreak{}constant, strongly-\allowbreak{}convex nonsmooth RRM nonconvergence alternative.\allowbreak{}
\par}
\reviewmatch{Matches source input}
\noindent\textbf{Reviewer annotation}\par
\reviewerbox
\clearpage
\hypertarget{source-claim-e1c86b604c5fef5c}{}
\section*{Review row 9}
\textbf{Source locator:} {\footnotesize\raggedright cited publication, lines 315-\allowbreak{}323\par}
\paragraph{Verbatim source input}
\begin{ReviewVerbatim}
if θ > 12 , and a point mass at 1 if θ < 21 . Clearly there is     Proposition 3.6. Suppose that the distribution map D(·) is
no performatively stable point, and RRM will simply result         ε-sensitive with ε > 0. RRM can fail to converge at all in
in the alternating sequence 1, 0, 1, 0, . . . . The performative   any of the following cases, for any choice of β, γ > 0:
optimum in this case is θPO = 12 .
                                                                    (a) The loss is β-jointly smooth and convex, but not
To show convergence of retraining schemes, it is hence                  strongly convex.
necessary to make a regularity assumption on D(·), such             (b) The loss is γ-strongly convex, but not jointly smooth.
as ε-sensitivity. We are now ready to state our main result         (c) The loss is β-jointly smooth and γ-strongly convex, but
regarding the convergence of repeated risk minimization.                ε > βγ .
\end{ReviewVerbatim}


\paragraph{Formalized review target}
The archival source above is not asserted equivalent to this formalized target.
\begin{ReviewVerbatim}
For every positive epsilon, beta, and gamma with gamma at most beta and epsilon at least gamma divided by beta, there is an epsilon-sensitive performative model with beta-jointly-smooth and gamma-strongly-convex loss whose repeated-risk-minimization trajectory diverges.
\end{ReviewVerbatim}

\textbf{Archival source anchor:} {\footnotesize\raggedright cited publication, lines 315-\allowbreak{}323\par}
\paragraph{Expanded PaperInterface specification}
\begin{ReviewVerbatim}
∀ (epsilon beta gamma : ℝ),
  0 < epsilon →
    0 < beta →
      0 < gamma →
        gamma ≤ beta →
          gamma / beta ≤ epsilon →
            AppliedModelingLib.IsMeasureRepeatedRiskMinimization
                (AppliedModelingLib.proposition36cRegularityMeasureModel epsilon beta gamma)
                (AppliedModelingLib.proposition36cRegularityRRMUpdate epsilon beta gamma) ∧
              AppliedModelingLib.IsJointlySmoothGradient
                  (AppliedModelingLib.proposition36cRegularityGradient beta gamma) beta.toNNReal ∧
                PZMH20PerformativePrediction.IsMeasurePointwiseGradientStronglyConvex
                    (AppliedModelingLib.proposition36cRegularityMeasureModel epsilon beta gamma)
                    (AppliedModelingLib.proposition36cRegularityGradient beta gamma) gamma ∧
                  (∀ (datum candidate : ℝ),
                      HasGradientAt
                        (fun (parameter : ℝ) =>
                          AppliedModelingLib.proposition36cRegularityLoss beta gamma datum parameter)
                        (AppliedModelingLib.proposition36cRegularityGradient beta gamma datum candidate) candidate) ∧
                    PZMH20PerformativePrediction.IsMeasureWassersteinSensitive
                        (AppliedModelingLib.proposition36cRegularityMeasureModel epsilon beta gamma) epsilon ∧
                      ∃ (initial : ℝ),
                        Filter.Tendsto
                          (fun (iteration : ℕ) =>
                            (AppliedModelingLib.proposition36cRegularityRRMUpdate epsilon beta gamma)^[iteration]
                              initial)
                          Filter.atTop Filter.atTop
\end{ReviewVerbatim}

\paragraph{Retained governing declarations}
\begin{itemize}\raggedright
\item \hyperlink{library-prerequisite-2e05b739f5b71f23}{Applied\allowbreak{}Modeling\allowbreak{}Lib.\allowbreak{}Is\allowbreak{}Jointly\allowbreak{}Smooth\allowbreak{}Gradient}
\item \hyperlink{governing-declaration-916f213e8c73a94a}{Applied\allowbreak{}Modeling\allowbreak{}Lib.\allowbreak{}Is\allowbreak{}Measure\allowbreak{}Repeated\allowbreak{}Risk\allowbreak{}Minimization}
\item \hyperlink{governing-declaration-954407b64b7a9ac2}{Applied\allowbreak{}Modeling\allowbreak{}Lib.\allowbreak{}proposition36c\allowbreak{}Regularity\allowbreak{}Gradient}
\item \hyperlink{governing-declaration-9a6823e6fdefa446}{Applied\allowbreak{}Modeling\allowbreak{}Lib.\allowbreak{}proposition36c\allowbreak{}Regularity\allowbreak{}Loss}
\item \hyperlink{governing-declaration-9b198d5b04ab29c5}{Applied\allowbreak{}Modeling\allowbreak{}Lib.\allowbreak{}proposition36c\allowbreak{}Regularity\allowbreak{}Measure\allowbreak{}Model}
\item \hyperlink{governing-declaration-39c6c11ddc030996}{Applied\allowbreak{}Modeling\allowbreak{}Lib.\allowbreak{}proposition36c\allowbreak{}Regularity\allowbreak{}RRM\allowbreak{}Update}
\item \hyperlink{governing-declaration-11c9ab6e3c6f528b}{PZMH20\allowbreak{}Performative\allowbreak{}Prediction.\allowbreak{}Is\allowbreak{}Measure\allowbreak{}Pointwise\allowbreak{}Gradient\allowbreak{}Strongly\allowbreak{}Convex}
\item \hyperlink{paper-prerequisite-2ea2f5769822c5c4}{PZMH20\allowbreak{}Performative\allowbreak{}Prediction.\allowbreak{}Is\allowbreak{}Measure\allowbreak{}Wasserstein\allowbreak{}Sensitive}
\end{itemize}
\paragraph{Lean-elaborated claim roles}\begin{ReviewVerbatim}
1. parameter: ℝ
2. parameter: ℝ
3. parameter: ℝ
4. assumption: 0 < epsilon
5. assumption: 0 < beta
6. assumption: 0 < gamma
7. assumption: gamma ≤ beta
8. assumption: gamma / beta ≤ epsilon
9. conclusion: AppliedModelingLib.IsMeasureRepeatedRiskMinimization
    (AppliedModelingLib.proposition36cRegularityMeasureModel epsilon beta gamma)
    (AppliedModelingLib.proposition36cRegularityRRMUpdate epsilon beta gamma) ∧
  AppliedModelingLib.IsJointlySmoothGradient (AppliedModelingLib.proposition36cRegularityGradient beta gamma)
      beta.toNNReal ∧
    PZMH20PerformativePrediction.IsMeasurePointwiseGradientStronglyConvex
        (AppliedModelingLib.proposition36cRegularityMeasureModel epsilon beta gamma)
        (AppliedModelingLib.proposition36cRegularityGradient beta gamma) gamma ∧
      (∀ (datum candidate : ℝ),
          HasGradientAt
            (fun (parameter : ℝ) => AppliedModelingLib.proposition36cRegularityLoss beta gamma datum parameter)
            (AppliedModelingLib.proposition36cRegularityGradient beta gamma datum candidate) candidate) ∧
        PZMH20PerformativePrediction.IsMeasureWassersteinSensitive
            (AppliedModelingLib.proposition36cRegularityMeasureModel epsilon beta gamma) epsilon ∧
          ∃ (initial : ℝ),
            Filter.Tendsto
              (fun (iteration : ℕ) =>
                (AppliedModelingLib.proposition36cRegularityRRMUpdate epsilon beta gamma)^[iteration] initial)
              Filter.atTop Filter.atTop
\end{ReviewVerbatim}

\paragraph{Lean proof endpoint (built separately)}\begin{ReviewVerbatim}
PZMH20PerformativePrediction.proposition36cFeasibleUniformCounterexample
\end{ReviewVerbatim}

\paragraph{Recorded source-to-Spec screening}
\textbf{Verdict:} \texttt{matches\_approved\_corrected\_target}
\\
{\raggedright
\textbf{Reason:} Compared the selected proposition\_\allowbreak{}3\_\allowbreak{}6\_\allowbreak{}counterexamples\_\allowbreak{}c source item, its identity-\allowbreak{}bound approved correction for proposition-\allowbreak{}3.\allowbreak{}6c-\allowbreak{}arbitrary-\allowbreak{}regularity-\allowbreak{}constants, the complete expanded Spec and claim roles, and all 17 displayed prerequisites.\allowbreak{} PMLR PDF page 5 visibly states epsilon >= gamma/\allowbreak{}beta;\allowbreak{} the extracted text loses the non-\allowbreak{}strict glyph and fraction layout.\allowbreak{} The approved correction expressly adds gamma <= beta and disclaims archival equivalence.\allowbreak{} For every positive epsilon, beta, gamma satisfying these two inequalities, the displayed real-\allowbreak{}valued model has D(theta) = delta\_\allowbreak{}(1 + epsilon*theta), loss gamma*x\textasciicircum{}2/\allowbreak{}2 -\allowbreak{} beta*z*x, and actual gradient gamma*x -\allowbreak{} beta*z.\allowbreak{} Its parameter-\allowbreak{}gradient and data-\allowbreak{}gradient Lipschitz constants are respectively gamma <= beta and beta, and its gamma-\allowbreak{}strong-\allowbreak{}convexity inequality holds with the required gamma/\allowbreak{}2 remainder.\allowbreak{} Dirac laws make every frozen loss integrable and give W1 = epsilon*|theta-\allowbreak{}theta\_\allowbreak{}prime| through the displayed finite-\allowbreak{}cost coupling definition.\allowbreak{} Frozen risk is minimized over all real candidates by G(theta) = (beta/\allowbreak{}gamma)*(1 + epsilon*theta), exactly the expanded RRM predicate.\allowbreak{} Writing G(theta) = a*theta + b gives a = epsilon*beta/\allowbreak{}gamma >= 1 and b = beta/\allowbreak{}gamma > 0;\allowbreak{} from initial 0, the iterates are at least n*b and tend to +infinity.\allowbreak{} Equality epsilon = gamma/\allowbreak{}beta is therefore covered by positive translation.\allowbreak{} The complete quantified claim matches the approved corrected target, with no additional restriction on its admissible constants and no archival-\allowbreak{}equivalence claim.\allowbreak{}
\par}
\reviewmatch{Matches formalized target}
\noindent\textbf{Reviewer annotation}\par
\reviewerbox
\clearpage
\hypertarget{source-claim-a27351d86a4be6d9}{}
\section*{Review row 10}
\textbf{Source locator:} {\footnotesize\raggedright cited publication, lines 373-\allowbreak{}415\par}
\paragraph{Verbatim source input}
\begin{ReviewVerbatim}
Theorem 3.10. Suppose that the loss `(z; θ) is β-jointly
Theorem 3.8. Suppose that the loss `(z; θ) is β-jointly                       smooth (A1) and γ-strongly convex (A2), and that there exist
smooth (A1) and γ-strongly convex (A2). If the distribution                                                def R        α
                                        γ                                     α > 1, µ > 0 such that ξα,µ = Rm eµ|x| dD(θ) is finite
map D(·) is ε-sensitive with ε < (β+γ)(1+1.5ηβ) , then RGD
                                                                              ∀θ ∈ Θ. Let δ ∈ (0, 1) be a radius of convergence.     Con-
                     2
with step size η 6 β+γ satisfies the following:
                                                                                                                                    � t
                                                                                                                            1
                                                                              sider running RERM or RGD with nt = O (εδ)m log p
(a) kGgd (θ) − Ggd (θ0 )k2                                                    samples at time t.
                                                                                                                            γ
                                                                              (a) If the map D(·) is ε-sensitive with ε < 2β  , then with
                         βγ                                          0
      6       1−η           − ε(1.5ηβ 2 + β)             kθ − θ k2 .              probability 1 − p, RERM satisfies,
                        β+γ
                                                                                                                  log 1δ kθ0 − θPS k2
                                                                                                                     �
(b) The iterates θt of RGD converge to a unique performa-                           kθt − θPS k2 6 δ, for all t >                       .
    tively stable point θPS at a linear rate,                                                                           1 − 2εβ
                                                                                                                              γ

                                                                                                                                 γ
                                               �1                             (b) If the map D(·) is ε-sensitive with ε < (β+γ)(1+1.5ηβ)  ,
                                         log    δ kθ0 − θPS k2
     kθt − θPS k2 6 δ for t >                                             .       then with probability 1 − p, REGD with satisfies,
                                          βγ          2
                                 η       β+γ − ε(1.5ηβ + β)
                                                                                                                    log 1δ kθ0 − θPS k2
                                                                                                                        �
                                                                                  kθt −θPS k2 6 δ, for all t >                              ,
The conclusion of Theorem 3.8 is a strict generalization of a                                                        βγ
                                                                                                                 η β+γ    − ε(3ηβ 2 + 2β)
classical optimization result which considers a static objec-

                                                        βγ
tive, in which case the rate of contraction is 1 − η β+γ                                                                   2
                                                                                   for a constant choice of step size η 6 β+γ .
(see for example Theorem 2.1.15 in Nesterov (2013) or
Lemma 3.7 in Hardt et al. (2016b)). Our rate exactly                          Proof sketch. The basic idea behind these results is the fol-
matches this standard result in the case that ε = 0. The                      lowing. While kθt − θPS k2 > δ, the sample size nt is
proof of Theorem 3.8 can be found in Appendix E.3.                            sufficiently large to ensure a behavior similar to that on a

[Next verbatim source excerpt]

Throughout our presentation, we focus on predictive models        Here, ε represents the performative aspect of the model. If
that are parametrized by a vector θ ∈ Θ, where the param-         ε = 0, outcomes are independent of the assigned scores
eter space Θ ⊆ Rd is a closed, convex set. We use capi-           and the problem reduces to standard supervised learning
tal letters to denote random variables and their lowercase        where the optimal predictor is the conditional expectation
counterparts to denote realizations of these variables. We        fθSL (x) = E[Y | X = x] = 12 + µx, with θSL = µ.
consider instances z = (x, y) defined as feature, outcome          In the performative setting with ε 6= 0, the optimal predic-
pairs, where x ∈ Rm−1 and y ∈ R. Whenever we define a              tion θPO balances between its predictive accuracy as well

[Next verbatim source excerpt]

formatively optimal classifiers need not be performatively            Definition 3.1 (ε-sensitivity). We say that a distribution
stable and performatively stable classifiers need not be per-         map D(·) is ε-sensitive if for all θ, θ0 ∈ Θ:
formatively optimal. We illustrate this point in the context
                                                                                   W1 D(θ), D(θ0 ) 6 εkθ − θ0 k2 ,
                                                                                     �
of our previous biased coin toss example.
Example 2.2 (continued). Consider again our model of
a biased coin toss. In order for a classifier fθ to be perfor-        where W1 denotes the earth mover’s distance.

[Next verbatim source excerpt]

In addition to this assumption on the distribution map, we          (b) If ε < βγ , the iterates θt of RRM converge to a unique
will often make standard assumptions on the loss function               performatively stable point θPS at a linear rate,
`(z; θ) which hold for broad classes of losses. Each tech-
nical result to follow will invoke some subset of them. To                                                log( 1δ kθ0 − θPS k2 )
                                 def                                         kθt − θPS k2 6 δ for t >                              .
simplify our presentation, let Z = ∪θ∈Θ supp(D(θ)).                                                             (1 − ε βγ )
(A1) (joint smoothness) A loss function `(z; θ) is β-jointly       The main message of this theorem is that in performative
     smooth if ∀θ, θ0 ∈ Θ and z, z 0 ∈ Z:                          prediction, if the loss function is sufficiently “nice” and the
                                                                   distribution map is sufficiently (in)sensitive, then one need
           krθ `(z; θ) − rθ `(z; θ0 )k2 6 β kθ − θ0 k2 ,
                                                                   only retrain a classifier a small number of times before it
           krθ `(z; θ) − rθ `(z 0 ; θ)k2 6 β kz − z 0 k2 .         converges to a unique stable point. Here, we provide a proof
                                                                   sketch illustrating the main ideas behind the theorem and
(A2) (strong convexity) A loss function `(z; θ) is γ-
                                                                   defer the full proof to Appendix E.1.
     strongly convex if ∀ θ, θ0 ∈ Θ and z ∈ Z:
                                                γ         2        Proof Sketch. Fix θ, θ0 ∈ Θ. Let f (ϕ) = EZ∼D(θ) `(Z; ϕ)
      `(z; θ) > `(z; θ0 )+rθ `(z; θ0 )> (θ−θ0 )+ kθ − θ0 k2 .
                                                2                  and f 0 (ϕ) = EZ∼D(θ0 ) `(Z; ϕ). By applying standard prop-
                                                                   erties of strong convexity and the fact that G(θ) is the unique
If γ = 0, this last assumption is equivalent to convexity. We

[Next verbatim source excerpt]

                                                                              We now extend our main results regarding the convergence
globally optimal. However, in our framework, strong con-
                                                                              of RRM and RGD to the finite-sample regime. To do so,
vexity is in fact necessary to guarantee convergence of re-
                                                                              we leverage the fact that, under mild regularity conditions,
peated risk minimization, even for arbitrarily smooth losses
                                                                              the empirical distribution Dn given by n samples drawn
and an arbitrarily small sensitivity parameter.
                                                                              i.i.d. from a true distribution D is with high probability
                                                                              close to D in earth mover’s distance (Fournier & Guillin,
3.3. Repeated Gradient Descent
                                                                              2015). We begin by defining the finite-sample version of
Theorem 3.5 demonstrates that repeated risk minimization                      these procedures.
converges to a unique stable point if the sensitivity parameter               Definition 3.9 (RERM & REGD). Define repeated empir-
ε is small enough. However, implementing RRM requires                         ical risk minimization (RERM) to be the procedure where
access to an exact optimization oracle. We now relax this                     starting from a classifier fθ0 at every iteration t > 0, we
requirement and demonstrate how a simple gradient descent                     collect nt samples from D(θt ) and perform the update:
algorithm also converges to a unique stable point.
                                                                                                       def
Definition 3.7 (RGD). Repeated gradient descent (RGD) is                            θt+1 = Gnt (θt ) = arg min              E       `(Z; θ).
                                                                                                             θ      Z∼D nt (θt )
the procedure where, starting from an initial model fθ0 we
perform the following sequence of updates for t > 0:                          Similarly, define repeated empirical gradient descent
                                                                              (REGD) to be the optimization procedure with update rule:
                        def
   θt+1 = Ggd (θt ) = ΠΘ (θt − η               E    rθ `(Z; θt )),                               def
                                         Z∼D(θt )                               θt+1 = Gngdt (θt ) = ΠΘ (θt − η        E         rθ `(Z; θt )).
                                                                                                                  Z∼D nt (θt )

where η > 0 is a fixed step size and ΠΘ denotes the Eu-
                                                                              Here, η > 0 is a step size and ΠΘ denotes the Euclidean
clidean projection operator onto Θ.
                                                                              projection operator onto Θ.

[Next verbatim source excerpt]

Lemma 3.7 in Hardt et al. (2016b)). Our rate exactly                          Proof sketch. The basic idea behind these results is the fol-
matches this standard result in the case that ε = 0. The                      lowing. While kθt − θPS k2 > δ, the sample size nt is
proof of Theorem 3.8 can be found in Appendix E.3.                            sufficiently large to ensure a behavior similar to that on a

[form-feed in source extraction]
                                                        Performative Prediction

population level: as in Theorems 3.5 and 3.8, the iterates θt         D(·) is ε-sensitive. Then, for every performatively stable
contract toward θPS . This implies that θt eventually enters a        point θPS and every performative optimum θPO :
δ-ball around θPS , for some large enough t. Once this hap-
                                                                                                            2Lz ε
pens, contrary to population-level results, a contraction is                            kθPO − θPS k2 6           .
no longer guaranteed due to the noise inherent in observing                                                  γ
only finite-sample approximations of D(θt ). Nevertheless,
the sample size nt is sufficiently large to ensure that θt             This result shows that in cases where repeated risk mini-
cannot escape a δ-ball around θPS either.                             mization converges to a stable point, the resulting classifier
\end{ReviewVerbatim}


\paragraph{Formalized review target}
The archival source above is not asserted equivalent to this formalized target.
\begin{ReviewVerbatim}
In data dimension greater than two, use the explicit supercritical batch-count schedule, evaluated at the canonical capped deviation, together with the inverse-square confidence schedule. Suppose the input data provide the stated smoothness, strong-convexity, Wasserstein-sensitivity, and small-sensitivity conditions; a convex parameter domain; a shared bounded exponential-moment and tail envelope along every adaptive history; measurable adaptive sampling and bad events; data-coordinate Lipschitzness, population-gradient integrability, local loss measurability, pointwise loss differentiation, and empirical-loss integrability; a checked RERM trajectory and population RRM update with a performatively stable point; and the stated numerical W1 budget, outer-contraction, absorbing-radius, and initial-distance conditions. Then there exists an entry iteration such that, under the adaptive conditionally IID sampling law, with probability at least one minus p, every RERM iterate at or after that entry iteration lies within the prescribed radius of the stable point. The model and expected-loss integrability are defined for every ambient parameter. Joint-gradient smoothness and pointwise strong convexity quantify over all ambient parameters and data points; Wasserstein sensitivity, gradient integrability, and loss differentiation quantify over all ambient parameters. RERM data-loss Lipschitzness also quantifies over every ambient parameter and data point. These global hypotheses are stronger in scope than the source assumptions on the parameter domain and the union of its induced supports; their necessity for the source conclusion is unresolved.
\end{ReviewVerbatim}

\textbf{Archival source anchor:} {\footnotesize\raggedright cited publication, lines 373-\allowbreak{}415\par}
\paragraph{Expanded PaperInterface specification}
\begin{ReviewVerbatim}
∀ {Parameter : Type u_1} [inst : MeasurableSpace Parameter] [inst_1 : NormedAddCommGroup Parameter]
  [inst_2 : InnerProductSpace ℝ Parameter] [inst_3 : CompleteSpace Parameter]
  (data : PZMH20PerformativePrediction.Theorem310CorrectedRERMAllRoundTrajectoryData Parameter),
  ∃ (entryIteration : ℕ),
    (AppliedModelingLib.measurePerformativeHeterogeneousBatchTraceInfiniteLaw data.model data.sampling
            (PZMH20PerformativePrediction.theorem310CorrectedRERMCountSchedule data.dimension data.cutoff data.eta
              data.alpha data.gamma data.momentBound data.tailBound data.scaledDeviation data.p data.headTolerance)
            data.deployedOfHistory data.hmeasurableDeployed).real
        {trace :
          (i : ℕ) →
            AppliedModelingLib.HeterogeneousBatchTraceCoordinate
              (fun (round : ℕ) =>
                Fin
                  (PZMH20PerformativePrediction.theorem310CorrectedRERMCountSchedule data.dimension data.cutoff data.eta
                    data.alpha data.gamma data.momentBound data.tailBound data.scaledDeviation data.p data.headTolerance
                    round))
              (EuclideanSpace ℝ (Fin data.dimension)) i |
          ∀ (iteration : ℕ),
            entryIteration ≤ iteration →
              dist
                  (data.deployedOfHistory iteration
                    (AppliedModelingLib.heterogeneousBatchTraceOfInfiniteTrace iteration trace))
                  data.stable ≤
                data.radius} ≥
      1 - data.p
\end{ReviewVerbatim}

\paragraph{Retained governing declarations}
\begin{itemize}\raggedright
\item \hyperlink{governing-declaration-37d7ae049de215a7}{Applied\allowbreak{}Modeling\allowbreak{}Lib.\allowbreak{}Heterogeneous\allowbreak{}Batch\allowbreak{}Trace\allowbreak{}Coordinate}
\item \hyperlink{governing-declaration-92d1582063c434f3}{Applied\allowbreak{}Modeling\allowbreak{}Lib.\allowbreak{}heterogeneous\allowbreak{}Batch\allowbreak{}Trace\allowbreak{}Coordinate\allowbreak{}Measurable\allowbreak{}Space}
\item \hyperlink{governing-declaration-644c195e65179f85}{Applied\allowbreak{}Modeling\allowbreak{}Lib.\allowbreak{}heterogeneous\allowbreak{}Batch\allowbreak{}Trace\allowbreak{}Of\allowbreak{}Infinite\allowbreak{}Trace}
\item \hyperlink{governing-declaration-2833c0091dd3eed5}{Applied\allowbreak{}Modeling\allowbreak{}Lib.\allowbreak{}measure\allowbreak{}Performative\allowbreak{}Heterogeneous\allowbreak{}Batch\allowbreak{}Trace\allowbreak{}Infinite\allowbreak{}Law}
\item \hyperlink{paper-prerequisite-613020783cc01a15}{PZMH20\allowbreak{}Performative\allowbreak{}Prediction.\allowbreak{}Theorem310\allowbreak{}Corrected\allowbreak{}RERM\allowbreak{}All\allowbreak{}Round\allowbreak{}Trajectory\allowbreak{}Data}
\item \hyperlink{governing-declaration-f4e329b046167e0c}{PZMH20\allowbreak{}Performative\allowbreak{}Prediction.\allowbreak{}theorem310\allowbreak{}Corrected\allowbreak{}RERM\allowbreak{}Count\allowbreak{}Schedule}
\end{itemize}
\paragraph{Lean-elaborated claim roles}\begin{ReviewVerbatim}
1. parameter: Type u_1
2. parameter: MeasurableSpace Parameter
3. parameter: NormedAddCommGroup Parameter
4. parameter: InnerProductSpace ℝ Parameter
5. assumption: CompleteSpace Parameter
6. parameter: PZMH20PerformativePrediction.Theorem310CorrectedRERMAllRoundTrajectoryData Parameter
7. conclusion: ∃ (entryIteration : ℕ),
  (AppliedModelingLib.measurePerformativeHeterogeneousBatchTraceInfiniteLaw data.model data.sampling
          (PZMH20PerformativePrediction.theorem310CorrectedRERMCountSchedule data.dimension data.cutoff data.eta
            data.alpha data.gamma data.momentBound data.tailBound data.scaledDeviation data.p data.headTolerance)
          data.deployedOfHistory data.hmeasurableDeployed).real
      {trace :
        (i : ℕ) →
          AppliedModelingLib.HeterogeneousBatchTraceCoordinate
            (fun (round : ℕ) =>
              Fin
                (PZMH20PerformativePrediction.theorem310CorrectedRERMCountSchedule data.dimension data.cutoff data.eta
                  data.alpha data.gamma data.momentBound data.tailBound data.scaledDeviation data.p data.headTolerance
                  round))
            (EuclideanSpace ℝ (Fin data.dimension)) i |
        ∀ (iteration : ℕ),
          entryIteration ≤ iteration →
            dist
                (data.deployedOfHistory iteration
                  (AppliedModelingLib.heterogeneousBatchTraceOfInfiniteTrace iteration trace))
                data.stable ≤
              data.radius} ≥
    1 - data.p
\end{ReviewVerbatim}

\paragraph{Lean proof endpoint (built separately)}\begin{ReviewVerbatim}
PZMH20PerformativePrediction.theorem310CorrectedRERMAllRoundTrajectoryReview
\end{ReviewVerbatim}

\paragraph{Recorded source-to-Spec screening}
\textbf{Verdict:} \texttt{matches\_approved\_corrected\_target}
\\
{\raggedright
\textbf{Reason:} The selected route displays the archival Theorem 3.\allowbreak{}10(a) text together with a complete approved corrected target, its three defect identifiers, scope note, and exact approval record, and it expressly disclaims archival equivalence.\allowbreak{} The carrier exposes dimension > 2, the capped supercritical count and inverse-\allowbreak{}square confidence schedules, a uniform adaptive moment/\allowbreak{}tail envelope, global smoothness, strong-\allowbreak{}convexity, Wasserstein, loss-\allowbreak{}gradient, measurability and integrability premises, the RERM and population-\allowbreak{}RRM trajectories, a stable point, and the numerical contraction, absorption, budget, radius, and initial-\allowbreak{}distance conditions named by that corrected target.\allowbreak{} The Spec then quantifies over every such carrier and chooses one entry iteration outside the event;\allowbreak{} under the displayed adaptive Markov-\allowbreak{}kernel law whose round batch is conditionally IID from the deployed data law, the event requires every iterate at or after entry to lie within radius, with measure at least 1 -\allowbreak{} p.\allowbreak{} The broader p > 0 input also includes every source-\allowbreak{}relevant 0 < p < 1 case without making those cases vacuous.\allowbreak{}
\par}
\reviewmatch{Matches formalized target}
\noindent\textbf{Reviewer annotation}\par
\reviewerbox
\clearpage
\hypertarget{source-claim-2e42d97b744932d5}{}
\section*{Review row 11}
\textbf{Source locator:} {\footnotesize\raggedright cited publication, lines 373-\allowbreak{}415\par}
\paragraph{Verbatim source input}
\begin{ReviewVerbatim}
Theorem 3.10. Suppose that the loss `(z; θ) is β-jointly
Theorem 3.8. Suppose that the loss `(z; θ) is β-jointly                       smooth (A1) and γ-strongly convex (A2), and that there exist
smooth (A1) and γ-strongly convex (A2). If the distribution                                                def R        α
                                        γ                                     α > 1, µ > 0 such that ξα,µ = Rm eµ|x| dD(θ) is finite
map D(·) is ε-sensitive with ε < (β+γ)(1+1.5ηβ) , then RGD
                                                                              ∀θ ∈ Θ. Let δ ∈ (0, 1) be a radius of convergence.     Con-
                     2
with step size η 6 β+γ satisfies the following:
                                                                                                                                    � t
                                                                                                                            1
                                                                              sider running RERM or RGD with nt = O (εδ)m log p
(a) kGgd (θ) − Ggd (θ0 )k2                                                    samples at time t.
                                                                                                                            γ
                                                                              (a) If the map D(·) is ε-sensitive with ε < 2β  , then with
                         βγ                                          0
      6       1−η           − ε(1.5ηβ 2 + β)             kθ − θ k2 .              probability 1 − p, RERM satisfies,
                        β+γ
                                                                                                                  log 1δ kθ0 − θPS k2
                                                                                                                     �
(b) The iterates θt of RGD converge to a unique performa-                           kθt − θPS k2 6 δ, for all t >                       .
    tively stable point θPS at a linear rate,                                                                           1 − 2εβ
                                                                                                                              γ

                                                                                                                                 γ
                                               �1                             (b) If the map D(·) is ε-sensitive with ε < (β+γ)(1+1.5ηβ)  ,
                                         log    δ kθ0 − θPS k2
     kθt − θPS k2 6 δ for t >                                             .       then with probability 1 − p, REGD with satisfies,
                                          βγ          2
                                 η       β+γ − ε(1.5ηβ + β)
                                                                                                                    log 1δ kθ0 − θPS k2
                                                                                                                        �
                                                                                  kθt −θPS k2 6 δ, for all t >                              ,
The conclusion of Theorem 3.8 is a strict generalization of a                                                        βγ
                                                                                                                 η β+γ    − ε(3ηβ 2 + 2β)
classical optimization result which considers a static objec-

                                                        βγ
tive, in which case the rate of contraction is 1 − η β+γ                                                                   2
                                                                                   for a constant choice of step size η 6 β+γ .
(see for example Theorem 2.1.15 in Nesterov (2013) or
Lemma 3.7 in Hardt et al. (2016b)). Our rate exactly                          Proof sketch. The basic idea behind these results is the fol-
matches this standard result in the case that ε = 0. The                      lowing. While kθt − θPS k2 > δ, the sample size nt is
proof of Theorem 3.8 can be found in Appendix E.3.                            sufficiently large to ensure a behavior similar to that on a

[Next verbatim source excerpt]

Throughout our presentation, we focus on predictive models        Here, ε represents the performative aspect of the model. If
that are parametrized by a vector θ ∈ Θ, where the param-         ε = 0, outcomes are independent of the assigned scores
eter space Θ ⊆ Rd is a closed, convex set. We use capi-           and the problem reduces to standard supervised learning
tal letters to denote random variables and their lowercase        where the optimal predictor is the conditional expectation
counterparts to denote realizations of these variables. We        fθSL (x) = E[Y | X = x] = 12 + µx, with θSL = µ.
consider instances z = (x, y) defined as feature, outcome          In the performative setting with ε 6= 0, the optimal predic-
pairs, where x ∈ Rm−1 and y ∈ R. Whenever we define a              tion θPO balances between its predictive accuracy as well

[Next verbatim source excerpt]

formatively optimal classifiers need not be performatively            Definition 3.1 (ε-sensitivity). We say that a distribution
stable and performatively stable classifiers need not be per-         map D(·) is ε-sensitive if for all θ, θ0 ∈ Θ:
formatively optimal. We illustrate this point in the context
                                                                                   W1 D(θ), D(θ0 ) 6 εkθ − θ0 k2 ,
                                                                                     �
of our previous biased coin toss example.
Example 2.2 (continued). Consider again our model of
a biased coin toss. In order for a classifier fθ to be perfor-        where W1 denotes the earth mover’s distance.

[Next verbatim source excerpt]

In addition to this assumption on the distribution map, we          (b) If ε < βγ , the iterates θt of RRM converge to a unique
will often make standard assumptions on the loss function               performatively stable point θPS at a linear rate,
`(z; θ) which hold for broad classes of losses. Each tech-
nical result to follow will invoke some subset of them. To                                                log( 1δ kθ0 − θPS k2 )
                                 def                                         kθt − θPS k2 6 δ for t >                              .
simplify our presentation, let Z = ∪θ∈Θ supp(D(θ)).                                                             (1 − ε βγ )
(A1) (joint smoothness) A loss function `(z; θ) is β-jointly       The main message of this theorem is that in performative
     smooth if ∀θ, θ0 ∈ Θ and z, z 0 ∈ Z:                          prediction, if the loss function is sufficiently “nice” and the
                                                                   distribution map is sufficiently (in)sensitive, then one need
           krθ `(z; θ) − rθ `(z; θ0 )k2 6 β kθ − θ0 k2 ,
                                                                   only retrain a classifier a small number of times before it
           krθ `(z; θ) − rθ `(z 0 ; θ)k2 6 β kz − z 0 k2 .         converges to a unique stable point. Here, we provide a proof
                                                                   sketch illustrating the main ideas behind the theorem and
(A2) (strong convexity) A loss function `(z; θ) is γ-
                                                                   defer the full proof to Appendix E.1.
     strongly convex if ∀ θ, θ0 ∈ Θ and z ∈ Z:
                                                γ         2        Proof Sketch. Fix θ, θ0 ∈ Θ. Let f (ϕ) = EZ∼D(θ) `(Z; ϕ)
      `(z; θ) > `(z; θ0 )+rθ `(z; θ0 )> (θ−θ0 )+ kθ − θ0 k2 .
                                                2                  and f 0 (ϕ) = EZ∼D(θ0 ) `(Z; ϕ). By applying standard prop-
                                                                   erties of strong convexity and the fact that G(θ) is the unique
If γ = 0, this last assumption is equivalent to convexity. We

[Next verbatim source excerpt]

                                                                              We now extend our main results regarding the convergence
globally optimal. However, in our framework, strong con-
                                                                              of RRM and RGD to the finite-sample regime. To do so,
vexity is in fact necessary to guarantee convergence of re-
                                                                              we leverage the fact that, under mild regularity conditions,
peated risk minimization, even for arbitrarily smooth losses
                                                                              the empirical distribution Dn given by n samples drawn
and an arbitrarily small sensitivity parameter.
                                                                              i.i.d. from a true distribution D is with high probability
                                                                              close to D in earth mover’s distance (Fournier & Guillin,
3.3. Repeated Gradient Descent
                                                                              2015). We begin by defining the finite-sample version of
Theorem 3.5 demonstrates that repeated risk minimization                      these procedures.
converges to a unique stable point if the sensitivity parameter               Definition 3.9 (RERM & REGD). Define repeated empir-
ε is small enough. However, implementing RRM requires                         ical risk minimization (RERM) to be the procedure where
access to an exact optimization oracle. We now relax this                     starting from a classifier fθ0 at every iteration t > 0, we
requirement and demonstrate how a simple gradient descent                     collect nt samples from D(θt ) and perform the update:
algorithm also converges to a unique stable point.
                                                                                                       def
Definition 3.7 (RGD). Repeated gradient descent (RGD) is                            θt+1 = Gnt (θt ) = arg min              E       `(Z; θ).
                                                                                                             θ      Z∼D nt (θt )
the procedure where, starting from an initial model fθ0 we
perform the following sequence of updates for t > 0:                          Similarly, define repeated empirical gradient descent
                                                                              (REGD) to be the optimization procedure with update rule:
                        def
   θt+1 = Ggd (θt ) = ΠΘ (θt − η               E    rθ `(Z; θt )),                               def
                                         Z∼D(θt )                               θt+1 = Gngdt (θt ) = ΠΘ (θt − η        E         rθ `(Z; θt )).
                                                                                                                  Z∼D nt (θt )

where η > 0 is a fixed step size and ΠΘ denotes the Eu-
                                                                              Here, η > 0 is a step size and ΠΘ denotes the Euclidean
clidean projection operator onto Θ.
                                                                              projection operator onto Θ.

[Next verbatim source excerpt]

Lemma 3.7 in Hardt et al. (2016b)). Our rate exactly                          Proof sketch. The basic idea behind these results is the fol-
matches this standard result in the case that ε = 0. The                      lowing. While kθt − θPS k2 > δ, the sample size nt is
proof of Theorem 3.8 can be found in Appendix E.3.                            sufficiently large to ensure a behavior similar to that on a

[form-feed in source extraction]
                                                        Performative Prediction

population level: as in Theorems 3.5 and 3.8, the iterates θt         D(·) is ε-sensitive. Then, for every performatively stable
contract toward θPS . This implies that θt eventually enters a        point θPS and every performative optimum θPO :
δ-ball around θPS , for some large enough t. Once this hap-
                                                                                                            2Lz ε
pens, contrary to population-level results, a contraction is                            kθPO − θPS k2 6           .
no longer guaranteed due to the noise inherent in observing                                                  γ
only finite-sample approximations of D(θt ). Nevertheless,
the sample size nt is sufficiently large to ensure that θt             This result shows that in cases where repeated risk mini-
cannot escape a δ-ball around θPS either.                             mization converges to a stable point, the resulting classifier
\end{ReviewVerbatim}


\paragraph{Formalized review target}
The archival source above is not asserted equivalent to this formalized target.
\begin{ReviewVerbatim}
In data dimension greater than two, use the explicit supercritical batch-count schedule, evaluated at the canonical capped deviation, together with the inverse-square confidence schedule. Suppose the input data provide the stated joint-gradient smoothness, pointwise strong convexity, Wasserstein sensitivity, step-size, and small-sensitivity conditions; a convex parameter domain with a variational nonexpansive projection; a shared bounded exponential-moment and tail envelope along every adaptive history; measurable adaptive sampling and bad events; population-gradient integrability, local loss measurability, pointwise loss differentiation, and empirical-gradient integrability; a checked REGD trajectory with a performatively stable point; and the stated numerical W1 budget, outer-contraction, absorbing-radius, and initial-distance conditions. Then there exists an entry iteration such that, under the adaptive conditionally IID sampling law, with probability at least one minus p, every REGD iterate at or after that entry iteration lies within the prescribed radius of the stable point. The model and expected-loss integrability are defined for every ambient parameter. Joint-gradient smoothness and pointwise strong convexity quantify over all ambient parameters and data points; Wasserstein sensitivity, gradient integrability, and loss differentiation quantify over all ambient parameters. These global hypotheses are stronger in scope than the source assumptions on the parameter domain and the union of its induced supports; their necessity for the source conclusion is unresolved.
\end{ReviewVerbatim}

\textbf{Archival source anchor:} {\footnotesize\raggedright cited publication, lines 373-\allowbreak{}415\par}
\paragraph{Expanded PaperInterface specification}
\begin{ReviewVerbatim}
∀ {Parameter : Type u_1} [inst : MeasurableSpace Parameter] [inst_1 : NormedAddCommGroup Parameter]
  [inst_2 : InnerProductSpace ℝ Parameter] [inst_3 : CompleteSpace Parameter]
  (data : PZMH20PerformativePrediction.Theorem310CorrectedRGDAllRoundTrajectoryData Parameter),
  ∃ (entryIteration : ℕ),
    (AppliedModelingLib.measurePerformativeHeterogeneousBatchTraceInfiniteLaw data.model data.sampling
            (AppliedModelingLib.pOneFournierGuillinTheorem310SupercriticalCappedConcreteCountSchedule data.dimension
              data.cutoff data.eta data.alpha data.gamma data.momentBound data.tailBound data.scaledDeviation data.p
              data.headTolerance)
            data.deployedOfHistory data.hmeasurableDeployed).real
        {trace :
          (i : ℕ) →
            AppliedModelingLib.HeterogeneousBatchTraceCoordinate
              (fun (round : ℕ) =>
                Fin
                  (AppliedModelingLib.pOneFournierGuillinTheorem310SupercriticalCappedConcreteCountSchedule
                    data.dimension data.cutoff data.eta data.alpha data.gamma data.momentBound data.tailBound
                    data.scaledDeviation data.p data.headTolerance round))
              (EuclideanSpace ℝ (Fin data.dimension)) i |
          ∀ (iteration : ℕ),
            entryIteration ≤ iteration →
              dist
                  (data.deployedOfHistory iteration
                    (AppliedModelingLib.heterogeneousBatchTraceOfInfiniteTrace iteration trace))
                  data.stable ≤
                data.radius} ≥
      1 - data.p
\end{ReviewVerbatim}

\paragraph{Retained governing declarations}
\begin{itemize}\raggedright
\item \hyperlink{governing-declaration-37d7ae049de215a7}{Applied\allowbreak{}Modeling\allowbreak{}Lib.\allowbreak{}Heterogeneous\allowbreak{}Batch\allowbreak{}Trace\allowbreak{}Coordinate}
\item \hyperlink{governing-declaration-92d1582063c434f3}{Applied\allowbreak{}Modeling\allowbreak{}Lib.\allowbreak{}heterogeneous\allowbreak{}Batch\allowbreak{}Trace\allowbreak{}Coordinate\allowbreak{}Measurable\allowbreak{}Space}
\item \hyperlink{governing-declaration-644c195e65179f85}{Applied\allowbreak{}Modeling\allowbreak{}Lib.\allowbreak{}heterogeneous\allowbreak{}Batch\allowbreak{}Trace\allowbreak{}Of\allowbreak{}Infinite\allowbreak{}Trace}
\item \hyperlink{governing-declaration-2833c0091dd3eed5}{Applied\allowbreak{}Modeling\allowbreak{}Lib.\allowbreak{}measure\allowbreak{}Performative\allowbreak{}Heterogeneous\allowbreak{}Batch\allowbreak{}Trace\allowbreak{}Infinite\allowbreak{}Law}
\item \hyperlink{governing-declaration-3cf6b14c7bcca80c}{Applied\allowbreak{}Modeling\allowbreak{}Lib.\allowbreak{}p\allowbreak{}One\allowbreak{}Fournier\allowbreak{}Guillin\allowbreak{}Theorem310\allowbreak{}Supercritical\allowbreak{}Capped\allowbreak{}Concrete\allowbreak{}Count\allowbreak{}Schedule}
\item \hyperlink{paper-prerequisite-5a443fd08b184a8b}{PZMH20\allowbreak{}Performative\allowbreak{}Prediction.\allowbreak{}Theorem310\allowbreak{}Corrected\allowbreak{}RGD\allowbreak{}All\allowbreak{}Round\allowbreak{}Trajectory\allowbreak{}Data}
\end{itemize}
\paragraph{Lean-elaborated claim roles}\begin{ReviewVerbatim}
1. parameter: Type u_1
2. parameter: MeasurableSpace Parameter
3. parameter: NormedAddCommGroup Parameter
4. parameter: InnerProductSpace ℝ Parameter
5. assumption: CompleteSpace Parameter
6. parameter: PZMH20PerformativePrediction.Theorem310CorrectedRGDAllRoundTrajectoryData Parameter
7. conclusion: ∃ (entryIteration : ℕ),
  (AppliedModelingLib.measurePerformativeHeterogeneousBatchTraceInfiniteLaw data.model data.sampling
          (AppliedModelingLib.pOneFournierGuillinTheorem310SupercriticalCappedConcreteCountSchedule data.dimension
            data.cutoff data.eta data.alpha data.gamma data.momentBound data.tailBound data.scaledDeviation data.p
            data.headTolerance)
          data.deployedOfHistory data.hmeasurableDeployed).real
      {trace :
        (i : ℕ) →
          AppliedModelingLib.HeterogeneousBatchTraceCoordinate
            (fun (round : ℕ) =>
              Fin
                (AppliedModelingLib.pOneFournierGuillinTheorem310SupercriticalCappedConcreteCountSchedule data.dimension
                  data.cutoff data.eta data.alpha data.gamma data.momentBound data.tailBound data.scaledDeviation data.p
                  data.headTolerance round))
            (EuclideanSpace ℝ (Fin data.dimension)) i |
        ∀ (iteration : ℕ),
          entryIteration ≤ iteration →
            dist
                (data.deployedOfHistory iteration
                  (AppliedModelingLib.heterogeneousBatchTraceOfInfiniteTrace iteration trace))
                data.stable ≤
              data.radius} ≥
    1 - data.p
\end{ReviewVerbatim}

\paragraph{Lean proof endpoint (built separately)}\begin{ReviewVerbatim}
PZMH20PerformativePrediction.theorem310CorrectedRGDAllRoundTrajectoryReview
\end{ReviewVerbatim}

\paragraph{Recorded source-to-Spec screening}
\textbf{Verdict:} \texttt{matches\_approved\_corrected\_target}
\\
{\raggedright
\textbf{Reason:} The selected route displays the archival Theorem 3.\allowbreak{}10(b) text together with a complete approved corrected target, its two defect identifiers, scope note, and exact approval record, and it expressly disclaims archival equivalence.\allowbreak{} The carrier exposes dimension > 2, the capped supercritical count and inverse-\allowbreak{}square confidence schedules, a uniform adaptive moment/\allowbreak{}tail envelope, global joint smoothness, strong convexity, Wasserstein sensitivity, measurability, loss differentiation and gradient-\allowbreak{}integrability premises, the variational nonexpansive projection, the source step-\allowbreak{}size and small-\allowbreak{}sensitivity inequalities, the exact empirical REGD recurrence, a stable point, and the numerical contraction, absorption, budget, radius, and initial-\allowbreak{}distance conditions named by that corrected target.\allowbreak{} The Spec quantifies over every such carrier and chooses one entry iteration outside the event;\allowbreak{} the displayed trace law samples each round's batch conditionally IID from the data law at the history-\allowbreak{}dependent deployment, and the event requires every iterate at or after entry to lie within radius with measure at least 1 -\allowbreak{} p.\allowbreak{} The broader p > 0 and radius domains retain all source-\allowbreak{}relevant nontrivial cases and only add cases to the universal result.\allowbreak{}
\par}
\reviewmatch{Matches formalized target}
\noindent\textbf{Reviewer annotation}\par
\reviewerbox

\end{document}
